% !TEX program = xelatex
% ============================================================================
% One Math Book — 大学数学，本科一年级
% Simplified Chinese edition of one_math_book_3_university_year_1.tex.
%
% Build:  latexmk one_math_book_3_university_year_1_zh.tex
% Output: build/one_math_book_3_university_year_1_zh.pdf
% ============================================================================
\documentclass[11pt,a4paper,twoside,openany]{book}

\newcommand{\booklang}{zh}
\usepackage{styles/onemath}
% Give TeX room around unbreakable Latin and math runs in Chinese lines.
\setlength{\emergencystretch}{3em}

% Brand the running header: rosette + ONE-COURSE.COM (linking to the site)
% above the chapter mark and the book title, on even pages; page number at the
% foot. See \ombrandheader in styles/onemath.sty.
\ombrandheader

% Print, at the end of every exercise and weekend problem, a hyperlinked
% "Solution p. N" pointer to its solution in the end-of-book appendix.
\omsolutionlinks

\newcommand{\bookdate}{\today}
\newcommand{\bookversion}{}
\IfFileExists{version.tex}{\input{version.tex}}{}

\newcommand{\bookline}{第 3 册：大学数学 -- 本科一年级}

\title{One Math Book}
\author{One Math Book 贡献者}
\date{\bookdate}

\begin{document}

\frontmatter
% One Course cover (mockup 81): pale blue-gray header band with the rosette
% sitting on its lower edge, ONE-COURSE.COM wordmark, four-color quad lines.
% Drawn as a full-page overlay: "remember picture" needs two LaTeX passes
% (latexmk handles the reruns). Coordinates in cm from the page's south west
% corner; the page is A4, 21 x 29.7.
%
% Text strings come from styles/lang/<lang>.tex (\bookmaintitle, …) and from
% the book entry file (\bookline).
\begin{titlepage}
  \begin{tikzpicture}[remember picture, overlay]
    \begin{scope}[shift={(current page.south west)}]
      % header band and rosette astride its lower edge
      \fill[ocPaleBlue] (0,24.6) rectangle (21,29.7);
      \ocQuadLine{0}{21}{24.36}{0.24}
      \ocRosette{10.5}{24.6}{2.4}{58}
      % wordmark, linking out to the site (\omSiteURL, styles/onemath.sty).
      % \href would recolor it (hyperref's urlcolor), so the ink is restated
      % inside the link.
      \node[text=ocInk, font=\sffamily\bfseries\LARGE] at (10.5,20.6)
        {\href{\omSiteURL}{\textcolor{ocInk}{\textls[210]{ONE-COURSE.COM}}}};
      % title block
      \node[text=ocDarkBlue, font=\fontsize{40}{44}\selectfont\bfseries]
        at (10.5,16.4) {\bookmaintitle};
      \node[text=ocInk, font=\Large\itshape] at (10.5,14.75)
        {\booktagline};
      \node[text=ocInk, font=\normalsize, align=center, text width=13cm]
        at (10.5,13.4)
        {\bookdescription};
      \ocQuadLine{6.6}{14.4}{12.0}{0.22}
      % this volume's name (defined by each book's entry file)
      \ifdefined\bookline
        \node[text=ocDarkBlue, font=\Large\bfseries] at (10.5,11.15)
          {\bookline};
      \fi
      % credits, version and date
      \node[text=ocInk, font=\normalsize] at (10.5,6.2)
        {\bookcontributors};
      \node[text=ocInk, font=\small] at (10.5,5.55)
        {Benjamin Virrion \;\textperiodcentered\; One Course};
      \node[text=ocInk, font=\small] at (10.5,4.7)
        {\ifdefempty{\bookversion}{}%
          {\bookversionlabel\ \bookversion\;\textperiodcentered\;}\bookdate};
      % bottom band
      \ocQuadLine{0}{21}{0}{1.0}
    \end{scope}
  \end{tikzpicture}
  \null
\end{titlepage}

% Reference page (verso of the title page): where the book lives and its licence.
\thispagestyle{empty}
\vspace*{\fill}
\begin{center}
\omcolophonsource\\[2pt]
\url{https://github.com/bvirrion/one-math-book}

\bigskip
\omcolophonlicense\\[2pt]
\url{https://creativecommons.org/licenses/by-nc-sa/4.0/}\\[6pt]
\omcolophoncredit
\end{center}
\vspace*{\fill}
\clearpage

\chapter{前言}

本书的愿望很简单：用一部前后连贯的数学书，涵盖从幼儿园到本科毕业的全部内容，
为世界各地的读者而写。

本书力求简洁而严谨：课程由定义、例子、命题和定理构成，只要在相应阶段可以理解，
就给出证明，随后是精心挑选的练习。不加证明给出的结果都明确标为“承认”。
所有练习的完整解答集中在书末。

每个学年构成本书的一个部分，每个部分再分为若干章。本项目目前涵盖一年级至十二年级
（6--18 岁），从小学一年级到中学最后一年，其中最后两年达到数学强化方向的水平。
学年越早，面向的读者越年幼：插图更多，演算步骤更详细，练习也更平缓。

\section*{如何阅读本书}

各陈述在每章内编号，并以颜色区分：定义为蓝色，定理为红色，命题及其同类为橙色，
方法（处理常见任务的简短步骤）为绿色。练习按难度分级，从直接应用的（$\star$）
到较难的（$\star\star\star$）。

\section*{参与贡献}

本书的源代码存放在一个公开的 git 仓库中：
\begin{center}
\url{https://github.com/bvirrion/one-math-book}
\end{center}
欢迎各种贡献，包括新的章节、勘误、更好的证明和更多的练习；
详见仓库根目录下的 \texttt{CONTRIBUTING.md}。

\tableofcontents

\mainmatter

% ============================================================================
% Bachelor Year 1 (age 18–19) — first year after secondary school.
% Contents follow the old French MPSI program (classes préparatoires,
% first year), written in English with international terminology. At this
% level essentially every theorem is proved; \admitted is reserved for
% genuinely out-of-scope results (construction of R, fundamental theorem
% of algebra, ...).
%
% One chapter per file. Solutions live in solutions/, one file per chapter,
% same numbering. Further chapters are added here as they are written.
% ============================================================================
\part{Bachelor Year 1 (Age 18--19)}

% Foundations and algebra
\chapter{Logic, Sets and Maps}\label{ch:b1:logic}

Up to now, proofs have been carried out with an informal but honest idea
of what ``proving'' means. This first chapter of undergraduate
mathematics makes the rules of the game explicit: what a mathematical
\omterm{def:b1:logic:statement}{statement} is, how the connectives and quantifiers combine \omterm{def:b1:logic:statement}{statements},
which moves are legal in a proof --- and then builds, on this footing,
the two universal languages of mathematics: \omterm{def:b1:logic:sets}{sets} and \omterm{def:b1:logic:map}{maps}.

\section{Statements and connectives}

\begin{definition}[Statement, connectives]\label{def:b1:logic:statement}
A \emph{statement}\index{statement} (or \emph{proposition}) is a
sentence which is either true (T) or false (F) --- exactly one of the
two. From statements $P$ and $Q$ one forms:
\begin{itemize}
  \item the \emph{negation} $\lnot P$ (``not $P$''), true exactly when
        $P$ is false;
  \item the \emph{conjunction} $P \land Q$ (``$P$ and $Q$''), true
        exactly when both are true;
  \item the \emph{disjunction} $P \lor Q$ (``$P$ or $Q$''), true
        exactly when at least one is true (this ``or'' is inclusive);
  \item the \emph{implication} $P \implies Q$, false exactly when $P$
        is true and $Q$ is false;
  \item the \emph{equivalence} $P \iff Q$, true exactly when $P$ and
        $Q$ have the same truth value.
\end{itemize}
\end{definition}

\begin{remark}
The truth table of $P \implies Q$ deserves a pause: when $P$ is false,
$P \implies Q$ is \emph{true}, whatever $Q$ is. ``If $2 < 1$ then
$0 = 5$'' is a true implication. An implication asserts nothing about
what happens when its hypothesis fails.
\end{remark}

\begin{proposition}[Rules of computation on statements]\label{prop:b1:logic:rules}
For all \omterm{def:b1:logic:statement}{statements} $P$, $Q$, $R$:
\begin{enumerate}
  \item $\lnot(\lnot P) \iff P$;
  \item De Morgan's laws\index{De Morgan's laws}:
        $\lnot(P \land Q) \iff (\lnot P) \lor (\lnot Q)$ and
        $\lnot(P \lor Q) \iff (\lnot P) \land (\lnot Q)$;
  \item $(P \implies Q) \iff \bigl((\lnot P) \lor Q\bigr)$, hence
        $\lnot(P \implies Q) \iff P \land (\lnot Q)$;
  \item contraposition\index{contraposition}:
        $(P \implies Q) \iff \bigl((\lnot Q) \implies (\lnot P)\bigr)$;
  \item $(P \iff Q) \iff \bigl((P \implies Q) \land (Q \implies
        P)\bigr)$;
  \item distributivity: $P \land (Q \lor R) \iff (P \land Q) \lor
        (P \land R)$ and $P \lor (Q \land R) \iff (P \lor Q) \land
        (P \lor R)$.
\end{enumerate}
\end{proposition}

\begin{proof}
Each equivalence is checked by comparing truth tables: two compound
\omterm{def:b1:logic:statement}{statements} built from $P$, $Q$, $R$ are equivalent exactly when they
take the same truth value in each of the (four or eight) cases. Let us
display one table in full, for the first De Morgan law:
\begin{center}
\begin{tabular}{cc|c|c|cc|c}
$P$ & $Q$ & $P \land Q$ & $\lnot(P \land Q)$ & $\lnot P$ & $\lnot Q$ &
$(\lnot P) \lor (\lnot Q)$ \\
\hline
T & T & T & F & F & F & F \\
T & F & F & T & F & T & T \\
F & T & F & T & T & F & T \\
F & F & F & T & T & T & T \\
\end{tabular}
\end{center}
Columns $4$ and $7$ coincide, which proves the law. For contraposition
a verbal shortcut is quicker: $P \implies Q$ is false exactly in the
case ($P$ true, $Q$ false), and $(\lnot Q) \implies (\lnot P)$ is false
exactly in the case ($\lnot Q$ true, $\lnot P$ false), i.e.\ ($Q$
false, $P$ true) --- the same single case, so the two implications have
identical tables. The remaining rules are checked the same way; note
that (3) reduces every implication to a disjunction, so that (2)
mechanically produces the negation rule $\lnot(P \implies Q) \iff P
\land (\lnot Q)$: to contradict an implication one must exhibit a case
where the hypothesis holds and the conclusion fails.
\end{proof}

\section{Quantifiers}

\begin{definition}[Quantifiers]\label{def:b1:logic:quantifiers}
Let $P(x)$ be a property of an element $x$ of a \omterm{def:b1:logic:sets}{set} $E$.
\begin{itemize}
  \item $\forall x \in E,\ P(x)$ (``for all $x$ in $E$, $P(x)$'') is
        true when every element of $E$ satisfies $P$;
  \item $\exists x \in E,\ P(x)$ (``there exists $x$ in $E$ such that
        $P(x)$'') is true when at least one element of $E$ satisfies
        $P$.
\end{itemize}
One writes $\exists!$ for ``there exists a unique''.
\end{definition}

\begin{proposition}[Negation of quantifiers]\label{prop:b1:logic:negquant}
\[
\lnot\bigl(\forall x \in E,\ P(x)\bigr) \iff
\exists x \in E,\ \lnot P(x),
\qquad
\lnot\bigl(\exists x \in E,\ P(x)\bigr) \iff
\forall x \in E,\ \lnot P(x).
\]
\end{proposition}

\begin{proof}
Let us argue the first equivalence in both directions; the second
is symmetric. If $\forall x \in E,\ P(x)$ is false, then not every
element satisfies $P$: the \omterm{def:b1:logic:sets}{set} $A = \{x \in E : \lnot P(x)\}$
cannot be empty, and any of its elements witnesses $\exists x \in
E,\ \lnot P(x)$. Conversely, if some $x_0 \in E$ satisfies $\lnot
P(x_0)$, then $x_0$ is a counterexample and the universal \omterm{def:b1:logic:statement}{statement}
fails. For the second rule: ``no $x$ satisfies $P$'' means the \omterm{def:b1:logic:sets}{set}
$\{x : P(x)\}$ is empty, i.e.\ every $x$ lies in its complement
$A$. Applied in cascade to a nested prefix of quantifiers, the two
rules give the mechanical procedure of \cref{ex:b1:logic:limit}:
the negation walks left to right, flipping every $\forall$ into
$\exists$ and every $\exists$ into $\forall$, and finally negates
the innermost predicate.
\end{proof}

\begin{example}[Negating everyday mathematical sentences]\label{ex:b1:logic:negpractice}
Let $f \colon \R \to \R$. The sentence ``$f$ is increasing'' reads
\[
\forall x \in \R,\ \forall y \in \R,\quad
x \leq y \implies f(x) \leq f(y) ,
\]
and its negation, by \cref{prop:b1:logic:negquant} plus the rule
$\lnot(P \implies Q) \iff P \land \lnot Q$:
\[
\exists x \in \R,\ \exists y \in \R,\quad
x \leq y \ \text{ and }\ f(x) > f(y) :
\]
a single witnessing pair suffices. Likewise ``$f$ is bounded'' is
$\exists M \in \R,\ \forall x \in \R,\ \abs{f(x)} \leq M$, with
negation
\[
\forall M \in \R,\ \exists x \in \R,\quad \abs{f(x)} > M :
\]
\emph{whatever} bound is proposed, some point beats it. The
insight: a correct negation never contains ``not'' applied to a
quantified block --- it is a new positive \omterm{def:b1:logic:statement}{statement}, in which the
roles are exchanged: one now produces the witnesses one previously
received.
\end{example}

\begin{example}[Order of quantifiers]\label{ex:b1:logic:order}
The \omterm{def:b1:logic:order}{order} of different quantifiers matters:
\[
\forall x \in \R,\ \exists y \in \R,\ y > x
\quad\text{is true (take } y = x+1\text{),}
\]
\[
\exists y \in \R,\ \forall x \in \R,\ y > x
\quad\text{is false (no real number exceeds all reals).}
\]
In the first \omterm{def:b1:logic:statement}{statement} $y$ may depend on $x$; in the second, one single
$y$ must work for all $x$. Two identical quantifiers, on the other
hand, always commute.
\end{example}

\begin{example}[Reading a definition with three quantifiers]\label{ex:b1:logic:limit}
The sentence ``the sequence $(u_n)$ converges to $\ell$'' will be
written in \cref{ch:b1:seq} as
\[
\forall \varepsilon > 0,\ \exists N \in \N,\ \forall n \geq N,\quad
\abs{u_n - \ell} \leq \varepsilon .
\]
Its negation, by \cref{prop:b1:logic:negquant} applied three times, is
\[
\exists \varepsilon > 0,\ \forall N \in \N,\ \exists n \geq N,\quad
\abs{u_n - \ell} > \varepsilon .
\]
Being able to negate such sentences mechanically, without thinking
about what they mean, is a genuine skill: it separates the logical work
from the mathematical work.
\end{example}

\section{Proof techniques}

\begin{method}[The standard proof patterns]\label{met:b1:logic:proofs}
To prove\dots
\begin{enumerate}
  \item \emph{an implication $P \implies Q$ directly}: assume $P$,
        deduce $Q$;
  \item \emph{by contraposition}: assume $\lnot Q$, deduce $\lnot P$
        --- valid by \cref{prop:b1:logic:rules}~(4);
  \item \emph{by contradiction}: assume the \omterm{def:b1:logic:statement}{statement} is false, derive
        a contradiction;
  \item \emph{an equivalence}: prove both implications separately (or
        chain known equivalences);
  \item \emph{a ``for all'' \omterm{def:b1:logic:statement}{statement}}: pick an \emph{arbitrary} $x$
        in $E$ (``let $x \in E$'') and prove $P(x)$;
  \item \emph{a ``there exists'' \omterm{def:b1:logic:statement}{statement}}: exhibit a witness, or
        prove existence indirectly;
  \item \emph{by induction}: see \cref{thm:b1:logic:induction}.
\end{enumerate}
When proving a \omterm{def:b1:logic:statement}{statement} about a well-chosen but arbitrary element,
never give the element extra properties: ``let $x \in \R$'' followed by
``since $x > 0$\dots'' proves nothing about negative $x$.
\end{method}

\begin{remark}[Common pitfalls in proofs]\label{rem:b1:logic:pitfalls}
Four classical traps, all worth naming once.
\begin{enumerate}
  \item \emph{Converse instead of contrapositive.} $Q \implies P$
        is \emph{not} equivalent to $P \implies Q$; only $\lnot Q
        \implies \lnot P$ is. ``If it rains, the street is wet''
        does not entitle one to conclude rain from a wet street.
  \item \emph{Proving an equivalence by one implication.} An
        ``iff'' claim is two theorems; announce which direction is
        being proved, and prove both. Chains of $\iff$ are legal
        only if \emph{every} link is genuinely reversible ---
        squaring an equation, for instance, is not.
  \item \emph{Backwards proofs.} Starting from the desired
        conclusion and deducing a true \omterm{def:b1:logic:statement}{statement} proves nothing
        (from $-1 = 1$ one derives the true $1 = 1$ by squaring).
        A computation may be \emph{discovered} backwards, but it
        must be \emph{written} forwards, or with explicit
        equivalences.
  \item \emph{Fixed witness versus arbitrary element.} To prove
        $\exists x,\ P(x)$, one may exhibit a single cleverly
        chosen $x$; to prove $\forall x,\ P(x)$, the chosen $x$
        must remain arbitrary. Mixing the two --- checking a
        universal claim on an example --- is the most common
        error in beginners' copies.
\end{enumerate}
\end{remark}

\begin{example}[Contraposition and contradiction at work]\label{ex:b1:logic:sqrt2}
\emph{For $n \in \N$: if $n^2$ is even then $n$ is even.} By
contraposition: if $n$ is odd, $n = 2k+1$, then
$n^2 = 4k^2 + 4k + 1$ is odd.

\emph{$\sqrt 2$ is irrational.} By contradiction: suppose
$\sqrt 2 = p/q$ with $p, q \in \N^*$ and the fraction in lowest terms.
Then $p^2 = 2q^2$ is even, so $p$ is even (previous point), $p = 2r$;
then $q^2 = 2r^2$ is even, so $q$ is even --- contradicting lowest
terms.
\end{example}

\begin{theorem}[Induction]\label{thm:b1:logic:induction}
Let $P(n)$ be a property of the integer $n$. If
\begin{enumerate}
  \item $P(0)$ is true, and
  \item for all $n \in \N$, $P(n) \implies P(n+1)$,
\end{enumerate}
then $P(n)$ is true for all $n \in \N$.

\emph{Strong induction:} conclusion unchanged if (2) is replaced by:
for all $n$, $\bigl(P(0) \land \dots \land P(n)\bigr) \implies
P(n+1)$.
\end{theorem}

\begin{proof}
This is a property of $\N$ itself, equivalent to: \emph{every nonempty
subset of $\N$ has a least element} (which we take as known). Indeed,
suppose (1) and (2) hold and let $A = \{n \in \N : P(n)
\text{ false}\}$. If $A \neq \emptyset$, it has a least element $m$;
$m \neq 0$ by (1); then $m - 1 \notin A$, so $P(m-1)$ holds, and (2)
gives $P(m)$ --- contradiction. So $A = \emptyset$. For strong
induction, apply the same argument: $P(0), \dots, P(m-1)$ all hold
since $m$ is least in $A$.
\end{proof}

\begin{example}[Proving unique existence]\label{ex:b1:logic:uniqueexistence}
A \omterm{def:b1:logic:statement}{statement} $\exists!\,x,\ P(x)$ is \emph{two} \omterm{def:b1:logic:statement}{statements}, proved
separately: existence (exhibit or construct some $x_0$ with
$P(x_0)$) and uniqueness (assume $P(x)$ and $P(x')$, deduce $x =
x'$). Sample: \emph{there is a unique real $x$ with $x^3 + x =
2$.} Existence: $x_0 = 1$ works, since $1 + 1 = 2$. Uniqueness:
if $x^3 + x = x'^3 + x'$, then
\[
0 = (x^3 - x'^3) + (x - x')
= (x - x')\,\bigl(x^2 + xx' + x'^2 + 1\bigr),
\]
and the second factor is positive (it equals $\bigl(x +
\tfrac{x'}2\bigr)^2 + \tfrac34 x'^2 + 1 \geq 1$), so $x = x'$.
Note the division of labor: existence used a lucky guess,
uniqueness used algebra valid for \emph{arbitrary} solutions ---
neither argument does the other's job, and forgetting the second
half is a standing temptation once a solution has been found.
\end{example}

\begin{example}\label{ex:b1:logic:induction}
For all $n \in \N^*$: $\;\sum_{k=1}^n k = \frac{n(n+1)}{2}$. Base case
$n = 1$: both sides equal $1$. Step: assuming the formula for $n$,
\[
\sum_{k=1}^{n+1} k = \frac{n(n+1)}{2} + (n+1)
= (n+1)\Bigl(\frac n2 + 1\Bigr) = \frac{(n+1)(n+2)}{2}. \qedhere
\]
\end{example}

\begin{example}[Strong induction at work]\label{ex:b1:logic:strongind}
\emph{Every integer $n \geq 2$ is a product of prime numbers} (a prime
being an integer $\geq 2$ whose only divisors $\geq 1$ are $1$ and
itself; primes are studied for their own sake in \cref{ch:b1:arith}).
Ordinary induction is helpless here: knowing that $95 = 5 \times 19$
factors says nothing about $96$. Strong induction fits exactly. Base
case: $2$ is prime, hence a (one-factor) product of primes. Step: let
$n \geq 2$ and assume every integer $m$ with $2 \leq m \leq n$ is a
product of primes. If $n + 1$ is prime, done. Otherwise $n + 1 = ab$
with $2 \leq a, b \leq n$; by the strong hypothesis both $a$ and $b$
are products of primes, hence so is $n + 1$. The insight: strong
induction is the right tool whenever the ``reason'' for $P(n+1)$ lives
at some unpredictable earlier rank, not at rank $n$.
\end{example}

\section{Sets}

\begin{definition}[Set operations]\label{def:b1:logic:sets}
We take the notion of \emph{set}\index{set} and the membership relation
$x \in E$ as primitive. For sets $A, B$ inside an ambient set $E$:
\begin{itemize}
  \item \emph{inclusion}: $A \subseteq B$ when $\forall x,\ x \in A
        \implies x \in B$; equality $A = B$ when $A \subseteq B$ and
        $B \subseteq A$;
  \item \emph{union} $A \cup B$, \emph{intersection} $A \cap B$,
        \emph{difference} $A \setminus B = \{x \in A : x \notin B\}$,
        \emph{complement} $\overline{A} = E \setminus A$;
  \item the \emph{empty set} $\emptyset$, contained in every set;
  \item the \emph{power set}\index{power set} $\mathcal{P}(E)$: the set
        of all subsets of $E$;
  \item the \emph{product} $E \times F$: the set of ordered pairs
        $(x, y)$ with $x \in E$, $y \in F$.
\end{itemize}
\end{definition}

\begin{example}[Getting used to the power set]\label{ex:b1:logic:powerset}
For $E = \{a, b\}$:
\[
\mathcal P(E) = \bigl\{\, \emptyset,\ \{a\},\ \{b\},\ \{a, b\}
\,\bigr\},
\]
four elements --- and note the type discipline: $a \in E$ but
$\{a\} \in \mathcal P(E)$; the \omterm{def:b1:logic:statement}{statements} $a \in \mathcal P(E)$
and $\{a\} \subseteq \mathcal P(E)$ are both false as written
(the second would require $a$ to be a \emph{subset} of $E$).
Iterating from nothing: $\mathcal P(\emptyset) = \{\emptyset\}$
has one element, $\mathcal P(\mathcal P(\emptyset)) = \{\emptyset,
\{\emptyset\}\}$ has two, the next has four --- \omterm{def:b1:logic:sets}{sets} of \omterm{def:b1:logic:sets}{sets} are
ordinary \omterm{def:b1:logic:sets}{sets}, and \cref{ch:b1:counting} will confirm the doubling
pattern: $\abs{\mathcal P(E)} = 2^{\abs E}$. Keeping the levels
($x$, $\{x\}$, $\{\{x\}\}$) straight is half the battle in
exercises like \cref{exo:b1:logic:11,exo:b1:logic:12}.
\end{example}

\begin{proposition}[Set algebra]\label{prop:b1:logic:setalgebra}
For subsets $A, B, C$ of $E$:
\begin{enumerate}
  \item $A \cap (B \cup C) = (A \cap B) \cup (A \cap C)$ and
        $A \cup (B \cap C) = (A \cup B) \cap (A \cup C)$;
  \item De Morgan: $\overline{A \cup B} = \overline{A} \cap
        \overline{B}$ and $\overline{A \cap B} = \overline{A} \cup
        \overline{B}$;
  \item $A \subseteq B \iff \overline{B} \subseteq \overline{A}$.
\end{enumerate}
\end{proposition}

\begin{proof}
Each identity translates a rule of \cref{prop:b1:logic:rules} through
the dictionary ($\in A$ or not) $\leftrightarrow$ (\omterm{def:b1:logic:statement}{statement} true or
false): e.g.\ $x \in \overline{A \cup B} \iff \lnot(x \in A \lor x \in
B) \iff (x \notin A) \land (x \notin B) \iff x \in \overline{A} \cap
\overline{B}$. Point (3) is contraposition. As a second sample, the
first distributivity law in full:
\[
x \in A \cap (B \cup C)
\iff (x \in A) \land \bigl(x \in B \lor x \in C\bigr)
\iff \bigl(x \in A \land x \in B\bigr) \lor
\bigl(x \in A \land x \in C\bigr),
\]
by the distributivity of \cref{prop:b1:logic:rules}~(6), and the
last \omterm{def:b1:logic:statement}{statement} reads $x \in (A \cap B) \cup (A \cap C)$. Every \omterm{def:b1:logic:sets}{set}
identity of this kind is provable by this one mechanical
translation --- which is why none of them needs memorizing.
\end{proof}

\begin{method}[Proving set equalities]\label{met:b1:logic:doubleinc}
To prove $A = B$, prove the two inclusions: let $x \in A$, show $x \in
B$; then let $x \in B$, show $x \in A$. Alternatively, chain
equivalences $x \in A \iff \dots \iff x \in B$ when each step really is
an equivalence.
\end{method}

\begin{omfigure}
\begin{tikzpicture}[scale=0.9]
  \draw[thick] (-2.6,-1.9) rectangle (2.6,1.9);
  \node[below right] at (-2.6,1.9) {$E$};
  \begin{scope}
    \clip (-2.6,-1.9) rectangle (2.6,1.9);
    \fill[omDef!20] (-2.6,-1.9) rectangle (2.6,1.9);
    \fill[white] (-0.7,0) circle (1.1);
    \fill[white] (0.7,0) circle (1.1);
  \end{scope}
  \draw[omDef, thick] (-0.7,0) circle (1.1);
  \draw[omDef, thick] (0.7,0) circle (1.1);
  \node[omDef] at (-1.35,0.5) {$A$};
  \node[omDef] at (1.35,0.5) {$B$};
  \begin{scope}[xshift=6.6cm]
    \draw[thick] (-2.6,-1.9) rectangle (2.6,1.9);
    \node[below right] at (-2.6,1.9) {$E$};
    \begin{scope}
      \clip (-2.6,-1.9) rectangle (2.6,1.9);
      \fill[omThm!20] (-2.6,-1.9) rectangle (2.6,1.9);
      \begin{scope}
        \clip (-0.7,0) circle (1.1);
        \fill[white] (0.7,0) circle (1.1);
      \end{scope}
    \end{scope}
    \draw[omThm, thick] (-0.7,0) circle (1.1);
    \draw[omThm, thick] (0.7,0) circle (1.1);
    \node[omThm] at (-1.35,0.5) {$A$};
    \node[omThm] at (1.35,0.5) {$B$};
  \end{scope}
\end{tikzpicture}

{\small De Morgan's laws in pictures: the shaded region on the left
is $\overline{A \cup B} = \overline A \cap \overline B$ (everything
outside both disks); on the right, $\overline{A \cap B} = \overline
A \cup \overline B$ (everything except the lens-shaped overlap). A
diagram is not a proof, but it makes the element-chasing proof of
\cref{prop:b1:logic:setalgebra} impossible to misremember.}
\end{omfigure}

\section{Maps}

\begin{definition}[Map, image, preimage]\label{def:b1:logic:map}
A \emph{map}\index{map} (or \emph{function}) $f \colon E \to F$
assigns to every element $x$ of the \omterm{def:b1:logic:sets}{set} $E$ (the \emph{domain}) exactly
one element $f(x)$ of the \omterm{def:b1:logic:sets}{set} $F$ (the \emph{codomain}). For $A
\subseteq E$ and $B \subseteq F$:
\[
f(A) = \{f(x) : x \in A\} \subseteq F,
\qquad
f^{-1}(B) = \{x \in E : f(x) \in B\} \subseteq E
\]
are the \emph{direct image} of $A$ and the \emph{preimage}
\index{preimage} of $B$. The \emph{composition} of $f \colon E \to F$
and $g \colon F \to G$ is $g \circ f \colon E \to G$, $x \mapsto
g(f(x))$.
\end{definition}

\begin{remark}
The notation $f^{-1}(B)$ does \emph{not} presuppose an inverse \omterm{def:b1:logic:map}{map}:
$f^{-1}(B)$ is defined for every $f$. \omterm{def:b1:logic:map}{Preimages} behave better than
images: $f^{-1}$ preserves unions, intersections and complements,
whereas $f(A \cap A') \subseteq f(A) \cap f(A')$ can be strict
(\cref{exo:b1:logic:8}).
\end{remark}

\begin{example}[Computing images and preimages]\label{ex:b1:logic:images}
Let $f \colon \R \to \R$, $x \mapsto x^2$. Then:
\[
f\bigl(\intcc{-1}{2}\bigr) = \intcc04, \qquad
f^{-1}\bigl(\intcc14\bigr) = \intcc{-2}{-1} \cup \intcc12, \qquad
f^{-1}(\{-1\}) = \emptyset .
\]
For the first: every $x \in \intcc{-1}2$ has $x^2 \in \intcc04$, and
every $y \in \intcc04$ is attained as $y = (\sqrt y)^2$ with $\sqrt y
\in \intcc02 \subseteq \intcc{-1}2$ --- note that the image is
\emph{not} $\intcc14 = \{(-1)^2, 2^2\}$: images of intervals are not
computed from endpoints alone. For the second: $1 \leq x^2 \leq 4 \iff
1 \leq \abs x \leq 2$, which splits into two pieces. The third
illustrates that a \omterm{def:b1:logic:map}{preimage} may be empty --- $f^{-1}(B)$ always makes
sense, however small the intersection of $B$ with the image. Finally
observe on this example the strictness phenomenon of the remark above:
with $A = \intcc{-1}0$ and $A' = \intcc01$, one has $f(A \cap A') =
f(\{0\}) = \{0\}$, while $f(A) \cap f(A') = \intcc01$.
\end{example}

\begin{definition}[Injective, surjective, bijective]\label{def:b1:logic:inj}
A \omterm{def:b1:logic:map}{map} $f \colon E \to F$ is:
\begin{itemize}
  \item \emph{injective}\index{injective} when distinct elements have
        distinct images: $\forall x, x' \in E,\ f(x) = f(x') \implies
        x = x'$;
  \item \emph{surjective}\index{surjective} when every element of $F$
        is attained: $\forall y \in F,\ \exists x \in E,\ f(x) = y$;
  \item \emph{bijective}\index{bijective} when it is both, i.e.\ every
        $y \in F$ has exactly one \omterm{def:b1:logic:map}{preimage}.
\end{itemize}
\end{definition}

\begin{theorem}[Inverse map]\label{thm:b1:logic:inverse}
A \omterm{def:b1:logic:map}{map} $f \colon E \to F$ is \omterm{def:b1:logic:inj}{bijective} if and only if there is a \omterm{def:b1:logic:map}{map}
$g \colon F \to E$ with $g \circ f = \mathrm{id}_E$ and $f \circ g =
\mathrm{id}_F$. In that case $g$ is unique; it is written $f^{-1}$ and
called the \emph{inverse} of $f$, and $f^{-1}$ is itself \omterm{def:b1:logic:inj}{bijective} with
$(f^{-1})^{-1} = f$.
\end{theorem}

\begin{proof}
($\Rightarrow$) If $f$ is \omterm{def:b1:logic:inj}{bijective}, every $y \in F$ has a unique
\omterm{def:b1:logic:map}{preimage}; define $g(y)$ to be that \omterm{def:b1:logic:map}{preimage}. Then $f(g(y)) = y$ by
construction, and $g(f(x)) = x$ because $x$ is \emph{the} \omterm{def:b1:logic:map}{preimage} of
$f(x)$.

($\Leftarrow$) Suppose such a $g$ exists. If $f(x) = f(x')$, applying
$g$ gives $x = x'$: $f$ is \omterm{def:b1:logic:inj}{injective}. For $y \in F$, $x = g(y)$
satisfies $f(x) = y$: $f$ is \omterm{def:b1:logic:inj}{surjective}.

Uniqueness: if $g$ and $h$ both work, then $g = g \circ \mathrm{id}_F
= g \circ (f \circ h) = (g \circ f) \circ h = h$. Finally the pair of
identities is symmetric in $f$ and $g$, so $g = f^{-1}$ is \omterm{def:b1:logic:inj}{bijective}
with inverse $f$.
\end{proof}

\begin{example}[Computing an inverse in practice]\label{ex:b1:logic:inversecomp}
Let $f \colon \R \to \intoo0{+\infty}$, $f(x) = \eu^{2x+1}$. To
invert, solve $y = f(x)$ for a given $y > 0$:
\[
y = \eu^{2x+1} \iff \ln y = 2x + 1 \iff x = \frac{\ln y - 1}2 ,
\]
each step being reversible on the announced domains. The
computation delivers everything at once: for each $y$ in the
codomain there is exactly one solution $x$, so $f$ is \omterm{def:b1:logic:inj}{bijective},
and
\[
f^{-1} \colon \intoo0{+\infty} \to \R,
\qquad
f^{-1}(y) = \frac{\ln y - 1}2 .
\]
A quick check of both compositions ($f^{-1}(f(x)) = \frac{(2x+1) -
1}2 = x$ and $f(f^{-1}(y)) = \eu^{\ln y} = y$) confirms
\cref{thm:b1:logic:inverse}'s criterion. The insight: ``solve for
$x$ and watch the equivalences'' is simultaneously the existence
proof, the uniqueness proof, and the formula --- but it only works
if the codomain was announced correctly ($f$ is \emph{not}
\omterm{def:b1:logic:inj}{surjective} onto $\R$).
\end{example}

\begin{proposition}[Composition and the three properties]\label{prop:b1:logic:comp}
Let $f \colon E \to F$ and $g \colon F \to G$.
\begin{enumerate}
  \item If $f$ and $g$ are \omterm{def:b1:logic:inj}{injective} (resp.\ \omterm{def:b1:logic:inj}{surjective}, \omterm{def:b1:logic:inj}{bijective}), so
        is $g \circ f$; and then $(g \circ f)^{-1} = f^{-1} \circ
        g^{-1}$ in the \omterm{def:b1:logic:inj}{bijective} case.
  \item If $g \circ f$ is \omterm{def:b1:logic:inj}{injective}, then $f$ is \omterm{def:b1:logic:inj}{injective}. If
        $g \circ f$ is \omterm{def:b1:logic:inj}{surjective}, then $g$ is \omterm{def:b1:logic:inj}{surjective}.
\end{enumerate}
\end{proposition}

\begin{proof}
(1) If $g(f(x)) = g(f(x'))$, injectivity of $g$ gives $f(x) = f(x')$,
then injectivity of $f$ gives $x = x'$. If $z \in G$, surjectivity of
$g$ gives $y$ with $g(y) = z$, then surjectivity of $f$ gives $x$ with
$f(x) = y$, so $g(f(x)) = z$. In the \omterm{def:b1:logic:inj}{bijective} case one checks directly
that $f^{-1} \circ g^{-1}$ is a two-sided inverse of $g \circ f$, and
uniqueness in \cref{thm:b1:logic:inverse} concludes.

(2) If $f(x) = f(x')$ then $g(f(x)) = g(f(x'))$, and injectivity of
$g \circ f$ gives $x = x'$. If $z \in G$, surjectivity of $g \circ f$
gives $x$ with $g(f(x)) = z$: then $y = f(x)$ satisfies $g(y) = z$.
\end{proof}

\begin{example}[Point (2) is sharp]\label{ex:b1:logic:compsharp}
In \cref{prop:b1:logic:comp}~(2), one cannot upgrade the
conclusions: $g \circ f$ \omterm{def:b1:logic:inj}{bijective} does \emph{not} force $f$
\omterm{def:b1:logic:inj}{surjective} or $g$ \omterm{def:b1:logic:inj}{injective}. Take $E = G = \{1\}$, $F = \{1,
2\}$, with $f(1) = 1$ and $g(1) = g(2) = 1$: then $g \circ f =
\mathrm{id}_E$ is \omterm{def:b1:logic:inj}{bijective}, yet $f$ misses the element $2$ and
$g$ collapses both elements. The moral is a precise bookkeeping
rule: composition information flows to the \emph{inner} \omterm{def:b1:logic:map}{map} for
injectivity and to the \emph{outer} \omterm{def:b1:logic:map}{map} for surjectivity, never
the other way. (\cref{exo:b1:logic:9} builds the same phenomenon
with infinite \omterm{def:b1:logic:sets}{sets}, where it is the engine behind one-sided
inverses.)
\end{example}

\begin{example}\label{ex:b1:logic:injexamples}
$f \colon \R \to \R$, $x \mapsto x^2$ is neither \omterm{def:b1:logic:inj}{injective} ($f(-1) =
f(1)$) nor \omterm{def:b1:logic:inj}{surjective} ($-1$ has no \omterm{def:b1:logic:map}{preimage}). Restricting the domain
and codomain, $f \colon \R_+ \to \R_+$, $x \mapsto x^2$ is \omterm{def:b1:logic:inj}{bijective},
with inverse $y \mapsto \sqrt y$. A \omterm{def:b1:logic:map}{map}'s injectivity or surjectivity
depends on the announced domain and codomain, not only on the formula.
\end{example}

\section{Relations}

\begin{definition}[Equivalence relation]\label{def:b1:logic:equiv}
A \emph{binary relation} $\mathcal{R}$ on a \omterm{def:b1:logic:sets}{set} $E$ is an
\emph{equivalence relation}\index{equivalence relation} when it is:
\emph{reflexive} ($x \mathbin{\mathcal{R}} x$ for all $x$),
\emph{symmetric} ($x \mathbin{\mathcal{R}} y \implies y
\mathbin{\mathcal{R}} x$) and \emph{transitive} ($x
\mathbin{\mathcal{R}} y$ and $y \mathbin{\mathcal{R}} z$ imply $x
\mathbin{\mathcal{R}} z$). The \emph{equivalence class} of $x$ is
$\mathrm{cl}(x) = \{y \in E : x \mathbin{\mathcal{R}} y\}$.
\end{definition}

\begin{example}[Checking the three axioms]\label{ex:b1:logic:fracpart}
On $\R$, declare $x \mathbin{\mathcal{R}} y$ when $x - y \in \Z$.
\emph{Reflexive:} $x - x = 0 \in \Z$. \emph{Symmetric:} if $x - y
\in \Z$ then $y - x = -(x - y) \in \Z$. \emph{Transitive:} if $x -
y \in \Z$ and $y - z \in \Z$, then $x - z = (x - y) + (y - z) \in
\Z$ (a sum of integers). So $\mathcal R$ is an \omterm{def:b1:logic:equiv}{equivalence
relation}, and $\mathrm{cl}(x) = x + \Z = \{x + k : k \in \Z\}$:
each class contains exactly one representative in $\intco01$, its
\emph{fractional part}. By contrast, the relation ``$\abs{x - y}
\leq 1$'' on $\R$ is reflexive and symmetric but \emph{not}
transitive ($0 \mathbin{\mathcal R} 1$ and $1 \mathbin{\mathcal R}
2$, yet $\abs{0 - 2} > 1$): closeness does not propagate, and no
\omterm{thm:b1:logic:partition}{partition} into classes exists --- a useful counterexample to keep
in mind when verifying the axioms feels routine.
\end{example}

\begin{theorem}[Classes form a partition]\label{thm:b1:logic:partition}
Let $\mathcal{R}$ be an \omterm{def:b1:logic:equiv}{equivalence relation} on $E$. Then the
equivalence classes are nonempty, pairwise disjoint or equal, and their
union is $E$: they form a \emph{partition}\index{partition} of $E$.
Conversely, every \omterm{thm:b1:logic:partition}{partition} of $E$ arises this way from exactly one
\omterm{def:b1:logic:equiv}{equivalence relation} (``being in the same piece'').
\end{theorem}

\begin{proof}
$x \in \mathrm{cl}(x)$ by reflexivity, so classes are nonempty with
union $E$. Suppose $\mathrm{cl}(x) \cap \mathrm{cl}(y) \neq \emptyset$,
say $z$ lies in both. Then $x \mathbin{\mathcal{R}} z$ and $y
\mathbin{\mathcal{R}} z$, so by symmetry and transitivity $x
\mathbin{\mathcal{R}} y$. Now for any $t \in \mathrm{cl}(y)$,
transitivity gives $t \in \mathrm{cl}(x)$, and symmetrically: the two
classes are equal. For the converse, let $(E_i)_{i \in I}$ be a
\omterm{thm:b1:logic:partition}{partition} of $E$ and define $x \mathbin{\mathcal S} y$ to mean
``some piece contains both $x$ and $y$''. \emph{Reflexive:} $x$
lies in some piece, which then contains $x$ twice over.
\emph{Symmetric:} the defining condition is symmetric in $x$ and
$y$. \emph{Transitive:} if $x, y \in E_i$ and $y, z \in E_j$, then
$y \in E_i \cap E_j$, so $E_i = E_j$ (distinct pieces are
disjoint) and $x, z$ share a piece. The $\mathcal S$-class of $x$
is exactly the piece containing $x$, so the classes are the given
pieces. Finally the relation is determined by its classes: two
\omterm{def:b1:logic:equiv}{equivalence relations} with the same classes relate the same pairs,
since each relates $x$ and $y$ exactly when $y$ belongs to the
class of $x$ --- whence the uniqueness claim.
\end{proof}

\begin{example}\label{ex:b1:logic:congruence}
On $\Z$, congruence modulo $n$ ($x \equiv y \pmod n$ when $n$ divides
$x - y$) is an \omterm{def:b1:logic:equiv}{equivalence relation}; its classes are the $n$ \omterm{def:b1:logic:sets}{sets} of
integers with a given remainder upon division by $n$. This example
becomes the ring $\Z/n\Z$ in \cref{ch:b1:structures}.
\end{example}

\begin{definition}[Order relation]\label{def:b1:logic:order}
A relation $\preceq$ on $E$ is an \emph{order}\index{order relation}
when it is reflexive, \emph{antisymmetric} ($x \preceq y$ and $y
\preceq x$ imply $x = y$) and transitive. The order is \emph{total}
when any two elements are comparable, \emph{partial} otherwise. An
element $M \in A \subseteq E$ is a \emph{greatest element} of $A$ when
$a \preceq M$ for all $a \in A$; greatest (and least) elements are
unique when they exist.
\end{definition}

\begin{example}\label{ex:b1:logic:orders}
$(\R, \leq)$ is totally ordered. $(\mathcal{P}(E), \subseteq)$ is
partially ordered as soon as $E$ has two elements: $\{a\}$ and
$\{b\}$ are not comparable. The subset $A = \{\{a\}, \{b\}\}$ of
$\mathcal{P}(\{a,b\})$ has no greatest element, yet has an upper bound
$\{a, b\}$: the distinction between greatest elements and upper bounds
returns, for $\R$, in \cref{ch:b1:reals}.
\end{example}

\begin{example}[Two orders on the grid $\N^2$]\label{ex:b1:logic:lexorder}
On pairs of naturals, compare componentwise: $(a, b) \preceq (a',
b')$ when $a \leq a'$ \emph{and} $b \leq b'$ (the \emph{product
\omterm{def:b1:logic:order}{order}}). This is an \omterm{def:b1:logic:order}{order} --- each axiom is inherited coordinate
by coordinate --- but a partial one: $(1, 3)$ and $(2, 0)$ are
incomparable. Now compare like a dictionary: $(a, b)
\preceq_{\mathrm{lex}} (a', b')$ when $a < a'$, or $a = a'$ and $b
\leq b'$ (the \emph{lexicographic \omterm{def:b1:logic:order}{order}}). Transitivity requires
a two-case check but holds, and any two pairs are now comparable:
the \omterm{def:b1:logic:order}{order} is total. The two \omterm{def:b1:logic:order}{orders} rank the same \omterm{def:b1:logic:sets}{set} differently
--- $(0, 100) \preceq_{\mathrm{lex}} (1, 0)$ though the product
\omterm{def:b1:logic:order}{order} says nothing --- a reminder that an \omterm{def:b1:logic:order}{order} is a structure
one \emph{chooses}, not a property of the \omterm{def:b1:logic:sets}{set}. Lexicographic
comparison is also the standard trick for turning several sorting
criteria into one.
\end{example}

\begin{remark}[Interlude: size as bijection]\label{rem:b1:logic:interlude}
A quiet theme of this chapter deserves a spotlight: bijections are
the mathematician's notion of ``same size''. For finite \omterm{def:b1:logic:sets}{sets} this
becomes the counting calculus of \cref{ch:b1:counting}, where
every formula is secretly a bijection; for infinite \omterm{def:b1:logic:sets}{sets} it
becomes the weekend problem below, where $\N$, $\Q$ and $\R$ turn
out to have genuinely different sizes. The same dictionary
reappears twice more in this volume in refined forms: sequences
(\cref{ch:b1:seq}) are nothing but \omterm{def:b1:logic:map}{maps} $\N \to \R$, so \omterm{def:b1:logic:statement}{statements}
about sequences are \omterm{def:b1:logic:statement}{statements} about a \omterm{def:b1:logic:sets}{set} of \omterm{def:b1:logic:map}{maps}; and linear
algebra will measure vector spaces not by bijections but by
\emph{linear} bijections, whose existence is governed by a single
number, the dimension (\cref{ch:b1:findim}). Whenever a new
``sameness'' appears --- equipotence, isomorphism of groups
(\cref{ch:b1:structures}), linear isomorphism --- the pattern of
\cref{thm:b1:logic:inverse} repeats: sameness is an invertible,
structure-respecting \omterm{def:b1:logic:map}{map}.
\end{remark}

\begin{remark}[Where this chapter is used]\label{rem:b1:logic:whereused}
Everywhere --- but a few places deserve flagging. The three-quantifier
gymnastics of \cref{ex:b1:logic:limit} is the daily bread of
\cref{ch:b1:seq,ch:b1:continuity}: every limit proof is a game played
against an arbitrary $\varepsilon$. Equivalence classes reappear as the
congruence classes of $\Z/n\Z$ in \cref{ch:b1:structures}, where the
\omterm{thm:b1:logic:partition}{partition} of \cref{thm:b1:logic:partition} acquires an algebraic
structure of its own. \omterm{def:b1:logic:order}{Order relations}, upper bounds and least upper
bounds become the axiomatic heart of $\R$ in \cref{ch:b1:reals}.
Injections, surjections and bijections return as the linear \omterm{def:b1:logic:map}{maps} of
\cref{ch:b1:linmaps}, where injectivity can be tested on a single
vector (the kernel); and the weekend problem below turns the bare
notion of bijection into a theory of the \emph{sizes of infinite
\omterm{def:b1:logic:sets}{sets}}, whose conclusions (countability of $\Q$, uncountability of
$\R$) resurface in \cref{ch:b1:reals,ch:b1:topology}.
\end{remark}

\section{Exercises}

\begin{exercise}[$\star$]\label{exo:b1:logic:1}
Write the negation of each \omterm{def:b1:logic:statement}{statement}, without using the word ``not'':
\begin{enumerate}
  \item $\forall x \in \R,\ \exists y \in \R,\ x + y > 0$;
  \item $\exists x \in \R,\ \forall y \in \R,\ xy = 0$;
  \item $\forall \varepsilon > 0,\ \exists \delta > 0,\ \forall x \in
        \R,\ \abs{x} \leq \delta \implies \abs{f(x)} \leq \varepsilon$
        (for a fixed \omterm{def:b1:logic:map}{map} $f \colon \R \to \R$).
\end{enumerate}
Then decide whether \omterm{def:b1:logic:statement}{statements} (1) and (2) are true.
\end{exercise}

\begin{exercise}[$\star$]\label{exo:b1:logic:2}
Let $P, Q$ be \omterm{def:b1:logic:statement}{statements}. Using truth tables, prove that
$\lnot(P \implies Q) \iff P \land (\lnot Q)$, and deduce the negation
of: ``if a function is differentiable then it is continuous''.
\end{exercise}

\begin{exercise}[$\star$]\label{exo:b1:logic:3}
Prove by contraposition: for $x \in \R$, if $x^3 + x \geq 2$ then
$x \geq 1$. Then prove by contradiction: there is no smallest strictly
positive real number.
\end{exercise}

\begin{exercise}[$\star$]\label{exo:b1:logic:4}
Prove by induction that for all $n \in \N$:
\begin{enumerate}
  \item $\sum_{k=0}^{n} 2^k = 2^{n+1} - 1$;
  \item $4^n + 5$ is divisible by $3$.
\end{enumerate}
\end{exercise}

\begin{exercise}[$\star$]\label{exo:b1:logic:5}
Find the flaw in the following ``proof'' that all pencils have the same
color. \emph{Let $P(n)$: ``in every \omterm{def:b1:logic:sets}{set} of $n$ pencils, all pencils
have the same color''. $P(1)$ is clear. Assume $P(n)$ and take $n+1$
pencils; removing the last one, the first $n$ share their color;
removing the first one, the last $n$ share their color; hence all
$n+1$ share their color.}
\end{exercise}

\begin{exercise}[$\star$]\label{exo:b1:logic:6}
Let $A, B, C$ be subsets of $E$. Prove:
\begin{enumerate}
  \item $A \setminus B = A \cap \overline{B}$;
  \item $(A \cup B) \setminus C = (A \setminus C) \cup (B \setminus
        C)$;
  \item $A \subseteq B \iff A \cup B = B \iff A \cap B = A$.
\end{enumerate}
\end{exercise}

\begin{exercise}[$\star\star$]\label{exo:b1:logic:7}
For each \omterm{def:b1:logic:map}{map}, decide (with proof) whether it is \omterm{def:b1:logic:inj}{injective}, \omterm{def:b1:logic:inj}{surjective},
\omterm{def:b1:logic:inj}{bijective}:
\begin{enumerate}
  \item $f \colon \N \to \N$, $n \mapsto n + 1$;
  \item $g \colon \Z \to \Z$, $n \mapsto n + 1$;
  \item $h \colon \R \setminus \{1\} \to \R$, $x \mapsto
        \frac{x+1}{x-1}$.
\end{enumerate}
For $h$, adjust the codomain to make it \omterm{def:b1:logic:inj}{bijective} and compute the
inverse.
\end{exercise}

\begin{exercise}[$\star\star$]\label{exo:b1:logic:8}
Let $f \colon E \to F$, let $A, A' \subseteq E$ and $B, B' \subseteq
F$.
\begin{enumerate}
  \item Prove $f^{-1}(B \cap B') = f^{-1}(B) \cap f^{-1}(B')$ and
        $f(A \cup A') = f(A) \cup f(A')$.
  \item Prove $f(A \cap A') \subseteq f(A) \cap f(A')$ and give an
        example where the inclusion is strict.
  \item Prove: $f$ is \omterm{def:b1:logic:inj}{injective} if and only if $f(A \cap A') = f(A)
        \cap f(A')$ for all $A, A'$.
\end{enumerate}
\end{exercise}

\begin{exercise}[$\star\star$]\label{exo:b1:logic:9}
Let $f \colon E \to F$ and $g \colon F \to E$ satisfy $g \circ f =
\mathrm{id}_E$. Prove that $f$ is \omterm{def:b1:logic:inj}{injective} and $g$ is \omterm{def:b1:logic:inj}{surjective}.
Give an example where neither $f$ nor $g$ is \omterm{def:b1:logic:inj}{bijective}.
\end{exercise}

\begin{exercise}[$\star\star$]\label{exo:b1:logic:10}
On $\R$, define $x \mathbin{\mathcal{R}} y \iff x^2 - y^2 = x - y$.
Prove that $\mathcal{R}$ is an \omterm{def:b1:logic:equiv}{equivalence relation} and describe the
\omterm{def:b1:logic:equiv}{equivalence class} of each real $x$. Which classes have exactly one
element?
\end{exercise}

\begin{exercise}[$\star\star\star$]\label{exo:b1:logic:11}
(Cantor) Let $E$ be a \omterm{def:b1:logic:sets}{set}. Prove that there is no surjection from $E$
onto $\mathcal{P}(E)$. \emph{Hint: given $f \colon E \to
\mathcal{P}(E)$, consider $D = \{x \in E : x \notin f(x)\}$.}
\end{exercise}

\begin{exercise}[$\star\star\star$]\label{exo:b1:logic:12}
Let $f \colon E \to F$ be a \omterm{def:b1:logic:map}{map}. Define $\Phi \colon \mathcal{P}(F)
\to \mathcal{P}(E)$ by $\Phi(B) = f^{-1}(B)$.
\begin{enumerate}
  \item Prove that $f$ is \omterm{def:b1:logic:inj}{surjective} if and only if $\Phi$ is
        \omterm{def:b1:logic:inj}{injective}.
  \item Prove that $f$ is \omterm{def:b1:logic:inj}{injective} if and only if $\Phi$ is
        \omterm{def:b1:logic:inj}{surjective}.
\end{enumerate}
\end{exercise}

\section{Problem: Comparing infinities}

\begin{problem}\label{pb:b1:logic:1}
When do two \omterm{def:b1:logic:sets}{sets} have ``the same number of elements''? Cantor's
answer --- when there is a bijection between them --- turns out to be
usable even for infinite \omterm{def:b1:logic:sets}{sets}, and it splits infinity into genuinely
different sizes. This problem builds the whole toolbox from the bare
definitions of this chapter: the Cantor--Schr\"oder--Bernstein theorem
(two injections manufacture a bijection), the countability of $\Q$,
the uncountability of $\R$ by the diagonal argument, and Cantor's
astonishing 1874 conclusion: \emph{\omterm{pb:b1:logic:1}{transcendental numbers} exist, and
massively so}, without exhibiting a single one.
\index{countable set}\index{Cantor--Schroder--Bernstein theorem}
Throughout, for \omterm{def:b1:logic:sets}{sets} $E$ and $F$, write $E \preceq F$ when there
exists an injection from $E$ into $F$, and $E \approx F$ (``$E$ and
$F$ are \emph{equipotent}'') when there exists a bijection from $E$
onto $F$.

\medskip
\noindent\textbf{Part I --- The vocabulary of comparison.}
\begin{enumerate}
  \item Show that $\approx$ behaves like an \omterm{def:b1:logic:equiv}{equivalence relation}:
        $E \approx E$; if $E \approx F$ then $F \approx E$; if
        $E \approx F$ and $F \approx G$ then $E \approx G$. (Quote
        precisely \cref{thm:b1:logic:inverse} and
        \cref{prop:b1:logic:comp}.)
  \item Show that $\preceq$ is transitive, and that an injection
        $f \colon E \to F$ always induces $E \approx f(E)$.
  \item Let $E \neq \emptyset$. Show that $E \preceq F$ if and only
        if there exists a surjection from $F$ onto $E$.
  \item Verify that $n \mapsto n + 1$ is a bijection from $\N$ onto
        $\N^* = \N \setminus \{0\}$, and that
        \[
        \sigma(n) = \frac n2 \ \ (n \text{ even}), \qquad
        \sigma(n) = -\frac{n+1}2 \ \ (n \text{ odd})
        \]
        is a bijection from $\N$ onto $\Z$. Hence removing a point,
        or doubling into negatives, does not change the size of
        $\N$.
\end{enumerate}

\medskip
\noindent\textbf{Part II --- The Cantor--Schr\"oder--Bernstein
theorem.} Let $f \colon E \to F$ and $g \colon F \to E$ be two
injections. Define
\[
C_0 = E \setminus g(F), \qquad C_{n+1} = g\bigl(f(C_n)\bigr)
\ \ (n \in \N), \qquad C = \bigcup_{n \in \N} C_n,
\]
and let $h \colon E \to F$ send $x \in C$ to $f(x)$, and $x \notin C$
to the unique $y \in F$ with $g(y) = x$.
\begin{enumerate}[resume]
  \item Check that $h$ is well defined: if $x \notin C$ then $x \in
        g(F)$, and the element $y$ with $g(y) = x$ is unique.
  \item Show that $g\bigl(f(C)\bigr) = \bigcup_{n \geq 1} C_n
        \subseteq C$. (Direct images commute with unions:
        \cref{exo:b1:logic:8}.)
  \item Show that $h$ is \omterm{def:b1:logic:inj}{injective}. (Three cases; in the mixed case
        $x \in C$, $x' \notin C$, show that $h(x) = h(x')$ would
        force $x' \in g(f(C)) \subseteq C$.)
  \item Show that $h$ is \omterm{def:b1:logic:inj}{surjective}: given $y \in F$, distinguish
        the cases $g(y) \notin C$ and $g(y) \in C_n$ for some $n
        \geq 1$ (why is $g(y) \in C_0$ impossible?), and exhibit a
        \omterm{def:b1:logic:map}{preimage} of $y$ in each case.
  \item Conclude with the \emph{Cantor--Schr\"oder--Bernstein
        theorem}: if $E \preceq F$ and $F \preceq E$, then $E
        \approx F$. Comment in one sentence on what makes this
        \omterm{def:b1:logic:statement}{statement} nontrivial.
  \item Two applications. (a) Show $\intcc01 \approx \intoo01$.
        (b) Show that $\varphi(p, q) = 2^p(2q + 1) - 1$ defines a
        bijection from $\N \times \N$ onto $\N$ --- injectivity by a
        parity argument, surjectivity by strong induction
        (\cref{thm:b1:logic:induction}). Hence $\N \times \N \approx
        \N$: the plane of integer points is no bigger than the
        line.
\end{enumerate}

\medskip
\noindent\textbf{Part III --- \omterm{pb:b1:logic:1}{Countable sets}.} Call a \omterm{def:b1:logic:sets}{set} $E$
\emph{at most countable} when $E \preceq \N$, and \emph{countable}
when $E \approx \N$.
\begin{enumerate}[resume]
  \item Show that every infinite subset $A \subseteq \N$ is
        countable. (Define $\varphi(n)$ recursively as the least
        element of $A \setminus \{\varphi(0), \dots,
        \varphi(n-1)\}$; show $\varphi$ is strictly increasing,
        satisfies $\varphi(n) \geq n$, and attains every element of
        $A$.)
  \item Deduce that a \omterm{def:b1:logic:sets}{set} is at most countable if and only if it is
        finite or countable, and observe that question 9 gives the
        shortcut: if $E \preceq \N$ and $\N \preceq E$, then $E$ is
        countable.
  \item Show that if $E$ and $F$ are at most countable, so is $E
        \times F$. Deduce that $\Z \times \N^*$ is countable.
  \item Show that $\Q$ is countable. (Inject $\Q$ into $\Z \times
        \N^*$ by writing each rational in lowest terms with positive
        denominator --- the uniqueness of that representation is
        proved in \cref{ch:b1:arith}; then apply question 12.)
  \item Show that a countable union of at most \omterm{pb:b1:logic:1}{countable sets} is at
        most countable: if each $E_n$ ($n \in \N$) is at most
        countable, so is $\bigcup_{n \in \N} E_n$. (Send $x$ to the
        pair $(n, f_n(x))$ where $n$ is the \emph{least} index with
        $x \in E_n$.)
  \item Show that the \omterm{def:b1:logic:sets}{set} of \emph{finite} subsets of $\N$ is
        countable. (Map a finite subset $F$ to $\sum_{i \in F} 2^i$;
        prove injectivity by comparing the largest element where two
        finite \omterm{def:b1:logic:sets}{sets} differ, using $\sum_{k=0}^{m-1} 2^k = 2^m - 1$
        from \cref{exo:b1:logic:4}.)
\end{enumerate}

\medskip
\noindent\textbf{Part IV --- Diagonalization.} Let $\{0,1\}^{\N}$
denote the \omterm{def:b1:logic:sets}{set} of all \omterm{def:b1:logic:map}{maps} $u \colon \N \to \{0, 1\}$, i.e.\ the \omterm{def:b1:logic:sets}{set}
of binary sequences.
\begin{enumerate}[resume]
  \item Construct a bijection between $\mathcal{P}(\N)$ and
        $\{0,1\}^{\N}$ (indicator functions).
  \item (The diagonal argument) Let $\Phi \colon \N \to
        \{0,1\}^{\N}$ be any \omterm{def:b1:logic:map}{map}. Consider the sequence $d$ defined
        by $d(n) = 1 - \Phi(n)(n)$. Show that $d$ is not in the
        image of $\Phi$, and conclude that $\{0,1\}^{\N}$ is
        \emph{not} at most countable. Explain in one sentence why,
        through question 17, this is exactly Cantor's theorem
        (\cref{exo:b1:logic:11}) for $E = \N$.
  \item Admit --- as familiar from school, and established
        rigorously in \cref{ch:b1:reals} --- that every $x \in
        \intco01$ has a unique \emph{proper} decimal expansion $x =
        0.d_1 d_2 d_3\dots$ (one not ending in an infinite string of
        $9$s). Given any sequence $(x_n)_{n \geq 1}$ of elements of
        $\intco01$, construct $x \in \intco01$ with $x \neq x_n$ for
        all $n$: choose its $n$-th digit to be $5$ if the $n$-th
        digit of $x_n$ differs from $5$, and $6$ otherwise. Justify
        carefully that $x$ is proper and avoids every $x_n$, and
        conclude that $\intco01$ is not at most countable.
  \item Deduce that $\R$ is uncountable, and that the \omterm{def:b1:logic:sets}{set} $\R
        \setminus \Q$ of irrational numbers is uncountable as well.
        In what precise sense are ``most'' real numbers irrational?
\end{enumerate}

\medskip
\noindent\textbf{Part V --- Cantor's 1874 theorem: \omterm{pb:b1:logic:1}{transcendental
numbers} exist.} A real number $x$ is \emph{algebraic}
\index{algebraic number} when $P(x) = 0$ for some nonzero polynomial
$P$ with integer coefficients, and \emph{transcendental}
\index{transcendental number} otherwise. Admit for this part --- it
is proved in \cref{ch:b1:poly} --- that a nonzero polynomial of
degree $n$ has at most $n$ real roots.
\begin{enumerate}[resume]
  \item Show that every rational number is algebraic, and find
        explicit polynomials with integer coefficients annihilating
        $\sqrt 2$ and $\sqrt 2 + \sqrt 3$.
  \item For fixed $n \in \N$, show that the \omterm{def:b1:logic:sets}{set} of polynomials of
        degree at most $n$ with integer coefficients is countable.
        (Inject it into $\Z^{n+1}$ and induct on $n$ with
        question 13.)
  \item Deduce that the \omterm{def:b1:logic:sets}{set} of \emph{all} polynomials with integer
        coefficients is countable.
  \item Prove \emph{Cantor's theorem on \omterm{pb:b1:logic:1}{algebraic numbers}}: the \omterm{def:b1:logic:sets}{set}
        $\mathcal{A}$ of algebraic real numbers is countable.
  \item Conclude: transcendental real numbers exist, and the \omterm{def:b1:logic:sets}{set} of
        \omterm{pb:b1:logic:1}{transcendental numbers} is uncountable. Then take stock of
        the whole problem in a few sentences: the chain $\N \approx
        \Z \approx \Q \approx \mathcal{A}$, the strict jump to $\R
        \approx$ (essentially) $\mathcal{P}(\N)$, where each tool
        (Cantor--Schr\"oder--Bernstein, countable unions, the
        diagonal) was decisive --- and the philosophical punch of
        proving that uncountably many \omterm{pb:b1:logic:1}{transcendental numbers} exist
        without naming even one. (Proving a \emph{specific} number
        such as $\pi$ transcendental is far harder and beyond this
        volume.)
\end{enumerate}
\end{problem}

\chapter{Counting}\label{ch:b1:counting}

Counting finite sets sounds elementary --- and quickly becomes subtle.
This chapter defines cardinality properly (through bijections, in the
spirit of \cref{ch:b1:logic}), establishes the handful of counting
principles from which everything follows, and derives the classical
counts: lists, permutations, subsets, binomial coefficients.

\section{Cardinality of finite sets}

\begin{definition}[Finite set, cardinality]\label{def:b1:counting:card}
For $n \in \N^*$, write $\intint{1}{n} = \{1, 2, \dots,
n\}$. A set $E$ is \emph{finite}\index{finite set} when $E = \emptyset$
or there is a bijection from $\intint{1}{n}$ onto $E$ for
some $n \in \N^*$; this $n$ is unique (\cref{thm:b1:counting:welldef})
and is the \emph{cardinality}\index{cardinality} of $E$, written
$\abs{E}$ (with $\abs{\emptyset} = 0$).
\end{definition}

\begin{theorem}[Cardinality is well defined]\label{thm:b1:counting:welldef}
If $m \neq n$, there is no bijection from $\intint{1}{m}$
onto $\intint{1}{n}$. More precisely, if $m > n$ there is
no injection from $\intint{1}{m}$ into $\intint{1}{n}$.
\end{theorem}

\begin{proof}
We prove by induction on $n$ the statement: \emph{for all $m > n$,
there is no injection $\intint{1}{m} \to \intint{1}{n}$}. For $n = 0$ the target is empty and $m \geq 1$: no map
exists at all. Assume the statement for $n$, and suppose $f \colon
\intint{1}{m} \to \intint{1}{n+1}$ is an
injection with $m > n + 1$. If the value $n + 1$ is not attained, $f$
is an injection into $\intint{1}{n}$, contradicting the
induction hypothesis. Otherwise $f(a) = n + 1$ for exactly one $a$;
swap $f(a)$ and $f(m)$ (formally: compose with the transposition of
the two values), so that the new injection $g$ has $g(m) = n + 1$.
Then the restriction of $g$ to $\intint{1}{m-1}$ is an
injection into $\intint{1}{n}$ with $m - 1 > n$ ---
contradiction again.
\end{proof}

\begin{corollary}[Pigeonhole principle]\label{cor:b1:counting:pigeonhole}
If $\abs{E} > \abs{F}$, no map $f \colon E \to F$ is injective: some
two elements of $E$ share their image\index{pigeonhole principle}.
\end{corollary}

\begin{proof}
Write $\abs E = m$, $\abs F = n$ with $m > n$, and choose bijections
$u \colon \intint1m \to E$ and $v \colon F \to \intint1n$. If $f$
were injective, $v \circ f \circ u$ would be an injection from
$\intint1m$ into $\intint1n$ (a composition of injections,
\cref{prop:b1:logic:comp}), contradicting
\cref{thm:b1:counting:welldef}.
\end{proof}

\begin{remark}[Interlude: why the swap in the theorem's proof?]\label{rem:b1:counting:swapinterlude}
The proof of \cref{thm:b1:counting:welldef} contains the
chapter's first genuinely clever move, worth replaying slowly.
The obstacle: to apply the induction hypothesis one wants to
delete the last point $m$ of the source \emph{and} the last point
$n+1$ of the target, but $f$ may send some other point $a$ to $n
+ 1$, and then deleting the target point damages the map
elsewhere. The cure: compose $f$ with the transposition of the
two \emph{values} $f(a)$ and $f(m)$ --- a bijection of the
target, so injectivity is preserved --- after which the
troublesome value $n + 1$ sits at the harmless position $m$, and
both deletions are clean. This ``normalize first, then cut''
pattern recurs: it is how the derangement recurrence redirects
$\sigma^{-1}(n+1)$ in the weekend problem of this chapter, and
how permutations are patched throughout
\cref{ch:b1:structures}'s problem on the symmetric group.
\end{remark}

\begin{proposition}[Injections, surjections and cardinality]\label{prop:b1:counting:injsur}
Let $E, F$ be finite sets with $\abs{E} = \abs{F}$, and $f \colon E
\to F$. Then
\[
f \text{ injective} \iff f \text{ surjective} \iff f \text{ bijective}.
\]
\end{proposition}

\begin{proof}
Suppose $f$ injective. Then $f$ is a bijection from $E$ onto $f(E)$,
so $\abs{f(E)} = \abs{E} = \abs{F}$. If $f(E)$ missed a point $y_0$ of
$F$, then $f$ would be an injection of $E$ into $F \setminus \{y_0\}$,
a set of cardinality $\abs{F} - 1 < \abs{E}$ --- impossible by the
pigeonhole principle. So $f(E) = F$: $f$ is surjective, hence
bijective.

Suppose $f$ surjective. Choose for each $y \in F$ one preimage $s(y)
\in E$; then $f \circ s = \mathrm{id}_F$, so $s$ is injective
(\cref{prop:b1:logic:comp}). By the previous paragraph applied to $s$
(the cardinalities are equal), $s$ is bijective. From $f \circ s =
\mathrm{id}_F$ we get $f = \mathrm{id}_F \circ s^{-1} = s^{-1}$, so
$f$ is bijective. Finally, a bijective map is by definition both
injective and surjective, which closes the cycle of implications.
\end{proof}

\begin{example}[Finiteness is essential]\label{ex:b1:counting:finiteonly}
On a \emph{finite} set, \cref{prop:b1:counting:injsur} is a powerful
shortcut: any injective map from $E$ to itself is automatically a
permutation of $E$ --- half of bijectivity comes for free. Both
implications collapse on infinite sets: $n \mapsto n + 1$ is
injective from $\N$ to $\N$ but misses $0$, and the map $\N \to \N$
sending $0 \mapsto 0$ and $n \mapsto n - 1$ for $n \geq 1$ is
surjective but not injective. Whenever this proposition is invoked,
the finiteness hypothesis is doing real work --- a theme that the
weekend problem of \cref{ch:b1:logic} explores from the other side,
where infinite sets are precisely those admitting such self-maps.
\end{example}

\begin{example}[Half the work, for free]\label{ex:b1:counting:mod7}
Consider the map $f$ on $\{0, 1, \dots, 6\}$ sending $k$ to the
remainder of $3k$ upon division by $7$; its table of values is
\[
0,\ 3,\ 6,\ 2,\ 5,\ 1,\ 4 .
\]
Is $f$ a bijection? Injectivity alone suffices
(\cref{prop:b1:counting:injsur}): if $3k$ and $3k'$ have the same
remainder, $7$ divides $3(k - k')$, and since $7$ is prime and
does not divide $3$, it divides $k - k'$ (Euclid's lemma, used at
High School level here and proved in \cref{ch:b1:arith}); with
$\abs{k - k'} \leq 6$ this forces $k = k'$. Surjectivity comes
free --- no need to solve $3k \equiv c$ for each $c$, though the
table confirms every value appears exactly once. The shortcut is
a workhorse: it proves the invertibility of modular
multiplication (\cref{ch:b1:arith}), powers the pairing in
Wilson's theorem, and returns in linear algebra as ``an
endomorphism of a finite-dimensional space is injective iff
surjective'' (\cref{ch:b1:findim}).
\end{example}

\section{The counting principles}

\begin{proposition}[Sum and product rules]\label{prop:b1:counting:rules}
Let $E, F$ be finite sets.
\begin{enumerate}
  \item If $E \cap F = \emptyset$, then $\abs{E \cup F} = \abs{E} +
        \abs{F}$; more generally, for a partition of $E$ into pieces
        $E_1, \dots, E_k$, $\abs{E} = \sum_i \abs{E_i}$.
  \item In general, $\abs{E \cup F} = \abs{E} + \abs{F} - \abs{E \cap
        F}$.
  \item $\abs{E \times F} = \abs{E} \times \abs{F}$.
  \item The set $F^E$ of all maps from $E$ to $F$ satisfies
        $\abs{F^E} = \abs{F}^{\abs{E}}$.
  \item $\abs{\mathcal{P}(E)} = 2^{\abs{E}}$.
\end{enumerate}
\end{proposition}

\begin{proof}
(1) Concatenate enumerations: if $E = \{x_1, \dots, x_m\}$ and $F =
\{y_1, \dots, y_n\}$ with no repetition, then $x_1, \dots, x_m, y_1,
\dots, y_n$ enumerates $E \cup F$ without repetition (disjointness).
Induction extends this to $k$ pieces.

(2) $E \cup F$ is the disjoint union of $E$ and $F \setminus E$, and
$F$ is the disjoint union of $F \cap E$ and $F \setminus E$; so
$\abs{E \cup F} = \abs{E} + \abs{F \setminus E} = \abs{E} + \abs{F} -
\abs{E \cap F}$.

(3) $E \times F$ is the disjoint union, over $x \in E$, of the sets
$\{x\} \times F$, each of cardinality $\abs{F}$; apply (1).

(4) A map from $E = \{x_1, \dots, x_m\}$ to $F$ is exactly the choice
of the $m$-tuple $(f(x_1), \dots, f(x_m)) \in F^m$; this
correspondence is a bijection, and $\abs{F^m} = \abs{F}^m$ by (3) and
induction.

(5) Subsets of $E$ correspond bijectively to maps $E \to \{0, 1\}$
(send $A$ to its indicator function); apply (4).
\end{proof}

\begin{example}[Complement counting]\label{ex:b1:counting:complement}
How many $4$-digit PIN codes (digits $0$--$9$, order matters,
repetition allowed) contain \emph{at least one} repeated digit?
Counting them directly means juggling the cases ``exactly one pair,
two pairs, a triple, a quadruple'' --- five overlapping
configurations. Count the complement instead: all codes number
$10^4 = 10\,000$ (product rule), the codes with four distinct
digits number $10 \times 9 \times 8 \times 7 = 5\,040$
($4$-arrangements), so the answer is
\[
10^4 - 10 \cdot 9 \cdot 8 \cdot 7 = 10\,000 - 5\,040 = 4\,960 .
\]
Nearly half of all PINs repeat a digit. The insight: whenever a
count is phrased with ``at least'' or ``not all'', try the
complement first --- the sum rule guarantees that
$\abs{A} = \abs{E} - \abs{\overline A}$, and the complement is
often a single clean configuration.
\end{example}

\begin{example}[Lattice paths]\label{ex:b1:counting:paths}
Count the shortest paths from the corner $(0,0)$ to the corner
$(4, 3)$ of a grid, moving only one step right (R) or one step up
(U) at a time. Every such path takes exactly $7$ steps, of which
$4$ are R and $3$ are U; conversely, any word of length $7$ in the
letters R, U with four R's describes exactly one path. The paths
therefore correspond bijectively to the choices of the positions
of the R's:
\[
\binom{7}{4} = 35 .
\]
The insight is the \emph{encoding}: the count became trivial the
moment each path was translated into a word, i.e.\ a subset of
positions --- one more instance of the slogan that a correct count
is a disguised bijection (\cref{met:b1:counting:which}).
\end{example}

\begin{omfigure}
\begin{tikzpicture}[scale=0.85]
  \draw[gray!60] (0,0) grid (4,3);
  \draw[->, thin] (-0.4,0) -- (4.6,0);
  \draw[->, thin] (0,-0.4) -- (0,3.6);
  \draw[omDef, very thick]
    (0,0) -- (1,0) -- (1,1) -- (3,1) -- (3,2) -- (4,2) -- (4,3);
  \fill (0,0) circle (0.07) node[below left] {$(0,0)$};
  \fill (4,3) circle (0.07) node[above right] {$(4,3)$};
  \node[omDef, below] at (0.5,0) {\small R};
  \node[omDef, left] at (1,0.5) {\small U};
  \node[omDef, below] at (1.5,1) {\small R};
  \node[omDef, below] at (2.5,1) {\small R};
  \node[omDef, left] at (3,1.5) {\small U};
  \node[omDef, below] at (3.5,2) {\small R};
  \node[omDef, left] at (4,2.5) {\small U};
\end{tikzpicture}

{\small One of the $\binom74 = 35$ shortest paths from $(0,0)$ to
$(4,3)$: the path shown encodes the word RURRURU, i.e.\ the choice
of positions $\{1,3,4,6\}$ for the letter R among the seven
steps.}
\end{omfigure}

\section{Lists, permutations, subsets}

\begin{definition}[Arrangements, permutations, combinations]\label{def:b1:counting:objects}
Let $E$ be a set with $\abs{E} = n$ and let $0 \leq k \leq n$.
\begin{itemize}
  \item A \emph{$k$-arrangement} of $E$ is an injective $k$-tuple of
        elements of $E$ (an ordered selection without repetition);
  \item a \emph{permutation}\index{permutation} of $E$ is a bijection
        from $E$ to itself --- equivalently, an $n$-arrangement;
  \item a \emph{$k$-combination} is a subset of $E$ with $k$ elements
        (an unordered selection without repetition). Their number is
        written $\binom{n}{k}$, read ``$n$ choose $k$''
        \index{binomial coefficient}.
\end{itemize}
\end{definition}

\begin{theorem}[The three counts]\label{thm:b1:counting:counts}
With $n = \abs{E}$ and $0 \leq k \leq n$:
\begin{enumerate}
  \item the number of $k$-arrangements of $E$ is $n (n-1) \cdots
        (n-k+1) = \dfrac{n!}{(n-k)!}$;
  \item the number of permutations of $E$ is $n!$;
  \item $\dbinom{n}{k} = \dfrac{n!}{k!\,(n-k)!}$.
\end{enumerate}
\end{theorem}

\begin{proof}
(1) Choose the first coordinate ($n$ ways), then the second ($n - 1$
remaining choices), \dots, then the $k$-th ($n - k + 1$ choices).
Formally, induct on $k$. For $k = 1$ there are $n$ one-term
injective tuples. Assume the count for $k - 1$. Each
$k$-arrangement $(x_1, \dots, x_k)$ is obtained from exactly one
$(k-1)$-arrangement --- its truncation $(x_1, \dots, x_{k-1})$
--- by appending a last coordinate outside $\{x_1, \dots,
x_{k-1}\}$, for which exactly $n - (k - 1)$ values are available.
The $k$-arrangements are thus partitioned, by truncation, into
classes of common size $n - k + 1$ indexed by the
$(k-1)$-arrangements, and the sum rule gives
\[
\frac{n!}{(n-k+1)!}\;(n - k + 1) = \frac{n!}{(n-k)!} .
\]

(2) is (1) with $k = n$.

(3) Each $k$-subset orders into $k!$ distinct $k$-arrangements, and
every $k$-arrangement arises from exactly one subset: so
$\frac{n!}{(n-k)!} = \binom nk \cdot k!$.
\end{proof}

\begin{example}[Round tables: quotienting by symmetry]\label{ex:b1:counting:roundtable}
In how many ways can $n$ guests sit around a round table, two
seatings being identical when each guest has the same left and
right neighbors --- i.e.\ up to rotation? Each circular seating
corresponds to exactly $n$ linear seatings (cut the circle at any
of the $n$ places), so the $n!$ linear orders collapse in groups
of $n$:
\[
\frac{n!}{n} = (n-1)! \quad\text{circular seatings.}
\]
Equivalently: seat one distinguished guest anywhere (killing the
rotational freedom), then order the remaining $n - 1$ guests
clockwise. For $n = 6$: $120$ tables. The two solutions illustrate
the two standard cures for overcounting: divide by the exact
number of repetitions, or \emph{break the symmetry} by pinning one
object down. Both require the group-of-repetitions size to be the
same for every configuration --- which the proof of the formula
$\binom nk = \frac{n!}{k!\,(n-k)!}$ above used as well, with
$k!$ in place of $n$.
\end{example}

\begin{example}[Adding a constraint]\label{ex:b1:counting:constraint}
Continuing the round table: among the $(n-1)!$ tables of $n \geq
3$ guests, how many seat two given guests $A$ and $B$ \emph{apart}
(not adjacent)? Count the complement. Tables where $A$ and $B$
sit together: glue them into a single block --- $n - 1$ objects
around the table, i.e.\ $(n-2)!$ circular arrangements --- then
order the pair inside its block ($2$ ways): $2\,(n-2)!$ adjacent
tables. Hence
\[
(n-1)! - 2\,(n-2)! = (n-2)!\,\bigl((n - 1) - 2\bigr)
= (n-3)\,(n-2)!
\]
tables keep them apart. Sanity checks: $n = 3$ gives $0$ (around
a triangle, everyone touches everyone) and $n = 4$ gives $2$,
easily listed by hand. The gluing trick --- treat a forced block
as one object, then count its internal arrangements --- is the
standard cure for adjacency constraints, linear or circular.
\end{example}

\begin{proposition}[Basic identities]\label{prop:b1:counting:identities}
For $0 \leq k \leq n$:
\[
\binom{n}{k} = \binom{n}{n-k},
\qquad
\binom{n}{k} = \binom{n-1}{k-1} + \binom{n-1}{k}
\quad (1 \leq k \leq n-1),
\qquad
\sum_{k=0}^{n} \binom{n}{k} = 2^n .
\]
\end{proposition}

\begin{proof}
First identity: $A \mapsto E \setminus A$ is a bijection between
$k$-subsets and $(n-k)$-subsets. Pascal's rule\index{Pascal's rule}:
fix an element $a \in E$; the $k$-subsets split into those containing
$a$ (choose the $k - 1$ others: $\binom{n-1}{k-1}$) and those avoiding
$a$ ($\binom{n-1}{k}$). Third identity: both sides count all subsets
of $E$, split by size on the left
(\cref{prop:b1:counting:rules}~(1) and (5)).
\end{proof}

\begin{theorem}[Binomial theorem]\label{thm:b1:counting:binomial}
For all $a, b$ in a commutative ring (say $\R$ or $\C$) and $n \in \N$:
\[
(a + b)^n = \sum_{k=0}^{n} \binom{n}{k} a^k b^{\,n-k} .
\]
\end{theorem}

\begin{proof}
Expanding $(a+b)(a+b)\cdots(a+b)$ distributively yields one term per
choice, in each factor, of either $a$ or $b$: the term $a^k b^{n-k}$
appears once for each way of choosing which $k$ of the $n$ factors
contribute $a$ --- that is, $\binom nk$ times. (Alternatively: induct
on $n$ using Pascal's rule.)
\end{proof}

\begin{example}\label{ex:b1:counting:binomial}
Two classical specializations: $a = b = 1$ recovers $\sum_k \binom nk
= 2^n$; $a = -1$, $b = 1$ gives $\sum_{k} (-1)^k \binom nk = 0$ for
$n \geq 1$: among the subsets of a nonempty set, exactly half have
even cardinality.
\end{example}

\begin{example}[One identity, two proofs]\label{ex:b1:counting:threen}
The specialization $a = 2$, $b = 1$ of the binomial theorem reads
\[
\sum_{k=0}^{n} \binom nk\,2^k = 3^n .
\]
Here is the same identity with no algebra at all. The right side
counts the words of length $n$ over the alphabet $\{0, 1, 2\}$
(product rule). Classify each word by the set $K$ of positions
carrying a nonzero letter: choosing $K$ with $\abs K = k$ costs
$\binom nk$, then each position of $K$ independently carries $1$
or $2$: $2^k$ ways. The sum rule over $k$ gives the left side.
Beyond the pleasure of agreement, the two proofs have different
virtues: the algebraic one generalizes to any value of $a$, the
combinatorial one \emph{explains} the formula and adapts to
constraints (forbid the letter $2$ in the last position, say) that
no substitution captures. Keeping both techniques active is the
practical skill this chapter trains.
\end{example}

\begin{method}[Which count applies?]\label{met:b1:counting:which}
Before computing, answer two questions about the selection: does the
\emph{order} matter, and are \emph{repetitions} allowed?
\begin{center}
\renewcommand{\arraystretch}{1.4}
\begin{tabular}{l|c|c}
 & order matters & order does not matter \\
\hline
no repetition & $\dfrac{n!}{(n-k)!}$ & $\dbinom{n}{k}$ \\[6pt]
repetition allowed & $n^k$ & (\cref{exo:b1:counting:10}) \\
\end{tabular}
\end{center}
Then hunt for a bijection or a partition reducing the problem to these
model counts; a correct count is a disguised bijection.
\end{method}

\begin{remark}[Common pitfalls in counting]\label{rem:b1:counting:pitfalls}
\begin{enumerate}
  \item \emph{Summing non-disjoint cases.} The sum rule requires a
        partition; if configurations can satisfy two cases at
        once, they are counted twice --- the cure is
        inclusion--exclusion (\cref{thm:b1:counting:inclexcl}) or
        a finer case split.
  \item \emph{Ordered versus unordered.} Choosing ``a committee of
        two'' is $\binom n2$, not $n(n-1)$: decide \emph{before
        computing} whether the selection carries an order, and if
        an ordered count is easier, divide by the number of
        orderings at the end --- but only when every unordered
        object arises from the \emph{same} number of ordered ones.
  \item \emph{Multi-stage choices that are not independent.}
        The product rule needs the number of options at each stage
        to be independent of the previous choices. ``Choose a
        captain, then a different vice-captain'' is fine ($n(n-1)$);
        ``choose two players who get along'' is not a two-stage
        product at all.
  \item \emph{Double-counting by construction.} Building each
        object twice --- e.g.\ counting hands with \emph{at least}
        one ace as (choose an ace) $\times$ (choose $4$ more
        cards) --- overcounts hands with two aces. ``At least''
        almost always calls for the complement
        (\cref{ex:b1:counting:complement}).
\end{enumerate}
\end{remark}

\begin{example}[Poker-style count]\label{ex:b1:counting:poker}
From a $52$-card deck, the number of $5$-card hands is
$\binom{52}{5} = 2\,598\,960$. Hands containing exactly one ace:
choose the ace ($4$ ways) then $4$ cards among the $48$ non-aces:
$4 \binom{48}{4} = 778\,320$. The product rule applies because the
choice splits into independent stages.
\end{example}

\begin{method}[Double counting]\label{met:b1:counting:doublecount}
To prove an identity between two counting expressions, find a single
finite set that both sides count --- typically a set of \emph{pairs}
--- and evaluate its cardinality in two different orders. The
prototype is the \emph{handshake lemma}: at a party, count the pairs
(person, hand shaken). Summing over people gives $\sum_p d_p$ (each
person $p$'s number of handshakes); summing over handshakes gives
twice the number of handshakes (each involves two people). Hence
$\sum_p d_p$ is even --- so the number of people who shook an odd
number of hands is always even, a nontrivial conclusion obtained
with no formula at all. The same engine runs
\cref{exo:b1:counting:12} and several questions of the weekend
problem below.
\end{method}

\begin{example}[The average subset]\label{ex:b1:counting:averagesubset}
What is the average cardinality of a subset of an $n$-element set
$E$, all $2^n$ subsets being equally likely? Double count the
pairs $(A, a)$ with $a \in A$: summing over subsets gives
$\sum_A \abs A$, the total we want; summing over elements gives
$n \cdot 2^{n-1}$ (each of the $n$ elements lies in exactly half
of the subsets --- pair each $A$ containing $a$ with $A
\setminus \{a\}$). Hence
\[
\frac{1}{2^n}\sum_{A \subseteq E} \abs A
= \frac{n\,2^{n-1}}{2^n} = \frac n2 :
\]
subsets are, on average, half-full --- as the symmetry $A
\leftrightarrow \overline A$ (which pairs sizes $k$ and $n - k$)
also predicts. Two proofs, one answer, and both avoid the direct
computation $\sum_k k\binom nk$ of \cref{exo:b1:counting:5}: a
well-chosen pairing often replaces an identity.
\end{example}

\section{Inclusion--exclusion}

\begin{theorem}[Inclusion--exclusion]\label{thm:b1:counting:inclexcl}
For finite sets $A_1, \dots, A_p$:
\[
\Bigl|\, \bigcup_{i=1}^{p} A_i \,\Bigr|
= \sum_{\emptyset \neq I \subseteq \intint{1}{p}}
(-1)^{\abs{I}+1} \Bigl|\, \bigcap_{i \in I} A_i \,\Bigr| .
\]
For $p = 3$:
$\abs{A \cup B \cup C} = \abs A + \abs B + \abs C - \abs{A \cap B} -
\abs{A \cap C} - \abs{B \cap C} + \abs{A \cap B \cap C}$.
\end{theorem}

\begin{proof}
Fix an element $x$ of the union and count its contribution to the
right-hand side. Let $J = \{i : x \in A_i\}$, of cardinality $m \geq
1$. The element $x$ is counted once in $\abs{\bigcap_{i \in I} A_i}$
exactly when $\emptyset \neq I \subseteq J$, with sign
$(-1)^{\abs I + 1}$; its total contribution is
\[
\sum_{k=1}^{m} \binom{m}{k} (-1)^{k+1}
= 1 - \sum_{k=0}^{m} \binom mk (-1)^k = 1 - 0 = 1
\]
by \cref{ex:b1:counting:binomial}. So each element of the union is
counted exactly once.
\end{proof}

\begin{example}[Counting coprime integers]\label{ex:b1:counting:coprime}
How many integers of $\intint1{120}$ are coprime to $120 = 2^3
\times 3 \times 5$? An integer shares a factor with $120$ exactly
when it is divisible by $2$, $3$ or $5$, so count the complement
of $A_2 \cup A_3 \cup A_5$, where $A_d$ collects the multiples of
$d$. Within $\intint1{120}$, the multiples of $d$ number $120/d$
whenever $d$ divides $120$ --- no floor functions needed --- and
$A_2 \cap A_3 = A_6$, etc. Inclusion--exclusion:
\[
\abs{A_2 \cup A_3 \cup A_5}
= 60 + 40 + 24 - 20 - 12 - 8 + 4 = 88 ,
\]
so $120 - 88 = 32$ integers are coprime to $120$. It is
instructive to regroup the computation as a product:
\[
120 - 88 = 120\Bigl(1 - \frac12\Bigr)\Bigl(1 -
\frac13\Bigr)\Bigl(1 - \frac15\Bigr) = 120 \cdot \frac12 \cdot
\frac23 \cdot \frac45 = 32 :
\]
expanding the three parentheses reproduces exactly the eight
signed terms of inclusion--exclusion, one per subset of $\{2, 3,
5\}$. This product form defines Euler's totient function, whose
arithmetic role appears with \cref{ch:b1:arith}'s congruences and
is developed in the Year 2 volume.
\end{example}

\begin{example}[Derangements]\label{ex:b1:counting:derangement}
A \emph{derangement} is a permutation with no fixed point. Let $A_i$
be the set of permutations of $\intint{1}{n}$ fixing $i$;
then $\abs{\bigcap_{i \in I} A_i} = (n - \abs I)!$, and
inclusion--exclusion counts the permutations with at least one fixed
point; the derangements number
\[
D_n = n! \sum_{k=0}^{n} \frac{(-1)^k}{k!} .
\]
Since $\sum (-1)^k / k! \to \eu^{-1}$ (see \cref{ch:b1:series}), about
$37\%$ of all permutations are derangements, whatever $n$ is.
\end{example}

\begin{remark}[Where this chapter is used]\label{rem:b1:counting:whereused}
Binomial coefficients are the most heavily reused objects of this
chapter: they drive the binomial theorem in \cref{ch:b1:poly}
(expansion of $(X + a)^n$), the Leibniz formula for the $n$-th
derivative of a product in \cref{ch:b1:derivative}, and the
coefficients of Taylor expansions in \cref{ch:b1:taylor}.
Permutations return as a group --- with the signature built from
counting inversions --- in \cref{ch:b1:structures}, and the
signature in turn defines determinants in \cref{ch:b1:det}.
Inclusion--exclusion and the counting principles are the finite
backbone of discrete probability, developed in the Year 2 volume;
the derangement numbers of \cref{ex:b1:counting:derangement} are
studied in depth in the weekend problem below.
\end{remark}

\section{Exercises}

\begin{exercise}[$\star$]\label{exo:b1:counting:1}
A license plate consists of two letters (A--Z), then three digits, then
two letters. How many plates are possible? How many with no repeated
letter among the four?
\end{exercise}

\begin{exercise}[$\star$]\label{exo:b1:counting:2}
How many anagrams (rearrangements of the letters, meaningful or not)
has the word \textsc{orange}? And \textsc{banana}?
\end{exercise}

\begin{exercise}[$\star$]\label{exo:b1:counting:3}
A committee of $4$ people is chosen from $7$ women and $5$ men. How
many committees: in total? with exactly $2$ women? with at least one
man?
\end{exercise}

\begin{exercise}[$\star$]\label{exo:b1:counting:4}
Prove that in any group of $13$ people, two share their birth month;
and that among any $n + 1$ integers chosen from $\intint{1}{2n}$, two are consecutive. \emph{(Pigeonhole both times: name
the boxes.)}
\end{exercise}

\begin{exercise}[$\star$]\label{exo:b1:counting:5}
Compute $\sum_{k=0}^{n} k \binom{n}{k}$. \emph{Hint: differentiate
$(1 + x)^n$, or use $k \binom nk = n \binom{n-1}{k-1}$ (prove it).}
\end{exercise}

\begin{exercise}[$\star\star$]\label{exo:b1:counting:6}
How many strictly increasing maps are there from $\intint{1}{k}$ to $\intint{1}{n}$? Deduce the number of
increasing (not necessarily strictly) maps. \emph{Hint for the second
count: $f$ increasing $\mapsto$ $g(i) = f(i) + i - 1$.}
\end{exercise}

\begin{exercise}[$\star\star$]\label{exo:b1:counting:7}
(Vandermonde) Prove, by counting $k$-subsets of a set split into two
blocks of sizes $m$ and $n$:
\[
\binom{m+n}{k} = \sum_{j=0}^{k} \binom{m}{j} \binom{n}{k-j} .
\]
Deduce $\sum_{j=0}^{n} \binom nj^2 = \binom{2n}{n}$.
\end{exercise}

\begin{exercise}[$\star\star$]\label{exo:b1:counting:8}
How many integers in $\intint{1}{1000}$ are divisible by
$2$ or $3$ or $5$? (Inclusion--exclusion; $\lfloor 1000/6 \rfloor$
counts the multiples of $6$, etc.)
\end{exercise}

\begin{exercise}[$\star\star$]\label{exo:b1:counting:9}
Count the surjections from a set of $4$ elements onto a set of $2$
elements; then onto a set of $3$ elements. \emph{Hint: count the
non-surjective maps with inclusion--exclusion on the missed values.}
\end{exercise}

\begin{exercise}[$\star\star$]\label{exo:b1:counting:10}
(Stars and bars) Prove that the number of $k$-selections of $n$
objects \emph{with} repetition, order ignored --- equivalently, the
number of $(x_1, \dots, x_n) \in \N^n$ with $x_1 + \dots + x_n = k$
--- is $\binom{n + k - 1}{k}$. \emph{Hint: encode a solution as a row
of $k$ stars and $n - 1$ bars.}
\end{exercise}

\begin{exercise}[$\star\star\star$]\label{exo:b1:counting:11}
Prove the formula of \cref{ex:b1:counting:derangement} for $D_n$ in
detail, and deduce $n! = \sum_{k=0}^{n} \binom{n}{k} D_{n-k}$ (also
prove this identity directly by classifying permutations by their
fixed-point set).
\end{exercise}

\begin{exercise}[$\star\star\star$]\label{exo:b1:counting:12}
For $n \in \N^*$, prove by a double count of pairs (subset, marked
element):
\[
\sum_{k=1}^{n} k \binom{n}{k} = n\, 2^{n-1},
\qquad\text{then}\qquad
\sum_{k=1}^{n} k^2 \binom{n}{k} = n(n+1)\, 2^{n-2} .
\]
\emph{For the second: count pairs of marked elements, equal or not.}
\end{exercise}

\section{Problem: Derangements, or the misaddressed letters}

\begin{problem}\label{pb:b1:counting:1}
A secretary puts $n$ letters into $n$ addressed envelopes at random:
what is the chance that \emph{nobody} receives the right letter?
This classical question (Montmort, 1708) leads to the derangement
numbers $D_n$ of \cref{ex:b1:counting:derangement}. The
inclusion--exclusion formula is only the opening move: this problem
develops the recurrences that compute $D_n$, two further independent
proofs of the formula, the striking theorem that $D_n$ is the
integer nearest to $n!/\eu$, the full distribution of fixed points
of a random permutation, and the curious arithmetic of the sequence
$(D_n)$.\index{derangement}
Throughout, $D_n$ denotes the number of derangements (fixed-point
free permutations) of $\intint1n$, with the convention $D_0 = 1$
(the empty permutation has no fixed point).

\medskip
\noindent\textbf{Part I --- Small cases and the fixed-point
census.}
\begin{enumerate}
  \item Compute $D_1, D_2, D_3$ directly, and $D_4$ by listing the
        derangements of $\{1, 2, 3, 4\}$ grouped by the value of
        $\sigma(1)$. (You should find $D_4 = 9$.)
  \item For $0 \leq k \leq n$, show that the number $P_k(n)$ of
        permutations of $\intint1n$ with \emph{exactly} $k$ fixed
        points is $\binom nk D_{n-k}$.
  \item Verify the census for $n = 4$: compute $P_0(4), \dots,
        P_4(4)$ and check that they sum to $4! = 24$. Which is more
        likely for four letters: no match, or exactly one match?
  \item By double counting (\cref{met:b1:counting:doublecount})
        the pairs $(\sigma, i)$ with $\sigma(i) = i$, show that
        \[
        \sum_{\sigma} \abs{\mathrm{Fix}(\sigma)} = n! :
        \]
        on average, a random permutation has \emph{exactly one}
        fixed point, whatever $n \geq 1$ is.
\end{enumerate}

\medskip
\noindent\textbf{Part II --- Two recurrences and two new proofs of
the formula.}
\begin{enumerate}[resume]
  \item Prove combinatorially, for $n \geq 1$:
        \[
        D_{n+1} = n\,(D_n + D_{n-1}) .
        \]
        (Classify the derangements $\sigma$ of $\intint1{n+1}$ by
        $j = \sigma(n+1)$, then by whether $\sigma(j) = n + 1$; in
        the case $\sigma(j) \neq n+1$, build a bijection with the
        derangements of $\intint1n$ by redirecting the preimage of
        $n + 1$ to $j$.) Check the recurrence numerically up to
        $D_6$.
  \item Setting $u_n = D_n - n D_{n-1}$, deduce from question 5
        that $u_{n+1} = -u_n$, and conclude the second recurrence:
        \[
        D_n = n D_{n-1} + (-1)^n \qquad (n \geq 1).
        \]
  \item From question 6, prove by induction the formula of
        \cref{ex:b1:counting:derangement},
        \[
        D_n = n! \sum_{k=0}^{n} \frac{(-1)^k}{k!},
        \]
        --- a proof wholly independent of inclusion--exclusion.
  \item (Binomial inversion) Let $(a_n)$ and $(b_n)$ be two
        sequences such that $a_n = \sum_{k=0}^n \binom nk b_k$ for
        all $n$. Prove that
        \[
        b_n = \sum_{k=0}^{n} (-1)^{n-k} \binom nk a_k
        \qquad (n \in \N).
        \]
        (First establish the \emph{trinomial revision}
        $\binom nk \binom kj = \binom nj \binom{n-j}{k-j}$, then
        use the alternating row sum of
        \cref{ex:b1:counting:binomial}.)
  \item Apply question 8 to the identity $n! = \sum_k \binom nk
        D_{n-k}$ of \cref{exo:b1:counting:11} to obtain a
        \emph{third} proof of the formula for $D_n$.
\end{enumerate}

\medskip
\noindent\textbf{Part III --- The nearest integer to
$n!/\eu$.} Admit for this part --- the theory is built in
\cref{ch:b1:series} --- that
$\eu^{-1} = \lim_{n \to \infty} s_n$ where $s_n = \sum_{k=0}^{n}
\frac{(-1)^k}{k!}$, with the strict alternating-series bound
$\abs{\eu^{-1} - s_n} < \frac1{(n+1)!}$ for every $n$.
\begin{enumerate}[resume]
  \item Show that $\bigl| D_n - n!/\eu \bigr| < \frac1{n+1}$ for
        all $n \in \N$.
  \item Deduce the headline theorem: \emph{for every $n \geq 1$,
        $D_n$ is the integer nearest to $n!/\eu$}. Why does the
        argument need $n \geq 1$?
  \item Determine the sign of the error: show that $D_n > n!/\eu$
        exactly when $n$ is even. (Locate the first neglected term
        of the alternating series.)
  \item Compute $D_7$ through $D_{10}$ with the recurrence of
        question 5, then check $D_{10}$ against $10!/\eu$
        ($10! = 3\,628\,800$, $\eu \approx 2.718281828$).
  \item (The hat-check probability) Let $p_n = D_n/n!$ be the
        probability that a uniformly random permutation is a
        derangement. Show $\abs{p_n - \eu^{-1}} < \frac1{(n+1)!}$
        and compute $p_6$ to five decimal places. Comment: why is
        the answer to Montmort's question essentially independent
        of $n$ --- already for a dozen letters?
\end{enumerate}

\medskip
\noindent\textbf{Part IV --- The distribution of fixed points.}
\begin{enumerate}[resume]
  \item Fix $k \in \N$. Show that the proportion of permutations
        of $\intint1n$ with exactly $k$ fixed points satisfies
        \[
        \frac{P_k(n)}{n!} = \frac{s_{n-k}}{k!}
        \;\xrightarrow[n \to \infty]{}\; \frac{\eu^{-1}}{k!} .
        \]
        (These limiting values, summing to $1$, form the
        \emph{Poisson distribution} of parameter $1$, a central
        object of the probability course in the Year 2 volume.)
  \item By double counting the triples $(\sigma, i, j)$ where $i
        \neq j$ are both fixed by $\sigma$, show that
        $\sum_\sigma \abs{\mathrm{Fix}(\sigma)}\,
        (\abs{\mathrm{Fix}(\sigma)} - 1) = n!$ for $n \geq 2$.
        Combined with question 4: the average of
        $\abs{\mathrm{Fix}}^2$ is $2$, so the ``spread'' (variance)
        of the number of fixed points equals $1$ --- again
        independent of $n$, again matching the Poisson law.
  \item Compute the proportion of permutations having at least one
        fixed point for $n = 4, 5, 6$ (as fractions and to four
        decimal places), and compare with $1 - \eu^{-1} \approx
        0.6321$.
  \item Show directly --- no limits needed --- that $s_{n+2} - s_n
        = (-1)^{n+1}\bigl(\frac1{(n+1)!} - \frac1{(n+2)!}\bigr)$,
        and deduce that the probabilities $p_n = s_n$ of
        question 14 oscillate: $p_0 > p_2 > p_4 > \dots$ and
        $p_1 < p_3 < p_5 < \dots$, the even (resp.\ odd) values
        decreasing (resp.\ increasing) toward the common limit
        $\eu^{-1}$.
  \item (Secret Santa) $n$ people each draw one name from a hat;
        if anyone draws their own name, the \emph{whole} drawing is
        restarted from scratch. Using the standard fact that an
        event of probability $p$ takes on average $1/p$ attempts,
        estimate the average number of complete drawings needed,
        and conclude that the procedure costs about $\eu \approx
        2.72$ drawings on average, essentially independently of
        $n$.
\end{enumerate}

\medskip
\noindent\textbf{Part V --- The arithmetic of $D_n$, and a
synthesis.}
\begin{enumerate}[resume]
  \item Refine question 5: show that for fixed $j \in
        \intint2n$, the derangements of $\intint1n$ with
        $\sigma(1) = j$ number exactly $D_{n-1} + D_{n-2}$,
        independently of $j$. Deduce that $n - 1$ divides $D_n$
        for every $n \geq 2$.
  \item Prove that $D_n$ is odd if and only if $n$ is even. (Work
        modulo $2$ in the recurrence of question 6.)
  \item Prove that $D_n \equiv (-1)^n \pmod n$ for $n \geq 1$, and
        check the congruence on the last digit of $D_{10}$.
  \item Show from question 6 that $\dfrac{D_n}{D_{n-1}} = n +
        \dfrac{(-1)^n}{D_{n-1}}$ for $n \geq 3$, so the ratio of
        consecutive derangement numbers is \emph{almost exactly}
        $n$; explain in one sentence why this is consistent with
        $D_n \approx n!/\eu$.
  \item Where exactly did this problem use: (i) the product and
        sum rules; (ii) double counting; (iii) the binomial
        theorem; (iv) the admitted alternating-series bound? One
        sentence each.
  \item Synthesis. The formula for $D_n$ now has three proofs
        (inclusion--exclusion, recurrence plus induction, binomial
        inversion). In a short paragraph, compare what each proof
        \emph{explains}: which one computes fastest, which one
        generalizes to other fixed-point counts, and which one
        reveals why $\eu$ appears in a problem about envelopes.
\end{enumerate}
\end{problem}

\chapter{Complex Numbers}\label{ch:b1:complex}

Complex numbers were met in the High School volume as a computational
device for quadratic equations. This chapter treats them as a central
object: the exponential form and its consequences (de Moivre, $n$-th
roots, \omterm{def:b1:complex:unity}{roots of unity}), the systematic translation between $\C$ and
plane geometry, and the use of $\eu^{\iu\theta}$ as a machine for
proving trigonometric identities.

\section{The field $\C$, modulus and conjugate}

\begin{definition}[The complex field]\label{def:b1:complex:field}
$\C = \{a + \iu b : a, b \in \R\}$ with the usual addition and the
multiplication determined by $\iu^2 = -1$. Every nonzero $z = a + \iu
b$ has an inverse: $z^{-1} = \frac{a - \iu b}{a^2 + b^2}$ --- in the
language of \cref{ch:b1:structures}, $\C$ is a field. One writes $a =
\Re(z)$, $b = \Im(z)$, $\conj{z} = a - \iu b$ (the
\emph{conjugate}\index{conjugate}) and $\abs{z} = \sqrt{a^2 + b^2}$
(the \emph{modulus}\index{modulus}).
\end{definition}

\begin{proposition}[Rules for conjugate and modulus]\label{prop:b1:complex:rules}
For $z, w \in \C$:
\begin{enumerate}
  \item $\conj{z + w} = \conj z + \conj w$, $\conj{zw} = \conj z\,
        \conj w$, $\conj{\conj z} = z$;
  \item $z \conj z = \abs z^2$; $\;\Re(z) = \frac{z + \conj z}{2}$,
        $\Im(z) = \frac{z - \conj z}{2\iu}$;
  \item $\abs{zw} = \abs z\, \abs w$, and $\abs{z^{-1}} =
        \abs{z}^{-1}$ for $z \neq 0$;
  \item (triangle inequality\index{triangle inequality}) $\abs{z + w}
        \leq \abs z + \abs w$, with equality if and only if $z$ and
        $w$ lie on a common ray from $0$ ($w = \lambda z$ or
        $z = \lambda w$ with $\lambda \geq 0$);
  \item (reverse triangle inequality) $\bigl|\abs z - \abs w\bigr|
        \leq \abs{z - w}$.
\end{enumerate}
\end{proposition}

\begin{proof}
(1) and (2) are direct computations on real and imaginary parts. (3):
$\abs{zw}^2 = zw\,\conj{zw} = z\conj z\, w \conj w = \abs z^2 \abs
w^2$, then take square roots; apply to $z \cdot z^{-1} = 1$ for the
inverse.

(4) Both sides are nonnegative, so compare squares:
\[
\abs{z+w}^2 = (z+w)(\conj z + \conj w)
= \abs z^2 + \abs w^2 + 2\,\Re(z \conj w),
\]
and $\Re(z\conj w) \leq \abs{z \conj w} = \abs z \abs w$ gives
$\abs{z+w}^2 \leq (\abs z + \abs w)^2$. Equality forces $\Re(z \conj
w) = \abs{z \conj w}$, i.e.\ $z \conj w \in \R_+$; if $w \neq 0$ this
gives $z = \frac{z\conj w}{\abs w^2}\, w = \lambda w$ with $\lambda
\geq 0$ (and the case $w = 0$ is trivial).

(5) $\abs z = \abs{(z - w) + w} \leq \abs{z-w} + \abs w$ gives $\abs z
- \abs w \leq \abs{z - w}$; swap $z$ and $w$ for the other sign.
\end{proof}

\begin{example}[A complete modulus-and-argument workout]\label{ex:b1:complex:workout}
Put $w = \dfrac{3 + 4\iu}{1 - 2\iu}$ in algebraic form and compute
its \omterm{def:b1:complex:field}{modulus} twice. Multiplying by the \omterm{def:b1:complex:field}{conjugate} of the
denominator:
\[
w = \frac{(3 + 4\iu)(1 + 2\iu)}{(1 - 2\iu)(1 + 2\iu)}
= \frac{3 + 6\iu + 4\iu - 8}{1 + 4}
= \frac{-5 + 10\iu}{5} = -1 + 2\iu .
\]
Directly: $\abs w = \sqrt{1 + 4} = \sqrt5$. Via the quotient rule
(\cref{prop:b1:complex:rules}~(3)): $\abs w = \frac{\abs{3 +
4\iu}}{\abs{1 - 2\iu}} = \frac5{\sqrt5} = \sqrt5$ --- same answer,
no algebraic form needed. The lesson generalizes: moduli and
arguments travel well through products and quotients, real and
imaginary parts travel well through sums. Choose the
representation that matches the operations at hand, and convert
only when forced.
\end{example}

\begin{example}[Equations involving the conjugate]\label{ex:b1:complex:conjeq}
Solve in $\C$: $\;z + 2\conj z = 6 + 2\iu$. An equation mixing $z$
and $\conj z$ is \emph{not} polynomial in $z$; the reliable move
is to split into real coordinates. With $z = x + \iu y$:
\[
z + 2\conj z = 3x - \iu y ,
\]
so the equation reads $3x = 6$ and $-y = 2$: the unique solution
is $z = 2 - 2\iu$. (Check: $(2 - 2\iu) + 2(2 + 2\iu) = 6 +
2\iu$.) Alternatively, \omterm{def:b1:complex:field}{conjugate} the whole equation to get
$\conj z + 2z = 6 - 2\iu$ and solve the linear system in the
unknowns $z, \conj z$ --- same answer, and a useful trick when
the coefficients are complex. Equations in $z$ and $\conj z$ are
really systems of two real equations; expecting ``degree $1$, one
solution'' is safe here, but $z\conj z = -1$ (no solution) shows
the polynomial intuition failing as soon as products appear.
\end{example}

\section{Exponential form}

\begin{definition}[Complex exponential of an imaginary argument]\label{def:b1:complex:expi}
For $\theta \in \R$ one defines
\[
\eu^{\iu\theta} = \cos\theta + \iu \sin\theta .
\]
Every $z \neq 0$ can be written $z = r\,\eu^{\iu\theta}$ with $r =
\abs z > 0$; $\theta$ is an \emph{argument}\index{argument} of $z$,
determined up to adding a multiple of $2\pi$. The value in
$\intoc{-\pi}{\pi}$ is the \emph{principal argument}, written
$\arg z$.
\end{definition}

\begin{example}[First values, and one famous identity]\label{ex:b1:complex:eulerid}
Reading the definition at the cardinal angles:
\[
\eu^{\iu\pi/2} = \iu, \qquad
\eu^{\iu\pi} = -1, \qquad
\eu^{2\iu\pi} = 1, \qquad
\eu^{\iu\pi/4} = \frac{\sqrt2}2\,(1 + \iu) .
\]
The second one, rearranged as $\eu^{\iu\pi} + 1 = 0$, is Euler's
celebrated identity linking $\eu$, $\iu$, $\pi$, $1$ and $0$; at
this stage of the book it is a definition unwinding rather than a
theorem, and its real content --- why the analytic exponential
function of \cref{ch:b1:functions}, extended to imaginary
arguments, deserves the same name --- is settled by the power
series of \cref{ch:b1:series}. Meanwhile the display above is
worth memorizing as a conversion table: it is used silently every
time an argument is read off a picture.
\end{example}

\begin{theorem}[The functional equation]\label{thm:b1:complex:funceq}
For all $\theta, \varphi \in \R$:
\[
\eu^{\iu\theta}\, \eu^{\iu\varphi} = \eu^{\iu(\theta + \varphi)},
\qquad
\abs{\eu^{\iu\theta}} = 1,
\qquad
\conj{\eu^{\iu\theta}} = \eu^{-\iu\theta} = (\eu^{\iu\theta})^{-1}.
\]
Consequently $\abs{zw} = \abs z \abs w$ carries the arguments too:
$\arg(zw) \equiv \arg z + \arg w \pmod{2\pi}$.
\end{theorem}

\begin{proof}
Expand the product and use the addition formulas:
\[
(\cos\theta + \iu\sin\theta)(\cos\varphi + \iu\sin\varphi)
= (\cos\theta\cos\varphi - \sin\theta\sin\varphi)
+ \iu\,(\sin\theta\cos\varphi + \cos\theta\sin\varphi),
\]
which is $\cos(\theta+\varphi) + \iu\sin(\theta+\varphi)$. The \omterm{def:b1:complex:field}{modulus}
is $\sqrt{\cos^2\theta + \sin^2\theta} = 1$, and the \omterm{def:b1:complex:field}{conjugate} formula
is the parity of cosine and sine; it inverts $\eu^{\iu\theta}$ since
$\eu^{\iu\theta}\eu^{-\iu\theta} = \eu^0 = 1$.
\end{proof}

\begin{corollary}[de Moivre's formula]\label{cor:b1:complex:demoivre}
For $\theta \in \R$ and $n \in \Z$:
$\;(\cos\theta + \iu\sin\theta)^n = \cos(n\theta) +
\iu\sin(n\theta)$.\index{de Moivre's formula}
\end{corollary}

\begin{proof}
For $n \geq 0$, induct. The case $n = 0$ reads $1 = 1$. Assuming
the formula for $n$, the functional equation
(\cref{thm:b1:complex:funceq}) gives
\[
\bigl(\eu^{\iu\theta}\bigr)^{n+1}
= \bigl(\eu^{\iu\theta}\bigr)^{n}\,\eu^{\iu\theta}
= \eu^{\iu n\theta}\,\eu^{\iu\theta}
= \eu^{\iu(n+1)\theta} ,
\]
which is the formula at rank $n + 1$. For $n < 0$, write $n = -m$
with $m > 0$: since $\eu^{\iu m\theta}$ has inverse
$\eu^{-\iu m\theta}$ (\cref{thm:b1:complex:funceq} again),
\[
\bigl(\eu^{\iu\theta}\bigr)^{-m}
= \bigl(\eu^{\iu m\theta}\bigr)^{-1}
= \eu^{\iu(-m)\theta} . \qedhere
\]
\end{proof}

\begin{example}[Expanding $\cos 3\theta$ by de Moivre]\label{ex:b1:complex:cos3}
Write $c = \cos\theta$, $s = \sin\theta$. De Moivre and the binomial
theorem give
\[
\cos 3\theta + \iu \sin 3\theta = (c + \iu s)^3
= c^3 - 3cs^2 + \iu\,(3c^2 s - s^3),
\]
and identifying real parts, then substituting $s^2 = 1 - c^2$:
\[
\cos 3\theta = c^3 - 3c(1 - c^2) = 4\cos^3\theta - 3\cos\theta .
\]
The imaginary part yields $\sin 3\theta = 3\sin\theta -
4\sin^3\theta$ for free: one complex identity always carries
\emph{two} real ones. Read backwards, the boxed identity is the
key to the classical trisection equation: constructing
$\cos(\theta/3)$ from $\cos\theta$ means solving the cubic $4x^3 -
3x = \cos\theta$, which is where the algebra of
\cref{ch:b1:poly} takes over.
\end{example}

\begin{proposition}[Euler's formulas]\label{prop:b1:complex:euler}
\[
\cos\theta = \frac{\eu^{\iu\theta} + \eu^{-\iu\theta}}{2},
\qquad
\sin\theta = \frac{\eu^{\iu\theta} - \eu^{-\iu\theta}}{2\iu}.
\]
\end{proposition}

\begin{proof}
Add, respectively subtract, $\eu^{\iu\theta} = \cos\theta +
\iu\sin\theta$ and $\eu^{-\iu\theta} = \cos\theta - \iu\sin\theta$.
\end{proof}

\begin{method}[Trigonometry via exponentials]\label{met:b1:complex:trig}
\begin{enumerate}
  \item \emph{Linearize} $\cos^p\theta \sin^q\theta$ (turn powers into
        a sum of $\cos k\theta$, $\sin k\theta$): substitute Euler's
        formulas, expand with the binomial theorem
        (\cref{thm:b1:counting:binomial}), regroup \omterm{def:b1:complex:field}{conjugate} terms.
  \item \emph{Expand} $\cos n\theta$ as a polynomial in $\cos\theta$:
        write $\cos n\theta = \Re\bigl((\cos\theta +
        \iu\sin\theta)^n\bigr)$, expand, and convert even powers of
        $\sin$ via $\sin^2 = 1 - \cos^2$.
  \item \emph{Sum trigonometric series} like $\sum_k \cos k\theta$:
        recognize the real part of a geometric sum $\sum_k
        (\eu^{\iu\theta})^k$.
  \item The \emph{half-angle factorization}: for all $p, q$,
        \[
        \eu^{\iu p} + \eu^{\iu q}
        = 2 \cos\tfrac{p - q}{2}\; \eu^{\iu \frac{p+q}{2}},
        \qquad
        \eu^{\iu p} - \eu^{\iu q}
        = 2\iu \sin\tfrac{p - q}{2}\; \eu^{\iu \frac{p+q}{2}} .
        \]
\end{enumerate}
\end{method}

\begin{example}[Linearization]\label{ex:b1:complex:linearize}
\[
\cos^3\theta
= \Bigl(\frac{\eu^{\iu\theta} + \eu^{-\iu\theta}}{2}\Bigr)^{\!3}
= \frac{\eu^{3\iu\theta} + 3\eu^{\iu\theta} + 3\eu^{-\iu\theta} +
\eu^{-3\iu\theta}}{8}
= \frac{\cos 3\theta + 3\cos\theta}{4}.
\]
This form integrates immediately --- the very reason linearization
matters in \cref{ch:b1:integration}. A mixed product works the
same way, only with both Euler formulas at once:
\[
\sin^2\theta\cos^2\theta
= \Bigl(\frac{\sin2\theta}2\Bigr)^{\!2}
= \frac{1}{4}\cdot
\Bigl(\frac{\eu^{2\iu\theta} - \eu^{-2\iu\theta}}{2\iu}
\Bigr)^{\!2}
= \frac{2 - \eu^{4\iu\theta} - \eu^{-4\iu\theta}}{16}
= \frac{1 - \cos4\theta}{8} ,
\]
where the double-angle shortcut in the first step saved one
binomial expansion --- always worth scanning for before
mechanizing.
\end{example}

\begin{example}[A binomial trigonometric sum]\label{ex:b1:complex:binomsum}
For $n \in \N$ and $\theta \in \R$, evaluate $S =
\sum_{k=0}^{n}\binom nk \cos k\theta$. Recognize the real part of
a binomial expansion:
\[
S = \Re\sum_{k=0}^n \binom nk \bigl(\eu^{\iu\theta}\bigr)^k
= \Re\bigl(1 + \eu^{\iu\theta}\bigr)^n ,
\]
then factor the half-angle (\cref{met:b1:complex:trig}~(4)): $1 +
\eu^{\iu\theta} = 2\cos\frac\theta2\,\eu^{\iu\theta/2}$, so
\[
S = \Re\Bigl(2^n\cos^n\frac\theta2\;\eu^{\iu n\theta/2}\Bigr)
= 2^n \cos^n\frac\theta2\,\cos\frac{n\theta}2 .
\]
The imaginary part yields $\sum_k\binom nk\sin k\theta =
2^n\cos^n\frac\theta2\sin\frac{n\theta}2$ for free. Sanity checks:
$\theta = 0$ recovers $\sum\binom nk = 2^n$, and $\theta = \pi$
gives $S = 0$ for $n \geq 1$ (each factor $\cos\frac\pi2$
vanishes), i.e.\ the alternating row sum of
\cref{ex:b1:counting:binomial}. The method --- ``see the cosine
sum as the shadow of a complex power, then factor half-angles''
--- is exactly that of \cref{exo:b1:complex:6}, with the binomial
theorem replacing the geometric series.
\end{example}

\section{Roots of complex numbers}

\begin{theorem}[$n$-th roots]\label{thm:b1:complex:roots}
Let $a = r\,\eu^{\iu\alpha} \neq 0$ and $n \in \N^*$. The equation
$z^n = a$ has exactly $n$ solutions:
\[
z_k = r^{1/n}\, \eu^{\iu\left(\frac{\alpha}{n} +
\frac{2k\pi}{n}\right)},
\qquad k = 0, 1, \dots, n - 1 .
\]
\end{theorem}

\begin{proof}
Write $z = \rho\,\eu^{\iu\theta}$ ($\rho > 0$). Then $z^n =
\rho^n \eu^{\iu n\theta} = r \eu^{\iu\alpha}$ if and only if $\rho^n
= r$ (moduli) and $n\theta \equiv \alpha \pmod{2\pi}$ (arguments),
i.e.\ $\rho = r^{1/n}$ and $\theta = \frac{\alpha}{n} +
\frac{2k\pi}{n}$ for some $k \in \Z$. It remains to see when two
integers $k, k'$ yield the same number: this happens exactly when
the angles differ by a multiple of $2\pi$,
\[
\frac{2k\pi}n - \frac{2k'\pi}n \in 2\pi\Z
\iff \frac{k - k'}n \in \Z
\iff n \mid k - k' .
\]
By Euclidean division, every $k \in \Z$ is congruent mod $n$ to
exactly one element of $\{0, 1, \dots, n-1\}$, so this range lists
each solution once and the count is exactly $n$. (The same
argument, run inside $\mathbb U_n$, shows the roots form a regular
$n$-gon: consecutive values of $k$ rotate by the fixed angle
$\frac{2\pi}n$.)
\end{proof}

\begin{example}[Cube roots of $-27$]\label{ex:b1:complex:cuberoots}
Solve $z^3 = -27$. Exponential form of the right side: $-27 =
27\,\eu^{\iu\pi}$, so the three roots are
\[
z_k = 3\,\eu^{\iu(\frac\pi3 + \frac{2k\pi}3)}, \quad k = 0, 1, 2 :
\qquad
z_0 = 3\eu^{\iu\pi/3} = \frac32 + \frac{3\sqrt3}2\,\iu,
\quad
z_1 = -3,
\quad
z_2 = \conj{z_0} .
\]
Two checks. First, the real root $-3$ is the obvious one, and the
other two are its rotations by $\pm\frac{2\pi}3$ --- equivalently
$-3j$ and $-3j^2$. Second, algebra confirms: $z^3 + 27 = (z +
3)(z^2 - 3z + 9)$, and the quadratic has discriminant $9 - 36 =
-27 < 0$ with roots $\frac{3 \pm 3\iu\sqrt3}2 = z_0, \conj{z_0}$.
The insight: for real right-hand sides, the non-real roots always
come in \omterm{def:b1:complex:field}{conjugate} pairs, so a picture of the solution \omterm{def:b1:logic:sets}{set} is
symmetric about the real axis --- a preview of the real
factorization theorem of \cref{ch:b1:poly}.
\end{example}

\begin{example}[A non-real right-hand side]\label{ex:b1:complex:fourthroots}
Solve $z^4 = -8 + 8\iu\sqrt3$. Exponential form of the right side:
\omterm{def:b1:complex:field}{modulus} $\sqrt{64 + 192} = 16$, argument $\theta$ with
$\cos\theta = -\frac12$, $\sin\theta = \frac{\sqrt3}2$, i.e.\
$\theta = \frac{2\pi}3$. The four roots are
\[
z_k = 2\,\eu^{\iu(\frac\pi6 + \frac{k\pi}2)}, \quad k = 0, 1, 2,
3 :
\qquad
z_0 = \sqrt3 + \iu,\quad z_1 = \iu z_0 = -1 + \iu\sqrt3,
\]
\[
z_2 = -z_0 = -\sqrt3 - \iu,\qquad z_3 = -\iu z_0 = 1 - \iu\sqrt3 .
\]
(Check: $z_0^2 = 2 + 2\iu\sqrt3$, so $z_0^4 = (2 +
2\iu\sqrt3)^2 = 4 - 12 + 8\iu\sqrt3 = -8 + 8\iu\sqrt3$.) Once
\emph{one} root is found, the other three come free: they are its
successive rotations by $\frac\pi2$, i.e.\ its products with the
fourth \omterm{def:b1:complex:unity}{roots of unity} --- the general structure behind
\cref{thm:b1:complex:roots}, worth exploiting before recomputing
each root from scratch. No \omterm{def:b1:complex:field}{conjugate} symmetry this time: the
right-hand side is not real.
\end{example}

\begin{definition}[Roots of unity]\label{def:b1:complex:unity}
The \emph{$n$-th roots of unity}\index{roots of unity} are the
solutions of $z^n = 1$:
\[
\mathbb{U}_n = \bigl\{\, \omega^k : k = 0, \dots, n-1 \,\bigr\},
\qquad \omega = \eu^{2\iu\pi/n}.
\]
They form a group under multiplication (\cref{ch:b1:structures}) and
sit at the vertices of a regular $n$-gon inscribed in the unit circle.
\end{definition}

\begin{omfigure}
\begin{tikzpicture}[scale=1.5]
  \draw[->, thin] (-1.35,0) -- (1.45,0) node[below] {$\Re$};
  \draw[->, thin] (0,-1.35) -- (0,1.45) node[left] {$\Im$};
  \draw[thin, gray] (0,0) circle (1);
  \foreach \k in {0,...,4} {
    \fill[omDef] ({cos(72*\k)},{sin(72*\k)}) circle (0.035);
  }
  \draw[omDef, thick] (1,0) -- ({cos(72)},{sin(72)}) -- ({cos(144)},{sin(144)})
    -- ({cos(216)},{sin(216)}) -- ({cos(288)},{sin(288)}) -- cycle;
  \node[right, font=\small, omDef] at (1.02,0.12) {$1$};
  \node[above right, font=\small, omDef] at ({cos(72)},{sin(72)})
    {$\omega$};
  \node[above left, font=\small, omDef] at ({cos(144)},{sin(144)})
    {$\omega^2$};
  \node[below left, font=\small, omDef] at ({cos(216)},{sin(216)})
    {$\omega^3$};
  \node[below right, font=\small, omDef] at ({cos(288)},{sin(288)})
    {$\omega^4$};
\end{tikzpicture}

{\small The fifth \omterm{def:b1:complex:unity}{roots of unity}, $\omega = \eu^{2\iu\pi/5}$: a
regular pentagon on the unit circle.}
\end{omfigure}

\begin{proposition}[Sum of the roots of unity]\label{prop:b1:complex:sumroots}
For $n \geq 2$, the $n$-th \omterm{def:b1:complex:unity}{roots of unity} sum to zero:
$\sum_{k=0}^{n-1} \omega^k = 0$.
\end{proposition}

\begin{proof}
Geometric sum with ratio $\omega \neq 1$:
$\sum_{k=0}^{n-1} \omega^k = \frac{\omega^n - 1}{\omega - 1} = 0$
since $\omega^n = 1$.
\end{proof}

\begin{example}[Reading real and imaginary parts]\label{ex:b1:complex:sumrootsparts}
Splitting $\sum_{k=0}^{n-1}\omega^k = 0$ into real and imaginary
parts yields two trigonometric identities for free:
\[
\sum_{k=0}^{n-1}\cos\frac{2k\pi}n = 0,
\qquad
\sum_{k=0}^{n-1}\sin\frac{2k\pi}n = 0
\qquad (n \geq 2).
\]
Geometrically: the centroid of a regular $n$-gon inscribed in the
unit circle is its center --- the vertices balance out exactly.
For $n = 5$ the first identity gives $1 + 2\cos\frac{2\pi}5 +
2\cos\frac{4\pi}5 = 0$ (pairing $k$ with $n - k$), the starting
point of \cref{exo:b1:complex:8}'s computation of
$\cos\frac{2\pi}5$.
\end{example}

\begin{example}[Square roots in algebraic form]\label{ex:b1:complex:sqrt}
To solve $z^2 = 3 + 4\iu$ without trigonometry, set $z = x + \iu y$:
\[
x^2 - y^2 = 3, \qquad 2xy = 4, \qquad x^2 + y^2 = \abs{3 + 4\iu} = 5 .
\]
Adding the first and last: $x^2 = 4$, so $x = \pm 2$, then $y = 2/x =
\pm 1$ with the \emph{same} sign pairing ($xy = 2 > 0$): $z = \pm(2 +
\iu)$. Combined with the usual formula, this solves every quadratic
equation with complex coefficients
(\cref{exo:b1:complex:7}).
\end{example}

\section{Complex numbers and plane geometry}

\begin{proposition}[Geometric dictionary]\label{prop:b1:complex:geometry}
Identify the point $M(x, y)$ of the plane with its \emph{affix} $z = x
+ \iu y$\index{affix}. For distinct points $A, B, C$ of affixes $a, b,
c$:
\begin{enumerate}
  \item $\abs{b - a}$ is the distance $AB$;
  \item $\arg\dfrac{c - a}{b - a}$ is the angle between the vectors
        $\vect{AB}$ and $\vect{AC}$ (mod $2\pi$);
  \item $A, B, C$ are aligned if and only if $\dfrac{c - a}{b - a} \in
        \R$; the lines $AB$ and $AC$ are perpendicular if and only if
        $\dfrac{c - a}{b - a} \in \iu\R$.
\end{enumerate}
\end{proposition}

\begin{proof}
(1) is the definition of the \omterm{def:b1:complex:field}{modulus} applied to $b - a$, the affix of
$\vect{AB}$. (2): write $b - a = r\eu^{\iu\theta}$, $c - a =
s\eu^{\iu\varphi}$; then $\frac{c-a}{b-a} = \frac sr
\eu^{\iu(\varphi - \theta)}$ has argument $\varphi - \theta$, the
angle from $\vect{AB}$ to $\vect{AC}$. (3): alignment means angle $0$
or $\pi$, i.e.\ argument in $\pi\Z$, i.e.\ the quotient is real;
perpendicularity means angle $\pm\frac\pi2$, i.e.\ the quotient is
purely imaginary. (The quotient is nonzero since $C \neq A$.)
\end{proof}

\begin{remark}[Interlude: $\C$ is the plane with a multiplication]\label{rem:b1:complex:interlude}
It is worth pausing on what makes this chapter possible at all:
the plane $\R^2$ carries additions in every direction, but no
God-given multiplication --- and $\C$ \emph{is} the plane
equipped with one, in which multiplying by a fixed number
rotates and scales. This one structure will be milked three more
times in this volume. In \cref{ch:b1:matrices}, multiplication by
$a + \iu b$ reappears as the $2 \times 2$ matrix with rows $(a,
-b)$ and $(b, a)$: complex arithmetic is a first, entirely
concrete family of matrix products. In \cref{ch:b1:euclid}, the
formula $\Re(\conj z\,w)$ turns out to be the dot product, and
$\abs z$ the Euclidean norm: the triangle inequality proved here
is the model for the Cauchy--Schwarz story there. And in
\cref{ch:b1:curves}, a moving point is best written $t \mapsto
z(t)$, so that velocity and acceleration become complex-valued
derivatives --- circular motion, for one, is simply $z(t) =
R\,\eu^{\iu\omega t}$. One good multiplication, four chapters of
dividends.
\end{remark}

\begin{proposition}[Maps $z \mapsto az + b$]\label{prop:b1:complex:similitude}
Let $a \in \C^*$, $b \in \C$. The transformation $f(z) = az + b$ of
the plane:
\begin{itemize}
  \item is a translation if $a = 1$;
  \item otherwise has a unique fixed point $\zeta = \frac{b}{1-a}$,
        and $f(z) - \zeta = a\,(z - \zeta)$: $f$ is the rotation of
        center $\zeta$ and angle $\arg a$, composed with the scaling
        of center $\zeta$ and ratio $\abs a$.
\end{itemize}
In particular $z \mapsto \eu^{\iu\theta} z$ is the rotation of angle
$\theta$ about the origin, and $\conj z$ is the reflection in the real
axis.
\end{proposition}

\begin{proof}
If $a = 1$, $f(z) = z + b$ translates by the vector of affix $b$. If
$a \neq 1$, the fixed-point equation $z = az + b$ has the unique
solution $\zeta = \frac{b}{1-a}$, and then $f(z) - \zeta = az + b -
(a\zeta + b) = a(z - \zeta)$. Writing $a = \abs a\, \eu^{\iu\arg a}$,
multiplication by $a$ scales distances to $\zeta$ by $\abs a$ and adds
$\arg a$ to angles at $\zeta$, which is the announced composition.
\end{proof}

\begin{example}[Classifying a map $z \mapsto az+b$]\label{ex:b1:complex:classify}
Take $f(z) = \iu z + 1$. Here $a = \iu \neq 1$: the fixed point is
\[
\zeta = \frac{b}{1 - a} = \frac1{1 - \iu}
= \frac{1 + \iu}{2},
\]
and since $\abs a = 1$, there is no scaling at all: $f$ is the pure
rotation of center $\frac{1+\iu}2$ and angle $\arg \iu = \frac\pi2$.
Check: $f(\zeta) = \iu\,\frac{1+\iu}2 + 1 = \frac{\iu - 1}2 + 1 =
\frac{1 + \iu}2 = \zeta$, and $f(0) = 1$, $f(1) = 1 + \iu$, $f(1 +
\iu) = \iu\,(1 + \iu) + 1 = \iu$: the four points $0, 1, 1+\iu, \iu$
of the unit square cycle around their center $\zeta$, one quarter
turn at a time --- exactly what a rotation by $\frac\pi2$ about the
square's center must do.
\end{example}

\begin{omfigure}
\begin{tikzpicture}[scale=1.6]
  \draw[->, thin] (-0.5,0) -- (1.6,0) node[below] {$\Re$};
  \draw[->, thin] (0,-0.35) -- (0,1.5) node[left] {$\Im$};
  \draw[gray, thin] (0,0) rectangle (1,1);
  \fill[omDef] (0.5,0.5) circle (0.03)
    node[below right, font=\small] {$\zeta$};
  \fill (0,0) circle (0.025) node[below left, font=\small] {$0$};
  \fill (1,0) circle (0.025) node[below right, font=\small] {$1$};
  \fill (1,1) circle (0.025)
    node[above right, font=\small] {$1{+}\iu$};
  \fill (0,1) circle (0.025) node[above left, font=\small] {$\iu$};
  \draw[omThm, ->, thick] (0.08,-0.12)
    .. controls (0.6,-0.28) and (1.12,0.4) .. (1.12,0.9);
  \draw[omThm, ->, thick] (1.12,1.08)
    .. controls (0.9,1.3) and (0.35,1.3) .. (0.05,1.12);
  \node[omThm, font=\small] at (0.95,-0.28) {$0 \mapsto 1$};
  \node[omThm, font=\small] at (1.62,0.75) {$1 \mapsto 1{+}\iu$};
\end{tikzpicture}

{\small The \omterm{def:b1:logic:map}{map} $f(z) = \iu z + 1$ of
\cref{ex:b1:complex:classify}: a quarter-turn rotation about the
fixed point $\zeta = \frac{1+\iu}2$. The vertices of the unit
square cycle $0 \to 1 \to 1{+}\iu \to \iu \to 0$; no point moves
along a straight line, yet the whole square is rigidly turned.}
\end{omfigure}

\begin{remark}[Common pitfalls with moduli and arguments]\label{rem:b1:complex:pitfalls}
\begin{enumerate}
  \item \emph{No inequalities in $\C$.} Writing $z \leq w$ for
        non-real numbers is meaningless; only moduli, real parts
        and imaginary parts can be compared.
  \item \emph{$\sqrt{\phantom z}$ is reserved for $\R_+$.} Every
        nonzero complex number has \emph{two} square roots and
        neither is privileged: write ``let $\delta$ be a square
        root of $\Delta$'' (computed as in
        \cref{ex:b1:complex:sqrt}), never $\sqrt\Delta$ ---
        the rule $\sqrt{ab} = \sqrt a\sqrt b$ already fails at
        $a = b = -1$.
  \item \emph{Arguments live modulo $2\pi$.} From
        $\eu^{\iu\alpha} = \eu^{\iu\beta}$ conclude $\alpha
        \equiv \beta \pmod{2\pi}$, not $\alpha = \beta$;
        forgetting the $2k\pi$ is how solution \omterm{def:b1:logic:sets}{sets} of $z^n = a$
        lose $n - 1$ of their $n$ elements.
  \item \emph{$\abs{z + w}$ is not $\abs z + \abs w$.} Equality
        in the triangle inequality is the exceptional aligned
        case (\cref{prop:b1:complex:rules}~(4)); in general the
        \omterm{def:b1:complex:field}{modulus} of a sum must be estimated, not computed.
\end{enumerate}
\end{remark}

\begin{example}\label{ex:b1:complex:equilateral}
$A, B, C$ (affixes $a, b, c$, pairwise distinct) form an equilateral
triangle with the vertices in direct (counterclockwise) order if and
only if $\frac{c - a}{b - a} = \eu^{\iu\pi/3}$: the rotation of
center $A$ and angle $\frac\pi3$ sends $B$ to $C$. Both orientations
together are captured by the symmetric equation $a^2 + b^2 + c^2 = ab
+ bc + ca$ (\cref{exo:b1:complex:10}).
\end{example}

\begin{remark}[Where this chapter is used]\label{rem:b1:complex:whereused}
The exponential form is the single most reused computational device
of the volume. \omterm{def:b1:complex:unity}{Roots of unity} become the standard example of a
cyclic group in \cref{ch:b1:structures} and drive the factorization
of $X^n - 1$ in \cref{ch:b1:poly}; grouping \omterm{def:b1:complex:field}{conjugate} roots there
produces the real factorizations used by partial fractions in
\cref{ch:b1:fractions}. Linearization (\cref{met:b1:complex:trig})
is the standard preparation for integrating trigonometric powers in
\cref{ch:b1:integration}, and the characteristic equations of
\cref{ch:b1:diffeq} have complex roots whose real and imaginary
parts produce the oscillating solutions $\eu^{\lambda t}\cos\omega
t$. The geometric dictionary returns in matrix clothing: rotations
and similarities become the orthogonal matrices of
\cref{ch:b1:matrices,ch:b1:euclid}, and parametrized curves in
\cref{ch:b1:curves} are often best written as \omterm{def:b1:logic:map}{maps} $t \mapsto z(t)
\in \C$.
\end{remark}

\section{Exercises}

\begin{exercise}[$\star$]\label{exo:b1:complex:1}
Put in exponential form: $1 + \iu$; $\;\sqrt 3 - \iu$; $\;-5$;
$\;\dfrac{1 + \iu}{\sqrt 3 - \iu}$. Deduce
$\cos\frac{5\pi}{12}$ and $\sin\frac{5\pi}{12}$.
\end{exercise}

\begin{exercise}[$\star$]\label{exo:b1:complex:2}
Compute $(1 + \iu)^{20}$. For which $n \in \N$ is $(1 + \iu)^n$ a
real number?
\end{exercise}

\begin{exercise}[$\star$]\label{exo:b1:complex:3}
Describe geometrically the \omterm{def:b1:logic:sets}{set} of $z \in \C$ such that:
$\abs{z - 2} = \abs{z + \iu}$; $\;\abs{z - 1} = 2$;
$\;\dfrac{z - 1}{z + 1} \in \iu\R$ (for $z \neq -1$).
\end{exercise}

\begin{exercise}[$\star$]\label{exo:b1:complex:4}
Linearize $\sin^4\theta$, and expand $\cos 4\theta$ as a polynomial in
$\cos\theta$.
\end{exercise}

\begin{exercise}[$\star$]\label{exo:b1:complex:5}
Solve $z^3 = 8\iu$ and place the solutions on a picture. Solve
$z^4 = -4$ and factor $X^4 + 4$ into two real quadratic polynomials.
\end{exercise}

\begin{exercise}[$\star\star$]\label{exo:b1:complex:6}
For $\theta \in \R$ with $\eu^{\iu\theta} \neq 1$ and $n \in \N$,
compute
\[
C_n = \sum_{k=0}^{n} \cos k\theta
\qquad\text{and}\qquad
S_n = \sum_{k=0}^{n} \sin k\theta
\]
by summing the geometric series $\sum_k \eu^{\iu k\theta}$ and using
the half-angle factorization.
\end{exercise}

\begin{exercise}[$\star\star$]\label{exo:b1:complex:7}
Solve in $\C$: $z^2 - (3 + 4\iu) z + (-1 + 5\iu) = 0$. \emph{(Compute
the discriminant, extract its square roots as in
\cref{ex:b1:complex:sqrt}.)}
\end{exercise}

\begin{exercise}[$\star\star$]\label{exo:b1:complex:8}
Let $\omega = \eu^{2\iu\pi/5}$.
\begin{enumerate}
  \item Justify $1 + \omega + \omega^2 + \omega^3 + \omega^4 = 0$.
  \item Set $u = \omega + \omega^4$ and $v = \omega^2 + \omega^3$.
        Compute $u + v$ and $uv$, and deduce that $u$ and $v$ are the
        roots of $X^2 + X - 1$.
  \item Conclude that $\cos\frac{2\pi}{5} = \frac{\sqrt 5 - 1}{4}$.
\end{enumerate}
\end{exercise}

\begin{exercise}[$\star\star$]\label{exo:b1:complex:9}
Prove that for all $z, w \in \C$ (parallelogram identity):
\[
\abs{z + w}^2 + \abs{z - w}^2 = 2\abs z^2 + 2\abs w^2 ,
\]
and interpret it geometrically in the parallelogram with vertices $0,
z, w, z + w$.
\end{exercise}

\begin{exercise}[$\star\star\star$]\label{exo:b1:complex:10}
Prove that three pairwise distinct points of affixes $a, b, c$ form an
equilateral triangle (either orientation) if and only if
\[
a^2 + b^2 + c^2 = ab + bc + ca .
\]
\emph{Hint: the direct case is $\frac{c-a}{b-a} = -j^2$ and the
indirect case $\frac{c-a}{b-a} = -j$, where $j = \eu^{2\iu\pi/3}$
satisfies $j^2 + j + 1 = 0$; or factor $a + jb + j^2c$ and $a + j^2 b
+ jc$.}
\end{exercise}

\begin{exercise}[$\star\star\star$]\label{exo:b1:complex:11}
For $n \in \N^*$, evaluate $P = \prod_{k=1}^{n-1}
\bigl(1 - \omega^k\bigr)$ where $\omega = \eu^{2\iu\pi/n}$.
\emph{Hint: $X^n - 1 = \prod_{k=0}^{n-1} (X - \omega^k)$; divide by
$X - 1$ and evaluate at $X = 1$.} Deduce
$\prod_{k=1}^{n-1} \sin\frac{k\pi}{n} = \dfrac{n}{2^{n-1}}$.
\end{exercise}

\begin{exercise}[$\star\star$]\label{exo:b1:complex:12}
Let $n \geq 2$. Solve the equation $(z + 1)^n = (z - 1)^n$ in $\C$:
show that it has exactly $n - 1$ solutions, all purely imaginary,
namely
\[
z_k = -\iu\,\frac{\cos\frac{k\pi}{n}}{\sin\frac{k\pi}{n}},
\qquad k = 1, \dots, n - 1 .
\]
\emph{Hint: $z = 1$ is not a solution, so divide and use the \omterm{def:b1:complex:unity}{roots of
unity}; then apply the half-angle factorization of
\cref{met:b1:complex:trig}.}
\end{exercise}

\section{Problem: Napoleon's theorem}

\begin{problem}\label{pb:b1:complex:1}
Erect an equilateral triangle outward on each side of an
\emph{arbitrary} triangle: the three centers of those triangles
always form an equilateral triangle. This is Napoleon's theorem
\index{Napoleon's theorem} --- a \omterm{def:b1:logic:statement}{statement} with no visible reason to
be true, which the algebra of $j = \eu^{2\iu\pi/3}$ proves in three
lines of computation. This problem builds the full toolkit
(rotations, direct similarities, the $j$-criterion of
\cref{exo:b1:complex:10}), proves Napoleon's theorem inner and
outer, locates the degenerate case, and closes with a second jewel
of the same school: van Schooten's theorem on equilateral triangles
and Ptolemy's inequality. Throughout, points of the plane are
identified with their affixes.

\medskip
\noindent\textbf{Part I --- The number $j$ and rotations.}
\begin{enumerate}
  \item Compute the algebraic form of $j$, then $j^2$, $j^3$, $1 + j
        + j^2$, $\conj j$ and $j^{-1}$. Place $1$, $j$, $j^2$ on a
        sketch of the unit circle.
  \item Show that the rotation of center $a$ and angle $\theta$ is
        $r(z) = a + \eu^{\iu\theta}(z - a)$, and that the
        composition of two rotations, of angles $\theta$ and
        $\theta'$, is a rotation of angle $\theta + \theta'$ if
        $\theta + \theta' \notin 2\pi\Z$, and a translation
        otherwise. (Use \cref{prop:b1:complex:similitude}.)
  \item Let $b \neq c$. Show that there are exactly two points $p$
        making $(b, c, p)$ equilateral, namely $p = b +
        \eu^{\pm\iu\pi/3}(c - b)$. For a \emph{direct} triangle
        $(a, b, c)$, check on the example $a = 0$, $b = 1$, $c =
        \iu$ that the choice $\eu^{-\iu\pi/3}$ is the one lying on
        the far side of the line $BC$ from $a$ --- the
        \emph{outward} apex.
  \item From the solution of \cref{exo:b1:complex:10}, $(a, b, c)$
        is direct equilateral iff $a + jb + j^2c = 0$. Prove the two
        rotation identities
        \[
        b + jc + j^2a = j^2\,(a + jb + j^2c),
        \qquad
        c + ja + j^2b = j\,(a + jb + j^2c),
        \]
        and deduce that the criterion is invariant under cyclic
        \omterm{def:b1:counting:objects}{permutation} of $(a, b, c)$.
  \item Show that if $a + jb + j^2c = 0$ and two of the three points
        coincide, then all three coincide. (So the criterion exactly
        characterizes: direct equilateral triangle, or a single
        point.)
\end{enumerate}

\medskip
\noindent\textbf{Part II --- Direct similarities.}
\begin{enumerate}[resume]
  \item Show that the \omterm{def:b1:logic:map}{maps} $f(z) = az + b$ with $a \in \C^*$
        (\emph{direct similarities}) are stable under composition
        and inversion: the composition of two of them, and the
        inverse of any of them, is again of this form. (In the
        language of \cref{ch:b1:structures}: they form a group.)
  \item Given $z_1 \neq z_2$ and $w_1 \neq w_2$, show there is a
        \emph{unique} direct similarity $f$ with $f(z_1) = w_1$ and
        $f(z_2) = w_2$, and give $a$ and $b$ explicitly.
  \item Show that a direct similarity $f(z) = az + b$ multiplies
        all distances by $\abs a$ and preserves the \emph{shape}
        $\frac{c - a'}{b' - a'}$ of any triangle $(a', b', c')$
        \emph{(numerator and denominator are both multiplied by
        $a$)}. Deduce that two triangles are directly similar
        exactly when their shapes are equal.
  \item Determine completely the direct similarity with $f(0) = 1$
        and $f(1) = \iu$: give $a$, $b$, the fixed point, the ratio
        and the angle.
  \item Let $\alpha + \beta = 1$. Show that the ``weighted point''
        $g(a', b') = \alpha a' + \beta b'$ commutes with every
        direct similarity: $f(\alpha a' + \beta b') = \alpha f(a')
        + \beta f(b')$. Deduce that centroids, midpoints, and the
        Napoleon centers below are all transported by similarities
        --- the algebraic license behind every ``without loss of
        generality, place the circumcircle at the unit circle''
        argument.
\end{enumerate}

\medskip
\noindent\textbf{Part III --- Napoleon's theorem.} Let $(a, b, c)$
be a direct triangle. On each side erect the outward equilateral
triangle (question 3) and let $n_a$, $n_b$, $n_c$ be the centers
(centroids) of the triangles erected on $[b, c]$, $[c, a]$,
$[a, b]$ respectively.
\begin{enumerate}[resume]
  \item Show that
        \[
        n_a = \alpha b + \beta c, \qquad
        n_b = \alpha c + \beta a, \qquad
        n_c = \alpha a + \beta b,
        \qquad\text{with}\quad
        \alpha = \frac{3 + \iu\sqrt3}{6},\ \beta = \conj\alpha .
        \]
  \item Check $\alpha + \beta = 1$, and deduce that the triangle
        $(n_a, n_b, n_c)$ has the \emph{same centroid} as
        $(a, b, c)$.
  \item Using the identities of question 4, show that
        \[
        n_a + j\,n_b + j^2 n_c
        = (\alpha j^2 + \beta j)\,(a + jb + j^2c) .
        \]
  \item Compute $\alpha j + \beta$ and conclude: $\alpha j^2 +
        \beta j = j(\alpha j + \beta) = 0$, so $n_a + j n_b + j^2
        n_c = 0$ for \emph{every} triangle: the outer Napoleon
        triangle is direct equilateral (or a point). Napoleon's
        theorem is proved.
  \item Erect the equilateral triangles \emph{inward} instead
        (choice $\eu^{+\iu\pi/3}$ in question 3) and let $m_a, m_b,
        m_c$ be their centers. Show $m_a = \beta b + \alpha c$ (and
        cyclically), then prove $m_a + j^2 m_b + j\,m_c = 0$: the
        inner Napoleon triangle is equilateral too, with the
        opposite orientation.
  \item Locate the degeneracy: show that $n_a = n_b = n_c$ happens
        exactly when $a + j^2 b + jc = 0$, i.e.\ when $(a, b, c)$
        is an \emph{indirect} equilateral triangle. \emph{(Two
        points among $n_a, n_b, n_c$ coincide iff all three do, by
        question 5; then use $n_a + j^2 n_b + j n_c = j(\alpha +
        \beta j)(a + j^2b + jc)$ and check $\alpha + \beta j \neq
        0$.)}
  \item Carry out the whole computation on the triangle $a = 0$,
        $b = 1$, $c = \iu$: give $n_a, n_b, n_c$ exactly, and
        verify by direct computation of the three squared side
        lengths that the triangle is equilateral, of squared side
        $\frac{2 + \sqrt3}{3}$.
\end{enumerate}

\medskip
\noindent\textbf{Part IV --- Ptolemy and van Schooten.}
\begin{enumerate}[resume]
  \item Prove the identity, valid for all complex $a, b, c, d$:
        \[
        (a - b)(c - d) + (a - d)(b - c) = (a - c)(b - d) .
        \]
  \item Deduce \emph{Ptolemy's inequality}\index{Ptolemy's
        inequality}: for any four points $A, B, C, D$,
        \[
        AC \cdot BD \;\leq\; AB \cdot CD + AD \cdot BC ,
        \]
        with equality if and only if $(a-b)(c-d)$ and
        $(a-d)(b-c)$ lie on a common ray from $0$.
  \item Show that for $\theta, \varphi \in \R$,
        $\abs{\eu^{\iu\theta} - \eu^{\iu\varphi}} =
        2\,\abs{\sin\frac{\theta - \varphi}2}$.
  \item (van Schooten's theorem\index{van Schooten's theorem}) Let
        $(a, b, c) = (1, j, j^2)$ --- by Part II this loses no
        generality among direct equilateral triangles --- and let
        $p = \eu^{\iu\theta}$ with $\theta \in
        \intoo{2\pi/3}{4\pi/3}$, a point of the circumcircle on the
        arc $BC$ not containing $A$. Prove
        \[
        PA = PB + PC .
        \]
        \emph{(Express the three distances with question 20 and use
        a sum-to-product formula.)}
  \item Verify at $p = -1$ that this is exactly the equality case
        of Ptolemy's inequality for the cyclic order $A, B, P, C$:
        compute $(a - b)(p - c)$ and $(a - c)(b - p)$ and check
        that their ratio is a positive real.
\end{enumerate}

\medskip
\noindent\textbf{Part V --- Synthesis.}
\begin{enumerate}[resume]
  \item Compute the outer Napoleon triangle of the equilateral
        triangle $(1, j, j^2)$ itself, and describe the result
        geometrically.
  \item Where exactly did the problem use: (i) multiplication as
        rotation; (ii) the group structure of similarities and the
        shape invariant; (iii) the half-angle factorization of
        \cref{met:b1:complex:trig}? One sentence each.
  \item In a short paragraph, articulate the moral of the problem:
        what does the dictionary between plane geometry and the
        algebra of $\C$ buy, what does it cost, and which of the
        two proofsteps --- the identity $\alpha j + \beta = 0$ or
        the classical picture --- \emph{explains} Napoleon's
        theorem better in your view? Mention one place where the
        dictionary will reappear in matrix form later in this
        volume.
\end{enumerate}
\end{problem}

\chapter{Standard Functions}\label{ch:b1:functions}

Analysis is only as useful as the stock of functions one masters. To
the collection inherited from secondary school --- powers, exponential,
logarithm, trigonometric functions --- this chapter adds their inverse
functions ($\arcsin$, $\arccos$, $\arctan$) and the hyperbolic family.
Derivatives are used freely at High School level; the
theory behind inverse functions is completed in
\cref{ch:b1:continuity,ch:b1:derivative}.

\section{Exponential, logarithm, powers}

\begin{proposition}[Recap and characterization]\label{prop:b1:functions:expln}
$\exp \colon \R \to \intoo{0}{+\infty}$ and $\ln \colon
\intoo{0}{+\infty} \to \R$ are reciprocal bijections, strictly
increasing, with
\[
\exp(x + y) = \exp x \exp y, \qquad \ln(xy) = \ln x + \ln y,
\qquad \exp' = \exp, \qquad \ln'(x) = \frac 1x .
\]
Moreover $\exp$ is the \emph{only} differentiable function $f \colon
\R \to \R$ with $f' = f$ and $f(0) = 1$.
\end{proposition}

\begin{proof}
The recap is High School material. For uniqueness, let $f' = f$,
$f(0) = 1$, and set $g(x) = f(x)\,\eu^{-x}$. Then $g' = f'\eu^{-x} -
f\eu^{-x} = 0$, so $g$ is constant equal to $g(0) = 1$: $f = \exp$.
\end{proof}

\begin{definition}[General powers]\label{def:b1:functions:powers}
For $x > 0$ and $\alpha \in \R$:
$\;x^\alpha = \eu^{\alpha \ln x}$\index{power function}. For $a > 0$,
$a \neq 1$, the \emph{logarithm to base $a$} is $\log_a x =
\frac{\ln x}{\ln a}$, the inverse of $x \mapsto a^x$.
\end{definition}

\begin{example}[Solving exponential equations]\label{ex:b1:functions:expequations}
Solve $2^x = 5^{\,x-1}$ in $\R$. Both sides are positive, so take
logarithms --- a reversible step:
\[
x\ln2 = (x - 1)\ln5
\iff x(\ln2 - \ln5) = -\ln5
\iff x = \frac{\ln5}{\ln5 - \ln2} = \frac{\ln 5}{\ln\frac52}
\approx 1.756 .
\]
Next, solve $x^{\sqrt2} = 3$ for $x > 0$: raise to the power
$\frac1{\sqrt2}$ (i.e.\ apply the reciprocal bijection):
$x = 3^{1/\sqrt2} = \eu^{(\ln 3)/\sqrt2} \approx 2.175$. The
insight: every equation mixing powers unwinds through $\ln$ and
$\exp$, because the definition $x^\alpha = \eu^{\alpha\ln x}$
reduces all power manipulations to arithmetic of exponents ---
but only on the domain $x > 0$ where that definition lives.
\end{example}

\begin{example}[Doubling times]\label{ex:b1:functions:doubling}
A quantity grows by $3\%$ per step: after $n$ steps it is
multiplied by $(1.03)^n$. When does it double? Solve $(1.03)^n
\geq 2$:
\[
n\ln(1.03) \geq \ln2
\iff n \geq \frac{\ln 2}{\ln 1.03}
= \frac{0.6931}{0.02956} \approx 23.4 ,
\]
so the first doubling occurs at step $24$. (The financiers'
``rule of $72$'', which estimates the doubling time as
$72$ divided by the rate in percent, is this computation with
the approximation $\ln(1 + x) \approx x$, quantified in
\cref{ch:b1:taylor}.) Exponential
processes are best reasoned about through their logarithms: on
that scale, growth is linear and questions become divisions.
\end{example}

\begin{example}[How long is $2^{2026}$?]\label{ex:b1:functions:digits}
The number of decimal digits of an integer $N \geq 1$ is
$\floor{\log_{10} N} + 1$ (indeed $N$ has $d$ digits exactly when
$10^{d-1} \leq N < 10^d$, i.e.\ $d - 1 \leq \log_{10}N < d$). For
$N = 2^{2026}$:
\[
\log_{10} 2^{2026} = 2026\,\log_{10}2
= 2026 \times 0.301030 = 609.887 ,
\]
so $2^{2026}$ has $610$ digits. The fractional part carries a
bonus: $10^{0.887} \approx 7.7$, so the number \emph{begins} with
$7$. One multiplication answered a question about a number no one
will ever write out --- logarithms compress multiplicative size
into additive size, which is their whole historical point
(\cref{ex:b1:structures:expmorphism} makes that sentence precise).
\end{example}

\begin{proposition}[Power rules and growth comparison]\label{prop:b1:functions:powerrules}
For $x, y > 0$ and $\alpha, \beta \in \R$:
\[
x^{\alpha+\beta} = x^\alpha x^\beta, \quad
(x^\alpha)^\beta = x^{\alpha\beta}, \quad
(xy)^\alpha = x^\alpha y^\alpha, \quad
(x^\alpha)' = \alpha\, x^{\alpha - 1}.
\]
Growth scale\index{growth comparison} as $x \to +\infty$, for every
$\alpha > 0$ and $\beta > 0$:
\[
\frac{(\ln x)^\beta}{x^\alpha} \longrightarrow 0,
\qquad
\frac{x^\alpha}{\eu^{\beta x}} \longrightarrow 0 .
\]
\end{proposition}

\begin{proof}
The identities transcribe those of $\exp$ and $\ln$ through the
definition. In detail for the second (the least obvious): $x^\alpha
> 0$ and $\ln(x^\alpha) = \alpha\ln x$ (apply $\ln$ to the
definition), so
\[
(x^\alpha)^\beta = \eu^{\beta\ln(x^\alpha)}
= \eu^{\beta\alpha\ln x} = x^{\alpha\beta} ;
\]
the first and third are the same two-line unwindings via
$\exp(u + v) = \exp u\exp v$ and $\ln(xy) = \ln x + \ln y$. The
derivative is the chain rule: $(\eu^{\alpha\ln x})' =
\frac{\alpha}{x} \eu^{\alpha\ln x}$.

Comparisons: from $\frac{\ln t}{t} \to 0$ (proved in the High
School volume), substitute $t = x^{\alpha/\beta}$:
$\frac{\ln x}{x^{\alpha/\beta}} \to 0$, then raise to the power
$\beta$. For the second, substitute $x = \ln u$ in the first.
\end{proof}

\begin{example}[Two limits every reader should own]\label{ex:b1:functions:xx}
What are $\lim_{x \to 0^+} x^x$ and $\lim_{x \to +\infty}
x^{1/x}$? Both powers are defined through the exponential, so
settle the exponent first:
\[
x^x = \eu^{x\ln x}, \qquad
x\ln x = -\frac{\ln(1/x)}{1/x} \longrightarrow 0
\quad (x \to 0^+),
\qquad\text{so } x^x \longrightarrow \eu^0 = 1 ;
\]
and $x^{1/x} = \eu^{(\ln x)/x} \to \eu^0 = 1$ as $x \to +\infty$,
directly by the growth comparison. The insight: an indeterminate
power (of shape $0^0$ or $\infty^0$) is \emph{always} handled by
rewriting $u^v = \eu^{v\ln u}$ and analyzing the product $v \ln u$
--- never by guessing from base and exponent separately. The
function $x \mapsto \frac{\ln x}{x}$ that decided both limits is
studied exhaustively in the weekend problem of this chapter.
\end{example}

\begin{omfigure}
\begin{tikzpicture}
\begin{axis}[omaxis, width=9.5cm, height=6cm,
    xmin=0, xmax=5.2, ymin=-1.8, ymax=8.5,
    xtick={1,2,3,4,5}, ytick={2,4,6,8}]
  \addplot[omMeth, thick, domain=0.02:5, samples=200] {ln(x)}
    node[pos=0.97, below, font=\small] {$\ln x$};
  \addplot[omDef, thick, domain=0:5, samples=100] {sqrt(x)}
    node[pos=0.9, above, font=\small] {$\sqrt x$};
  \addplot[omProp, thick, domain=0:5] {x}
    node[pos=0.72, left, font=\small] {$x$};
  \addplot[omThm, thick, domain=0:2.14, samples=120] {exp(x)}
    node[pos=0.85, left, font=\small] {$\eu^x$};
\end{axis}
\end{tikzpicture}

{\small The growth scale of
\cref{prop:b1:functions:powerrules}: near the right edge, $\ln x$
has barely passed $1.6$ while $\eu^x$ has left the frame. Every
ratio (log over power, power over exponential) tends to $0$ ---
the picture only hints at what the substitutions in the proof make
exact.}
\end{omfigure}

\begin{remark}[Common pitfalls with powers and logarithms]\label{rem:b1:functions:pitfalls}
\begin{enumerate}
  \item \emph{Domains.} $x^\alpha$ for irrational $\alpha$ needs
        $x > 0$; $(-8)^{1/3}$ is best avoided in favor of ``the
        real cube root of $-8$'', because the rule
        $(x^\alpha)^\beta = x^{\alpha\beta}$ silently fails on
        negatives: $\bigl((-8)^2\bigr)^{1/6} = 2 \neq -2$.
  \item \emph{$\sqrt{x^2} = \abs x$}, not $x$: forgetting the
        absolute value is the classic source of lost negative
        solutions.
  \item \emph{$\ln(xy) = \ln x + \ln y$ requires $x, y > 0$.} For
        negative arguments, $\ln(xy)$ can be defined while the
        right-hand side is not.
  \item \emph{Which function grows?} In $x^\alpha$ the
        \emph{base} varies; in $a^x$ the \emph{exponent} does.
        Their growths are wildly different
        (\cref{prop:b1:functions:powerrules}), and hybrid
        expressions like $x^{\ln x}$ or $x^{1/x}$ must be
        rewritten as $\eu^{(\cdot)}$ before any reasoning ---
        as in \cref{ex:b1:functions:xx}.
\end{enumerate}
\end{remark}

\section{Inverse trigonometric functions}

\begin{definition}[$\arcsin$, $\arccos$, $\arctan$]\label{def:b1:functions:arc}
The restrictions
\[
\sin \colon \intcc{-\tfrac\pi2}{\tfrac\pi2} \to \intcc{-1}{1},
\qquad
\cos \colon \intcc{0}{\pi} \to \intcc{-1}{1},
\qquad
\tan \colon \intoo{-\tfrac\pi2}{\tfrac\pi2} \to \R
\]
are strictly monotonic bijections. Their inverse maps are written
$\arcsin$\index{arcsine}, $\arccos$\index{arccosine} and
$\arctan$\index{arctangent}. Thus, for instance, $y = \arcsin x$
($x \in \intcc{-1}{1}$) is \emph{the} angle in
$\intcc{-\frac\pi2}{\frac\pi2}$ whose sine is $x$.
\end{definition}

\begin{omfigure}
\begin{tikzpicture}
\begin{axis}[omaxis, width=5.6cm, height=5.6cm,
    xmin=-1.25, xmax=1.25, ymin=-1.8, ymax=1.8,
    xtick={-1,1}, ytick={-1.5708,1.5708},
    yticklabels={$-\frac{\pi}{2}$,$\frac{\pi}{2}$}]
  \addplot[omDef, very thick, domain=-1:1, samples=200]
    {rad(asin(x))};
\end{axis}
\begin{scope}[xshift=5.3cm]
\begin{axis}[omaxis, width=5.6cm, height=5.6cm,
    xmin=-1.25, xmax=1.25, ymin=-0.4, ymax=3.4,
    xtick={-1,1}, ytick={1.5708,3.1416},
    yticklabels={$\frac{\pi}{2}$,$\pi$}]
  \addplot[omThm, very thick, domain=-1:1, samples=200]
    {rad(acos(x))};
\end{axis}
\end{scope}
\begin{scope}[xshift=10.6cm]
\begin{axis}[omaxis, width=5.6cm, height=5.6cm,
    xmin=-4, xmax=4, ymin=-1.9, ymax=1.9,
    xtick={-2,2}, ytick={-1.5708,1.5708},
    yticklabels={$-\frac{\pi}{2}$,$\frac{\pi}{2}$}]
  \addplot[omMeth, very thick, domain=-4:4, samples=200]
    {rad(atan(x))};
  \addplot[dashed, gray, domain=-4:4] {1.5708};
  \addplot[dashed, gray, domain=-4:4] {-1.5708};
\end{axis}
\end{scope}
\end{tikzpicture}

{\small From left to right: $\arcsin$, $\arccos$, $\arctan$. Each
inherits its graph from the restricted direct function by reflection
in the line $y = x$.}
\end{omfigure}

\begin{proposition}[Derivatives]\label{prop:b1:functions:arcderiv}
On the interior of their domains:
\[
\arcsin' x = \frac{1}{\sqrt{1 - x^2}},
\qquad
\arccos' x = \frac{-1}{\sqrt{1 - x^2}},
\qquad
\arctan' x = \frac{1}{1 + x^2}.
\]
\end{proposition}

\begin{proof}
Anticipating the inverse-function rule proved in
\cref{ch:b1:derivative}: if $f$ is a differentiable bijection and
$f' \neq 0$, then $(f^{-1})'(x) = \frac{1}{f'(f^{-1}(x))}$. For
$\arcsin$: with $y = \arcsin x$,
\[
\arcsin' x = \frac{1}{\cos y}
= \frac{1}{\sqrt{1 - \sin^2 y}} = \frac{1}{\sqrt{1 - x^2}},
\]
where $\cos y = +\sqrt{1 - \sin^2 y}$ because $y \in
\intoo{-\frac\pi2}{\frac\pi2}$ forces $\cos y > 0$. Similarly
$\arccos' x = \frac{1}{-\sin y}$ with $\sin y > 0$ on
$\intoo{0}{\pi}$, and $\arctan' x = \frac{1}{1 + \tan^2 y} =
\frac{1}{1 + x^2}$ using $\tan' = 1 + \tan^2$.
\end{proof}

\begin{example}[Recognizing a hidden constant]\label{ex:b1:functions:hiddenconst}
Study $g(x) = \arctan\dfrac{1 - x}{1 + x}$ on $\intoo{-1}
{+\infty}$. Chain rule and a short computation:
\[
g'(x) = \frac{1}{1 + \bigl(\frac{1-x}{1+x}\bigr)^2}\cdot
\frac{-(1+x) - (1-x)}{(1+x)^2}
= \frac{-2}{(1+x)^2 + (1-x)^2}
= \frac{-1}{1 + x^2} ,
\]
since $(1+x)^2 + (1-x)^2 = 2 + 2x^2$. So $g' = -\arctan'$: the
function $g + \arctan$ has zero derivative on the \emph{interval}
$\intoo{-1}{+\infty}$, hence is constant there; its value at $x =
0$ is $\arctan 1 + \arctan 0 = \frac\pi4$. Conclusion:
\[
\arctan\frac{1 - x}{1 + x} = \frac\pi4 - \arctan x
\qquad (x > -1) .
\]
On $\intoo{-\infty}{-1}$ the same derivative computation holds but
the constant is different ($-\frac{3\pi}4$: evaluate the limit as
$x \to -\infty$). The insight: ``zero derivative implies
constant'' is an interval-by-interval statement --- exactly the
subtlety exploited in \cref{exo:b1:functions:6} and
\cref{exo:b1:functions:10}.
\end{example}

\begin{proposition}[Standard identities]\label{prop:b1:functions:arcidentities}
\begin{enumerate}
  \item For $x \in \intcc{-1}{1}$: $\arcsin x + \arccos x =
        \dfrac{\pi}{2}$.
  \item For $x > 0$: $\arctan x + \arctan\dfrac 1x = \dfrac{\pi}{2}$
        (and $-\frac\pi2$ for $x < 0$).
  \item $\sin(\arccos x) = \cos(\arcsin x) = \sqrt{1 - x^2}$;
        $\;\tan(\arcsin x) = \dfrac{x}{\sqrt{1 - x^2}}$ for
        $\abs x < 1$.
\end{enumerate}
\end{proposition}

\begin{proof}
(1) The derivative of $x \mapsto \arcsin x + \arccos x$ is zero on
$\intoo{-1}{1}$ (\cref{prop:b1:functions:arcderiv}), so the function
is constant there, equal to its value $\frac\pi2$ at $0$; the endpoint
values $\pm 1$ are checked directly ($\frac\pi2 + 0$ and $-\frac\pi2 +
\pi$).

(2) Same method on $\intoo{0}{+\infty}$: the derivative is
$\frac{1}{1+x^2} + \frac{-1/x^2}{1 + 1/x^2} = \frac{1}{1+x^2} -
\frac{1}{x^2 + 1} = 0$, and at $x = 1$ the sum is $2\arctan 1 =
\frac\pi2$. For $x < 0$, use oddness of $\arctan$.

(3) With $y = \arccos x \in \intcc{0}{\pi}$: $\sin y \geq 0$, so
$\sin y = \sqrt{1 - \cos^2 y} = \sqrt{1 - x^2}$; likewise for
$\cos(\arcsin x)$, and the tangent formula is the quotient.
\end{proof}

\begin{remark}
$\arcsin(\sin\theta) = \theta$ holds \emph{only} for $\theta \in
\intcc{-\frac\pi2}{\frac\pi2}$: for instance
$\arcsin(\sin\pi) = 0 \neq \pi$. The composition in the other
direction, $\sin(\arcsin x) = x$, is valid on all of
$\intcc{-1}{1}$.
\end{remark}

\begin{example}[Wrapping back to the principal interval]
\label{ex:b1:functions:wrap}
Compute
\[
\arctan\Bigl(\tan\frac{3\pi}4\Bigr)
\qquad\text{and}\qquad
\arctan\Bigl(\tan\frac{17\pi}5\Bigr).
\]
The recipe: replace the
angle by the unique angle of $\intoo{-\frac\pi2}{\frac\pi2}$ with
the same tangent, i.e.\ subtract the right multiple of $\pi$
(the period of $\tan$). First: $\frac{3\pi}4 - \pi =
-\frac\pi4$, so the answer is $-\frac\pi4$. Second: $\frac{17\pi}
5 - 3\pi = \frac{2\pi}5 \in \intoo{-\frac\pi2}{\frac\pi2}$, so
the answer is $\frac{2\pi}5$. The computation is a Euclidean
division of the angle by $\pi$ in disguise --- and the analogous
recipes for $\arcsin$ (reflect into
$\intcc{-\frac\pi2}{\frac\pi2}$, period $2\pi$) and $\arccos$
(reflect into $\intcc0\pi$) drive \cref{exo:b1:functions:1} and
the piecewise answer of \cref{exo:b1:functions:10}.
\end{example}

\begin{example}[Adding arctangents safely]\label{ex:b1:functions:arctansum}
Let us prove
\[
\arctan\frac12 + \arctan\frac15 + \arctan\frac18 = \frac\pi4 .
\]
Two ingredients: the tangent addition formula, and --- the step
beginners forget --- a \emph{localization} of the sum. First,
$\tan(u + v) = \frac{\tan u + \tan v}{1 - \tan u\tan v}$ with
$u = \arctan\frac12$, $v = \arctan\frac15$ gives
\[
\tan(u + v) = \frac{\frac12 + \frac15}{1 - \frac1{10}}
= \frac{7/10}{9/10} = \frac79,
\qquad\text{then}\qquad
\tan\Bigl(u + v + \arctan\frac18\Bigr)
= \frac{\frac79 + \frac18}{1 - \frac7{72}}
= \frac{65/72}{65/72} = 1 .
\]
Second, each of the three angles lies in $\intoo0{\frac\pi4}$
(their arguments are less than $1$), so the sum lies in
$\intoo0{\frac{3\pi}4}$; the only angle there with tangent $1$ is
$\frac\pi4$. Without the localization the conclusion would be
``$\frac\pi4$ up to a multiple of $\pi$'' --- half a proof. The
same two-step discipline runs \cref{exo:b1:functions:8} and
\cref{exo:b1:functions:12}.
\end{example}

\begin{example}[An arctangent equation]\label{ex:b1:functions:arctaneq}
Solve $\arctan x + \arctan 2x = \dfrac\pi4$. First localize: the
left side has the sign of $x$ (both terms do), so any solution has
$x > 0$, and then the sum lies in $\intoo0\pi$. Take tangents
(injective on no interval of length $\pi$, but combined with the
localization this will suffice): the addition formula gives
\[
\tan\bigl(\arctan x + \arctan 2x\bigr)
= \frac{3x}{1 - 2x^2} = 1
\iff 2x^2 + 3x - 1 = 0
\iff x = \frac{-3 \pm \sqrt{17}}4 .
\]
The negative root is discarded by the localization; for the
positive candidate $x_0 = \frac{\sqrt{17} - 3}4 \approx 0.28$,
the sum lies in $\intoo0\pi$ with tangent $1$, and the only such
angle is $\frac\pi4$ (in $\intoo{\frac\pi2}\pi$ the tangent is
negative): $x_0$ is the unique solution. Note the shape of the
argument: taking tangents may create solutions, never lose them,
so one solves the polynomial equation \emph{then} filters by
localization --- the same forward-check discipline as squaring
an equation.
\end{example}

\section{Hyperbolic functions}

\begin{definition}[$\cosh$, $\sinh$, $\tanh$]\label{def:b1:functions:hyperbolic}
For $x \in \R$:
\[
\cosh x = \frac{\eu^x + \eu^{-x}}{2},
\qquad
\sinh x = \frac{\eu^x - \eu^{-x}}{2},
\qquad
\tanh x = \frac{\sinh x}{\cosh x}
\]
(\emph{hyperbolic cosine, sine, tangent}\index{hyperbolic
functions}). $\cosh$ is even, $\sinh$ and $\tanh$ are odd.
\end{definition}

\begin{proposition}[Basic properties]\label{prop:b1:functions:hyprules}
For all $x, y \in \R$:
\begin{enumerate}
  \item $\cosh^2 x - \sinh^2 x = 1$;
  \item $\cosh' = \sinh$, $\;\sinh' = \cosh$, $\;\tanh' = 1 - \tanh^2
        = \dfrac{1}{\cosh^2}$;
  \item $\cosh(x+y) = \cosh x\cosh y + \sinh x \sinh y$ and
        $\sinh(x+y) = \sinh x \cosh y + \cosh x \sinh y$;
  \item $\sinh$ is a strictly increasing bijection of $\R$ onto $\R$;
        $\cosh$ restricted to $\R_+$ is a strictly increasing
        bijection onto $\intco{1}{+\infty}$; $\tanh$ is a strictly
        increasing bijection of $\R$ onto $\intoo{-1}{1}$.
\end{enumerate}
\end{proposition}

\begin{proof}
(1) $\bigl(\frac{\eu^x + \eu^{-x}}{2}\bigr)^2 - \bigl(\frac{\eu^x -
\eu^{-x}}{2}\bigr)^2 = \frac{(\eu^{2x} + 2 + \eu^{-2x}) - (\eu^{2x} -
2 + \eu^{-2x})}{4} = 1$.

(2) Differentiate the defining formulas; for $\tanh$, the quotient
rule gives $\frac{\cosh^2 - \sinh^2}{\cosh^2}$, which is both
$1 - \tanh^2$ and $\frac{1}{\cosh^2}$ by (1).

(3) Expand the right-hand sides using the definitions. In detail
for the first formula:
\[
\cosh x\cosh y + \sinh x\sinh y
= \frac{(\eu^x + \eu^{-x})(\eu^y + \eu^{-y})
+ (\eu^x - \eu^{-x})(\eu^y - \eu^{-y})}{4} ;
\]
of the eight products, the four ``mixed'' ones
($\eu^{x-y}$, $\eu^{y-x}$) cancel in pairs, while
$\eu^{x+y}$ and $\eu^{-(x+y)}$ each appear twice:
the total is $\frac{2\eu^{x+y} + 2\eu^{-(x+y)}}{4} = \cosh(x+y)$.
The sine formula is identical with the mixed terms surviving
instead.

(4) $\sinh' = \cosh \geq 1 > 0$, so $\sinh$ is strictly increasing,
with limits $\pm\infty$ (dominant term $\pm\frac{\eu^{\abs
x}}{2}$); the bijection statement then follows from the intermediate
value theorem (used at High School level; systematic treatment in
\cref{ch:b1:continuity}). On $\R_+$, $\cosh' = \sinh \geq 0$
vanishing only at $0$: strictly increasing from $\cosh 0 = 1$ to
$+\infty$. And $\tanh' > 0$ with limits $\pm 1$ at $\pm\infty$: in
detail,
\[
\tanh x = \frac{\eu^x - \eu^{-x}}{\eu^x + \eu^{-x}}
= \frac{1 - \eu^{-2x}}{1 + \eu^{-2x}}
\xrightarrow[x \to +\infty]{} 1 ,
\]
after dividing numerator and denominator by $\eu^x$, and oddness
gives the limit $-1$ at $-\infty$; the strictly increasing
function $\tanh$ therefore maps $\R$ onto $\intoo{-1}1$.
\end{proof}

\begin{omfigure}
\begin{tikzpicture}
\begin{axis}[omaxis, width=6.4cm, height=6cm,
    xmin=-2.6, xmax=2.6, ymin=-3.6, ymax=3.9,
    xtick={-2,-1,1,2}, ytick={-3,-2,-1,1,2,3}]
  \addplot[omDef, very thick, domain=-2.05:2.05, samples=120]
    {cosh(x)} node[pos=0.1, right, font=\small] {$\cosh$};
  \addplot[omThm, very thick, domain=-2.05:2.05, samples=120]
    {sinh(x)} node[pos=0.93, above left, font=\small] {$\sinh$};
  \addplot[dashed, gray, domain=0:2.5] {exp(x)/2};
\end{axis}
\begin{scope}[xshift=7.4cm]
\begin{axis}[omaxis, width=6.4cm, height=6cm,
    xmin=-3.2, xmax=3.2, ymin=-1.5, ymax=1.5,
    xtick={-2,-1,1,2}, ytick={-1,1}]
  \addplot[omMeth, very thick, domain=-3:3, samples=150]
    {tanh(x)} node[pos=0.9, above left, font=\small] {$\tanh$};
  \addplot[dashed, gray, domain=-3:3] {1};
  \addplot[dashed, gray, domain=-3:3] {-1};
\end{axis}
\end{scope}
\end{tikzpicture}

{\small Left: $\cosh$ and $\sinh$, asymptotically glued to
$\frac{\eu^x}{2}$ (dashed). Right: $\tanh$, increasing from $-1$ to
$1$.}
\end{omfigure}

\begin{remark}[Why ``hyperbolic''?]
The point $(\cosh t, \sinh t)$ runs along the branch $x > 0$ of the
hyperbola $x^2 - y^2 = 1$ (by
\cref{prop:b1:functions:hyprules}~(1)), exactly as $(\cos t, \sin
t)$ runs along the circle $x^2 + y^2 = 1$. Every circular identity has
a hyperbolic sibling, with sign changes governed by
$\sinh^2 \leftrightarrow -\sin^2$.
\end{remark}

\begin{example}[The addition formula for $\tanh$]\label{ex:b1:functions:tanhadd}
Dividing the two addition formulas of
\cref{prop:b1:functions:hyprules}~(3) by $\cosh x\cosh y$:
\[
\tanh(x + y)
= \frac{\sinh x\cosh y + \cosh x\sinh y}
{\cosh x\cosh y + \sinh x\sinh y}
= \frac{\tanh x + \tanh y}{1 + \tanh x\,\tanh y} ,
\]
the hyperbolic sibling of the tangent addition formula --- with a
$+$ where trigonometry has a $-$. A dividend: since $\abs{\tanh}
< 1$, the right-hand side is a ``velocity addition'' rule that
never leaves $\intoo{-1}1$: if $u, v \in \intoo{-1}1$ then
$\frac{u + v}{1 + uv} \in \intoo{-1}1$ as well (write $u = \tanh
a$, $v = \tanh b$, possible by bijectivity, and read the formula
backwards). Checking that algebraically, without hyperbolic
functions, is a slightly painful exercise; parametrizing by
$\tanh$ makes it one line --- the same strategy that circular
functions provide for the unit circle.
\end{example}

\begin{proposition}[Inverse hyperbolic functions]\label{prop:b1:functions:invhyp}
The inverses granted by \cref{prop:b1:functions:hyprules}~(4) have
closed forms:
\[
\operatorname{arsinh} x = \ln\bigl(x + \sqrt{x^2 + 1}\bigr)
\ (x \in \R),
\qquad
\operatorname{arcosh} x = \ln\bigl(x + \sqrt{x^2 - 1}\bigr)
\ (x \geq 1),
\]
\[
\operatorname{artanh} x = \frac 12 \ln\frac{1 + x}{1 - x}
\ (\abs x < 1),
\]
with derivatives $\frac{1}{\sqrt{x^2+1}}$, $\frac{1}{\sqrt{x^2-1}}$
($x > 1$) and $\frac{1}{1 - x^2}$ respectively.
\end{proposition}

\begin{proof}
For $\operatorname{arsinh}$: solve $x = \sinh y = \frac{\eu^y -
\eu^{-y}}{2}$. Setting $u = \eu^y > 0$: $u^2 - 2xu - 1 = 0$, so $u = x
+ \sqrt{x^2 + 1}$ (the root $x - \sqrt{x^2+1}$ is negative), and $y =
\ln(x + \sqrt{x^2+1})$. The other two are identical computations
($u^2 - 2xu + 1 = 0$ for $\operatorname{arcosh}$, keeping the root
$\geq 1$; a two-line rearrangement for $\operatorname{artanh}$).
Derivatives: differentiate the logarithmic expressions, e.g.
\[
\bigl(\ln(x + \sqrt{x^2+1})\bigr)'
= \frac{1 + \frac{x}{\sqrt{x^2+1}}}{x + \sqrt{x^2+1}}
= \frac{1}{\sqrt{x^2 + 1}} . \qedhere
\]
\end{proof}

\begin{example}[Closed forms at work]\label{ex:b1:functions:arcoshtwo}
The solution of $\cosh t = 2$ with $t \geq 0$ is, by the closed
form, $t = \operatorname{arcosh} 2 = \ln(2 + \sqrt3) \approx
1.317$; the other solution is $-t$, by evenness --- and indeed
$\ln(2 - \sqrt3) = \ln\frac1{2 + \sqrt3} = -\ln(2 + \sqrt3)$: the
two roots $u = \eu^t$ of the quadratic $u^2 - 4u + 1 = 0$ are
reciprocals, as their product $1$ (Vieta) demands. This tiny
computation displays the general pattern: hyperbolic equations
convert to quadratics in $\eu^t$, and the symmetry $t \mapsto -t$
appears as the symmetry $u \mapsto 1/u$ of the quadratic ---
worth remembering when solving \cref{exo:b1:functions:7}.
\end{example}

\begin{method}[Choosing the right primitive form]\label{met:b1:functions:primitive}
The three derivative patterns to memorize for integration
(\cref{ch:b1:integration}):
\[
\int \frac{\dd x}{1 + x^2} = \arctan x + C,
\qquad
\int \frac{\dd x}{\sqrt{1 - x^2}} = \arcsin x + C,
\qquad
\int \frac{\dd x}{\sqrt{x^2 + 1}} = \operatorname{arsinh} x + C,
\]
and for a general quadratic, reduce to these by completing the square
and rescaling.
\end{method}

\begin{remark}[Where this chapter is used]\label{rem:b1:functions:whereused}
This chapter is the working vocabulary of all the analysis to come.
The growth scale of \cref{prop:b1:functions:powerrules} decides
convergence questions throughout
\cref{ch:b1:seq,ch:b1:series}; the derivative formulas of
\cref{prop:b1:functions:arcderiv} and
\cref{prop:b1:functions:invhyp} are the primitives most often
needed in \cref{ch:b1:integration}, via
\cref{met:b1:functions:primitive}; hyperbolic functions
parametrize the solutions of the equation $y'' = y$ in
\cref{ch:b1:diffeq} exactly as circular functions parametrize
those of $y'' = -y$. The characterization of $\exp$ by $f' = f$,
$f(0) = 1$ (\cref{prop:b1:functions:expln}) is the one-dimensional
seed of the theory of linear differential equations, and the
catenary curve $y = \cosh x$ returns among the plane curves of
\cref{ch:b1:curves}.
\end{remark}

\begin{remark}[Interlude: the inverse-function program]\label{rem:b1:functions:interlude}
This chapter ran one program four times: restrict a function
until it becomes a strictly monotonic bijection, name the inverse,
transport every formula through it. What was \emph{used} at each
step --- that a continuous strictly monotonic function on an
interval is a bijection onto an interval, and that its inverse is
continuous, then differentiable away from critical points --- was
borrowed on High School credit. The debt is repaid inside this
volume: \cref{ch:b1:continuity} proves the bijection statement
(the monotone inverse theorem, resting on the intermediate value
theorem), and \cref{ch:b1:derivative} proves the derivative rule
$(f^{-1})' = 1/(f' \circ f^{-1})$ that silently produced every
formula of \cref{prop:b1:functions:arcderiv} and
\cref{prop:b1:functions:invhyp}. Reading those chapters with
$\arcsin$ and $\operatorname{arsinh}$ in mind --- as the worked
examples the theory was built to justify --- is the intended way
around the apparent circularity.
\end{remark}

\section{Exercises}

\begin{exercise}[$\star$]\label{exo:b1:functions:1}
Compute without calculator: $\arcsin\frac{\sqrt 3}{2}$;
$\;\arccos\bigl(-\frac 12\bigr)$; $\;\arctan(-1)$;
$\;\arcsin\bigl(\sin\frac{5\pi}{6}\bigr)$;
$\;\arccos\bigl(\cos\frac{7\pi}{4}\bigr)$.
\end{exercise}

\begin{exercise}[$\star$]\label{exo:b1:functions:2}
Give the domain of definition of $f(x) = \arcsin(2x - 1)$ and compute
$f'$ where defined. Same questions for $g(x) = \arctan\sqrt{x}$.
\end{exercise}

\begin{exercise}[$\star$]\label{exo:b1:functions:3}
Prove that for all $x \in \R$: $\cosh x + \sinh x = \eu^x$ and
$(\cosh x + \sinh x)^n = \cosh nx + \sinh nx$ for every $n \in \Z$
(a hyperbolic de Moivre formula).
\end{exercise}

\begin{exercise}[$\star$]\label{exo:b1:functions:4}
Simplify $\cos(2\arcsin x)$ and $\sin(2\arctan x)$ into algebraic
expressions of $x$.
\end{exercise}

\begin{exercise}[$\star$]\label{exo:b1:functions:5}
Order, for large $x$, the functions
$x^{100}$, $\;\eu^{\sqrt x}$, $\;(\ln x)^{1000}$, $\;\eu^{x/100}$,
$\;x^{\ln x}$, from slowest to fastest growth, with justifications
based on \cref{prop:b1:functions:powerrules}.
\end{exercise}

\begin{exercise}[$\star\star$]\label{exo:b1:functions:6}
Study the function $f(x) = \arctan\dfrac{2x}{1 - x^2}$ on its domain:
compute $f'$, compare with $(2\arctan x)'$, and express $f(x)$ in
terms of $\arctan x$ on each of the three intervals of the domain.
\end{exercise}

\begin{exercise}[$\star\star$]\label{exo:b1:functions:7}
Solve in $\R$: $\;\cosh x = 2$; then $5\cosh x - 4 \sinh x = 3$
(express the solutions with logarithms). \emph{Hint for the second:
write everything with $u = \eu^x$.}
\end{exercise}

\begin{exercise}[$\star\star$]\label{exo:b1:functions:8}
Prove the identity $\arctan\frac{1}{2} + \arctan\frac{1}{3} =
\frac{\pi}{4}$, then Machin's formula\index{Machin's formula}
\[
4\arctan\frac 15 - \arctan\frac{1}{239} = \frac{\pi}{4}.
\]
\emph{Hint: compute the tangent of both sides using the addition
formula, and control in which interval each side lies.}
\end{exercise}

\begin{exercise}[$\star\star$]\label{exo:b1:functions:9}
Prove that for all $x \geq 0$: $\;x - \dfrac{x^3}{6} \leq \sin x \leq
x$ \emph{(study the successive derivatives of the differences)}, and
deduce $\lim_{x \to 0^+} \frac{x - \sin x}{x^3} = \frac 16$ is
plausible --- the limit itself is established in
\cref{ch:b1:taylor}.
\end{exercise}

\begin{exercise}[$\star\star\star$]\label{exo:b1:functions:10}
For $x \in \intcc{-1}{1}$, set
$h(x) = \arcsin\bigl(2x\sqrt{1 - x^2}\,\bigr) - 2\arcsin x$. One might
expect $h = 0$ from the identity $\sin 2y = 2\sin y\cos y$ --- but
$h$ is \emph{not} identically zero. Compute $h'$ on the open intervals
where it exists, evaluate $h$ at well-chosen points, and give the full
piecewise-constant description of $h$ on $\intcc{-1}{1}$.
\end{exercise}

\begin{exercise}[$\star\star\star$]\label{exo:b1:functions:11}
(Gudermannian) Let $g(x) = \arctan(\sinh x)$ for $x \in \R$. Prove
that $g$ is an odd, strictly increasing bijection from $\R$ onto
$\intoo{-\frac\pi2}{\frac\pi2}$, that $g'(x) = \frac{1}{\cosh x}$,
and that $\tan g(x) = \sinh x$, $\;\sin g(x) = \tanh x$,
$\;\cos g(x) = \frac{1}{\cosh x}$: the function $g$ links circular
and hyperbolic trigonometry without complex numbers.
\end{exercise}

\begin{exercise}[$\star\star$]\label{exo:b1:functions:12}
Prove that
\[
\arctan 1 + \arctan 2 + \arctan 3 = \pi .
\]
\emph{Hint: compute $\arctan 2 + \arctan 3$ first, localizing the
sum as in \cref{ex:b1:functions:arctansum}; beware that the
addition formula's denominator is negative here.}
\end{exercise}

\section{Problem: The equation $x^y = y^x$}

\begin{problem}\label{pb:b1:functions:1}
Which pairs of positive numbers satisfy $x^y = y^x$? Everyone knows
one accidental-looking example, $2^4 = 4^2 = 16$; this problem shows
that nothing about it is accidental. The whole equation is governed
by the variations of the single function
\[
f(t) = \frac{\ln t}{t} \qquad (t > 0),
\]
whose study yields: the complete solution set (a diagonal plus one
curved branch through $(\eu, \eu)$), a rational parametrization of
the branch, the fact that $(2, 4)$ is its \emph{only} integer point
and the classification of all its rational points --- plus, as
dividends, the comparison of $\eu^\pi$ with $\pi^\eu$ and the
monotone convergence of $\bigl(1 + \frac1n\bigr)^n$ to $\eu$.
\index{power function}

\medskip
\noindent\textbf{Part I --- The function $f(t) = \ln t / t$.}
\begin{enumerate}
  \item Justify that $f$ is differentiable on $\intoo0{+\infty}$,
        compute $f'$, and draw up the variation table: $f$
        increases strictly on $\intoc0\eu$, decreases strictly on
        $\intco\eu{+\infty}$, with maximum $f(\eu) = \frac1\eu$.
  \item Determine the limits of $f$ at $0^+$ and $+\infty$
        (\cref{prop:b1:functions:powerrules}), the sign of $f$
        (negative on $\intoo01$, zero at $1$, positive beyond),
        and sketch the graph.
  \item Check by direct computation that $f(2) = f(4)$. (Keep this
        equality in mind: the entire problem grows out of it.)
  \item Let $c \in \R$. Discuss, according to the value of $c$, the
        number of solutions of $f(t) = c$: exactly one for $c \leq
        0$; exactly two (one in $\intoo1\eu$, one in
        $\intoo\eu{+\infty}$) for $0 < c < \frac1\eu$; exactly one
        for $c = \frac1\eu$; none for $c > \frac1\eu$.
  \item Deduce that $\eu^x > x^\eu$ for every $x > 0$ with $x \neq
        \eu$, and in particular settle which of $\eu^\pi$ and
        $\pi^\eu$ is larger.
\end{enumerate}

\medskip
\noindent\textbf{Part II --- The equation and its curve.} In this
part $x, y > 0$.
\begin{enumerate}[resume]
  \item Show that $x^y = y^x \iff f(x) = f(y)$.
  \item Deduce the structure of the solution set: all diagonal
        pairs $(x, x)$; and the \emph{nontrivial} pairs ($x \neq
        y$), which satisfy: both coordinates are $> 1$, and one
        lies in $\intoo1\eu$ while the other lies in
        $\intoo\eu{+\infty}$.
  \item Parametrize the nontrivial pairs: writing $y = tx$ with $t
        > 0$, $t \neq 1$, show that $x^y = y^x$ forces
        \[
        x(t) = t^{\frac1{t-1}},
        \qquad
        y(t) = t\,x(t) = t^{\frac{t}{t-1}},
        \]
        and that conversely every such pair is a solution.
  \item Verify that $t = 2$ gives $(2, 4)$, and prove the symmetry
        $x(1/t) = y(t)$, $y(1/t) = x(t)$: inverting the parameter
        swaps the two coordinates.
  \item Determine the limits of $x(t)$ and $y(t)$ as $t \to 1$
        (both tend to $\eu$), as $t \to +\infty$ ($x \to 1$, $y
        \to +\infty$) and as $t \to 0^+$ ($x \to +\infty$, $y \to
        1$). Describe the resulting branch: a curve asymptotic to
        the lines $x = 1$ and $y = 1$, crossing the diagonal at
        $(\eu, \eu)$.
  \item Show that $t \mapsto x(t)$ is strictly decreasing on
        $\intoo1{+\infty}$. \emph{(Study $h(t) = 1 - \frac1t - \ln
        t$, whose sign controls the derivative of $\ln x(t) =
        \frac{\ln t}{t - 1}$.)}
  \item Conclude that the nontrivial solutions define a strictly
        decreasing bijection $\varphi \colon \intoo1\eu \to
        \intoo\eu{+\infty}$ with $x^{\varphi(x)} = \varphi(x)^x$,
        and that $\varphi$ is an involution of the branch:
        $\varphi(\varphi(x)) = x$ wherever both sides are defined.
\end{enumerate}

\medskip
\noindent\textbf{Part III --- Integer and rational points.}
\begin{enumerate}[resume]
  \item Prove that $(2, 4)$ and $(4, 2)$ are the only nontrivial
        \emph{integer} solutions of $x^y = y^x$.
  \item For $n \in \N^*$, apply the parametrization with $t = 1 +
        \frac1n$ and show that
        \[
        x_n = \Bigl(1 + \frac1n\Bigr)^{n},
        \qquad
        y_n = \Bigl(1 + \frac1n\Bigr)^{n+1}
        \]
        is a nontrivial \emph{rational} solution for every $n$,
        with $(x_1, y_1) = (2, 4)$.
  \item Conversely, let $(x, y)$ be a nontrivial rational solution
        with $y > x$, and write $t = y/x = 1 + \frac rs$ in lowest
        terms ($r, s \in \N^*$). Using $x = t^{s/r}$ and admitting
        the uniqueness of prime factorization (familiar from
        school; proved in \cref{ch:b1:arith}), show that $x$
        rational forces both $s$ and $s + r$ to be $r$-th powers
        of integers.
  \item Show with the binomial theorem that $a^r - b^r = r$ has no
        integer solutions $a > b \geq 1$ when $r \geq 2$, and
        conclude: the rational solutions of $x^y = y^x$ are
        exactly the pairs $(x_n, y_n)$ of question 14 (and their
        swaps).
  \item Verify the case $n = 2$ numerically: compute $f(9/4)$ and
        $f(27/8)$ to four decimal places and check they agree.
\end{enumerate}

\medskip
\noindent\textbf{Part IV --- Dividends.}
\begin{enumerate}[resume]
  \item Use questions 7 and 11 to prove, with no further
        computation, that the sequence $x_n = \bigl(1 +
        \frac1n\bigr)^n$ is strictly increasing with $x_n < \eu$,
        that $y_n = \bigl(1 + \frac1n\bigr)^{n+1}$ is strictly
        decreasing with $y_n > \eu$, and that both converge to
        $\eu$. (The parameter $t_n = 1 + \frac1n$ decreases to
        $1$.)
  \item Establish the general comparison rule for $1 < a < b$:
        if $\eu \leq a$ then $a^b > b^a$; if $b \leq \eu$ then
        $a^b < b^a$; and show by the two examples $(2, 3)$ and
        $(2, 5)$ that in the mixed case $a < \eu < b$ both
        outcomes really occur.
  \item Suppose the involution $\varphi$ of question 12 is
        differentiable at $\eu$ (it is). Differentiate the
        identity $\varphi(\varphi(x)) = x$ at $x = \eu$ and deduce
        $\varphi'(\eu) = -1$: the branch crosses the diagonal at
        right angles to it.
  \item Find the unique solution pair with $y = 3x$, in closed
        form, and check it numerically to four decimals via $f$.
  \item Among the numbers $\sqrt2$, $\sqrt[3]3$, $\sqrt[4]4$,
        $\sqrt[5]5$, determine the largest and identify the two
        that are equal. \emph{(Compare $n^{1/n} =
        \eu^{f(n)}$.)}
\end{enumerate}

\medskip
\noindent\textbf{Part V --- Synthesis.}
\begin{enumerate}[resume]
  \item Describe the complete solution set of $x^y = y^x$ in the
        quarter-plane $x, y > 0$ --- diagonal plus branch, their
        intersection $(\eu, \eu)$, the asymptotes, the integer
        point $(2, 4)$, the rational points accumulating at
        $(\eu, \eu)$ --- in a form you could sketch from memory.
  \item Where exactly did the problem use: (i) the growth
        comparisons of \cref{prop:b1:functions:powerrules}; (ii)
        the intermediate value theorem (through the bijection
        statements); (iii) the admitted uniqueness of prime
        factorization? One sentence each.
  \item Moral, in a short paragraph: a single variation table
        solved an equation in two unknowns, classified its
        rational points, and proved the monotone convergence of
        $\bigl(1 + \frac1n\bigr)^n$ --- comment on this economy,
        and name where each thread is industrialized later in the
        volume (\cref{ch:b1:seq} for the sequence,
        \cref{ch:b1:derivative} for variation tables,
        \cref{ch:b1:taylor} for the precision the table lacks).
\end{enumerate}
\end{problem}

\chapter{Linear Differential Equations}\label{ch:b1:diffeq}

First met in the High School volume, differential equations are treated
here with full proofs and in greater generality: first-order linear
equations with variable coefficients (solved completely by the
variation-of-constants method) and second-order linear equations with
constant coefficients, the model for oscillations. Both cases display
the same structure: \emph{general solution $=$ one particular solution
$+$ general solution of the homogeneous equation}.

\section{First-order linear equations}

\begin{definition}\label{def:b1:diffeq:linear1}
Let $I$ be an interval and $a, b \colon I \to \R$ (or $\C$) be
continuous. The equation
\[
(E)\colon\quad y' + a(x)\,y = b(x),
\]
in the unknown differentiable function $y \colon I \to \R$ (or $\C$),
is a \emph{first-order linear differential
equation}\index{differential equation!first-order linear}. The
equation $(H)\colon y' + a(x) y = 0$ is its \emph{homogeneous}
equation.
\end{definition}

\begin{theorem}[Solving the homogeneous equation]\label{thm:b1:diffeq:homogeneous1}
Let $A$ be a primitive of $a$ on $I$ (it exists:
\cref{ch:b1:integration}). The solutions of $(H)$ on $I$ are exactly
the functions
\[
y(x) = \lambda\, \eu^{-A(x)}, \qquad \lambda \in \R \text{ (or } \C).
\]
\end{theorem}

\begin{proof}
These functions are solutions: $y' = -\lambda A' \eu^{-A} = -a y$.
Conversely, let $y$ solve $(H)$ and set $z(x) = y(x)\, \eu^{A(x)}$.
Then
\[
z' = y' \eu^{A} + y\, a\, \eu^{A} = (y' + ay)\, \eu^{A} = 0 ,
\]
so $z$ is constant on the interval $I$, say $z = \lambda$: $y =
\lambda \eu^{-A}$. (Note the logic: no solution is lost, because
\emph{every} solution has been written in the announced form.)
\end{proof}

\begin{example}[A homogeneous equation with variable coefficient]\label{ex:b1:diffeq:varcoeff}
Solve $y' + (\cos x)\,y = 0$ on $\R$. A primitive of $a(x) =
\cos x$ is $A(x) = \sin x$, so the solutions are
\[
y(x) = \lambda\,\eu^{-\sin x}, \qquad \lambda \in \R .
\]
Two readings. Every solution is periodic (period $2\pi$) and never
vanishes unless $\lambda = 0$ --- the sign of $\lambda$ is the
sign of $y$ forever, since an exponential cannot cross zero. And
the solution through $y(0) = y_0$ is $y_0\eu^{-\sin x}$: exactly
one curve of the family through each initial point, the
one-dimensional picture of \cref{thm:b1:diffeq:voc}~(2).
\end{example}

\begin{theorem}[Variation of constants; Cauchy problem]\label{thm:b1:diffeq:voc}
With the notation above:
\begin{enumerate}
  \item The solutions of $(E)$ on $I$ are exactly
        \[
        y(x) = \Bigl(\lambda + \int_{x_0}^{x} b(t)\, \eu^{A(t)} \dd t
        \Bigr)\, \eu^{-A(x)},
        \qquad \lambda \in \R,
        \]
        where $x_0 \in I$ is fixed. Equivalently: general solution of
        $(H)$ plus one particular solution of $(E)$.
  \item For every $x_0 \in I$ and $y_0$, the \emph{Cauchy
        problem}\index{Cauchy problem} ``$(E)$ and $y(x_0) = y_0$''
        has exactly one solution on $I$.
\end{enumerate}
\end{theorem}

\begin{proof}
(1) Following the method named \emph{variation of
constants}\index{variation of constants}, look for solutions of the
form $y = \mu(x)\, \eu^{-A(x)}$ with $\mu$ differentiable --- no
generality is lost, since every function on $I$ can be so written
($\mu = y\,\eu^{A}$). Substituting,
\[
y' + ay = \mu' \eu^{-A} - \mu a \eu^{-A} + a\mu\eu^{-A}
= \mu'\, \eu^{-A},
\]
so $y$ solves $(E)$ if and only if $\mu'(x) = b(x)\,\eu^{A(x)}$, if
and only if $\mu(x) = \lambda + \int_{x_0}^x b(t)\eu^{A(t)}\dd t$ for
some constant $\lambda$ (two primitives of the same continuous
function on an interval differ by a constant).

(2) In the formula, $y(x_0) = \lambda\,\eu^{-A(x_0)}$: the condition
$y(x_0) = y_0$ determines $\lambda = y_0 \eu^{A(x_0)}$ uniquely.
\end{proof}

\begin{example}\label{ex:b1:diffeq:order1}
Solve $y' + \dfrac{y}{x} = x^2$ on $I = \intoo{0}{+\infty}$. Here
$a(x) = \frac 1x$, $A(x) = \ln x$, $\eu^{-A(x)} = \frac 1x$.
Homogeneous solutions: $\frac{\lambda}{x}$. \omterm{thm:b1:diffeq:voc}{Variation of constants}:
$\mu'(x) = x^2 \cdot x = x^3$, so $\mu = \frac{x^4}{4} + \lambda$, and
\[
y(x) = \frac{x^3}{4} + \frac{\lambda}{x}, \qquad \lambda \in \R .
\]
With the initial condition $y(1) = 0$: $\lambda = -\frac14$.
\emph{Check:} $y' + \frac yx = \frac{3x^2}{4} - \frac{\lambda}{x^2} +
\frac{x^2}{4} + \frac{\lambda}{x^2} = x^2$.
\end{example}

\begin{example}[Guessing beats integrating]\label{ex:b1:diffeq:guess}
Solve $y' + 2x\,y = x$ on $\R$. \omterm{thm:b1:diffeq:voc}{Variation of constants} works ($A =
x^2$, $\mu' = x\,\eu^{x^2}$, $\mu = \frac12\eu^{x^2} + \lambda$),
but observing that the \emph{constant} $y_p = \frac12$ solves the
equation ($0 + 2x\cdot\frac12 = x$) is faster. With the
homogeneous solutions $\lambda\,\eu^{-x^2}$:
\[
y(x) = \frac12 + \lambda\,\eu^{-x^2}, \qquad \lambda \in \R .
\]
Every solution converges to $\frac12$ extremely fast as $x \to
\pm\infty$: the constant particular solution is an
\emph{equilibrium} that all solutions join. The insight: before
launching the general method, spend ten seconds looking for an
obvious particular solution (constant, monomial, multiple of the
right-hand side); the structure theorem then finishes the job.
\end{example}

\begin{remark}[Intervals matter]
The theorem lives on an \emph{interval} where $a$ and $b$ are
continuous. For $y' + \frac yx = 0$ on $\R^*$, the solutions are
$\frac{\lambda}{x}$ on $\intoo{0}{+\infty}$ and
$\frac{\mu}{x}$ on $\intoo{-\infty}{0}$ with \emph{independent}
constants: there is no reason for a single formula to glue across the
singularity at $0$.
\end{remark}

\begin{example}[A complex right-hand side, two real answers]\label{ex:b1:diffeq:complexrhs}
Solve $y' - y = \cos x$ and $y' - y = \sin x$ in one stroke. Work
in $\C$ with the right-hand side $\eu^{\iu x}$: trying $y_p =
c\,\eu^{\iu x}$ gives $c(\iu - 1)\eu^{\iu x} = \eu^{\iu x}$, so
\[
c = \frac1{\iu - 1} = \frac{-1 - \iu}2,
\qquad
y_p = -\frac{(1 + \iu)(\cos x + \iu\sin x)}2
= \frac{\sin x - \cos x}2 + \iu\,\frac{-\sin x - \cos x}2 .
\]
Since the equation has real coefficients, real and imaginary
parts split: $\frac{\sin x - \cos x}2$ solves $y' - y = \cos x$,
and $-\frac{\sin x + \cos x}2$ solves $y' - y = \sin x$ (check
the first: derivative $\frac{\cos x + \sin x}2$, minus the
function, gives $\cos x$). One complex line replaced two runs of
\omterm{thm:b1:diffeq:voc}{variation of constants} --- the same economy that
\cref{met:b1:diffeq:particular} systematizes for second order,
and a recurring dividend of \cref{ch:b1:complex}.
\end{example}

\section{Second-order linear equations with constant coefficients}

\begin{definition}\label{def:b1:diffeq:linear2}
Let $a, b \in \R$ and $f \colon I \to \R$ continuous. The equation
\[
(E)\colon\quad y'' + a\,y' + b\,y = f(x)
\]
is a \emph{second-order linear equation with constant
coefficients}\index{differential equation!second-order linear};
$(H)\colon y'' + ay' + by = 0$ is its homogeneous equation, and
$\chi(r) = r^2 + ar + b$ its \emph{characteristic
polynomial}\index{characteristic polynomial}.
\end{definition}

\begin{theorem}[Homogeneous solutions]\label{thm:b1:diffeq:homogeneous2}
Let $\Delta = a^2 - 4b$ be the discriminant of $\chi$. The real
solutions of $(H)$ on $\R$ are:
\begin{enumerate}
  \item if $\Delta > 0$, with $r_1 \neq r_2$ the two real roots:
        $\;y = \lambda\, \eu^{r_1 x} + \mu\, \eu^{r_2 x}$;
  \item if $\Delta = 0$, with $r_0$ the double root:
        $\;y = (\lambda + \mu x)\, \eu^{r_0 x}$;
  \item if $\Delta < 0$, with roots $\alpha \pm \iu\omega$
        ($\omega > 0$):
        $\;y = \eu^{\alpha x} (\lambda \cos\omega x + \mu
        \sin\omega x)$;
\end{enumerate}
in each case with $(\lambda, \mu)$ running over $\R^2$.
\end{theorem}

\begin{proof}
First observe that for $r \in \C$, $x \mapsto \eu^{rx}$ solves $(H)$
if and only if $\chi(r) = 0$ (substitute: $(r^2 + ar + b)\eu^{rx} =
0$). This is why exponentials are the natural first guess:
differentiation acts on $\eu^{rx}$ as multiplication by the number
$r$, so the differential equation becomes the numerical equation
$\chi(r) = 0$ --- the whole analytic problem is compressed into
finding the roots of one quadratic.

The key step is a change of unknown that reduces the order. Let $r$ be
a (possibly complex) root of $\chi$ and write $y = z\, \eu^{rx}$,
which loses no generality. Then
\[
y'' + ay' + by
= \bigl(z'' + (2r + a) z' + \chi(r) z\bigr)\eu^{rx}
= \bigl(z'' + (2r + a) z'\bigr)\eu^{rx},
\]
so $(H)$ becomes the \emph{first-order} equation $u' + (2r + a) u = 0$
for $u = z'$.

\emph{Case $\Delta \neq 0$:} choose $r = r_1$; then $2r_1 + a = r_1 -
r_2$ (since $r_1 + r_2 = -a$). By
\cref{thm:b1:diffeq:homogeneous1}, $z' = c\,\eu^{(r_2 - r_1)x}$ for
some constant $c$; integrating on $\R$, $z = \mu\, \eu^{(r_2 - r_1)x}
+ \lambda$ with $\mu = \frac{c}{r_2 - r_1}$, and therefore $y = z\,
\eu^{r_1 x} = \lambda\,\eu^{r_1x} + \mu\,\eu^{r_2x}$.
When $\Delta < 0$, the roots are $\alpha \pm \iu\omega$ and the
complex solutions are $y = c_1\eu^{(\alpha+\iu\omega)x} +
c_2\eu^{(\alpha-\iu\omega)x}$ with $c_1, c_2 \in \C$. Which of
these are real-valued? Since $\conj{\eu^{(\alpha+\iu\omega)x}} =
\eu^{(\alpha-\iu\omega)x}$, the conjugate of $y$ is $\conj{c_2}\,
\eu^{(\alpha+\iu\omega)x} + \conj{c_1}\,\eu^{(\alpha-\iu\omega)x}$,
and $y = \conj y$ for all $x$ forces $c_2 = \conj{c_1}$ (the two
exponentials are linearly independent: evaluate at two points, or
compare at $x=0$ after dividing by $\eu^{\alpha x}$). Writing $c_1
= \frac{\lambda - \iu\mu}2$ with $\lambda, \mu$ real:
\[
y = 2\,\Re\Bigl(c_1\,\eu^{\alpha x}(\cos\omega x +
\iu\sin\omega x)\Bigr)
= \eu^{\alpha x}\bigl(\lambda\cos\omega x + \mu\sin\omega
x\bigr),
\]
and conversely every such function is a solution (real part of a
complex solution of a real equation): the real solution space is
as announced.

\emph{Case $\Delta = 0$:} $r = r_0$, $2r_0 + a = 0$, so $z'' = 0$:
$z = \lambda + \mu x$ and $y = (\lambda + \mu x)\eu^{r_0 x}$.
\end{proof}

\begin{example}[A Cauchy problem, start to finish]\label{ex:b1:diffeq:cauchy2}
Solve $y'' - 3y' + 2y = 0$ with $y(0) = 0$, $y'(0) = 1$. The
\omterm{def:b1:diffeq:linear2}{characteristic polynomial} $r^2 - 3r + 2 = (r - 1)(r - 2)$ has the
real roots $1$ and $2$: general solution $y = \lambda\eu^{x} +
\mu\eu^{2x}$. The two conditions give the linear system
\[
\lambda + \mu = 0, \qquad \lambda + 2\mu = 1 ,
\]
so $\mu = 1$, $\lambda = -1$:
\[
y(x) = \eu^{2x} - \eu^{x} .
\]
Check: $y(0) = 0$; $y' = 2\eu^{2x} - \eu^x$ has $y'(0) = 1$; and
$y'' - 3y' + 2y = (4 - 6 + 2)\eu^{2x} + (-1 + 3 - 2)\eu^x = 0$.
Note the shape of the answer: near $-\infty$ the slow mode
$-\eu^x$ dominates; near $+\infty$ the fast mode $\eu^{2x}$ does.
Reading solutions as superpositions of modes with different decay
or growth rates is the profitable habit --- it is how the
transient/steady-state split of the weekend problem is organized.
\end{example}

\begin{omfigure}
\begin{tikzpicture}[xscale=1.35, yscale=1.55]
  \draw[->, thin] (0,0) -- (6.4,0) node[below] {$x$};
  \draw[->, thin] (0,-0.85) -- (0,1.25) node[left] {$y$};
  \draw[omDef, thick, domain=0:6.2, samples=400, smooth]
    plot (\x, {exp(-0.35*\x)*cos(4*\x r)});
  \draw[omThm, thick, domain=0:6.2, samples=120, smooth]
    plot (\x, {(1+1.8*\x)*exp(-1.1*\x)});
  \draw[omProp, thick, domain=0:6.2, samples=120, smooth]
    plot (\x, {1.15*exp(-0.45*\x) - 0.15*exp(-2.5*\x)});
  \node[omDef, font=\small] at (2.1,-0.62)
    {$\Delta < 0$: oscillates};
  \node[omThm, font=\small] at (0.95,1.1) {$\Delta = 0$};
  \node[omProp, font=\small] at (4.9,0.38)
    {$\Delta > 0$: no oscillation};
\end{tikzpicture}

{\small The three regimes of $y'' + ay' + by = 0$ with decaying
solutions: damped oscillation (complex roots), critical return
(double root), overdamped decay (two real roots). Which regime
occurs is read off the sign of $\Delta = a^2 - 4b$ alone ---
before solving anything.}
\end{omfigure}

\begin{theorem}[Structure and Cauchy problem]\label{thm:b1:diffeq:structure2}
\begin{enumerate}
  \item If $y_p$ is one particular solution of $(E)$, the solutions of
        $(E)$ are exactly $y_p + y_h$, $y_h$ running over the
        solutions of $(H)$.
  \item (Superposition) If $y_1$ solves $y'' + ay' + by = f_1$ and
        $y_2$ solves $y'' + ay' + by = f_2$, then $y_1 + y_2$ solves
        the equation with right-hand side $f_1 + f_2$.
  \item For all $x_0 \in I$ and $(y_0, y_0')$, the \omterm{thm:b1:diffeq:voc}{Cauchy problem}
        ``$(E)$, $y(x_0) = y_0$, $y'(x_0) = y_0'$'' has exactly one
        solution on $I$. \emph{(Existence granted a particular
        solution; uniqueness in full.)}
\end{enumerate}
\end{theorem}

\begin{proof}
(1) $y$ solves $(E)$ iff $y - y_p$ solves $(H)$, by linearity of $y
\mapsto y'' + ay' + by$. (2) is the same linearity.

(3) By (1) it suffices to prove that the constants $(\lambda,
\mu)$ can always be adjusted, uniquely, to any data $(y_0, y_0')$.
Translating the variable, assume $x_0 = 0$. In case (1) of
\cref{thm:b1:diffeq:homogeneous2}, $y = \lambda\eu^{r_1x} +
\mu\eu^{r_2x}$ gives
\[
y(0) = \lambda + \mu, \qquad y'(0) = r_1\lambda + r_2\mu :
\]
a linear system in $(\lambda, \mu)$ whose determinant is $r_2 -
r_1 \neq 0$; solving it explicitly, $\mu = \frac{y_0' -
r_1y_0}{r_2 - r_1}$ and $\lambda = y_0 - \mu$: exactly one
solution. In case (2), $y(0) = \lambda$ and $y'(0) = r_0\lambda +
\mu$: the system is triangular with determinant $1$, solved by
$\lambda = y_0$, $\mu = y_0' - r_0y_0$. In case (3), $y(0) =
\lambda$ and $y'(0) = \alpha\lambda + \omega\mu$: determinant
$\omega \neq 0$, solved by $\lambda = y_0$, $\mu = \frac{y_0' -
\alpha y_0}\omega$. In each case the \omterm{def:b1:logic:map}{map} $(\lambda, \mu) \mapsto
(y(x_0), y'(x_0))$ is a linear bijection --- the language of
\cref{ch:b1:linmaps} will compress this case check into one
sentence.
\end{proof}

\begin{method}[Particular solution for $f(x) = P(x)\,\eu^{\gamma x}$]\label{met:b1:diffeq:particular}
When the right-hand side is $P(x)\,\eu^{\gamma x}$ with $P$ a
polynomial and $\gamma \in \R$ (this covers polynomials, exponentials,
and via complex $\gamma$ or superposition, $\cos$ and $\sin$): look
for a particular solution of the form
\[
y_p(x) = x^{m}\, Q(x)\, \eu^{\gamma x},
\qquad
m = \text{multiplicity of } \gamma \text{ as a root of } \chi
\ (m = 0, 1 \text{ or } 2),
\]
with $Q$ a polynomial of the same degree as $P$, whose coefficients
are found by substitution and identification. For $f = K\cos\omega x$
(or $\sin$), solve with right-hand side $K\eu^{\iu\omega x}$ and take
the real (resp.\ imaginary) part.
\end{method}

\begin{example}[Superposition in action]\label{ex:b1:diffeq:superposition}
Solve $y'' - y = \eu^{x} + 4$ on $\R$. Homogeneous: $\chi(r) = r^2
- 1$, roots $\pm1$, so $y_h = \lambda\eu^x + \mu\eu^{-x}$. Split
the right-hand side and treat each piece by the method box.
\emph{Piece $\eu^x$:} here $\gamma = 1$ is a simple root of
$\chi$, so try $y_1 = c\,x\,\eu^x$: then $y_1'' - y_1 = c(x +
2)\eu^x - cx\eu^x = 2c\,\eu^x$, giving $c = \frac12$.
\emph{Piece $4$:} $\gamma = 0$ is not a root; the constant $y_2 =
-4$ works. By superposition
(\cref{thm:b1:diffeq:structure2}~(2)):
\[
y = \frac{x\,\eu^x}2 - 4 + \lambda\,\eu^x + \mu\,\eu^{-x},
\qquad (\lambda, \mu) \in \R^2 .
\]
Note how the two pieces demanded \emph{different} shapes ($m = 1$
versus $m = 0$): the multiplicity test is applied to each
exponent separately, which is the whole point of splitting the
right-hand side before guessing.
\end{example}

\begin{example}[The multiplicity rule at work]\label{ex:b1:diffeq:multiplicity}
Solve $y'' + y' = x$ on $\R$. The right-hand side is $P(x)\eu^{0
\cdot x}$ with $P(x) = x$, and $\gamma = 0$ is a \emph{simple} root
of $\chi(r) = r^2 + r = r(r + 1)$: so $m = 1$, and the correct
guess is $y_p = x\,(\alpha x + \beta) = \alpha x^2 + \beta x$, one
degree higher than $P$. Substituting:
\[
y_p'' + y_p' = 2\alpha + (2\alpha x + \beta)
= 2\alpha x + (2\alpha + \beta) ,
\]
and identification with $x$ gives $\alpha = \frac12$, $\beta = -1$:
$y_p = \frac{x^2}2 - x$. General solution: $y = \frac{x^2}2 - x +
\lambda + \mu\,\eu^{-x}$. Had we guessed $y_p = \alpha x + \beta$
(ignoring the multiplicity), substitution would give $y_p'' + y_p'
= \alpha$, a constant --- no choice of $\alpha, \beta$ can match
$x$, and the failure is structural: constants already solve the
homogeneous equation, so they are invisible to the left-hand side.
The factor $x^m$ exists precisely to climb out of the homogeneous
solution space.
\end{example}

\begin{remark}[The thirty-second insurance policy]\label{rem:b1:diffeq:check}
Every solved equation in this chapter ends with a substitution
check, and this is not decorative. A differential-equation
computation chains many small steps (a primitive, a product rule,
two constants), and a single sign error propagates invisibly;
substituting the final formula back into the equation catches
essentially all of them at the cost of one differentiation.
Cultivate the reflex in three layers: check the \emph{particular
solution} alone (the homogeneous part cancels anyway), check the
\emph{initial conditions} on the full solution, and when a
parameter is present, check a \emph{degenerate value} (does the
formula for general $\Omega$ reproduce the known answer at
$\Omega = 0$?). The habit costs half a minute; it converts
``probably right'' into ``verified''.
\end{remark}

\begin{remark}[Common pitfalls]\label{rem:b1:diffeq:pitfalls}
\begin{enumerate}
  \item \emph{Normalize first.} The formulas assume the equation
        reads $y' + a(x)y = b(x)$ --- coefficient $1$ on $y'$.
        For $xy' - 2y = x^3$, divide by $x$ (on an interval
        avoiding $0$) before identifying $a$ and $b$, as in
        \cref{exo:b1:diffeq:2}.
  \item \emph{One constant per dimension, fixed at the end.} The
        general first-order solution carries one constant, the
        second-order one two; initial conditions are imposed on
        the \emph{complete} solution $y_p + y_h$, never on $y_h$
        alone --- imposing them before adding $y_p$ is the most
        frequent structural error.
  \item \emph{Mind the multiplicity.} A particular-solution guess
        that solves the homogeneous equation is invisible to the
        left-hand side; the factor $x^m$ of
        \cref{met:b1:diffeq:particular} is not optional
        (\cref{ex:b1:diffeq:multiplicity}).
  \item \emph{Intervals are part of the answer.} Solutions live
        on intervals where the coefficients are continuous;
        gluing across a singularity can create spurious constants
        (\cref{exo:b1:diffeq:12}) or destroy uniqueness. ``Solve
        on $\R^*$'' means two independent problems.
\end{enumerate}
\end{remark}

\begin{example}[Off-resonance forcing]\label{ex:b1:diffeq:offresonance}
Solve $y'' + 4y = \sin x$ on $\R$. Natural frequency $2$, forcing
frequency $1$: since $\iu$ is \emph{not} a root of $\chi(r) = r^2
+ 4$, the multiplicity is $m = 0$ and a plain sinusoid suffices.
Trying $y_p = \alpha\sin x$ (no cosine needed: the equation has
no $y'$ term, and $\sin$ regenerates $\sin$):
\[
y_p'' + 4y_p = -\alpha\sin x + 4\alpha\sin x = 3\alpha\sin x ,
\]
so $\alpha = \frac13$ and the general solution is
\[
y = \frac{\sin x}3 + \lambda\cos 2x + \mu\sin 2x .
\]
Every solution stays bounded: a superposition of two oscillations
at the frequencies $1$ (forced) and $2$ (natural). Compare with
the next example, where forcing \emph{at} the natural frequency
changes the answer's very shape.
\end{example}

\begin{example}[A forced oscillation]\label{ex:b1:diffeq:oscillation}
Solve $y'' + y = \cos x$, $y(0) = 0$, $y'(0) = 0$.

\emph{Homogeneous:} $\chi(r) = r^2 + 1$, roots $\pm\iu$: $y_h =
\lambda\cos x + \mu \sin x$.

\emph{Particular:} right-hand side $\Re(\eu^{\iu x})$ with $\gamma =
\iu$ a simple root of $\chi$: try $z_p = c\, x\, \eu^{\iu x}$
($c \in \C$). Then $z_p'' + z_p = c\,(2\iu)\eu^{\iu x}$, which equals
$\eu^{\iu x}$ for $c = \frac{1}{2\iu} = -\frac\iu2$. So $z_p =
-\frac{\iu}{2} x (\cos x + \iu \sin x)$ and $y_p = \Re(z_p) =
\frac{x \sin x}{2}$.

\emph{General solution:} $y = \frac{x\sin x}{2} + \lambda\cos x +
\mu\sin x$. Conditions: $y(0) = \lambda = 0$; $y' = \frac{\sin x +
x\cos x}{2} + \mu\cos x$, so $y'(0) = \mu = 0$. Answer: $y =
\frac{x\sin x}{2}$ --- an oscillation whose amplitude grows linearly:
the \emph{resonance}\index{resonance} phenomenon, caused by forcing
the system at its natural frequency.
\end{example}

\begin{omfigure}
\begin{tikzpicture}
\begin{axis}[omaxis, width=11cm, height=5.2cm,
    xmin=0, xmax=25, ymin=-13, ymax=13,
    xtick={0,5,10,15,20,25}, ytick=\empty]
  \addplot[omThm, thick, domain=0:25, samples=400] {x*sin(deg(x))/2};
  \addplot[dashed, gray, domain=0:25] {x/2};
  \addplot[dashed, gray, domain=0:25] {-x/2};
\end{axis}
\end{tikzpicture}

{\small \omterm{ex:b1:diffeq:oscillation}{Resonance}: the solution $y = \frac{x \sin x}{2}$ of $y'' + y =
\cos x$ oscillates between the lines $y = \pm\frac x2$ (dashed), with
ever-growing amplitude.}
\end{omfigure}

\begin{remark}[Interlude: linearity is a geometry]\label{rem:b1:diffeq:interlude}
Look back at the shape of every solution \omterm{def:b1:logic:sets}{set} in this chapter: a
special solution plus a space of homogeneous solutions with one
free constant (first order) or two (second order). The chapters on
linear algebra (\cref{ch:b1:vspaces,ch:b1:findim,ch:b1:linmaps})
will supply the exact
vocabulary: the \omterm{def:b1:logic:map}{map} $L(y) = y'' + ay' + by$ is \emph{linear}, its
homogeneous solutions form the \emph{kernel} of $L$, a vector
space whose \emph{dimension} equals the order of the equation ---
that is the honest content of ``one constant per order'' --- and
the solution \omterm{def:b1:logic:sets}{set} of $L(y) = f$ is an \emph{affine subspace}, a
translate of the kernel. Even the Cauchy \omterm{def:b1:logic:map}{map} $(\lambda, \mu)
\mapsto (y(x_0), y'(x_0))$ of \cref{thm:b1:diffeq:structure2} is a
linear bijection between two planes, i.e.\ an invertible $2
\times 2$ system (\cref{ch:b1:matrices}). Nothing in this chapter
will need to be redone --- it will only need to be renamed, and
the renaming is the best possible warm-up for linear algebra:
every abstract definition there has already earned its living
here.
\end{remark}

\begin{remark}[Where this chapter is used]\label{rem:b1:diffeq:whereused}
The structure theorem --- solutions of $(E)$ form ``a particular
solution plus the solutions of $(H)$'' --- is the first appearance
of a pattern that \cref{ch:b1:vspaces,ch:b1:linmaps} will name: the
solution \omterm{def:b1:logic:sets}{set} of $(H)$ is the \emph{kernel} of the linear \omterm{def:b1:logic:map}{map} $y
\mapsto y'' + ay' + by$, and the solution \omterm{def:b1:logic:sets}{set} of $(E)$ is an
affine translate of it. The \omterm{def:b1:diffeq:linear2}{characteristic polynomial} reappears as
the \omterm{def:b1:diffeq:linear2}{characteristic polynomial} of a matrix in \cref{ch:b1:matrices}:
a second-order equation is a $2 \times 2$ first-order system in
disguise, a viewpoint the Year 2 volume systematizes. The
integrals demanded by \omterm{thm:b1:diffeq:voc}{variation of constants} are supplied by
\cref{ch:b1:integration}, and the weekend problem below --- the
driven damped oscillator --- is the model case for every
oscillation question in the sciences, from circuits to suspension
bridges.
\end{remark}

\section{Exercises}

\begin{exercise}[$\star$]\label{exo:b1:diffeq:1}
Solve on $\R$: $\;y' + 2y = \eu^{3x}$; then the \omterm{thm:b1:diffeq:voc}{Cauchy problem} $y(0) =
1$.
\end{exercise}

\begin{exercise}[$\star$]\label{exo:b1:diffeq:2}
Solve on $\intoo{0}{+\infty}$: $\;x y' - 2y = x^3$ \emph{(put the
equation in the normalized form first)}.
\end{exercise}

\begin{exercise}[$\star$]\label{exo:b1:diffeq:3}
Solve on $\R$, giving the real general solution:
$\;y'' - 3y' + 2y = 0$; $\;y'' + 4y' + 4y = 0$; $\;y'' - 2y' + 5y =
0$.
\end{exercise}

\begin{exercise}[$\star$]\label{exo:b1:diffeq:4}
Solve $y'' - y = x^2$ on $\R$, then the \omterm{thm:b1:diffeq:voc}{Cauchy problem} $y(0) = 0$,
$y'(0) = 1$.
\end{exercise}

\begin{exercise}[$\star\star$]\label{exo:b1:diffeq:5}
Solve on $\intoo{-\frac\pi2}{\frac\pi2}$:
$\;y' + y\tan x = \sin 2x$.
\end{exercise}

\begin{exercise}[$\star\star$]\label{exo:b1:diffeq:6}
Solve $y'' - 4y' + 3y = (2x + 1)\,\eu^{x}$ on $\R$. \emph{(Mind the
multiplicity: is $1$ a root of the \omterm{def:b1:diffeq:linear2}{characteristic polynomial}?)}
\end{exercise}

\begin{exercise}[$\star\star$]\label{exo:b1:diffeq:7}
Solve $y'' + 4y = \sin 2x + x$ on $\R$ \emph{(superposition; treat
each right-hand side separately)}.
\end{exercise}

\begin{exercise}[$\star\star$]\label{exo:b1:diffeq:8}
A cup of coffee at temperature $T_0 = 80\,^\circ$C sits in a room at
$20\,^\circ$C. Newton's law of cooling states $T' = -k\,(T - 20)$ with
$k > 0$. Solve for $T(t)$, and given that the coffee is at
$50\,^\circ$C after $10$ minutes, find when it reaches
$25\,^\circ$C.
\end{exercise}

\begin{exercise}[$\star\star\star$]\label{exo:b1:diffeq:9}
(Damped oscillator) For $\varepsilon \geq 0$, consider
$y'' + 2\varepsilon y' + y = 0$.
\begin{enumerate}
  \item Solve for $\varepsilon \in \intco{0}{1}$, $\varepsilon = 1$,
        and $\varepsilon > 1$.
  \item Show that for $\varepsilon > 0$ every solution tends to $0$
        at $+\infty$, and that for $\varepsilon = 0$ the nonzero
        solutions do not.
  \item For $\varepsilon \in \intoo{0}{1}$, show that the zeros of a
        nonzero solution are regularly spaced, with gap
        $\frac{\pi}{\sqrt{1 - \varepsilon^2}}$.
\end{enumerate}
\end{exercise}

\begin{exercise}[$\star\star\star$]\label{exo:b1:diffeq:10}
Find all functions $f \colon \R \to \R$, twice differentiable, such
that
\[
\forall x, y \in \R, \qquad f(x + y) + f(x - y) = 2 f(x) f(y),
\]
with $f(0) \ne 0$ and $f$ not constant.
\emph{Hint: fix $y$, differentiate twice with respect to $x$ at $0$;
show $f(0) = 1$ and $f'' = c f$ for some constant $c$; then solve
according to the sign of $c$ and check which solutions satisfy the
functional equation.}
\end{exercise}

\begin{exercise}[$\star\star$]\label{exo:b1:diffeq:11}
(Euler equation) Solve $x^2 y'' - x y' + y = 0$ on
$\intoo{0}{+\infty}$. \emph{Hint: set $z(t) = y(\eu^t)$, i.e.\
substitute $x = \eu^t$, and show that $z$ satisfies a linear
equation with constant coefficients.}
\end{exercise}

\begin{exercise}[$\star\star\star$]\label{exo:b1:diffeq:12}
Consider the equation $x\,y' = 2y$ on the whole real line, in the
unknown differentiable function $y \colon \R \to \R$.
\begin{enumerate}
  \item Solve on $\intoo{0}{+\infty}$ and on $\intoo{-\infty}{0}$.
  \item Show that for \emph{any} constants $a, b \in \R$, the
        function equal to $ax^2$ for $x \geq 0$ and to $bx^2$ for
        $x < 0$ is differentiable on $\R$ and solves the equation
        everywhere.
  \item Conclude that the solution \omterm{def:b1:logic:sets}{set} on $\R$ is a two-parameter
        family, and explain why this does not contradict the
        uniqueness in \cref{thm:b1:diffeq:voc}.
\end{enumerate}
\end{exercise}

\section{Problem: The driven damped oscillator}

\begin{problem}\label{pb:b1:diffeq:1}
One equation governs a mass on a spring in a viscous medium, the
charge in an RLC circuit, and a building swaying in the wind:
\[
(E_\Omega)\colon\quad
x'' + 2\lambda x' + \omega_0^2\,x = A\cos(\Omega t),
\]
with $\lambda \geq 0$ the damping, $\omega_0 > 0$ the natural
frequency, and $A > 0$, $\Omega > 0$ the amplitude and frequency of
the forcing. This problem extracts its complete behavior: the decay
of transients, the unique periodic steady state, the \omterm{ex:b1:diffeq:oscillation}{resonance}
curve and its sharpness (the \emph{\omterm{pb:b1:diffeq:1}{quality factor}}), the beats of
the undamped case, and the energy balance that sustains the
oscillation.\index{resonance}\index{oscillator}
Unless stated otherwise, $0 < \lambda < \omega_0$ (underdamped
regime) and we write $\omega_d = \sqrt{\omega_0^2 - \lambda^2}$.

\medskip
\noindent\textbf{Part I --- The free oscillator.} Here $A = 0$.
\begin{enumerate}
  \item Solve the homogeneous equation $(H)$ for $0 < \lambda <
        \omega_0$, and for $\lambda = 0$. (The regimes $\lambda
        \geq \omega_0$ were treated in \cref{exo:b1:diffeq:9};
        quote them.)
  \item Show that for every $\lambda > 0$, all solutions of $(H)$
        tend to $0$ at $+\infty$ --- in all three regimes.
  \item Define the energy $\mathcal E(t) = \frac12 x'(t)^2 +
        \frac12\omega_0^2\,x(t)^2$ along a solution of $(H)$. Show
        $\mathcal E'(t) = -2\lambda\,x'(t)^2 \leq 0$, and deduce
        (without solving anything) that the \omterm{thm:b1:diffeq:voc}{Cauchy problem}
        ``$(H)$, $x(t_0) = x'(t_0) = 0$'' has only the zero
        solution, for every $\lambda \geq 0$.
  \item For $0 < \lambda < \omega_0$, write the nonzero solution as
        $x(t) = R\,\eu^{-\lambda t}\cos(\omega_d t - \varphi)$ and
        let $T_d = \frac{2\pi}{\omega_d}$ be the pseudo-period.
        Show that $x(t + T_d) = \eu^{-\lambda T_d}\,x(t)$: each
        swing is the previous one shrunk by the constant factor
        $\eu^{-\delta}$, $\delta = \frac{2\pi\lambda}{\omega_d}$
        (the \emph{logarithmic decrement}). Compute $\delta$ for
        $\omega_0 = 1$, $\lambda = 0.1$.
  \item Define the \emph{\omterm{pb:b1:diffeq:1}{quality factor}} $Q =
        \dfrac{\omega_0}{2\lambda}$\index{quality factor}. Show
        that after the time $\frac1\lambda$ (one amplitude
        $\eu$-folding), the oscillator has completed
        $\frac{\omega_d}{2\pi\lambda}$ pseudo-periods, which for
        weak damping ($\lambda \ll \omega_0$) is approximately
        $\frac Q\pi$: the \omterm{pb:b1:diffeq:1}{quality factor} counts, up to $\pi$, the
        oscillations survived before the amplitude decays by
        $\eu$.
\end{enumerate}

\medskip
\noindent\textbf{Part II --- The steady state.} Now $A > 0$ and
$\lambda > 0$.
\begin{enumerate}[resume]
  \item Look for a particular solution as the real part of
        $z\,\eu^{\iu\Omega t}$ with $z \in \C$
        (\cref{met:b1:diffeq:particular}). Show that this works
        with
        \[
        z = \frac{A}{\omega_0^2 - \Omega^2 + 2\iu\lambda\Omega} .
        \]
  \item Deduce the steady state in amplitude--phase form:
        $x_p(t) = R(\Omega)\cos\bigl(\Omega t -
        \varphi(\Omega)\bigr)$ with
        \[
        R(\Omega) = \frac{A}{\sqrt{(\omega_0^2 - \Omega^2)^2 +
        4\lambda^2\Omega^2}},
        \qquad
        \tan\varphi = \frac{2\lambda\Omega}{\omega_0^2 -
        \Omega^2},
        \quad \varphi \in \intoo0\pi .
        \]
  \item Interpret the two extreme regimes: compute the limits of
        $R$ and $\varphi$ as $\Omega \to 0^+$ (quasi-static
        response $A/\omega_0^2$, phase $0$) and as $\Omega \to
        +\infty$ ($R \sim A/\Omega^2 \to 0$, phase $\to \pi$: the
        mass moves opposite to a too-fast forcing).
  \item Show that \emph{every} solution of $(E_\Omega)$ is $x_p$
        plus a solution of $(H)$, hence converges to the steady
        state $x_p$ as $t \to +\infty$, whatever the initial
        conditions: after the transient dies, the oscillator has
        no memory of how it started.
  \item Show that $x_p$ is the \emph{only} periodic solution of
        $(E_\Omega)$.
  \item Work one \omterm{thm:b1:diffeq:voc}{Cauchy problem} to the end: for $x'' + 2x' + 2x =
        \cos t$ with $x(0) = x'(0) = 0$, show that the solution is
        \[
        x(t) = \frac{\cos t + 2\sin t}5
        - \eu^{-t}\,\frac{\cos t + 3\sin t}5 ,
        \]
        and identify transient and steady parts.
\end{enumerate}

\medskip
\noindent\textbf{Part III --- The \omterm{ex:b1:diffeq:oscillation}{resonance} curve.} Study of
$\Omega \mapsto R(\Omega)$ on $\intoo0{+\infty}$.
\begin{enumerate}[resume]
  \item Setting $u = \Omega^2$ and $g(u) = (\omega_0^2 - u)^2 +
        4\lambda^2 u$, show: if $2\lambda^2 < \omega_0^2$, then $R$
        attains a strict maximum at the \emph{\omterm{ex:b1:diffeq:oscillation}{resonance} frequency}
        $\Omega_r = \sqrt{\omega_0^2 - 2\lambda^2}$, with
        \[
        R_{\max} = R(\Omega_r)
        = \frac{A}{2\lambda\sqrt{\omega_0^2 - \lambda^2}} .
        \]
  \item Show that $\dfrac{R_{\max}}{R(0)} = Q\,\bigl(1 -
        \tfrac{\lambda^2}{\omega_0^2}\bigr)^{-1/2} \geq Q$: at
        \omterm{ex:b1:diffeq:oscillation}{resonance}, the forcing is amplified by (essentially) the
        \omterm{pb:b1:diffeq:1}{quality factor}.
  \item Prove that the \emph{velocity} amplitude $V(\Omega) =
        \Omega\,R(\Omega)$ is maximal exactly at $\Omega =
        \omega_0$ (not at $\Omega_r$), and that the phase there is
        $\varphi(\omega_0) = \frac\pi2$: at $\Omega = \omega_0$
        the velocity is exactly in phase with the force.
  \item (Bandwidth) Solve $g(u) = 2\,g(u_r)$ exactly, where $u_r =
        \omega_0^2 - 2\lambda^2$, and deduce that the two
        frequencies $\Omega_\pm$ where $R = R_{\max}/\sqrt2$
        satisfy $\Omega_+^2 - \Omega_-^2 =
        4\lambda\sqrt{\omega_0^2 - \lambda^2}$; conclude that for
        weak damping the bandwidth is $\Omega_+ - \Omega_-
        \approx 2\lambda$, i.e.\ $Q \approx
        \frac{\omega_0}{\Omega_+ - \Omega_-}$: sharp \omterm{ex:b1:diffeq:oscillation}{resonance}
        peaks are high-$Q$ systems.
  \item Numerical portrait for $\omega_0 = 1$, $\lambda = 0.05$
        ($Q = 10$), $A = 1$: compute $\Omega_r$, $R_{\max}$, the
        static response $R(0)$, and the approximate bandwidth.
  \item Show that if $2\lambda^2 \geq \omega_0^2$, then $R$ is
        strictly decreasing on $\intoo0{+\infty}$: heavily damped
        systems have no \omterm{ex:b1:diffeq:oscillation}{resonance} peak at all.
\end{enumerate}

\medskip
\noindent\textbf{Part IV --- No damping: beats and \omterm{ex:b1:diffeq:oscillation}{resonance}.}
Here $\lambda = 0$.
\begin{enumerate}[resume]
  \item For $\Omega \neq \omega_0$, find the general solution of
        $x'' + \omega_0^2 x = A\cos(\Omega t)$.
  \item Solve the \omterm{thm:b1:diffeq:voc}{Cauchy problem} $x(0) = x'(0) = 0$ and transform
        the answer into the product form
        \[
        x(t) = \frac{2A}{\omega_0^2 - \Omega^2}\,
        \sin\Bigl(\frac{(\omega_0 - \Omega)t}2\Bigr)
        \sin\Bigl(\frac{(\omega_0 + \Omega)t}2\Bigr) .
        \]
  \item For $\Omega$ close to $\omega_0$, read the product as a
        fast oscillation at frequency $\frac{\omega_0 + \Omega}2$
        modulated by a slow envelope at frequency
        $\frac{\abs{\omega_0 - \Omega}}2$: the \emph{beats}
        \index{beats}. Give the period of the envelope and the
        maximal amplitude, and note how both blow up as $\Omega
        \to \omega_0$.
  \item Fix $t$ and let $\Omega \to \omega_0$ in question 19's
        formula: show the limit is
        \[
        x_\infty(t) = \frac{A\,t\,\sin(\omega_0 t)}{2\omega_0} ,
        \]
        and check directly that $x_\infty$ solves the resonant
        equation $x'' + \omega_0^2 x = A\cos(\omega_0 t)$ with
        $x(0) = x'(0) = 0$ (compare
        \cref{ex:b1:diffeq:oscillation}): \omterm{ex:b1:diffeq:oscillation}{resonance} is the limit
        of ever-slower, ever-larger beats.
  \item Contrast the two fates of \omterm{ex:b1:diffeq:oscillation}{resonance}: linear growth
        $\frac{At}{2\omega_0}$ without damping, versus saturation
        at $R_{\max} \approx Q\,\frac{A}{\omega_0^2}$ with weak
        damping. In one sentence: what physical mechanism converts
        the first into the second?
\end{enumerate}

\medskip
\noindent\textbf{Part V --- Energy balance and synthesis.}
\begin{enumerate}[resume]
  \item In the steady state of Part II, compute the average over
        one period $\frac{2\pi}\Omega$ of (a) the power injected
        by the forcing, $P_{\mathrm{in}}(t) = A\cos(\Omega t)
        \cdot x_p'(t)$, and (b) the power dissipated by the
        damping, $P_{\mathrm{diss}}(t) = 2\lambda\,x_p'(t)^2$.
        Show both averages equal $\lambda\,R^2\Omega^2$: the
        forcing feeds in exactly what the damping burns --- this
        is why the steady state is steady.
  \item Where exactly did the problem use: (i) the structure
        theorem \cref{thm:b1:diffeq:structure2}; (ii) the complex
        exponential method; (iii) a real-variable function study
        in the style of \cref{ch:b1:functions}? One sentence
        each.
  \item Synthesis: describe the full behavior \omterm{def:b1:logic:map}{map} of $(E_\Omega)$
        --- free versus forced, damped versus undamped, the role
        of $Q$ as the single dimensionless dial tuning peak
        height, bandwidth and transient lifetime --- and mention
        where the story continues: first-order $2 \times 2$
        systems (\cref{ch:b1:matrices} and the Year 2 volume) and
        the decomposition of a general periodic forcing into
        sinusoids (Fourier series, in the Year 3 volume), for
        which this problem's sinusoidal case is the fundamental
        building block.
\end{enumerate}
\end{problem}



\appendix
\part{\omsolutionspart}
% Solutions to the Bachelor Year 1 exercises — one file per chapter.
\chapter{Bachelor Year 1 (Age 18--19)}

\section*{Chapter \ref{ch:b1:logic} --- Logic, Sets and Maps}

\begin{solution}{exo:b1:logic:1}
Negations, pushing $\lnot$ through each quantifier
(\cref{prop:b1:logic:negquant}) and using $\lnot(P \implies Q) \iff
P \land \lnot Q$:
\begin{enumerate}
  \item $\exists x \in \R,\ \forall y \in \R,\ x + y \leq 0$;
  \item $\forall x \in \R,\ \exists y \in \R,\ xy \neq 0$;
  \item $\exists \varepsilon > 0,\ \forall \delta > 0,\ \exists x \in
        \R,\ \abs{x} \leq \delta \text{ and } \abs{f(x)} >
        \varepsilon$.
\end{enumerate}
Statement (1) is true: given $x$, take $y = -x + 1$; then $x + y = 1 >
0$. Statement (2) is true: $x = 0$ satisfies $xy = 0$ for all $y$.
\end{solution}

\begin{solution}{exo:b1:logic:2}
Truth table, writing T/F for the four cases $(P, Q)$:
\begin{center}
\begin{tabular}{cc|c|c|c|c}
$P$ & $Q$ & $P \implies Q$ & $\lnot(P \implies Q)$ & $\lnot Q$ &
$P \land \lnot Q$ \\
\hline
T & T & T & F & F & F \\
T & F & F & T & T & T \\
F & T & T & F & F & F \\
F & F & T & F & T & F \\
\end{tabular}
\end{center}
Columns $4$ and $6$ coincide, proving the equivalence. The negation of
``if a function is differentiable then it is continuous'' is therefore:
``there is a function which is differentiable and not continuous''
(a false statement, as it happens: the original implication is true,
see \cref{ch:b1:derivative}).
\end{solution}

\begin{solution}{exo:b1:logic:3}
\emph{Contraposition.} Assume $x < 1$. Then $x^3 < 1$ (the cube
function is increasing) and $x < 1$, so $x^3 + x < 2$. This proves the
contrapositive, hence the statement.

\emph{Contradiction.} Suppose $a > 0$ is the smallest strictly positive
real. Then $a/2$ is strictly positive and $a/2 < a$ (since $a > 0$),
contradicting minimality. Hence no such $a$ exists.
\end{solution}

\begin{solution}{exo:b1:logic:4}
\begin{enumerate}
  \item Base case $n = 0$: $2^0 = 1 = 2^1 - 1$. Step: assuming the
        identity for $n$,
        \[
        \sum_{k=0}^{n+1} 2^k = (2^{n+1} - 1) + 2^{n+1}
        = 2 \cdot 2^{n+1} - 1 = 2^{n+2} - 1 .
        \]
  \item Base case $n = 0$: $4^0 + 5 = 6 = 3 \times 2$. Step: if
        $4^n + 5 = 3m$, then
        \[
        4^{n+1} + 5 = 4(4^n + 5) - 15 = 3(4m - 5),
        \]
        divisible by $3$.
\end{enumerate}
\end{solution}

\begin{solution}{exo:b1:logic:5}
The inductive step silently assumes that the two groups (``the first
$n$'' and ``the last $n$'') overlap, so that the shared pencils carry
the color from one group to the other. For $n + 1 = 2$ the two groups
are $\{$first pencil$\}$ and $\{$second pencil$\}$: they are disjoint,
and the argument breaks. So $P(1) \implies P(2)$ was never proved, and
the induction collapses --- even though $P(n) \implies P(n+1)$ is
valid for every $n \geq 2$.
\end{solution}

\begin{solution}{exo:b1:logic:6}
\begin{enumerate}
  \item $x \in A \setminus B \iff x \in A \land x \notin B \iff x \in
        A \land x \in \overline{B} \iff x \in A \cap \overline{B}$.
  \item Using (1) and distributivity
        (\cref{prop:b1:logic:setalgebra}):
        $(A \cup B) \cap \overline{C} = (A \cap \overline{C}) \cup
        (B \cap \overline{C})$.
  \item Suppose $A \subseteq B$. Then $A \cup B \subseteq B$ (both
        pieces lie in $B$) and $B \subseteq A \cup B$ always, so
        $A \cup B = B$. Suppose $A \cup B = B$: then $A \cap B
        \subseteq A$ always, and $A \subseteq A \cup B = B$ gives
        $A \subseteq A \cap B$, so $A \cap B = A$. Suppose $A \cap B =
        A$: then $A = A \cap B \subseteq B$. The three conditions are
        thus equivalent (we proved a cycle of implications).
\end{enumerate}
\end{solution}

\begin{solution}{exo:b1:logic:7}
\begin{enumerate}
  \item $f$ is injective ($n + 1 = m + 1 \implies n = m$) but not
        surjective: $0$ has no preimage in $\N$.
  \item $g$ is bijective: $n \mapsto n - 1$ is a two-sided inverse on
        $\Z$.
  \item $h$ is injective: $\frac{x+1}{x-1} = \frac{x'+1}{x'-1}$ gives
        $(x+1)(x'-1) = (x'+1)(x-1)$, i.e.\ $xx' - x + x' - 1 = xx' -
        x' + x - 1$, so $2x' = 2x$. It is not surjective onto $\R$:
        solving $y = \frac{x+1}{x-1}$ gives $x(y - 1) = y + 1$, which
        has no solution when $y = 1$ (the equation reads $0 = 2$).
        With codomain $\R \setminus \{1\}$, the same computation gives
        the unique preimage $x = \frac{y+1}{y-1}$, so
        $h \colon \R \setminus \{1\} \to \R \setminus \{1\}$ is
        bijective and $h^{-1}(y) = \frac{y+1}{y-1} = h(y)$: $h$ is its
        own inverse.
\end{enumerate}
\end{solution}

\begin{solution}{exo:b1:logic:8}
\begin{enumerate}
  \item $x \in f^{-1}(B \cap B') \iff f(x) \in B \cap B' \iff f(x) \in
        B \land f(x) \in B' \iff x \in f^{-1}(B) \cap f^{-1}(B')$.
        For images: $y \in f(A \cup A')$ iff $y = f(x)$ for some $x$
        in $A$ or in $A'$, iff $y \in f(A)$ or $y \in f(A')$.
  \item If $y \in f(A \cap A')$, then $y = f(x)$ with $x \in A$ and
        $x \in A'$, so $y \in f(A)$ and $y \in f(A')$. Strictness:
        take $f \colon \R \to \R$, $x \mapsto x^2$, $A = \{-1\}$,
        $A' = \{1\}$: then $f(A \cap A') = f(\emptyset) = \emptyset$
        but $f(A) \cap f(A') = \{1\}$.
  \item ($\Leftarrow$) With $A = \{x\}$, $A' = \{x'\}$ for $x \neq
        x'$: if $f(x) = f(x')$, then $f(A) \cap f(A') = \{f(x)\}$
        while $f(A \cap A') = \emptyset$, contradicting the assumed
        equality; so $f$ is injective. ($\Rightarrow$) Let $f$ be
        injective and $y \in f(A) \cap f(A')$: $y = f(x) = f(x')$ with
        $x \in A$, $x' \in A'$; injectivity gives $x = x' \in A \cap
        A'$, so $y \in f(A \cap A')$. With (2), equality holds.
\end{enumerate}
\end{solution}

\begin{solution}{exo:b1:logic:9}
$g \circ f = \mathrm{id}_E$ is injective and surjective, so by
\cref{prop:b1:logic:comp}~(2), $f$ is injective and $g$ is surjective.
Example: $E = \N$, $F = \Z$, $f$ the inclusion $n \mapsto n$, and
$g \colon \Z \to \N$, $g(n) = n$ for $n \geq 0$ and $g(n) = 0$ for
$n < 0$. Then $g(f(n)) = n$ for all $n \in \N$, but $f$ is not
surjective and $g$ is not injective.
\end{solution}

\begin{solution}{exo:b1:logic:10}
$x^2 - y^2 = x - y \iff (x - y)(x + y) = x - y \iff (x - y)(x + y - 1)
= 0 \iff y = x$ or $y = 1 - x$. \emph{Reflexive:} $y = x$ works.
\emph{Symmetric:} the condition ``$y = x$ or $y = 1 - x$'' is
symmetric in $x$ and $y$ (if $y = 1 - x$ then $x = 1 - y$).
\emph{Transitive:} suppose $x \mathbin{\mathcal{R}} y$ and $y
\mathbin{\mathcal{R}} z$; going through the four cases, $z$ equals $x$
or $1 - x$ each time (e.g.\ $y = 1 - x$ and $z = 1 - y$ give $z = x$).
So $\mathcal{R}$ is an equivalence relation and $\mathrm{cl}(x) =
\{x,\, 1 - x\}$. This class has one element exactly when $x = 1 - x$,
i.e.\ for $x = \frac12$.
\end{solution}

\begin{solution}{exo:b1:logic:11}
Let $f \colon E \to \mathcal{P}(E)$ be any map and set
$D = \{x \in E : x \notin f(x)\} \in \mathcal{P}(E)$. Suppose $D =
f(a)$ for some $a \in E$. If $a \in D$, then by definition of $D$,
$a \notin f(a) = D$: contradiction. If $a \notin D$, then $a \notin
f(a)$, so by definition of $D$, $a \in D$: contradiction. Hence $D$ is
not in the image of $f$, and $f$ is not surjective. (In particular no
set is in bijection with its power set: there are ``more'' subsets of
$\N$ than integers.)
\end{solution}

\begin{solution}{exo:b1:logic:12}
\begin{enumerate}
  \item ($\Rightarrow$) Let $f$ be surjective and $\Phi(B) =
        \Phi(B')$. For $y \in B$, pick $x$ with $f(x) = y$; then $x
        \in f^{-1}(B) = f^{-1}(B')$, so $y = f(x) \in B'$. Hence $B
        \subseteq B'$, and symmetrically $B' \subseteq B$: $\Phi$ is
        injective. ($\Leftarrow$) If $f$ is not surjective, pick $y_0
        \in F$ outside the image; then $f^{-1}(\{y_0\}) = \emptyset =
        f^{-1}(\emptyset)$ with $\{y_0\} \neq \emptyset$, so $\Phi$ is
        not injective.
  \item ($\Rightarrow$) Let $f$ be injective and $A \subseteq E$. Set
        $B = f(A)$; then $f^{-1}(B) = \{x : f(x) \in f(A)\}$, and
        injectivity gives $f(x) \in f(A) \iff x \in A$, so $\Phi(B) =
        A$: $\Phi$ is surjective. ($\Leftarrow$) If $f$ is not
        injective, take $x \neq x'$ with $f(x) = f(x')$. Every
        preimage set $f^{-1}(B)$ contains $x$ if and only if it
        contains $x'$; hence $\{x\}$ is not of the form $\Phi(B)$, and
        $\Phi$ is not surjective.
\end{enumerate}
\end{solution}

\section*{Chapter \ref{ch:b1:counting} --- Counting}

\begin{solution}{exo:b1:counting:1}
Independent stages and the product rule: $26^2 \times 10^3 \times 26^2
= 26^4 \times 1000 = 456\,976\,000$ plates. With the four letters
pairwise distinct, the letter stages form a $4$-arrangement of the
alphabet: $26 \times 25 \times 24 \times 23 = 358\,800$ ways, so
$358\,800 \times 1000 = 358\,800\,000$ plates.
\end{solution}

\begin{solution}{exo:b1:counting:2}
\textsc{orange} has $6$ distinct letters: $6! = 720$ anagrams.
\textsc{banana} has $6$ letters with repetitions ($3$ a's, $2$ n's,
$1$ b): each anagram is determined by the positions of the a's
($\binom 63$ choices), then of the n's among the remaining $3$ spots
($\binom 32$), the b taking the last spot:
$\binom{6}{3}\binom{3}{2} = 20 \times 3 = 60$ anagrams
(equivalently $6!/(3!\,2!\,1!) = 60$).
\end{solution}

\begin{solution}{exo:b1:counting:3}
Total: $\binom{12}{4} = 495$. Exactly $2$ women: choose them
($\binom 72 = 21$) and $2$ men ($\binom 52 = 10$): $210$ committees.
At least one man: complement of ``no man'',
$\binom{12}{4} - \binom{7}{4} = 495 - 35 = 460$.
\end{solution}

\begin{solution}{exo:b1:counting:4}
\emph{Birthdays:} the boxes are the $12$ months; $13$ people into $12$
boxes force two in the same box
(\cref{cor:b1:counting:pigeonhole}).

\emph{Consecutive integers:} the boxes are the $n$ pairs $\{1,2\},
\{3,4\}, \dots, \{2n-1, 2n\}$, which partition $\intint{1}{2n}$.
Choosing $n + 1$ integers puts two in the same pair, and the two
elements of a pair are consecutive.
\end{solution}

\begin{solution}{exo:b1:counting:5}
For $1 \leq k \leq n$,
\[
k \binom nk = k\,\frac{n!}{k!\,(n-k)!}
= n\,\frac{(n-1)!}{(k-1)!\,(n-k)!} = n \binom{n-1}{k-1}.
\]
Summing and reindexing with $j = k - 1$:
\[
\sum_{k=0}^{n} k \binom nk = n \sum_{j=0}^{n-1} \binom{n-1}{j}
= n\, 2^{n-1}
\]
by \cref{prop:b1:counting:identities}. (Alternative: differentiate
$(1+x)^n = \sum_k \binom nk x^k$ and set $x = 1$.)
\end{solution}

\begin{solution}{exo:b1:counting:6}
A strictly increasing map $f \colon \intint{1}{k} \to \intint{1}{n}$
is determined by its image, a $k$-subset of $\intint{1}{n}$ (list the
subset in increasing order); conversely every $k$-subset gives exactly
one such map. Hence $\binom nk$ strictly increasing maps.

If $f$ is merely increasing, set $g(i) = f(i) + i - 1$. Then $g$ is
strictly increasing (between consecutive arguments, $f$ gains $\geq 0$
and $i - 1$ gains $1$) with values in $\intint{1}{n + k - 1}$; and
$f(i) = g(i) - i + 1$ recovers $f$ from any strictly increasing $g$
into $\intint{1}{n+k-1}$. This is a bijection, so there are
$\binom{n + k - 1}{k}$ increasing maps.
\end{solution}

\begin{solution}{exo:b1:counting:7}
Split a set $E$ with $m + n$ elements into blocks $M$ ($m$ elements)
and $N$ ($n$ elements). A $k$-subset of $E$ contains some $j$ elements
of $M$ ($0 \leq j \leq k$) and $k - j$ of $N$; for fixed $j$ there are
$\binom mj \binom{n}{k-j}$ such subsets, and the cases $j = 0, \dots,
k$ partition the $k$-subsets. The sum rule gives Vandermonde's
identity\index{Vandermonde's identity}.

With $m = n = k$: $\binom{2n}{n} = \sum_{j=0}^{n} \binom nj
\binom{n}{n-j} = \sum_{j=0}^{n} \binom nj^2$, using $\binom{n}{n-j} =
\binom nj$.
\end{solution}

\begin{solution}{exo:b1:counting:8}
Let $A_d$ be the multiples of $d$ in $\intint{1}{1000}$, so
$\abs{A_d} = \lfloor 1000/d \rfloor$. Inclusion--exclusion
(\cref{thm:b1:counting:inclexcl}) with $A_2, A_3, A_5$, noting
$A_2 \cap A_3 = A_6$ etc.:
\[
500 + 333 + 200 - 166 - 100 - 66 + 33 = 734 .
\]
So $734$ integers are divisible by $2$, $3$ or $5$.
\end{solution}

\begin{solution}{exo:b1:counting:9}
Onto $2$ elements: all $2^4 = 16$ maps except the $2$ constant maps:
$14$ surjections.

Onto $3$ elements: by inclusion--exclusion on the missed values, the
number of maps from a $4$-set to a $3$-set missing at least one value
is $\binom 31 2^4 - \binom 32 1^4 = 48 - 3 = 45$; total maps $3^4 =
81$; surjections: $81 - 45 = 36$. (Check: a surjection from $4$ onto
$3$ elements doubles up exactly one value: choose the doubled value
($3$), the pair mapping to it ($\binom 42 = 6$), and a bijection for
the rest ($2$): $3 \times 6 \times 2 = 36$.)
\end{solution}

\begin{solution}{exo:b1:counting:10}
A solution of $x_1 + \dots + x_n = k$ in $\N^n$ encodes as a row of
$k$ stars and $n - 1$ bars: write $x_1$ stars, a bar, $x_2$ stars, a
bar, \dots, ending with $x_n$ stars. This is a bijection onto the
words of length $k + n - 1$ using $k$ stars and $n - 1$ bars, and
those words are determined by the positions of the stars:
$\binom{n + k - 1}{k}$. Selections with repetition correspond to
solutions of the equation ($x_i$ = number of copies of object $i$), so
the count is the same.
\end{solution}

\begin{solution}{exo:b1:counting:11}
With $A_i = \{\sigma : \sigma(i) = i\}$, a permutation in
$\bigcap_{i \in I} A_i$ fixes every $i \in I$ and permutes the other
$n - \abs I$ points freely: $\abs{\bigcap_{i \in I} A_i} =
(n - \abs I)!$. Inclusion--exclusion:
\[
\Bigl|\bigcup_i A_i\Bigr|
= \sum_{k=1}^{n} (-1)^{k+1} \binom nk (n-k)!
= \sum_{k=1}^{n} (-1)^{k+1} \frac{n!}{k!} ,
\]
since there are $\binom nk$ subsets $I$ of size $k$. Hence
\[
D_n = n! - \Bigl|\bigcup_i A_i\Bigr|
= n!\Bigl(1 - \sum_{k=1}^{n} \frac{(-1)^{k+1}}{k!}\Bigr)
= n! \sum_{k=0}^{n} \frac{(-1)^k}{k!} .
\]
For the second identity: classify the permutations $\sigma$ of
$\intint{1}{n}$ by their fixed-point set $F(\sigma)$. For a fixed
$k$-subset $F$, the permutations with $F(\sigma) = F$ are exactly the
derangements of the complement: $D_{n-k}$ of them. Summing over the
$\binom nk$ choices of $F$ for each $k$:
$n! = \sum_{k=0}^{n} \binom nk D_{n-k}$.
\end{solution}

\begin{solution}{exo:b1:counting:12}
\emph{First identity.} Count the pairs $(A, a)$ where $A \subseteq E$
($\abs E = n$) and $a \in A$. By size of $A$: $\sum_k \binom nk k$
pairs. By choosing the marked element first: $n$ choices for $a$, then
any subset of the remaining $n - 1$ elements to complete $A$:
$n\,2^{n-1}$ pairs.

\emph{Second identity.} Count the triples $(A, a, b)$ with $a, b \in
A$ (possibly $a = b$). By size: $\sum_k k^2 \binom nk$. Directly:
either $a = b$ ($n\,2^{n-1}$ triples, previous count) or $a \neq b$
($n(n-1)$ ordered choices, then any subset of the other $n - 2$
elements: $n(n-1)\,2^{n-2}$). Total
\[
n\,2^{n-1} + n(n-1)\,2^{n-2} = n\,2^{n-2}\,(2 + n - 1)
= n(n+1)\,2^{n-2} .
\]
\end{solution}

\begin{solution}{pb:b1:counting:1}
\textbf{1.} $D_1 = 0$ (the only permutation fixes $1$), $D_2 = 1$
(the swap), $D_3 = 2$ (in one-line notation: $231$ and $312$). For
$n = 4$, group by $\sigma(1)$: with $\sigma(1) = 2$ the derangements
are $2143$, $2341$, $2413$; with $\sigma(1) = 3$: $3142$, $3412$,
$3421$; with $\sigma(1) = 4$: $4123$, $4312$, $4321$. Three in each
group: $D_4 = 9$.

\textbf{2.} A permutation with exactly $k$ fixed points is
determined by the choice of its fixed-point set $F$ ($\binom nk$
ways) together with its restriction to the complement, which must
be a permutation of $n - k$ points with \emph{no} fixed point
($D_{n-k}$ ways). The two choices are independent and the
correspondence is bijective: $P_k(n) = \binom nk D_{n-k}$.

\textbf{3.} $P_0(4) = D_4 = 9$; $P_1(4) = \binom41 D_3 = 4 \times 2
= 8$; $P_2(4) = \binom42 D_2 = 6$; $P_3(4) = \binom43 D_1 = 0$
(three fixed points force a fourth); $P_4(4) = 1$. Sum: $9 + 8 + 6
+ 0 + 1 = 24 = 4!$. No match ($9$ cases) beats exactly one match
($8$ cases) --- narrowly.

\textbf{4.} Count the pairs $(\sigma, i)$ with $\sigma(i) = i$. For
fixed $i$, the permutations fixing $i$ are the permutations of the
other $n - 1$ points: $(n-1)!$ of them. Hence the number of pairs
is $n \cdot (n-1)! = n!$, and this number is also
$\sum_\sigma \abs{\mathrm{Fix}(\sigma)}$. Dividing by the number
$n!$ of permutations: the average number of fixed points is
exactly $1$, for every $n \geq 1$.

\textbf{5.} Let $\sigma$ be a derangement of $\intint1{n+1}$ and
$j = \sigma(n+1) \in \intint1n$: $n$ possible values. \emph{Case
$\sigma(j) = n+1$:} the points $j$ and $n+1$ swap, and $\sigma$
restricted to the remaining $n - 1$ points is an arbitrary
derangement of them: $D_{n-1}$ possibilities. \emph{Case $\sigma(j)
\neq n+1$:} let $i_0 = \sigma^{-1}(n+1)$; here $i_0 \neq j$ and
$i_0 \leq n$. Define $\tau$ on $\intint1n$ by $\tau(i) = \sigma(i)$
for $i \neq i_0$ and $\tau(i_0) = j$. Then $\tau$ is a permutation
of $\intint1n$ (the value $n+1$ has been replaced by the missing
value $j$), and it is a derangement: $\tau(i_0) = j \neq i_0$, and
$\tau(i) = \sigma(i) \neq i$ elsewhere. Conversely, from a
derangement $\tau$ of $\intint1n$ and the value $j$, one recovers
$\sigma$ by setting $\sigma(n+1) = j$, $\sigma(\tau^{-1}(j)) = n+1$
and $\sigma = \tau$ elsewhere: a bijection, giving $D_n$
possibilities. Summing over $j$: $D_{n+1} = n(D_n + D_{n-1})$.
Numerically: $D_5 = 4(9 + 2) = 44$, $D_6 = 5(44 + 9) = 265$.

\textbf{6.} From question 5, $D_{n+1} = nD_n + nD_{n-1}$, so
\[
u_{n+1} = D_{n+1} - (n+1)D_n = nD_n + nD_{n-1} - (n+1)D_n
= -(D_n - nD_{n-1}) = -u_n .
\]
Since $u_1 = D_1 - 1 \cdot D_0 = -1$, induction gives $u_n =
(-1)^n$, i.e.\ $D_n = nD_{n-1} + (-1)^n$ for $n \geq 1$.

\textbf{7.} Induction on $n$. Base: $D_0 = 1 = 0!\,s_0$. Step:
assuming $D_{n-1} = (n-1)!\,s_{n-1}$,
\[
D_n = nD_{n-1} + (-1)^n = n!\,s_{n-1} + (-1)^n
= n!\Bigl(s_{n-1} + \frac{(-1)^n}{n!}\Bigr) = n!\,s_n ,
\]
which is the formula. No inclusion--exclusion was used: only the
combinatorial recurrence of question 5.

\textbf{8.} Trinomial revision, by factorials:
\[
\binom nk \binom kj
= \frac{n!}{k!\,(n-k)!} \cdot \frac{k!}{j!\,(k-j)!}
= \frac{n!}{j!\,(n-j)!} \cdot \frac{(n-j)!}{(k-j)!\,(n-k)!}
= \binom nj \binom{n-j}{k-j} .
\]
Now substitute $a_k = \sum_j \binom kj b_j$ and exchange the two
finite sums:
\[
\sum_{k=0}^{n} (-1)^{n-k} \binom nk a_k
= \sum_{j=0}^{n} b_j \binom nj
\sum_{k=j}^{n} (-1)^{n-k} \binom{n-j}{k-j}
= \sum_{j=0}^{n} b_j \binom nj
\sum_{i=0}^{n-j} (-1)^{(n-j)-i} \binom{n-j}{i} .
\]
The inner sum is the expansion of $(1 + (-1))^{n-j} = 0^{n-j}$
(binomial theorem, \cref{thm:b1:counting:binomial}): it vanishes
for $j < n$ and equals $1$ for $j = n$. Only $j = n$ survives, and
the right-hand side is $b_n$, as claimed.

\textbf{9.} By the symmetry $\binom nk = \binom n{n-k}$, the
identity of \cref{exo:b1:counting:11} rewrites as $n! =
\sum_{k=0}^n \binom nk D_k$. Apply question 8 with $a_n = n!$ and
$b_k = D_k$:
\[
D_n = \sum_{k=0}^{n} (-1)^{n-k} \binom nk k!
= \sum_{k=0}^{n} (-1)^{n-k} \frac{n!}{(n-k)!}
= n! \sum_{j=0}^{n} \frac{(-1)^j}{j!} ,
\]
reindexing by $j = n - k$: the formula a third time.

\textbf{10.} $D_n = n!\,s_n$ (question 7), so
\[
\Bigl| D_n - \frac{n!}{\eu} \Bigr|
= n!\,\abs{s_n - \eu^{-1}} < \frac{n!}{(n+1)!} = \frac1{n+1} .
\]

\textbf{11.} For $n \geq 1$, $\frac1{n+1} \leq \frac12$, and the
inequality of question 10 is strict: $D_n$ lies at distance
$< \frac12$ from $n!/\eu$, so it is the unique nearest integer. For
$n = 0$ the bound only gives distance $< 1$, and indeed the claim
fails there: $0!/\eu \approx 0.368$ has nearest integer $0$, while
$D_0 = 1$.

\textbf{12.} $\eu^{-1} - s_n = \sum_{k \geq n+1} (-1)^k/k!$ is an
alternating series with strictly decreasing terms, so its sign is
the sign of its first term $(-1)^{n+1}/(n+1)!$. Hence $s_n -
\eu^{-1}$ has the sign of $(-1)^n$: for $n$ even, $s_n > \eu^{-1}$
and $D_n = n!\,s_n > n!/\eu$; for $n$ odd, $D_n < n!/\eu$.

\textbf{13.} $D_7 = 6(265 + 44) = 6 \times 309 = 1854$; $D_8 =
7(1854 + 265) = 7 \times 2119 = 14\,833$; $D_9 = 8(14\,833 + 1854)
= 8 \times 16\,687 = 133\,496$; $D_{10} = 9(133\,496 + 14\,833) =
9 \times 148\,329 = 1\,334\,961$. Check: $10!/\eu = 3\,628\,800 /
2.718281828 \approx 1\,334\,960.92$, whose nearest integer is
$1\,334\,961$ --- and $D_{10} > 10!/\eu$, as question 12 predicts
for even $n$.

\textbf{14.} $\abs{p_n - \eu^{-1}} = \abs{s_n - \eu^{-1}} <
\frac1{(n+1)!}$. For $n = 6$: $p_6 = 265/720 = 0.36806$ (five
decimals), against $\eu^{-1} = 0.36788$; the gap is below $1/7! =
1/5040 < 2 \times 10^{-4}$. The bound $1/(n+1)!$ collapses so fast
that the probability is fixed to many decimal places already for a
dozen letters: the answer ``about $36.8\%$'' is, for every
practical purpose, independent of $n$ --- the famous surprise of
the problem.

\textbf{15.} By question 2 and $D_m = m!\,s_m$:
\[
\frac{P_k(n)}{n!} = \frac{\binom nk D_{n-k}}{n!}
= \frac{D_{n-k}}{k!\,(n-k)!} = \frac{s_{n-k}}{k!}
\;\longrightarrow\; \frac{\eu^{-1}}{k!}
\]
as $n \to \infty$ with $k$ fixed, since $s_{n-k} \to \eu^{-1}$.
The limit values $\eu^{-1}/k!$ ($k \in \N$) are the weights of the
Poisson distribution of parameter $1$.

\textbf{16.} Count the triples $(\sigma, i, j)$ with $i \neq j$,
$\sigma(i) = i$, $\sigma(j) = j$. Choosing the ordered pair first:
$n(n-1)$ ways; the permutations fixing both $i$ and $j$ are the
permutations of the remaining $n - 2$ points: $(n-2)!$ of them.
Total: $n(n-1)(n-2)! = n!$. Summing instead over $\sigma$ first
counts, for each $\sigma$, the ordered pairs of distinct fixed
points: $\abs{\mathrm{Fix}(\sigma)}(\abs{\mathrm{Fix}(\sigma)}-1)$.
Hence the stated identity; dividing by $n!$, the average of
$\abs{\mathrm{Fix}}(\abs{\mathrm{Fix}} - 1)$ is $1$, so the average
of $\abs{\mathrm{Fix}}^2$ is $1 + 1 = 2$ and the variance is $2 -
1^2 = 1$.

\textbf{17.} The proportions $1 - p_n$: for $n = 4$, $1 - \frac
9{24} = \frac{15}{24} = 0.6250$; for $n = 5$, $1 - \frac{44}{120} =
\frac{76}{120} = 0.6333$; for $n = 6$, $1 - \frac{265}{720} =
\frac{455}{720} = 0.6319$. All within a percent of $1 - \eu^{-1}
\approx 0.6321$, oscillating around it.

\textbf{18.} Directly:
\[
s_{n+2} - s_n = \frac{(-1)^{n+1}}{(n+1)!} +
\frac{(-1)^{n+2}}{(n+2)!}
= (-1)^{n+1}\Bigl(\frac1{(n+1)!} - \frac1{(n+2)!}\Bigr),
\]
and the parenthesis is $> 0$. For $n$ even the difference is
negative: $s_{n+2} < s_n$, so $p_0 > p_2 > p_4 > \dots$; for $n$
odd it is positive: $p_1 < p_3 < p_5 < \dots$ Combined with
question 12 (evens above $\eu^{-1}$, odds below) and question 14
(distance to $\eu^{-1}$ tends to $0$): the two staircases squeeze
$\eu^{-1}$ between them.

\textbf{19.} One complete drawing is a uniform random permutation,
valid when it is a derangement: probability $p_n \approx \eu^{-1}$.
By the quoted fact, the average number of drawings until success is
$1/p_n$, and question 14 gives $1/p_n \approx \eu$ up to an error
that is negligible already for small $n$. So a Secret Santa with
restarts costs on average about $\eu \approx 2.72$ complete
drawings --- whether the office has $6$ people or $600$.

\textbf{20.} Fix $j \geq 2$ and run question 5's classification on
the value $\sigma(1) = j$. If $\sigma(j) = 1$: the remaining
$n - 2$ points carry an arbitrary derangement, $D_{n-2}$ ways. If
$\sigma(j) \neq 1$: redirect the preimage $i_0 = \sigma^{-1}(1)$ to
$j$ exactly as in question 5; this is a bijection with the
derangements of the $n - 1$ points $\{2, \dots, n\}$: $D_{n-1}$
ways. Total $D_{n-1} + D_{n-2}$, the same for every $j$. Summing
over the $n - 1$ values of $j$: $D_n = (n-1)(D_{n-1} + D_{n-2})$,
which displays the factor $n - 1$: $(n-1) \mid D_n$.

\textbf{21.} Claim: $D_n$ is odd iff $n$ is even. Induction using
$D_n = nD_{n-1} + (-1)^n$, i.e.\ $D_n \equiv nD_{n-1} + 1 \pmod 2$.
Base: $D_1 = 0$ is even, $n = 1$ odd: claim holds. If $n$ is even,
$nD_{n-1}$ is even and $D_n \equiv 1$: odd, as claimed. If $n$ is
odd, then $n - 1$ is even, so $D_{n-1}$ is odd by hypothesis, and
$D_n \equiv D_{n-1} + 1 \equiv 0$: even. The induction closes.

\textbf{22.} Reducing $D_n = nD_{n-1} + (-1)^n$ modulo $n$ kills
the first term: $D_n \equiv (-1)^n \pmod n$. For $n = 10$:
$(-1)^{10} = 1$, and indeed $D_{10} = 1\,334\,961$ ends in the
digit $1$.

\textbf{23.} For $n \geq 3$, $D_{n-1} \geq 1$ and division of
question 6's recurrence by $D_{n-1}$ gives $D_n/D_{n-1} = n +
(-1)^n/D_{n-1}$, with $\abs{(-1)^n/D_{n-1}} \leq 1$ and rapidly
tending to $0$. Consistency: if $D_n \approx n!/\eu$, then
$D_n/D_{n-1} \approx n!/(n-1)! = n$ --- the factor $\eu$ cancels
in the ratio, and the recurrence confirms it to accuracy
$1/D_{n-1}$.

\textbf{24.} (i) The product and sum rules underlie every count:
questions 2 and 5 partition sets of permutations into independent
stages. (ii) Double counting gave the mean (question 4) and the
variance (question 16) of the number of fixed points with no
formula for $D_n$ at all. (iii) The binomial theorem evaluated the
alternating inner sum $(1-1)^{n-j}$ that makes binomial inversion
work (question 8). (iv) The alternating-series bound converted the
exact but opaque sum $n!\,s_n$ into the transparent statement
``nearest integer to $n!/\eu$'' (questions 10--14).

\textbf{25.} Inclusion--exclusion
(\cref{ex:b1:counting:derangement} and \cref{exo:b1:counting:11})
is the conceptual proof: it explains the alternating sum as
overcounting corrections, and it generalizes verbatim to counting
elements avoiding any family of ``bad'' sets. The recurrence route
(questions 5--7) computes fastest --- linear time, exact integer
arithmetic, no factorials --- and is the source of the arithmetic
facts of Part V. Binomial inversion (questions 8--9) places the
formula inside a general transform that will reappear wherever
two triangular systems of identities face each other. And the
appearance of $\eu$ is best explained by the formula itself: the
proportion of derangements is the partial sum $s_n$ of the series
for $\eu^{-1}$, so Montmort's envelopes were, three decades before
Euler's notation, already computing the number $\eu$.
\end{solution}

\section*{Chapter \ref{ch:b1:complex} --- Complex Numbers}

\begin{solution}{exo:b1:complex:1}
$1 + \iu = \sqrt 2\, \eu^{\iu\pi/4}$;
$\;\sqrt 3 - \iu = 2\, \eu^{-\iu\pi/6}$;
$\;-5 = 5\, \eu^{\iu\pi}$;
\[
\frac{1 + \iu}{\sqrt 3 - \iu}
= \frac{\sqrt 2}{2}\, \eu^{\iu(\pi/4 + \pi/6)}
= \frac{\sqrt 2}{2}\, \eu^{5\iu\pi/12}.
\]
Computing the same quotient algebraically (multiply by the conjugate):
\[
\frac{(1 + \iu)(\sqrt 3 + \iu)}{4}
= \frac{(\sqrt 3 - 1) + \iu(\sqrt 3 + 1)}{4}.
\]
Identifying with $\frac{\sqrt 2}{2}(\cos\frac{5\pi}{12} +
\iu\sin\frac{5\pi}{12})$:
\[
\cos\frac{5\pi}{12} = \frac{\sqrt 6 - \sqrt 2}{4},
\qquad
\sin\frac{5\pi}{12} = \frac{\sqrt 6 + \sqrt 2}{4}.
\]
\end{solution}

\begin{solution}{exo:b1:complex:2}
$(1+\iu)^2 = 2\iu$, so $(1+\iu)^{20} = (2\iu)^{10} = 2^{10}\,\iu^{10}
= 1024 \times (-1) = -1024$.

In exponential form $(1+\iu)^n = 2^{n/2}\, \eu^{\iu n\pi/4}$, real if
and only if $\sin\frac{n\pi}{4} = 0$, i.e.\ $4 \mid n$. So $(1+\iu)^n
\in \R$ exactly for the multiples of $4$ (value $(-4)^{n/4}$).
\end{solution}

\begin{solution}{exo:b1:complex:3}
$\abs{z - 2} = \abs{z + \iu}$: equidistance from the points $2$ and
$-\iu$ --- the perpendicular bisector of the segment joining $(2, 0)$
and $(0, -1)$.

$\abs{z - 1} = 2$: the circle of center $1$ and radius $2$.

$\frac{z-1}{z+1} \in \iu\R$: by
\cref{prop:b1:complex:geometry}~(3), the points $M(z)$, $A(1)$,
$B(-1)$ satisfy: lines $MA$ and $MB$ perpendicular (or $z = 1$, where
the quotient is $0 \in \iu\R$). The set is the circle of diameter
$[-1, 1]$ (the unit circle), minus the point $-1$ where the quotient
is undefined. \emph{Check by computation:} $z = \eu^{\iu\theta}$ gives
$\frac{z-1}{z+1} = \frac{\eu^{\iu\theta/2}(\eu^{\iu\theta/2} -
\eu^{-\iu\theta/2})}{\eu^{\iu\theta/2}(\eu^{\iu\theta/2} +
\eu^{-\iu\theta/2})} = \iu\tan\frac\theta2 \in \iu\R$.
\end{solution}

\begin{solution}{exo:b1:complex:4}
Linearization:
\[
\sin^4\theta
= \Bigl(\frac{\eu^{\iu\theta} - \eu^{-\iu\theta}}{2\iu}\Bigr)^{\!4}
= \frac{\eu^{4\iu\theta} - 4\eu^{2\iu\theta} + 6 - 4\eu^{-2\iu\theta}
+ \eu^{-4\iu\theta}}{16}
= \frac{\cos 4\theta - 4\cos 2\theta + 3}{8}.
\]
Expansion: by de Moivre, $\cos 4\theta = \Re\bigl((c +
\iu s)^4\bigr) = c^4 - 6c^2 s^2 + s^4$ with $c = \cos\theta$, $s =
\sin\theta$; substituting $s^2 = 1 - c^2$:
\[
\cos 4\theta = c^4 - 6c^2(1 - c^2) + (1 - c^2)^2
= 8c^4 - 8c^2 + 1 .
\]
\end{solution}

\begin{solution}{exo:b1:complex:5}
$8\iu = 8\,\eu^{\iu\pi/2}$, so the cube roots are $2\,\eu^{\iu(\pi/6 +
2k\pi/3)}$, $k = 0, 1, 2$:
\[
z_0 = 2\eu^{\iu\pi/6} = \sqrt 3 + \iu,\quad
z_1 = 2\eu^{5\iu\pi/6} = -\sqrt 3 + \iu,\quad
z_2 = 2\eu^{3\iu\pi/2} = -2\iu :
\]
an equilateral triangle on the circle of radius $2$.

$-4 = 4\eu^{\iu\pi}$, so $z^4 = -4$ has solutions $\sqrt 2\,
\eu^{\iu(\pi/4 + k\pi/2)}$: $\;1 + \iu$, $-1 + \iu$, $-1 - \iu$,
$1 - \iu$. Pairing conjugate roots:
\[
X^4 + 4 = \bigl(X^2 - 2X + 2\bigr)\bigl(X^2 + 2X + 2\bigr),
\]
since $(X - (1+\iu))(X - (1-\iu)) = X^2 - 2X + 2$ and similarly for
the other pair. (Expand to check.)
\end{solution}

\begin{solution}{exo:b1:complex:6}
$C_n + \iu S_n = \sum_{k=0}^{n} \eu^{\iu k\theta} =
\frac{\eu^{\iu(n+1)\theta} - 1}{\eu^{\iu\theta} - 1}$ (geometric sum,
ratio $\eu^{\iu\theta} \neq 1$). Half-angle factorization
(\cref{met:b1:complex:trig}~(4)) on numerator and denominator:
\[
\frac{2\iu\sin\frac{(n+1)\theta}{2}\, \eu^{\iu(n+1)\theta/2}}
{2\iu\sin\frac{\theta}{2}\, \eu^{\iu\theta/2}}
= \frac{\sin\frac{(n+1)\theta}{2}}{\sin\frac{\theta}{2}}\,
\eu^{\iu n\theta/2}.
\]
Taking real and imaginary parts:
\[
C_n = \frac{\sin\frac{(n+1)\theta}{2}}{\sin\frac{\theta}{2}}
\cos\frac{n\theta}{2},
\qquad
S_n = \frac{\sin\frac{(n+1)\theta}{2}}{\sin\frac{\theta}{2}}
\sin\frac{n\theta}{2}.
\]
\end{solution}

\begin{solution}{exo:b1:complex:7}
Discriminant: $\Delta = (3 + 4\iu)^2 - 4(-1 + 5\iu) = 9 + 24\iu - 16 +
4 - 20\iu = -3 + 4\iu$. Square roots of $-3 + 4\iu$: solve $x^2 - y^2
= -3$, $2xy = 4$, $x^2 + y^2 = 5$; then $x^2 = 1$, $y = 2/x$, same
signs: $\delta = \pm(1 + 2\iu)$. Hence
\[
z = \frac{(3 + 4\iu) \pm (1 + 2\iu)}{2}
\in \{\, 2 + 3\iu,\; 1 + \iu \,\}.
\]
\emph{Check:} sum $= 3 + 4\iu$ and product $(2+3\iu)(1+\iu) = -1 +
5\iu$, as required by the coefficients.
\end{solution}

\begin{solution}{exo:b1:complex:8}
\begin{enumerate}
  \item \cref{prop:b1:complex:sumroots} with $n = 5$.
  \item $u + v = \omega + \omega^2 + \omega^3 + \omega^4 = -1$ by (1).
        For the product, expand and reduce exponents modulo $5$:
        \[
        uv = (\omega + \omega^4)(\omega^2 + \omega^3)
        = \omega^3 + \omega^4 + \omega^6 + \omega^7
        = \omega^3 + \omega^4 + \omega + \omega^2 = -1 .
        \]
        So $u$ and $v$ have sum $-1$ and product $-1$: they are the
        two roots of $X^2 + X - 1$.
  \item $u = \omega + \conj\omega = 2\cos\frac{2\pi}{5} > 0$ (the
        angle is acute), and the positive root of $X^2 + X - 1$ is
        $\frac{-1 + \sqrt 5}{2}$. Hence $\cos\frac{2\pi}{5} =
        \frac{\sqrt 5 - 1}{4}$.
\end{enumerate}
\end{solution}

\begin{solution}{exo:b1:complex:9}
Expand both squares as in the proof of
\cref{prop:b1:complex:rules}~(4):
\[
\abs{z + w}^2 = \abs z^2 + \abs w^2 + 2\Re(z\conj w),
\qquad
\abs{z - w}^2 = \abs z^2 + \abs w^2 - 2\Re(z\conj w),
\]
and add. Geometrically, $\abs{z+w}$ and $\abs{z-w}$ are the lengths of
the two diagonals of the parallelogram with vertices $0, z, z+w, w$,
while $\abs z$ and $\abs w$ are the side lengths: the sum of the
squares of the diagonals equals the sum of the squares of the four
sides.
\end{solution}

\begin{solution}{exo:b1:complex:10}
Let $j = \eu^{2\iu\pi/3}$, so $j^2 + j + 1 = 0$ and $\eu^{\iu\pi/3} =
-j^2$, $\eu^{-\iu\pi/3} = -j$. The triangle is equilateral direct iff
$c - a = -j^2 (b - a)$, indirect iff $c - a = -j(b - a)$
(\cref{ex:b1:complex:equilateral}).

Consider $P = a + jb + j^2 c$ and $Q = a + j^2 b + jc$. Using $1 +
j^2 = -j$ and $j^3 = 1$:
\[
c - a + j^2(b - a) = c + j^2 b - (1 + j^2)\,a = c + j^2 b + ja
= j\,(a + jb + j^2 c) = jP,
\]
so the direct case reads $jP = 0 \iff P = 0$; the same computation
with $j$ in place of $j^2$ gives $c - a + j(b - a) = j^2 Q$, so the
indirect case reads $Q = 0$. Hence: equilateral (either orientation)
$\iff PQ = 0$. Expanding, with $j + j^2 = -1$:
\[
PQ = a^2 + b^2 + c^2 + (j + j^2)(ab + bc + ca)
= a^2 + b^2 + c^2 - (ab + bc + ca).
\]
So the triangle is equilateral if and only if $a^2 + b^2 + c^2 = ab +
bc + ca$.
\end{solution}

\begin{solution}{exo:b1:complex:11}
Since the $\omega^k$, $k = 0, \dots, n-1$, are exactly the $n$ roots
of $X^n - 1$ (\cref{thm:b1:complex:roots}), and the polynomial is
monic:
\[
X^n - 1 = \prod_{k=0}^{n-1} (X - \omega^k)
= (X - 1) \prod_{k=1}^{n-1} (X - \omega^k).
\]
Dividing by $X - 1$: $\;1 + X + \dots + X^{n-1} = \prod_{k=1}^{n-1}
(X - \omega^k)$. Evaluating at $X = 1$ gives $P = n$.

Now $1 - \omega^k = -\eu^{\iu k\pi/n}\bigl(\eu^{\iu k\pi/n} -
\eu^{-\iu k\pi/n}\bigr) = -2\iu\,\eu^{\iu k\pi/n} \sin\frac{k\pi}{n}$,
so taking moduli in $P = n$ (each $\sin\frac{k\pi}n > 0$ for $1 \leq k
\leq n-1$):
\[
n = \abs P = \prod_{k=1}^{n-1} 2\sin\frac{k\pi}{n}
= 2^{n-1} \prod_{k=1}^{n-1} \sin\frac{k\pi}{n},
\qquad\text{hence}\qquad
\prod_{k=1}^{n-1} \sin\frac{k\pi}{n} = \frac{n}{2^{n-1}} .
\]
\end{solution}

\section*{Chapter \ref{ch:b1:functions} --- Standard Functions}

\begin{solution}{exo:b1:functions:1}
$\arcsin\frac{\sqrt 3}{2} = \frac\pi3$;
$\;\arccos\bigl(-\frac12\bigr) = \frac{2\pi}{3}$;
$\;\arctan(-1) = -\frac\pi4$.

$\arcsin\bigl(\sin\frac{5\pi}{6}\bigr)$: $\sin\frac{5\pi}{6} =
\frac12$, and the angle of $\intcc{-\frac\pi2}{\frac\pi2}$ with sine
$\frac12$ is $\frac\pi6$ (not $\frac{5\pi}{6}$).

$\arccos\bigl(\cos\frac{7\pi}{4}\bigr)$: $\cos\frac{7\pi}{4} =
\frac{\sqrt 2}{2}$, and the angle of $\intcc{0}{\pi}$ with that cosine
is $\frac\pi4$.
\end{solution}

\begin{solution}{exo:b1:functions:2}
$f(x) = \arcsin(2x - 1)$ requires $-1 \leq 2x - 1 \leq 1$, i.e.\ $x
\in \intcc{0}{1}$. On $\intoo{0}{1}$, the chain rule and
\cref{prop:b1:functions:arcderiv} give
\[
f'(x) = \frac{2}{\sqrt{1 - (2x-1)^2}} = \frac{2}{\sqrt{4x - 4x^2}}
= \frac{1}{\sqrt{x(1 - x)}} .
\]
$g(x) = \arctan\sqrt x$ is defined for $x \geq 0$, and for $x > 0$:
\[
g'(x) = \frac{1}{1 + x} \cdot \frac{1}{2\sqrt x}
= \frac{1}{2\sqrt x\,(1 + x)} .
\]
\end{solution}

\begin{solution}{exo:b1:functions:3}
$\cosh x + \sinh x = \frac{\eu^x + \eu^{-x}}{2} + \frac{\eu^x -
\eu^{-x}}{2} = \eu^x$. Hence for $n \in \Z$:
\[
(\cosh x + \sinh x)^n = (\eu^x)^n = \eu^{nx}
= \cosh nx + \sinh nx .
\]
\end{solution}

\begin{solution}{exo:b1:functions:4}
With $y = \arcsin x$: $\cos 2y = 1 - 2\sin^2 y = 1 - 2x^2$, so
$\cos(2\arcsin x) = 1 - 2x^2$.

With $y = \arctan x$: $\sin 2y = 2 \sin y \cos y = 2 \tan y \cos^2 y =
\frac{2\tan y}{1 + \tan^2 y}$, so $\sin(2\arctan x) =
\dfrac{2x}{1 + x^2}$.
\end{solution}

\begin{solution}{exo:b1:functions:5}
From slowest to fastest:
\[
(\ln x)^{1000} \;\ll\; x^{100} \;\ll\; x^{\ln x} \;\ll\;
\eu^{\sqrt x} \;\ll\; \eu^{x/100} .
\]
Justifications: $(\ln x)^{1000}/x^{100} \to 0$ by
\cref{prop:b1:functions:powerrules} (logs lose to powers).
$x^{100} = \eu^{100\ln x}$ and $x^{\ln x} = \eu^{(\ln x)^2}$: since
$(\ln x)^2 - 100\ln x \to +\infty$, the second wins. $x^{\ln x} =
\eu^{(\ln x)^2} \ll \eu^{\sqrt x}$ because $(\ln x)^2/\sqrt x \to 0$
(logs lose to the power $x^{1/4}$, squared). Finally $\sqrt x - x/100
\to -\infty$, so $\eu^{\sqrt x} \ll \eu^{x/100}$.
\end{solution}

\begin{solution}{exo:b1:functions:6}
Domain: $x \neq \pm 1$, three intervals. On each,
\[
f'(x) = \frac{\bigl(\frac{2x}{1-x^2}\bigr)'}{1 +
\frac{4x^2}{(1-x^2)^2}}
= \frac{\frac{2(1-x^2) + 4x^2}{(1-x^2)^2}}
{\frac{(1-x^2)^2 + 4x^2}{(1-x^2)^2}}
= \frac{2(1 + x^2)}{(1 + x^2)^2} = \frac{2}{1 + x^2}
= (2\arctan x)' .
\]
So $f(x) - 2\arctan x$ is constant on each interval. Values: at $x =
0$, $f(0) = 0$: the constant is $0$ on $\intoo{-1}{1}$. As $x \to
+\infty$, $\frac{2x}{1-x^2} \to 0^-$ so $f \to 0$, while $2\arctan x
\to \pi$: the constant is $-\pi$ on $\intoo{1}{+\infty}$. By oddness,
it is $+\pi$ on $\intoo{-\infty}{-1}$. Summary: $f = 2\arctan x$ on
$\intoo{-1}{1}$, $f = 2\arctan x - \pi$ for $x > 1$, $f = 2\arctan x +
\pi$ for $x < -1$. (This is the double-angle formula for the tangent,
read through $\arctan$.)
\end{solution}

\begin{solution}{exo:b1:functions:7}
$\cosh x = 2$: with $u = \eu^x > 0$, $u + u^{-1} = 4$, so $u^2 - 4u +
1 = 0$, $u = 2 \pm \sqrt 3$: $x = \ln(2 + \sqrt 3)$ or $x = \ln(2 -
\sqrt 3) = -\ln(2 + \sqrt 3)$ (the two solutions are opposite, as
$\cosh$ is even; both are valid). Equivalently $x = \pm
\operatorname{arcosh} 2$.

$5\cosh x - 4\sinh x = 3$: substituting the exponential definitions,
$\frac{5(u + u^{-1}) - 4(u - u^{-1})}{2} = 3$, i.e.\ $u + 9u^{-1} =
6$, i.e.\ $u^2 - 6u + 9 = (u - 3)^2 = 0$: $u = 3$, one solution $x =
\ln 3$.
\end{solution}

\begin{solution}{exo:b1:functions:8}
Let $\alpha = \arctan\frac12 + \arctan\frac13$. Addition formula:
\[
\tan\alpha = \frac{\frac12 + \frac13}{1 - \frac12\cdot\frac13}
= \frac{5/6}{5/6} = 1 .
\]
Both arctangents lie in $\intoo{0}{\frac\pi4}$ (their arguments are in
$\intoo{0}{1}$), so $\alpha \in \intoo{0}{\frac\pi2}$; the only angle
there with tangent $1$ is $\frac\pi4$.

Machin: let $\beta = \arctan\frac15$. Double angle twice:
\[
\tan 2\beta = \frac{2/5}{1 - 1/25} = \frac{5}{12},
\qquad
\tan 4\beta = \frac{2 \cdot 5/12}{1 - 25/144} = \frac{120}{119}.
\]
Then
\[
\tan\Bigl(4\beta - \frac\pi4\Bigr)
= \frac{\frac{120}{119} - 1}{1 + \frac{120}{119}}
= \frac{1/119}{239/119} = \frac{1}{239}.
\]
Location: $\beta < \arctan 1 = \frac\pi4$ and in fact $\tan 4\beta =
\frac{120}{119}$ close to $1$ with $4\beta \in \intoo{0}{\frac\pi2}$
(since $\beta < \frac\pi8$, as $\tan\frac\pi8 = \sqrt 2 - 1 > \frac15$);
so $4\beta - \frac\pi4 \in \intoo{-\frac\pi4}{\frac\pi4}$, where
$\arctan$ inverts $\tan$: $4\beta - \frac\pi4 =
\arctan\frac{1}{239}$, which is Machin's formula.
\end{solution}

\begin{solution}{exo:b1:functions:9}
Let $f(x) = x - \sin x$: $f(0) = 0$ and $f'(x) = 1 - \cos x \geq 0$,
so $f \geq 0$ on $\R_+$: $\sin x \leq x$.

Let $g(x) = \sin x - x + \frac{x^3}{6}$: $g(0) = 0$, $g'(x) = \cos x -
1 + \frac{x^2}{2}$, $g'(0) = 0$, $g''(x) = -\sin x + x = f(x) \geq 0$
on $\R_+$. So $g'$ is increasing with $g'(0) = 0$, hence $g' \geq 0$,
hence $g$ is increasing with $g(0) = 0$: $g \geq 0$, i.e.\ $x -
\frac{x^3}{6} \leq \sin x$ on $\R_+$.

Consequently $0 \leq \frac{x - \sin x}{x^3} \leq \frac 16$ for $x >
0$: the ratio is trapped near $\frac16$'s scale, and
\cref{ch:b1:taylor} shows its limit is exactly $\frac 16$.
\end{solution}

\begin{solution}{exo:b1:functions:10}
Write $y = \arcsin x$, so $2x\sqrt{1 - x^2} = 2\sin y\cos y = \sin
2y$ and $h(x) = \arcsin(\sin 2y) - 2y$.

For $\abs x \leq \frac{\sqrt 2}{2}$: $2y \in
\intcc{-\frac\pi2}{\frac\pi2}$, so $\arcsin(\sin 2y) = 2y$ and $h =
0$.

For $x \in \intoc{\frac{\sqrt 2}{2}}{1}$: $2y \in
\intoc{\frac\pi2}{\pi}$, and the angle of
$\intcc{-\frac\pi2}{\frac\pi2}$ with sine $\sin 2y$ is $\pi - 2y$:
so $h(x) = \pi - 4\arcsin x$. (Check by derivative: there
$h'(x) = \frac{2(1 - 2x^2)}{\abs{1 - 2x^2}\sqrt{1 - x^2}} -
\frac{2}{\sqrt{1-x^2}} = \frac{-4}{\sqrt{1 - x^2}}$, the derivative
of $-4\arcsin x$; and at $x = 1$, $h(1) = \arcsin 0 - 2\cdot\frac\pi2
= -\pi = \pi - 4\cdot\frac\pi2$.)

For $x \in \intco{-1}{-\frac{\sqrt 2}{2}}$, by oddness of $h$:
$h(x) = -\pi - 4\arcsin x$.
\end{solution}

\begin{solution}{exo:b1:functions:11}
$g = \arctan \circ \sinh$ is a composition of odd, strictly
increasing functions, so it is odd and strictly increasing; as $x \to
+\infty$, $\sinh x \to +\infty$ so $g(x) \to \frac\pi2$, and $g$ is a
bijection onto $\intoo{-\frac\pi2}{\frac\pi2}$ (continuity plus the
intermediate value theorem). Chain rule with
\cref{prop:b1:functions:hyprules}~(1):
\[
g'(x) = \frac{\cosh x}{1 + \sinh^2 x} = \frac{\cosh x}{\cosh^2 x}
= \frac{1}{\cosh x} .
\]
By construction $\tan g(x) = \sinh x$. Then, since $g(x) \in
\intoo{-\frac\pi2}{\frac\pi2}$ has positive cosine,
\begin{align*}
\cos g(x) &= \frac{1}{\sqrt{1 + \tan^2 g(x)}}
= \frac{1}{\sqrt{1 + \sinh^2 x}} = \frac{1}{\cosh x},\\
\sin g(x) &= \tan g(x) \cos g(x) = \frac{\sinh x}{\cosh x} = \tanh x .
\end{align*}
\end{solution}

\begin{solution}{exo:b1:functions:12}
Set $u = \arctan 2$ and $v = \arctan 3$; both lie in
$\intoo{\frac\pi4}{\frac\pi2}$ (their arguments exceed $1$), so $u
+ v \in \intoo{\frac\pi2}{\pi}$. The addition formula gives
\[
\tan(u + v) = \frac{2 + 3}{1 - 6} = -1 ,
\]
and the unique angle of $\intoo{\frac\pi2}{\pi}$ with tangent $-1$
is $\frac{3\pi}4$: so $\arctan 2 + \arctan 3 = \frac{3\pi}4$.
(Blindly applying $\arctan$ to the tangent would give
$-\frac\pi4$, off by $\pi$ --- the localization is what saves the
computation, and the negative denominator $1 - 6$ is precisely the
signal that the sum has left
$\intoo{-\frac\pi2}{\frac\pi2}$.) Adding $\arctan 1 = \frac\pi4$:
\[
\arctan 1 + \arctan 2 + \arctan 3 = \frac\pi4 + \frac{3\pi}4 = \pi .
\]
\end{solution}

\begin{solution}{pb:b1:functions:1}
\textbf{1.} $f$ is a quotient of differentiable functions with
nonvanishing denominator on $\intoo0{+\infty}$, and
\[
f'(t) = \frac{\frac1t \cdot t - \ln t}{t^2} = \frac{1 - \ln
t}{t^2} ,
\]
positive for $t < \eu$, zero at $t = \eu$, negative for $t > \eu$:
$f$ increases strictly on $\intoc0\eu$, decreases strictly on
$\intco\eu{+\infty}$, with maximum $f(\eu) = \frac1\eu$.

\textbf{2.} As $t \to 0^+$: $\ln t \to -\infty$ and $\frac1t \to
+\infty$, so $f(t) \to -\infty$. As $t \to +\infty$: $f(t) \to 0$
by the growth comparison of \cref{prop:b1:functions:powerrules}
($\beta = \alpha = 1$). Sign: that of $\ln t$, so $f < 0$ on
$\intoo01$, $f(1) = 0$, $f > 0$ on $\intoo1{+\infty}$. The graph
climbs from $-\infty$, crosses zero at $1$, peaks at $(\eu,
\frac1\eu)$, then decays to $0^+$.

\textbf{3.} $f(4) = \frac{\ln 4}4 = \frac{2\ln 2}4 = \frac{\ln 2}2
= f(2)$.

\textbf{4.} By the variation table and the intermediate value
theorem (used at High School level; formalized in
\cref{ch:b1:continuity}). For $c < 0$: solutions exist only where
$f < 0$, i.e.\ in $\intoo01$ where $f$ is a strictly increasing
bijection onto $\intoo{-\infty}0$: exactly one. For $c = 0$: only
$t = 1$. For $0 < c < \frac1\eu$: on $\intoo1\eu$, $f$ increases
from $0$ to $\frac1\eu$: one solution; on $\intoo\eu{+\infty}$,
$f$ decreases from $\frac1\eu$ to $0$: one more; total two. For
$c = \frac1\eu$: only the maximum point $t = \eu$. For $c >
\frac1\eu$: none.

\textbf{5.} For $x \neq \eu$, strict maximality gives $f(x) <
f(\eu)$, i.e.\ $\frac{\ln x}x < \frac1\eu$, i.e.\ $\eu\ln x < x$,
i.e.\ $\ln(x^\eu) < \ln(\eu^x)$: $x^\eu < \eu^x$. With $x = \pi$:
$\eu^\pi > \pi^\eu$ (numerically $23.14 > 22.46$).

\textbf{6.} For $x, y > 0$, both sides are positive, so
\[
x^y = y^x \iff y\ln x = x\ln y \iff \frac{\ln x}x = \frac{\ln y}y
\iff f(x) = f(y) ,
\]
dividing by $xy > 0$.

\textbf{7.} Let $f(x) = f(y) = c$ with $x \neq y$. If $c \leq 0$ or
$c = \frac1\eu$, question 4 says the equation $f = c$ has a single
solution: impossible. So $0 < c < \frac1\eu$, and again by
question 4 the two solutions are one point of $\intoo1\eu$ and one
of $\intoo\eu{+\infty}$: both coordinates exceed $1$ and they
straddle $\eu$.

\textbf{8.} Substituting $y = tx$ in $y\ln x = x\ln y$:
\[
tx\ln x = x(\ln t + \ln x)
\iff (t - 1)\ln x = \ln t
\iff \ln x = \frac{\ln t}{t - 1} ,
\]
so $x = t^{1/(t-1)}$ and $y = tx = t^{1 + \frac1{t-1}} =
t^{t/(t-1)}$. Conversely, for these values, $\ln y = \frac{t\ln
t}{t-1} = t\ln x$ and $y = tx$, so $\frac{\ln y}y = \frac{t\ln
x}{tx} = \frac{\ln x}x$: a solution. Every nontrivial solution has
some ratio $t = y/x \in \intoo0{+\infty} \setminus \{1\}$, so the
parametrization is complete.

\textbf{9.} $t = 2$: $x = 2^{1/1} = 2$, $y = 2^{2/1} = 4$. And
\[
x(1/t) = (1/t)^{\frac1{\frac1t - 1}}
= (t^{-1})^{\frac{t}{1 - t}}
= t^{\frac{t}{t - 1}} = y(t) ,
\]
then $y(1/t) = \frac1t\,x(1/t) = \frac{y(t)}t = x(t)$: the
parameter change $t \mapsto 1/t$ swaps the coordinates, as the
symmetry of the equation demands.

\textbf{10.} As $t \to 1$: $\frac{\ln t}{t-1} \to 1$ (it is the
difference quotient of $\ln$ at $1$), so $x \to \eu^1 = \eu$ and
$y = tx \to \eu$. As $t \to +\infty$: $\ln x = \frac{\ln t}{t-1}
\to 0$ so $x \to 1$, while $\ln y = \frac{t\ln t}{t-1} \to
+\infty$ so $y \to +\infty$. As $t \to 0^+$: $\ln x = \frac{\ln
t}{t-1} \to \frac{-\infty}{-1} = +\infty$ so $x \to +\infty$,
while $\ln y = \frac{t\ln t}{t-1} \to \frac{0}{-1} = 0$ so $y \to
1$ (using $t\ln t \to 0$, \cref{ex:b1:functions:xx}). The branch
therefore runs from the asymptote $y = 1$ (far right), up through
$(\eu, \eu)$ on the diagonal, and off along the asymptote $x = 1$
(far top) --- symmetric about the diagonal by question 9.

\textbf{11.} $\ln x(t) = g(t) = \frac{\ln t}{t-1}$, and
\[
g'(t) = \frac{\frac{t-1}t - \ln t}{(t-1)^2}
= \frac{h(t)}{(t-1)^2},
\qquad h(t) = 1 - \frac1t - \ln t .
\]
$h(1) = 0$ and $h'(t) = \frac1{t^2} - \frac1t = \frac{1 -
t}{t^2} < 0$ for $t > 1$: $h < 0$ on $\intoo1{+\infty}$, so $g' <
0$ and $x = \eu^g$ is strictly decreasing there (from $\eu$ to
$1$, by question 10).

\textbf{12.} By question 11, $t \mapsto x(t)$ is a strictly
decreasing bijection from $\intoo1{+\infty}$ onto $\intoo1\eu$;
write $t(x)$ for its inverse (also strictly decreasing) and set
$\varphi(x) = y(t(x))$. Note $h < 0$ also on $\intoo01$ (there
$h' > 0$ and $h(1) = 0$), so $x(\cdot)$ is decreasing on
$\intoo01$ as well; hence $y(t) = x(1/t)$ (question 9) is
\emph{increasing} in $t$ on $\intoo1{+\infty}$, from $\eu$ to
$+\infty$. Composing: $\varphi = y \circ t(\cdot)$ is strictly
decreasing from $\intoo1\eu$ onto $\intoo\eu{+\infty}$, and
$x^{\varphi(x)} = \varphi(x)^x$ by question 8. Finally, for a
common value $c = f(x) \in \intoo0{\frac1\eu}$, the pair $\{x,
\varphi(x)\}$ is \emph{the} two-element solution set of $f = c$
(question 4); extending $\varphi$ to $\intoo\eu{+\infty}$ as the
inverse map, $\varphi(\varphi(x))$ returns to the other (i.e.\
original) element: $\varphi \circ \varphi = \mathrm{id}$.

\textbf{13.} If $(x, y)$ is a nontrivial integer solution with $x
< y$, question 7 puts $x \in \intoo1\eu$: the only integer there
is $x = 2$. Then $f(y) = f(2) = \frac{\ln2}2$ with $y > \eu$; by
question 4 the equation $f = \frac{\ln2}2$ has exactly one
solution beyond $\eu$, and question 3 exhibits it: $y = 4$. Hence
$(2, 4)$, and its swap, are the only ones.

\textbf{14.} $t = 1 + \frac1n = \frac{n+1}n$ gives $\frac1{t - 1}
= n$ and $\frac t{t-1} = n + 1$, so
\[
x_n = \Bigl(\frac{n+1}n\Bigr)^{n},
\qquad
y_n = \Bigl(\frac{n+1}n\Bigr)^{n+1},
\]
manifestly rational, distinct ($y_n = t\,x_n \neq x_n$), and $n =
1$ gives $x_1 = 2$, $y_1 = 4$.

\textbf{15.} $t = y/x$ is rational and $> 1$; write $t - 1 =
\frac rs$ in lowest terms, so $t = \frac{s + r}s$ with $\gcd(s +
r, s) = \gcd(r, s) = 1$. Question 8 gives $x = t^{s/r}$, hence
\[
x^r = t^s = \frac{(s+r)^s}{s^s} ,
\]
a fraction in lowest terms (no prime divides both $s + r$ and
$s$). Writing $x = \frac pq$ in lowest terms, $x^r = \frac{p^r}
{q^r}$ is also in lowest terms, and by uniqueness of the reduced
representation: $p^r = (s+r)^s$ and $q^r = s^s$. Compare the
exponent of any prime $\ell$ in $q^r = s^s$: $r \cdot
v_\ell(q) = s \cdot v_\ell(s)$, so $r$ divides $s\,v_\ell(s)$;
since $\gcd(r, s) = 1$, $r$ divides $v_\ell(s)$ for every prime
$\ell$ (unique factorization, admitted; proved in
\cref{ch:b1:arith}), so $s = b^r$ for an integer $b$. The same
argument on $p^r = (s+r)^s$ gives $s + r = a^r$.

\textbf{16.} Suppose $a^r - b^r = r$ with integers $a > b \geq 1$
and $r \geq 2$. Then $a \geq b + 1$, so by the binomial theorem
\[
a^r - b^r \geq (b+1)^r - b^r
= \sum_{k=0}^{r-1}\binom rk b^k
\geq \binom r0 + \binom r{r-1} b^{r-1}
\geq 1 + r > r ,
\]
a contradiction. From question 15, $a^r - b^r = (s + r) - s = r$
forces $r = 1$: $t = 1 + \frac1s$, and the solution is $(x_s,
y_s)$ of question 14. Together with the swapped pairs, the
classification is complete.

\textbf{17.} $f\bigl(\tfrac94\bigr) = \frac{\ln 2.25}{2.25} =
\frac{0.81093}{2.25} = 0.3604$ and $f\bigl(\tfrac{27}8\bigr) =
\frac{\ln 3.375}{3.375} = \frac{1.21640}{3.375} = 0.3604$ (four
decimals): equal, as the construction promises --- so
$\bigl(\tfrac94\bigr)^{27/8} = \bigl(\tfrac{27}8\bigr)^{9/4}$.

\textbf{18.} The parameters $t_n = 1 + \frac1n$ decrease strictly
to $1$. Since $x(\cdot)$ is strictly decreasing on
$\intoo1{+\infty}$ (question 11), $x_n = x(t_n)$ is strictly
\emph{increasing}; since $y(\cdot)$ is strictly increasing there
(question 12), $y_n = y(t_n)$ is strictly \emph{decreasing}. By
question 7, $x_n \in \intoo1\eu$ and $y_n \in
\intoo\eu{+\infty}$: so $x_n < \eu < y_n$ for every $n$. Finally
$t_n \to 1$ and question 10 give $x_n \to \eu$ and $y_n \to \eu$.
The classical monotone convergence of $\bigl(1 +
\frac1n\bigr)^n$ falls out of the geometry of the branch, with no
new inequality.

\textbf{19.} If $\eu \leq a < b$: $f$ strictly decreasing on
$\intco\eu{+\infty}$ gives $f(a) > f(b)$, i.e.\ $b\ln a > a\ln
b$, i.e.\ $a^b > b^a$. If $1 < a < b \leq \eu$: $f$ strictly
increasing gives $f(a) < f(b)$, hence $a^b < b^a$. Mixed case:
$2 < \eu < 3$ and $2^3 = 8 < 9 = 3^2$; but $2 < \eu < 5$ and
$2^5 = 32 > 25 = 5^2$: both outcomes occur, decided by which side
of the branch the point $(a, b)$ falls.

\textbf{20.} The branch meets the diagonal at $\eu$ and $\varphi$
extends continuously there with $\varphi(\eu) = \eu$.
Differentiating $\varphi(\varphi(x)) = x$ at $x = \eu$ by the
chain rule: $\varphi'(\varphi(\eu))\,\varphi'(\eu) =
\varphi'(\eu)^2 = 1$, so $\varphi'(\eu) = \pm1$; $\varphi$ is
decreasing, so $\varphi'(\eu) = -1$. The branch crosses the
diagonal with slope $-1$: perpendicularly.

\textbf{21.} $y = 3x$ is the case $t = 3$: $x = 3^{1/2} = \sqrt3$
and $y = 3^{3/2} = 3\sqrt3$. Check: $f(\sqrt3) =
\frac{0.54931}{1.73205} = 0.3171$ and $f(3\sqrt3) =
\frac{\ln 5.19615}{5.19615} = \frac{1.64792}{5.19615} = 0.3171$:
equal to four decimals, so $(\sqrt3)^{3\sqrt3} =
(3\sqrt3)^{\sqrt3}$.

\textbf{22.} $n^{1/n} = \eu^{f(n)}$, and $\exp$ is increasing, so
the ordering is that of $f(n)$: $f(3) = \frac{\ln3}3 \approx
0.3662$ beats $f(2) = f(4) = \frac{\ln2}2 \approx 0.3466$ (question
3) and $f(5) \approx 0.3219$. The largest is $\sqrt[3]3$, and the
equal pair is $\sqrt2 = \sqrt[4]4$ --- the integer pair $(2, 4)$
in yet another disguise.

\textbf{23.} The solution set is the union of the diagonal $\{(x,
x) : x > 0\}$ and a single branch, symmetric about the diagonal,
strictly decreasing from the asymptote $x = 1$ (as $y \to
+\infty$) to the asymptote $y = 1$ (as $x \to +\infty$), crossing
the diagonal exactly once, at $(\eu, \eu)$, with slope $-1$
there. On the branch sit the integer points $(2, 4)$ and $(4, 2)$
--- the only ones --- and the rational points $(x_n, y_n)$,
which march monotonically along the branch toward $(\eu, \eu)$
without ever reaching it ($\eu$ is irrational; the branch's
rational points accumulate at an irrational corner).

\textbf{24.} (i) Growth comparisons gave the limits of $f$ at
$+\infty$ (question 2) and $t\ln t \to 0$ (question 10),
shaping both the variation table and the asymptotes. (ii) The
intermediate value theorem, through the bijection statements,
converted the variation table into exact solution counts
(question 4) and into the existence of the inverse map $t(x)$
(question 12). (iii) Unique factorization powered the two
lowest-terms identifications of question 15, the arithmetic heart
of the rational-point classification.

\textbf{25.} One derivative computation --- the sign of $1 - \ln
t$ --- generated everything: the solution count for each level
$c$, the shape and asymptotes of the branch, the extremal
inequality $\eu^x \geq x^\eu$, the classification of integer and
rational solutions, and the monotone convergence of $\bigl(1 +
\frac1n\bigr)^n$. This is the economy of thinking with variation
tables: a one-dimensional study resolves a two-variable equation
because the equation factors through a single function.
\cref{ch:b1:seq} will redo the convergence of $(x_n)$ with the
general theory of monotone sequences; \cref{ch:b1:derivative}
founds variation tables rigorously on the mean value theorem; and
\cref{ch:b1:taylor} supplies what the table cannot --- the
\emph{speed} of the convergence $x_n \to \eu$ (it is of order
$1/n$).
\end{solution}

\section*{Chapter \ref{ch:b1:diffeq} --- Linear Differential Equations}

\begin{solution}{exo:b1:diffeq:1}
Homogeneous: $y_h = \lambda\,\eu^{-2x}$. Particular: try $y_p =
c\,\eu^{3x}$ ($3$ is not a root of $r + 2$): $3c + 2c = 1$, $c =
\frac15$. General solution: $y = \frac{\eu^{3x}}{5} +
\lambda\,\eu^{-2x}$. With $y(0) = 1$: $\frac15 + \lambda = 1$,
$\lambda = \frac45$, so $y = \frac{\eu^{3x} + 4\,\eu^{-2x}}{5}$.
\end{solution}

\begin{solution}{exo:b1:diffeq:2}
On $\intoo{0}{+\infty}$, divide by $x$: $y' - \frac{2}{x}\,y = x^2$.
Here $A(x) = -2\ln x$, $\eu^{-A(x)} = x^2$: homogeneous solutions
$\lambda x^2$. \omterm{thm:b1:diffeq:voc}{Variation of constants}: $\mu'(x) = x^2 \cdot x^{-2} =
1$, so $\mu = x + \lambda$ and
\[
y(x) = x^3 + \lambda x^2, \qquad \lambda \in \R .
\]
\emph{Check:} $x(3x^2 + 2\lambda x) - 2(x^3 + \lambda x^2) = x^3$.
\end{solution}

\begin{solution}{exo:b1:diffeq:3}
$y'' - 3y' + 2y = 0$: roots $1$ and $2$; $\;y = \lambda\,\eu^{x} +
\mu\,\eu^{2x}$.

$y'' + 4y' + 4y = 0$: double root $-2$; $\;y = (\lambda + \mu
x)\,\eu^{-2x}$.

$y'' - 2y' + 5y = 0$: roots $1 \pm 2\iu$; $\;y = \eu^{x}(\lambda\cos
2x + \mu\sin 2x)$.
\end{solution}

\begin{solution}{exo:b1:diffeq:4}
Homogeneous: roots $\pm 1$, $y_h = \lambda\,\eu^x + \mu\,\eu^{-x}$.
Particular with polynomial right-hand side ($\gamma = 0$ not a root):
$y_p = ax^2 + bx + c$; substituting, $2a - (ax^2 + bx + c) = x^2$
gives $a = -1$, $b = 0$, $c = 2a = -2$: $y_p = -x^2 - 2$. General
solution $y = -x^2 - 2 + \lambda\eu^x + \mu\eu^{-x}$.

Cauchy: $y(0) = -2 + \lambda + \mu = 0$ and $y'(0) = \lambda - \mu =
1$: $\lambda = \frac32$, $\mu = \frac12$. So
$y = -x^2 - 2 + \frac{3\eu^x + \eu^{-x}}{2}$.
\end{solution}

\begin{solution}{exo:b1:diffeq:5}
$a(x) = \tan x$, $A(x) = -\ln(\cos x)$ (valid: $\cos > 0$ on the
interval), $\eu^{-A} = \cos x$: homogeneous solutions $\lambda\cos x$.
\omterm{thm:b1:diffeq:voc}{Variation of constants}: $\mu'(x) = \sin 2x \cdot \frac{1}{\cos x} =
2\sin x$, so $\mu = -2\cos x + \lambda$ and
\[
y(x) = -2\cos^2 x + \lambda \cos x .
\]
\emph{Check:} $y' = 4\cos x \sin x - \lambda\sin x$ and $y\tan x =
-2\cos x\sin x + \lambda \sin x$; their sum is $2\cos x\sin x =
\sin 2x$, as required.
\end{solution}

\begin{solution}{exo:b1:diffeq:6}
$\chi(r) = r^2 - 4r + 3 = (r-1)(r-3)$: $\gamma = 1$ is a simple root
($m = 1$). Try $y_p = x(ax + b)\,\eu^x$. With $u = ax^2 + bx$,
\[
y_p'' - 4y_p' + 3y_p = \bigl(u'' + (2 - 4)u' + \chi(1) u\bigr)\eu^x
= \bigl(2a - 2(2ax + b)\bigr)\eu^x .
\]
Identify with $(2x + 1)\eu^x$: $-4a = 2$ and $2a - 2b = 1$, so $a =
-\frac12$, $b = -1$. General solution:
\[
y = -\Bigl(\frac{x^2}{2} + x\Bigr)\eu^{x} + \lambda\,\eu^{x} +
\mu\,\eu^{3x}, \qquad (\lambda, \mu) \in \R^2 .
\]
\end{solution}

\begin{solution}{exo:b1:diffeq:7}
Homogeneous: $y_h = \lambda\cos 2x + \mu\sin 2x$.

Right-hand side $x$ ($\gamma = 0$ not a root): $y_1 = ax + b$ with
$4(ax + b) = x$: $y_1 = \frac x4$.

Right-hand side $\sin 2x = \Im(\eu^{2\iu x})$, $2\iu$ simple root of
$r^2 + 4$: try $z = c\,x\,\eu^{2\iu x}$; then $z'' + 4z =
4\iu c\,\eu^{2\iu x}$, equal to $\eu^{2\iu x}$ for $c =
\frac{1}{4\iu} = -\frac{\iu}{4}$. So $z = -\frac{\iu x}{4}(\cos 2x +
\iu \sin 2x)$ and $y_2 = \Im(z) = -\frac{x\cos 2x}{4}$.

By superposition:
\[
y = \frac{x}{4} - \frac{x\cos 2x}{4} + \lambda\cos 2x + \mu\sin 2x .
\]
\end{solution}

\begin{solution}{exo:b1:diffeq:8}
The equation $T' + kT = 20k$ has constant particular solution $20$ and
homogeneous solutions $\lambda\eu^{-kt}$: $T(t) = 20 +
\lambda\,\eu^{-kt}$, and $T(0) = 80$ gives $\lambda = 60$:
\[
T(t) = 20 + 60\,\eu^{-kt} .
\]
$T(10) = 50$: $\eu^{-10k} = \frac12$, so $k = \frac{\ln 2}{10}$. Then
$T(t) = 25$ requires $\eu^{-kt} = \frac{5}{60} = \frac{1}{12}$, i.e.
\[
t = \frac{\ln 12}{k} = 10\,\frac{\ln 12}{\ln 2} \approx 35.8
\text{ minutes.}
\]
\end{solution}

\begin{solution}{exo:b1:diffeq:9}
\begin{enumerate}
  \item $\chi(r) = r^2 + 2\varepsilon r + 1$, $\Delta = 4(\varepsilon^2
        - 1)$.
        For $\varepsilon \in \intco{0}{1}$: roots $-\varepsilon \pm
        \iu\sqrt{1 - \varepsilon^2}$, so
        $y = \eu^{-\varepsilon t}\bigl(\lambda\cos\omega t +
        \mu\sin\omega t\bigr)$ with $\omega = \sqrt{1 -
        \varepsilon^2}$.
        For $\varepsilon = 1$: double root $-1$, $y = (\lambda + \mu
        t)\,\eu^{-t}$.
        For $\varepsilon > 1$: real roots $r_\pm = -\varepsilon \pm
        \sqrt{\varepsilon^2 - 1}$, both $< 0$, and $y =
        \lambda\eu^{r_+t} + \mu\eu^{r_-t}$.
  \item For $\varepsilon \in \intoo{0}{1}$: $\abs y \leq
        \eu^{-\varepsilon t}(\abs\lambda + \abs\mu) \to 0$. For
        $\varepsilon = 1$: $(\lambda + \mu t)\eu^{-t} \to 0$
        (exponential beats polynomial,
        \cref{prop:b1:functions:powerrules}). For $\varepsilon > 1$:
        both exponentials decay since $r_\pm < 0$ (indeed
        $\sqrt{\varepsilon^2 - 1} < \varepsilon$). For $\varepsilon =
        0$: $y = \lambda\cos t + \mu\sin t$ has constant amplitude
        $\sqrt{\lambda^2 + \mu^2} \neq 0$ unless $y = 0$.
  \item Write $\lambda\cos\omega t + \mu\sin\omega t =
        R\cos(\omega t - \varphi)$ with $R = \sqrt{\lambda^2 + \mu^2}
        > 0$. The zeros of $y$ are those of $\cos(\omega t -
        \varphi)$ (the factor $\eu^{-\varepsilon t}$ never
        vanishes): $\omega t - \varphi \equiv \frac\pi2 \pmod \pi$,
        an arithmetic progression with gap $\frac{\pi}{\omega} =
        \frac{\pi}{\sqrt{1 - \varepsilon^2}}$.
\end{enumerate}
\end{solution}

\begin{solution}{exo:b1:diffeq:10}
Set $x = y = 0$: $2f(0) = 2f(0)^2$, and $f(0) \neq 0$ gives $f(0) =
1$. Fix $x$ and differentiate the equation twice with respect to $y$:
\[
f''(x+y) + f''(x-y) = 2 f(x) f''(y) .
\]
Setting $y = 0$: $\;2f''(x) = 2 f(x) f''(0)$, that is
\[
f''(x) = c\,f(x), \qquad c = f''(0).
\]
\emph{Case $c = \omega^2 > 0$:} $f(x) = \lambda\cosh\omega x +
\mu\sinh\omega x$; $f(0) = 1$ gives $\lambda = 1$. Plugging into the
functional equation and using the addition formulas
(\cref{prop:b1:functions:hyprules}), the equation forces $\mu = 0$
(compare the coefficients of $\sinh\omega x \sinh\omega y$ or
evaluate at $x = y$): $f = \cosh\omega x$, which does satisfy
$\cosh(x+y) + \cosh(x-y) = 2\cosh x\cosh y$.

\emph{Case $c = -\omega^2 < 0$:} similarly $f(x) = \cos\omega x$
($\omega \neq 0$), which satisfies the equation.

\emph{Case $c = 0$:} $f$ affine with $f(0) = 1$: $f(x) = 1 + \mu x$;
the equation forces $\mu = 0$, excluded ($f$ not constant).

Conclusion: the solutions are $f(x) = \cos\omega x$ and $f(x) =
\cosh\omega x$, $\omega > 0$.
\end{solution}

\begin{solution}{exo:b1:diffeq:11}
Set $z(t) = y(\eu^t)$, so that $y(x) = z(\ln x)$ for $x > 0$. Then
\[
y'(x) = \frac{z'(\ln x)}x,
\qquad
y''(x) = \frac{z''(\ln x) - z'(\ln x)}{x^2} ,
\]
and substituting into the equation:
\[
x^2y'' - xy' + y = \bigl(z'' - z'\bigr) - z' + z
= z'' - 2z' + z = 0 .
\]
\omterm{def:b1:diffeq:linear2}{Characteristic polynomial} $(r - 1)^2$: double root $1$, so $z(t) =
(\lambda + \mu t)\,\eu^t$ and, back in the variable $x = \eu^t$:
\[
y(x) = (\lambda + \mu\ln x)\,x,
\qquad \lambda, \mu \in \R .
\]
\end{solution}

\begin{solution}{exo:b1:diffeq:12}
\begin{enumerate}
  \item In normalized form $y' - \frac2x\,y = 0$ on each interval:
        $A(x) = -2\ln\abs x$, so the solutions are $y = a x^2$ on
        $\intoo0{+\infty}$ and $y = b x^2$ on $\intoo{-\infty}0$,
        with independent constants
        (\cref{thm:b1:diffeq:homogeneous1}).
  \item Let $y = ax^2$ for $x \geq 0$ and $bx^2$ for $x < 0$. On
        each open half-line $y$ is differentiable with $xy' = 2y$.
        At $0$: the difference quotients $\frac{y(h) - y(0)}h =
        ah$ or $bh$ tend to $0$, so $y'(0) = 0$ exists, and the
        equation at $x = 0$ reads $0 \cdot y'(0) = 2y(0) = 0$:
        satisfied. So $y$ solves the equation on all of $\R$.
  \item The solutions on $\R$ are exactly these glued functions: a
        \emph{two}-parameter family for a first-order equation.
        There is no contradiction with \cref{thm:b1:diffeq:voc},
        whose hypotheses fail here: written as $y' + a(x)y = 0$,
        the coefficient $a(x) = -\frac2x$ is not continuous at $0$
        --- indeed not defined --- so $\R$ is not an interval on
        which the theorem applies. The singularity at $0$
        disconnects the two half-lines, and the value $y(0) = 0$
        is forced, carrying no information across. Every Cauchy
        datum at $x_0 \neq 0$ determines the solution only on the
        half-line containing $x_0$.
\end{enumerate}
\end{solution}

\begin{solution}{pb:b1:diffeq:1}
\textbf{1.} $\chi(r) = r^2 + 2\lambda r + \omega_0^2$, $\Delta =
4(\lambda^2 - \omega_0^2) < 0$ for $0 < \lambda < \omega_0$: roots
$-\lambda \pm \iu\omega_d$, so by
\cref{thm:b1:diffeq:homogeneous2}
\[
x(t) = \eu^{-\lambda t}\bigl(\lambda_1\cos\omega_d t +
\mu_1\sin\omega_d t\bigr), \qquad (\lambda_1, \mu_1) \in \R^2 .
\]
For $\lambda = 0$: $x = \lambda_1\cos\omega_0 t +
\mu_1\sin\omega_0 t$. The critical ($\lambda = \omega_0$) and
overdamped ($\lambda > \omega_0$) regimes are those of
\cref{exo:b1:diffeq:9} (after rescaling time): $(\lambda_1 +
\mu_1 t)\eu^{-\lambda t}$, resp.\ combinations of $\eu^{r_\pm t}$
with $r_\pm = -\lambda \pm \sqrt{\lambda^2 - \omega_0^2}$.

\textbf{2.} Underdamped: $\abs x \leq \eu^{-\lambda
t}(\abs{\lambda_1} + \abs{\mu_1}) \to 0$. Critical: $(\lambda_1 +
\mu_1 t)\eu^{-\lambda t} \to 0$ since exponentials beat
polynomials (\cref{prop:b1:functions:powerrules}). Overdamped:
$r_- < r_+ = -\lambda + \sqrt{\lambda^2 - \omega_0^2} < 0$
because $\sqrt{\lambda^2 - \omega_0^2} < \lambda$; both
exponentials decay.

\textbf{3.} Along a solution of $(H)$, using $x'' = -2\lambda x' -
\omega_0^2 x$:
\[
\mathcal E'(t) = x'x'' + \omega_0^2 x x'
= x'\bigl(-2\lambda x' - \omega_0^2 x\bigr) + \omega_0^2 xx'
= -2\lambda\,x'^2 \leq 0 .
\]
If $x(t_0) = x'(t_0) = 0$ then $\mathcal E(t_0) = 0$; $\mathcal E$
is nonnegative and nonincreasing, so $\mathcal E \equiv 0$ on
$\intco{t_0}{+\infty}$, forcing $x \equiv 0$ there; for $t \leq
t_0$, run the same argument on $\tilde x(t) = x(2t_0 - t)$, which
solves the equation with damping $-\lambda$ but still has
$\tilde{\mathcal E}(t_0) = 0$ and $\tilde{\mathcal E}' =
+2\lambda\tilde x'^2 \geq 0$ with $\tilde{\mathcal E} \geq 0$;
nonnegative, nondecreasing and zero at the right end of
$\intoc{-\infty}{t_0}$ means zero throughout. So $x \equiv 0$ on
$\R$ --- an energy proof of uniqueness, valid for all $\lambda
\geq 0$.

\textbf{4.} $x(t + T_d) = R\,\eu^{-\lambda t}\eu^{-\lambda
T_d}\cos(\omega_d t + 2\pi - \varphi) = \eu^{-\lambda T_d}x(t)$.
The shrink factor per pseudo-period is $\eu^{-\delta}$ with
$\delta = \lambda T_d = \frac{2\pi\lambda}{\omega_d}$. For
$\omega_0 = 1$, $\lambda = 0.1$: $\omega_d = \sqrt{0.99} =
0.99499$, so $\delta = \frac{0.62832}{0.99499} = 0.6315$: each
swing keeps $\eu^{-0.63} \approx 53\%$ of its amplitude.

\textbf{5.} The amplitude factor is $\eu^{-\lambda t}$, which
decays by $\eu$ over $t = \frac1\lambda$. That interval contains
$\frac{1/\lambda}{T_d} = \frac{\omega_d}{2\pi\lambda}$
pseudo-periods. For $\lambda \ll \omega_0$, $\omega_d \approx
\omega_0$ and this is $\approx \frac{\omega_0}{2\pi\lambda} =
\frac Q\pi$. A $Q = 300$ guitar string rings for about a hundred
periods; a $Q = 1$ door damper does not complete one.

\textbf{6.} Substituting $z\,\eu^{\iu\Omega t}$ into the left-hand
side gives $z\,(-\Omega^2 + 2\iu\lambda\Omega +
\omega_0^2)\,\eu^{\iu\Omega t}$, which equals $A\,\eu^{\iu\Omega
t}$ exactly for $z = \frac{A}{\omega_0^2 - \Omega^2 +
2\iu\lambda\Omega}$ (the denominator is nonzero: its imaginary
part is $2\lambda\Omega > 0$). Since the coefficients are real,
the real part $x_p = \Re\bigl(z\eu^{\iu\Omega t}\bigr)$ solves the
equation with right-hand side $\Re\bigl(A\eu^{\iu\Omega t}\bigr) =
A\cos\Omega t$.

\textbf{7.} Write $\omega_0^2 - \Omega^2 + 2\iu\lambda\Omega =
\sqrt{D}\,\eu^{\iu\varphi}$ with $D = (\omega_0^2 - \Omega^2)^2 +
4\lambda^2\Omega^2$ and $\varphi \in \intoo0\pi$ (the imaginary
part $2\lambda\Omega$ is positive), so $\tan\varphi =
\frac{2\lambda\Omega}{\omega_0^2 - \Omega^2}$. Then $z =
\frac{A}{\sqrt D}\eu^{-\iu\varphi}$ and
\[
x_p(t) = \Re\Bigl(\frac A{\sqrt D}\,\eu^{\iu(\Omega t -
\varphi)}\Bigr) = R\cos(\Omega t - \varphi),
\qquad R = \frac A{\sqrt D} .
\]

\textbf{8.} As $\Omega \to 0^+$: $D \to \omega_0^4$, so $R \to
A/\omega_0^2$ and $\tan\varphi \to 0^+$ with $\varphi \in
\intoo0{\frac\pi2}$: $\varphi \to 0$. The mass follows the force
quasi-statically, displaced by force/stiffness. As $\Omega \to
+\infty$: $D \sim \Omega^4$, so $R \sim A/\Omega^2 \to 0$, and
$\varphi \to \pi$ (the complex number $\omega_0^2 - \Omega^2 +
2\iu\lambda\Omega$ goes to the second quadrant with argument $\to
\pi$): the mass barely moves, and in opposition of phase ---
inertia dominates.

\textbf{9.} By \cref{thm:b1:diffeq:structure2}~(1), every solution
is $x = x_p + x_h$ with $x_h$ solving $(H)$; by question 2, $x_h(t)
\to 0$, so $x(t) - x_p(t) \to 0$: all solutions converge to the
same steady state. The initial conditions only shape the transient.

\textbf{10.} If $x$ is a periodic solution, $x - x_p = x_h$ is a
periodic solution of $(H)$ which tends to $0$ at $+\infty$; a
periodic function with limit $0$ is identically $0$ (its values on
one period repeat forever, so every value is a limit of a
subsequence tending to $0$). Hence $x = x_p$.

\textbf{11.} Here $\lambda = 1$, $\omega_0^2 = 2$, $\Omega = 1$,
$A = 1$: $z = \frac1{2 - 1 + 2\iu} = \frac{1 - 2\iu}5$, so
\[
x_p = \Re\Bigl(\frac{(1 - 2\iu)(\cos t + \iu\sin t)}5\Bigr)
= \frac{\cos t + 2\sin t}5 .
\]
Homogeneous: roots of $r^2 + 2r + 2$ are $-1 \pm \iu$: $x_h =
\eu^{-t}(C\cos t + S\sin t)$. Conditions: $x(0) = \frac15 + C = 0$
gives $C = -\frac15$; differentiating, $x'(0) = \frac25 - C + S =
0$ gives $S = C - \frac25 = -\frac35$. Hence
\[
x(t) = \underbrace{\frac{\cos t + 2\sin t}5}_{\text{steady}}
- \underbrace{\eu^{-t}\,\frac{\cos t + 3\sin
t}5}_{\text{transient}} ,
\]
the transient dying like $\eu^{-t}$.

\textbf{12.} Expanding, $g(u) = u^2 - 2(\omega_0^2 - 2\lambda^2)u
+ \omega_0^4 = (u - u_r)^2 + g(u_r)$ with $u_r = \omega_0^2 -
2\lambda^2$ and
\[
g(u_r) = \omega_0^4 - u_r^2 = (\omega_0^2 - u_r)(\omega_0^2 +
u_r) = 2\lambda^2\bigl(2\omega_0^2 - 2\lambda^2\bigr) =
4\lambda^2(\omega_0^2 - \lambda^2) .
\]
If $2\lambda^2 < \omega_0^2$, then $u_r > 0$ is an admissible
squared frequency: $g$ has a strict minimum there, so $R = A/\sqrt
g$ has a strict maximum at $\Omega_r = \sqrt{u_r} =
\sqrt{\omega_0^2 - 2\lambda^2}$, with $R_{\max} = A/\sqrt{g(u_r)}
= \frac{A}{2\lambda\sqrt{\omega_0^2 - \lambda^2}}$.

\textbf{13.} $R(0) = A/\omega_0^2$, so
\[
\frac{R_{\max}}{R(0)}
= \frac{\omega_0^2}{2\lambda\sqrt{\omega_0^2 - \lambda^2}}
= \frac{\omega_0}{2\lambda}\cdot
\frac{\omega_0}{\sqrt{\omega_0^2 - \lambda^2}}
= Q\,\Bigl(1 - \frac{\lambda^2}{\omega_0^2}\Bigr)^{-1/2}
\geq Q .
\]
For weak damping the correction factor is close to $1$: \omterm{ex:b1:diffeq:oscillation}{resonance}
multiplies the static displacement by essentially $Q$.

\textbf{14.} $V(\Omega)^2 = \frac{A^2 u}{g(u)}$ with $u =
\Omega^2$. Its derivative has the sign of $g(u) - u\,g'(u) =
(\omega_0^2 - u)^2 + 4\lambda^2 u - u\bigl(2(u - \omega_0^2) +
4\lambda^2\bigr) = (\omega_0^2 - u)^2 + 2u(\omega_0^2 - u) =
(\omega_0^2 - u)(\omega_0^2 + u)$, positive for $u < \omega_0^2$
and negative beyond: strict maximum exactly at $\Omega =
\omega_0$, for every damping. There $\tan\varphi$ blows up with
$\varphi \in \intoo0\pi$: $\varphi = \frac\pi2$, and $x_p'(t) =
-R\Omega\sin(\Omega t - \frac\pi2) = R\Omega\cos(\Omega t)$ is
exactly in phase with the force: optimal power transfer.

\textbf{15.} $R = R_{\max}/\sqrt2 \iff g(u) = 2g(u_r) \iff (u -
u_r)^2 = g(u_r) = 4\lambda^2(\omega_0^2 - \lambda^2)$, giving
\[
u_\pm = u_r \pm 2\lambda\sqrt{\omega_0^2 - \lambda^2},
\qquad
\Omega_+^2 - \Omega_-^2 = 4\lambda\sqrt{\omega_0^2 - \lambda^2} .
\]
Then $\Omega_+ - \Omega_- = \frac{\Omega_+^2 -
\Omega_-^2}{\Omega_+ + \Omega_-}$, and for $\lambda \ll \omega_0$
both $\Omega_\pm \approx \omega_0$: $\Omega_+ - \Omega_- \approx
\frac{4\lambda\omega_0}{2\omega_0} = 2\lambda$, so
$\frac{\omega_0}{\Omega_+ - \Omega_-} \approx
\frac{\omega_0}{2\lambda} = Q$. Measuring a \omterm{ex:b1:diffeq:oscillation}{resonance} peak's width
measures its \omterm{pb:b1:diffeq:1}{quality factor}.

\textbf{16.} $Q = 10$; $\Omega_r = \sqrt{1 - 2(0.05)^2} =
\sqrt{0.995} = 0.9975$; $R_{\max} = \frac1{2 \times 0.05
\sqrt{1 - 0.0025}} = \frac1{0.1 \times 0.99875} = 10.01$; static
response $R(0) = 1$; bandwidth $\approx 2\lambda = 0.1$. A tall
thin spike of height $\approx Q$ over a plateau of height $1$.

\textbf{17.} If $2\lambda^2 \geq \omega_0^2$ then $u_r \leq 0$ and
$g'(u) = 2(u - u_r) > 0$ for all $u > 0$: $g$ increases strictly
on $\intoo0{+\infty}$, so $R = A/\sqrt g$ decreases strictly from
$R(0) = A/\omega_0^2$: the response is largest at zero frequency
and there is no peak.

\textbf{18.} $\gamma = \iu\Omega$ is not a root of $r^2 +
\omega_0^2$ (as $\Omega \neq \omega_0$), so
\cref{met:b1:diffeq:particular} with $m = 0$ gives $x_p =
\frac{A\cos\Omega t}{\omega_0^2 - \Omega^2}$ (substitute and
check: $-\Omega^2 + \omega_0^2$ times the cosine). General
solution:
\[
x(t) = \frac{A\cos\Omega t}{\omega_0^2 - \Omega^2}
+ \lambda_1\cos\omega_0 t + \mu_1\sin\omega_0 t .
\]

\textbf{19.} $x(0) = 0$ forces $\lambda_1 = -\frac A{\omega_0^2 -
\Omega^2}$, and $x'(0) = 0$ forces $\mu_1 = 0$:
\[
x(t) = \frac{A}{\omega_0^2 - \Omega^2}\,
\bigl(\cos\Omega t - \cos\omega_0 t\bigr)
= \frac{2A}{\omega_0^2 - \Omega^2}\,
\sin\Bigl(\frac{(\omega_0 - \Omega)t}2\Bigr)
\sin\Bigl(\frac{(\omega_0 + \Omega)t}2\Bigr) ,
\]
by the product formula $\cos a - \cos b =
2\sin\frac{b - a}2\sin\frac{b + a}2$ applied with $a = \Omega t$,
$b = \omega_0 t$.

\textbf{20.} For $\Omega$ near $\omega_0$, the second sine
oscillates at the fast frequency $\frac{\omega_0 + \Omega}2
\approx \omega_0$, while the first is a slow envelope of frequency
$\frac{\abs{\omega_0 - \Omega}}2$: the amplitude of the fast
oscillation waxes and wanes with envelope period
$\frac{2\pi}{\abs{\omega_0 - \Omega}}$ (two beats per envelope
period), reaching maxima $\frac{2A}{\abs{\omega_0^2 - \Omega^2}}$.
As $\Omega \to \omega_0$, the beats become both slower (period
$\to \infty$) and taller (amplitude $\to \infty$).

\textbf{21.} Fix $t$. As $\Omega \to \omega_0$:
\[
\frac{2A}{\omega_0^2 - \Omega^2}\sin\Bigl(\frac{(\omega_0 -
\Omega)t}2\Bigr)
= \frac{2A}{\omega_0 + \Omega}\cdot
\frac{\sin\bigl(\frac{(\omega_0 - \Omega)t}2\bigr)}
{\omega_0 - \Omega}
\longrightarrow \frac{2A}{2\omega_0}\cdot\frac t2
= \frac{At}{2\omega_0} ,
\]
while $\sin\bigl(\frac{(\omega_0 + \Omega)t}2\bigr) \to
\sin(\omega_0 t)$: the limit is $x_\infty(t) =
\frac{At\sin\omega_0 t}{2\omega_0}$. Direct check: with $C =
\frac A{2\omega_0}$, $x_\infty = Ct\sin\omega_0 t$ has
$x_\infty'' = 2C\omega_0\cos\omega_0 t - C\omega_0^2 t
\sin\omega_0 t$, so $x_\infty'' + \omega_0^2 x_\infty =
2C\omega_0\cos\omega_0 t = A\cos\omega_0 t$, with $x_\infty(0) =
0$ and $x_\infty'(0) = C\sin 0 + C\omega_0 \cdot 0 \cdot \cos 0 =
0$ --- for $\omega_0 = A = 1$ this is
exactly \cref{ex:b1:diffeq:oscillation}. \omterm{ex:b1:diffeq:oscillation}{Resonance} is the
degeneration of beats: the first swell of the envelope, stretched
to infinite length.

\textbf{22.} Without damping the resonant amplitude grows linearly
and without bound; with damping $\lambda > 0$ the growth saturates
at $R_{\max} \approx Q\,\frac A{\omega_0^2}$. The mechanism:
dissipation removes energy at a rate growing with the amplitude
(question 23), so the build-up stops exactly when the damping
burns energy as fast as the forcing supplies it.

\textbf{23.} With $x_p' = -R\Omega\sin(\Omega t - \varphi)$, over
a period the averages $\langle\cos^2\rangle =
\langle\sin^2\rangle = \frac12$ and $\langle\sin\cos\rangle = 0$
give:
\[
\langle P_{\mathrm{diss}}\rangle
= 2\lambda\,R^2\Omega^2\,\langle\sin^2(\Omega t -
\varphi)\rangle = \lambda R^2\Omega^2 ;
\]
and expanding $\sin(\Omega t - \varphi) = \sin\Omega t\cos\varphi
- \cos\Omega t\sin\varphi$:
\[
\langle P_{\mathrm{in}}\rangle
= -AR\Omega\,\bigl\langle\cos\Omega t\,\sin(\Omega t -
\varphi)\bigr\rangle
= AR\Omega\,\frac{\sin\varphi}2 .
\]
Since $\sin\varphi = \frac{2\lambda\Omega}{\sqrt D} =
\frac{2\lambda\Omega R}A$, this is $\frac{AR\Omega}2 \cdot
\frac{2\lambda\Omega R}A = \lambda R^2\Omega^2$: injected and
dissipated powers balance exactly --- the defining property of a
steady regime.

\textbf{24.} (i) The structure theorem split every solution into
steady state plus transient (questions 9--11) and reduced
uniqueness to the homogeneous problem. (ii) The complex method
turned the search for a particular solution into one division of
complex numbers (question 6), with amplitude and phase read off a
\omterm{def:b1:complex:field}{modulus} and an argument. (iii) The \omterm{ex:b1:diffeq:oscillation}{resonance} curve is a pure
function study --- a quadratic in $u = \Omega^2$, its minimum,
its level \omterm{def:b1:logic:sets}{sets} --- in the style of \cref{ch:b1:functions}
(questions 12--17).

\textbf{25.} Free and damped: decaying pseudo-oscillations,
lifetime $\frac1\lambda$, about $\frac Q\pi$ swings. Forced and
damped: transients die, and a unique sinusoidal steady state
survives at the forcing frequency, with amplitude peaking near
$\omega_0$ (height $\approx Q \times$ static, width $\approx
\frac{\omega_0}Q$) and phase sweeping from $0$ to $\pi$ through
$\frac\pi2$ at $\omega_0$. Free and undamped: perpetual
oscillation. Forced and undamped: beats, degenerating into
linearly growing \omterm{ex:b1:diffeq:oscillation}{resonance} at exact tuning. One dimensionless
number, $Q = \frac{\omega_0}{2\lambda}$, tunes everything ---
peak height, bandwidth, and transient lifetime are three readings
of the same dial. The sequel: rewriting $x'' + 2\lambda x' +
\omega_0^2x$ as a first-order system opens the matrix methods of
\cref{ch:b1:matrices} and the Year 2 volume, and decomposing an
arbitrary periodic forcing into sinusoids (Fourier series, Year 3
volume) makes this problem's single-frequency analysis the
universal building block: solve for each frequency, superpose.
\end{solution}

\section*{Chapter \ref{ch:b1:arith} --- Integer Arithmetic}

\begin{solution}{exo:b1:arith:1}
$1001 = 1 \times 777 + 224$; $777 = 3 \times 224 + 105$; $224 = 2
\times 105 + 14$; $105 = 7 \times 14 + 7$; $14 = 2 \times 7 + 0$. So
$\gcd(1001, 777) = 7$. Backwards:
\[
7 = 105 - 7 \times 14
= 105 - 7(224 - 2\times 105) = 15 \times 105 - 7 \times 224
\]
\[
= 15(777 - 3\times 224) - 7\times 224 = 15 \times 777 - 52 \times 224
= 15 \times 777 - 52(1001 - 777) = 67 \times 777 - 52 \times 1001 .
\]
Check: $67 \times 777 = 52\,059$ and $52 \times 1001 = 52\,052$;
difference $7$. Bézout pair: $(u, v) = (-52, 67)$ for $1001u + 777v =
7$.
\end{solution}

\begin{solution}{exo:b1:arith:2}
Since $10 \equiv 1 \pmod 9$: $10^k \equiv 1$, so $\sum_k d_k 10^k
\equiv \sum_k d_k \pmod 9$. Since $10 \equiv -1 \pmod{11}$: $10^k
\equiv (-1)^k$, so the integer is congruent to the alternating sum
$\sum_k (-1)^k d_k$ mod $11$ (starting from the \emph{units} digit
with sign $+$).

$123\,456\,789$: digit sum $45 \equiv 0 \pmod 9$. Alternating sum from
the units: $9 - 8 + 7 - 6 + 5 - 4 + 3 - 2 + 1 = 5$, so the number is
$\equiv 5 \pmod{11}$.
\end{solution}

\begin{solution}{exo:b1:arith:3}
Euclid: $237 = 2 \times 91 + 55$; $91 = 1 \times 55 + 36$; $55 = 1
\times 36 + 19$; $36 = 1 \times 19 + 17$; $19 = 1 \times 17 + 2$; $17
= 8 \times 2 + 1$. Backwards:
\[
1 = 17 - 8\times 2 = 17 - 8(19 - 17) = 9\times 17 - 8\times 19
= 9(36 - 19) - 8\times 19 = 9\times 36 - 17\times 19
\]
\[
= 9\times 36 - 17(55 - 36) = 26\times 36 - 17\times 55
= 26(91 - 55) - 17\times 55 = 26\times 91 - 43\times 55
\]
\[
= 26\times 91 - 43(237 - 2\times 91) = 112 \times 91 - 43 \times 237.
\]
So $91 \times 112 \equiv 1 \pmod{237}$: the solutions are $x \equiv
112 \pmod{237}$. (Check: $91 \times 112 = 10\,192 = 43 \times 237 +
1$.)
\end{solution}

\begin{solution}{exo:b1:arith:4}
$\gcd(17, 39) = 1$: Euclid gives $39 = 2\times 17 + 5$, $17 = 3\times
5 + 2$, $5 = 2\times 2 + 1$, and backwards
\[
1 = 5 - 2\times 2 = 5 - 2(17 - 3\times 5) = 7\times 5 - 2\times 17
= 7(39 - 2\times 17) - 2\times 17 = 7\times 39 - 16\times 17 .
\]
Particular solution $(x_0, y_0) = (-16, 7)$. General solution of the
homogeneous equation $17x + 39y = 0$: $x = 39k$, $y = -17k$ (since
$17 \mid 39y$ and $\gcd(17,39) = 1$ force $17 \mid y$ ---
Gauss's lemma). Hence
\[
(x, y) = (-16 + 39k,\; 7 - 17k), \qquad k \in \Z .
\]
For the right-hand side $5$, multiply the particular solution by $5$:
$(x, y) = (-80 + 39k,\; 35 - 17k)$, $k \in \Z$.
\end{solution}

\begin{solution}{exo:b1:arith:5}
For every \omterm{def:b1:arith:prime}{prime} $p$, with $\alpha = v_p(a)$ and $\beta = v_p(b)$:
\[
v_p\bigl(\gcd(a,b)\bigr) + v_p\bigl(\operatorname{lcm}(a,b)\bigr)
= \min(\alpha, \beta) + \max(\alpha, \beta)
= \alpha + \beta = v_p(ab) .
\]
Two positive integers with the same valuation at every \omterm{def:b1:arith:prime}{prime} are equal
(\cref{prop:b1:arith:valuation}), so $\gcd(a,b)\operatorname{lcm}(a,b)
= ab$.
\end{solution}

\begin{solution}{exo:b1:arith:6}
Valuations: $\gcd(a, b) = 2^{\min(10,6)} 3^{\min(4,7)} 5^{\min(2,0)}
7^{\min(0,1)} = 2^6\, 3^4 = 5184$;
$\operatorname{lcm}(a,b) = 2^{10}\, 3^7\, 5^2\, 7$.

Divisor count: a positive divisor of $n = \prod p_i^{\alpha_i}$ is
exactly a choice $\prod p_i^{\beta_i}$ with $0 \leq \beta_i \leq
\alpha_i$ (\cref{prop:b1:arith:valuation}); the choices are
independent, so there are $\prod_i (\alpha_i + 1)$ divisors. For $a$:
$(10+1)(4+1)(2+1) = 165$.
\end{solution}

\begin{solution}{exo:b1:arith:7}
Suppose $\sqrt p = \frac rq$ with $r, q \in \N^*$, i.e.\ $p q^2 =
r^2$. Apply $v_p$: $v_p(pq^2) = 1 + 2v_p(q)$ is odd, while $v_p(r^2)
= 2 v_p(r)$ is even. An integer cannot have an odd and an even
$p$-valuation at once: contradiction. So $\sqrt p \notin \Q$.
\end{solution}

\begin{solution}{exo:b1:arith:8}
\emph{General fact.} With $\gcd(m,n) = 1$, Bézout gives $mu + nv = 1$.
Set $x_0 = b\,mu + a\,nv$. Then $x_0 \equiv a\,nv \equiv a(1 - mu)
\equiv a \pmod m$ and similarly $x_0 \equiv b \pmod n$: existence. If
$x$ and $x'$ are two solutions, $m$ and $n$ divide $x - x'$, so
$mn \mid x - x'$ (\cref{thm:b1:arith:gauss}~(2)): uniqueness mod
$mn$.

\emph{Numerically:} $m = 7$, $n = 11$: $7 \times (-3) + 11 \times 2 =
1$. So $x_0 = 5 \times 7 \times (-3) + 2 \times 11 \times 2 = -105 +
44 = -61 \equiv 16 \pmod{77}$. Check: $16 = 2\times 7 + 2 \equiv 2
\pmod 7$; $16 = 11 + 5 \equiv 5 \pmod{11}$. Solutions: $x \equiv 16
\pmod{77}$.
\end{solution}

\begin{solution}{exo:b1:arith:9}
Mod $7$: Fermat gives $3^6 \equiv 1$, and $1000 = 6 \times 166 + 4$,
so $3^{1000} \equiv 3^4 = 81 \equiv 4 \pmod 7$.

Last two digits of $7^{100}$: work mod $4$ and mod $25$. Mod $4$: $7
\equiv -1$, so $7^{100} \equiv 1$. Mod $25$: $7^2 = 49 \equiv -1$, so
$7^4 \equiv 1$ and $7^{100} = (7^4)^{25} \equiv 1$. By the Chinese
remainder theorem (\cref{exo:b1:arith:8}), $7^{100} \equiv 1
\pmod{100}$: the last two digits are $01$.
\end{solution}

\begin{solution}{exo:b1:arith:10}
Write $m = nq + r$, $0 \leq r < n$. Then
\[
2^m - 1 = 2^r\bigl(2^{nq} - 1\bigr) + 2^r - 1,
\]
and $2^n - 1$ \omterm{def:b1:arith:divides}{divides} $2^{nq} - 1 = (2^n - 1)(2^{n(q-1)} + \dots +
1)$. So mod $2^n - 1$, $\;2^m - 1 \equiv 2^r - 1$, and since $0 \leq
2^r - 1 < 2^n - 1$, this \emph{is} the Euclidean remainder.

Therefore the \omterm{met:b1:arith:euclid}{Euclidean algorithm} on the pair $(2^m - 1, 2^n - 1)$
mirrors, exponent by exponent, the algorithm on $(m, n)$: each
division step replaces $(m, n)$ by $(n, r)$ upstairs and $(2^m - 1,
2^n - 1)$ by $(2^n - 1, 2^r - 1)$ downstairs. The algorithm upstairs
terminates at $\gcd(m,n)$, so downstairs it terminates at
$2^{\gcd(m,n)} - 1$.
\end{solution}

\begin{solution}{exo:b1:arith:11}
First solve $x^2 \equiv 1 \pmod p$: $p \mid (x-1)(x+1)$, so by
Euclid's lemma $x \equiv 1$ or $x \equiv -1 \pmod p$.

In the product $(p-1)! = 1 \times 2 \times \dots \times (p-1)$, every
factor $a$ is invertible mod $p$, and its inverse $a^{-1}$ is again
one of the factors (\cref{prop:b1:arith:invmod}). Pair each $a$ with
$a^{-1}$: the pairs multiply to $1$, except that self-paired factors
($a = a^{-1}$, i.e.\ $a^2 \equiv 1$) stand alone --- and these are
exactly $1$ and $p - 1$. Hence
\[
(p-1)! \equiv 1 \times (p - 1) \equiv -1 \pmod p .
\]
(For $p = 2$: $1! = 1 \equiv -1 \pmod 2$; the pairing argument
degenerates but the result holds.)

\emph{Converse.} Let $n \geq 2$ be composite, $n = ab$ with $1 < a
\leq b < n$. If $a < b$, both appear as distinct factors of $(n-1)!$,
so $n \mid (n-1)!$ and $(n-1)! \equiv 0 \not\equiv -1$. If $a = b$
(i.e.\ $n = a^2$): for $a \geq 3$, both $a$ and $2a$ are $< n$, so
$n = a^2 \mid a \times 2a \mid (n-1)!$, same conclusion; for $n = 4$,
$(n-1)! = 6 \equiv 2 \not\equiv -1 \pmod 4$.
\end{solution}

\begin{solution}{exo:b1:arith:12}
\begin{enumerate}
  \item Induction. For $n = 1$: $F_0 = 3 = F_1 - 2 = 5 - 2$.
        Assuming $F_0\cdots F_{n-1} = F_n - 2$:
        \[
        F_0 \cdots F_n = (F_n - 2)F_n
        = \bigl(2^{2^n} - 1\bigr)\bigl(2^{2^n} + 1\bigr)
        = 2^{2^{n+1}} - 1 = F_{n+1} - 2 .
        \]
  \item Let $m < n$ and $d = \gcd(F_m, F_n)$. By (1), $F_m$
        \omterm{def:b1:arith:divides}{divides} $F_n - 2$, so $d$ \omterm{def:b1:arith:divides}{divides} both $F_n$ and $F_n -
        2$, hence \omterm{def:b1:arith:divides}{divides} $2$. But every Fermat number is odd, so
        $d = 1$.
  \item Each $F_n \geq 3$ has a \omterm{def:b1:arith:prime}{prime} divisor $p_n$
        (\cref{thm:b1:arith:euclidprimes}'s first step). If $m
        \neq n$, then $p_m \neq p_n$, since a common \omterm{def:b1:arith:prime}{prime} would
        divide $\gcd(F_m, F_n) = 1$. The \omterm{def:b1:logic:map}{map} $n \mapsto p_n$ is
        therefore \omterm{def:b1:logic:inj}{injective} from $\N$ into the \omterm{def:b1:arith:prime}{primes}: there are
        infinitely many \omterm{def:b1:arith:prime}{primes}.
\end{enumerate}
\end{solution}

\begin{solution}{pb:b1:arith:1}
\textbf{1.} $10! = 3\,628\,800$: two trailing zeros. Valuations
factor by factor: powers of $2$ come from $2, 4 = 2^2, 6, 8 = 2^3,
10$, totaling $v_2(10!) = 1 + 2 + 1 + 3 + 1 = 8$; powers of $5$
from $5$ and $10$: $v_5(10!) = 2$. Trailing zeros $=
\min(v_2, v_5) = 2$, consistent.

\textbf{2.} Write the Euclidean division $\lfloor x\rfloor = nq +
r$, $0 \leq r \leq n - 1$. Then $x = nq + r + \{x\}$ with $0 \leq
r + \{x\} < n$, so $\lfloor x/n \rfloor = q = \bigl\lfloor \lfloor
x \rfloor / n \bigr\rfloor$.

\textbf{3.} Let $\alpha = v_p(a) \leq \beta = v_p(b)$ (swap if
needed) and write $a = p^\alpha a'$, $b = p^\beta b'$ with $p
\nmid a', b'$. Then $a + b = p^\alpha\bigl(a' + p^{\beta -
\alpha}b'\bigr)$, so $v_p(a + b) \geq \alpha = \min$. If $\alpha <
\beta$, the parenthesis is $a' + p^{\beta-\alpha}b' \equiv a'
\not\equiv 0 \pmod p$: the valuation is exactly $\alpha$.

\textbf{4.} The multiples of $m$ in $\intint1n$ are $m, 2m, \dots,
qm$ where $q$ is the largest integer with $qm \leq n$, i.e.\ $q =
\lfloor n/m \rfloor$.

\textbf{5.} By unique factorization, $v_p(n!) = \sum_{j=1}^{n}
v_p(j)$. Count differently: each $j$ contributes $v_p(j) =
\#\{k \geq 1 : p^k \mid j\}$, so
\[
v_p(n!) = \sum_{j=1}^n \#\{k : p^k \mid j\}
= \sum_{k\geq1} \#\{j \leq n : p^k \mid j\}
= \sum_{k\geq1} \Bigl\lfloor \frac n{p^k} \Bigr\rfloor
\]
by question 4 --- Legendre's formula. The sum is finite: terms
with $p^k > n$ vanish.

\textbf{6.} $v_5(1000!) = 200 + 40 + 8 + 1 = 249$ (divisions by
$5, 25, 125, 625$); $v_2(1000!) = 500 + 250 + 125 + 62 + 31 + 15 +
7 + 3 + 1 = 994$. Trailing zeros of $1000!$: each zero consumes
one $2$ and one $5$, so there are $\min(994, 249) = 249$ of them.

\textbf{7.} With $n = \sum_i a_ip^i$, question 2 gives $\lfloor
n/p^k \rfloor = \sum_{i \geq k} a_ip^{i-k}$ (truncate the base-$p$
expansion). Summing over $k \geq 1$ and exchanging the two finite
sums:
\[
v_p(n!) = \sum_{i\geq1} a_i \sum_{k=1}^{i} p^{i-k}
= \sum_{i\geq0} a_i\,\frac{p^i - 1}{p - 1}
= \frac{n - s_p(n)}{p - 1} .
\]

\textbf{8.} For $p = 2$: $v_2(n!) = n - s_2(n)$. Since $n \geq 1$
has $s_2(n) \geq 1$, always $v_2(n!) \leq n - 1 < n$: $2^n \nmid
n!$. And $v_2(n!) = n - 1$ iff $s_2(n) = 1$ iff $n$ is a power of
$2$.

\textbf{9.} $n$ has $\lfloor \log_p n \rfloor + 1$ base-$p$
digits, each at most $p - 1$, so $1 \leq s_p(n) \leq
(p-1)\bigl(\log_p(n) + 1\bigr)$. Substituting in question 7:
\[
\frac n{p-1} - \log_p(n) - 1 \;\leq\; v_p(n!) \;<\; \frac n{p-1},
\]
and dividing by $n$: $\frac{v_p(n!)}n \to \frac1{p-1}$.

\textbf{10.} $Z(n) - Z(n-1) = v_5(n!/(n-1)!) = v_5(n)$: the count
of trailing zeros jumps by $v_5(n)$ at each multiple of $5$ and is
constant in between. $Z(24) = \lfloor24/5\rfloor = 4$ and $Z(25) =
5 + 1 = 6$: at $n = 25$ the count jumps from $4$ straight to $6$
($v_5(25) = 2$), and since $Z$ is nondecreasing with $Z \leq 4$
before and $Z \geq 6$ after, the value $5$ is never attained: no
factorial ends in exactly five zeros.

\textbf{11.} Write $x = \lfloor x\rfloor + \{x\}$: $\lfloor x +
y\rfloor = \lfloor x\rfloor + \lfloor y\rfloor + \lfloor \{x\} +
\{y\}\rfloor$, and $0 \leq \{x\} + \{y\} < 2$ makes the last floor
$0$ or $1$. Then, by Legendre applied three times,
\[
v_p\binom{m+n}m = v_p\bigl((m{+}n)!\bigr) - v_p(m!) - v_p(n!)
= \sum_{k\geq1}\Bigl(
\Bigl\lfloor\frac{m+n}{p^k}\Bigr\rfloor
- \Bigl\lfloor\frac{m}{p^k}\Bigr\rfloor
- \Bigl\lfloor\frac{n}{p^k}\Bigr\rfloor\Bigr),
\]
a finite sum of $0$s and $1$s (apply the first claim to $x =
m/p^k$, $y = n/p^k$).

\textbf{12.} Fix $k \geq 1$ and write $m = p^km_1 + m_0$, $n =
p^kn_1 + n_0$ with $0 \leq m_0, n_0 < p^k$ (Euclidean division:
$m_0$ is the number formed by the $k$ low digits of $m$). Then
\[
\Bigl\lfloor\frac{m+n}{p^k}\Bigr\rfloor
- \Bigl\lfloor\frac m{p^k}\Bigr\rfloor
- \Bigl\lfloor\frac n{p^k}\Bigr\rfloor
= \Bigl\lfloor\frac{m_0 + n_0}{p^k}\Bigr\rfloor ,
\]
which is $1$ if $m_0 + n_0 \geq p^k$ and $0$ otherwise. But $m_0 +
n_0 \geq p^k$ says precisely that adding the $k$ low digits of $m$
and $n$ overflows into position $k$ --- a carry into position $k$
in the schoolbook addition algorithm. Summing over $k$:
$v_p\binom{m+n}m$ is the number of carries in the base-$p$
addition $m + n$. (Kummer, 1852.)

\textbf{13.} Apply Kummer to $m = j$, $n = p^k - j$, sum $p^k =
(1\underbrace{0\cdots0}_{k})_p$. Let $a = v_p(j)$, so the base-$p$
digits of $j$ at positions $0, \dots, a-1$ are $0$ and the digit
at position $a$ is nonzero. The digits of $p^k - j$ below position
$a$ are $0$ too ($p^k - j = p^a(p^{k-a} - j/p^a)$). At position
$a$, the two nonzero digits must sum to $p$ (result digit $0$):
one carry; at each position $a+1, \dots, k-1$, digits plus the
incoming carry sum to $p$ (result digit $0$ again): the carry
propagates. Total: $k - a$ carries, so $v_p\binom{p^k}j = k -
v_p(j)$. For $k = 1$: $v_p\binom pj = 1$ for $0 < j < p$, the
\omterm{def:b1:arith:divides}{divisibility} used in \cref{thm:b1:arith:fermat}.

\textbf{14.} By the digital form (question 7), using $s_2(2n) =
s_2(n)$ (appending a zero digit):
\[
v_2\binom{2n}n = \bigl(2n - s_2(2n)\bigr) - 2\bigl(n -
s_2(n)\bigr) = 2s_2(n) - s_2(2n) = s_2(n) \geq 1 :
\]
$\binom{2n}n$ is always even, and $v_2 = 1$ (i.e.\ $\binom{2n}n
\equiv 2 \pmod 4$) exactly when $s_2(n) = 1$, i.e.\ when $n$ is a
power of $2$.

\textbf{15.} Vandermonde with $m = n = k = p$: $\binom{2p}p =
\sum_{j=0}^p \binom pj\binom p{p-j} = \sum_{j=0}^p \binom pj^2$.
For $0 < j < p$, $p \mid \binom pj$ (question 13), so $\binom
pj^2 \equiv 0 \pmod p$; the end terms give $1 + 1$:
$\binom{2p}p \equiv 2 \pmod p$.

\textbf{16.} Base $3$: $500 = 486 + 9 + 3 + 2$, digits (low to
high) $(2, 1, 1, 0, 0, 2)$, so $s_3(500) = 6$; and $1000 = 729 +
243 + 27 + 1$, digits $(1, 0, 0, 1, 0, 1, 1)$, so $s_3(1000) = 4$.
\emph{Kummer:} add $500 + 500$ in base $3$: position $0$: $2 + 2 =
4$, digit $1$ carry $1$; position $1$: $1 + 1 + 1 = 3$, digit $0$
carry $1$; position $2$: $1 + 1 + 1 = 3$, digit $0$ carry $1$;
position $3$: $0 + 0 + 1 = 1$, no carry; position $4$: $0$;
position $5$: $2 + 2 = 4$, digit $1$ carry $1$; position $6$:
carry lands: digit $1$. Four carries: $v_3\binom{1000}{500} = 4$.
\emph{Legendre:} $v_3(1000!) = \frac{1000 - 4}2 = 498$ and
$v_3(500!) = \frac{500 - 6}2 = 247$, so $v_3\binom{1000}{500} =
498 - 2\times247 = 4$. The two computations agree --- and the
addition digits $(1, 0, 0, 1, 0, 1, 1)$ reproduce $1000$, as they
must.

\textbf{17.} By Kummer ($p = 2$, $m = k$, $n' = n - k$):
$\binom nk$ is odd iff the addition $k + (n - k)$ in base $2$ has
no carry, iff at every position the digits satisfy $k_i + (n -
k)_i = n_i$; in that case $k_i \leq n_i$ for all $i$. Conversely,
if $k_i \leq n_i$ for all $i$, then the number with digits $n_i -
k_i$ is $n - k$ and the addition is carry-free. Same proof in base
$p$: $p \nmid \binom nk$ iff every base-$p$ digit of $k$ is at
most the corresponding digit of $n$.

\textbf{18.} Counting the $k \in \intint0n$ whose digits obey
$k_i \leq n_i$: each digit of $k$ is chosen independently among
$n_i + 1$ values, giving $\prod_i (n_i + 1)$ choices; in base $2$
this is $2^{\#\{i : n_i = 1\}} = 2^{s_2(n)}$. Row $4 =
(100)_2$: $2^1 = 2$ odd entries --- indeed $1, 4, 6, 4, 1$ has odd
entries only at the ends. Row $5 = (101)_2$: $2^2 = 4$ --- indeed
$1, 5, 10, 10, 5, 1$.

\textbf{19.} All interior entries even $\iff$ the row has exactly
$2$ odd entries (the two ends always are odd) $\iff 2^{s_2(n)} =
2 \iff s_2(n) = 1 \iff n$ is a power of $2$.

\textbf{20.} In question 11's sum, the $k$-th term vanishes as
soon as $p^k > m + n$ (all three floors are then equal, indeed the
first is $0$ when $p^k > m+n$; more simply each term is $0$).
Hence at most $\lfloor \log_p(m+n)\rfloor$ terms are nonzero, each
worth $1$: $a = v_p\binom{m+n}m \leq \log_p(m+n)$, i.e.\ $p^a \leq
m + n$.

\textbf{21.} For every \omterm{def:b1:arith:prime}{prime} $p$, $v_p\bigl(\operatorname{lcm}(1,
\dots, 2n)\bigr) = \lfloor\log_p(2n)\rfloor$ (the largest power of
$p$ not exceeding $2n$ appears among $1, \dots, 2n$). Question 20
with $m = n$ gives $v_p\binom{2n}n \leq \lfloor\log_p(2n)\rfloor$
for every $p$: by \cref{prop:b1:arith:valuation}, $\binom{2n}n
\mid \operatorname{lcm}(1, \dots, 2n)$. For the size: the ratio
$\binom{2n}{k+1}/\binom{2n}k = \frac{2n-k}{k+1} \geq 1$ exactly
for $k < n$, so the central entry is the largest of the $2n + 1$
entries of row $2n$, whence $4^n = \sum_k \binom{2n}k \leq
(2n+1)\binom{2n}n$. Combining:
\[
\operatorname{lcm}(1, \dots, 2n) \geq \binom{2n}n \geq
\frac{4^n}{2n + 1} .
\]
If there were few \omterm{def:b1:arith:prime}{primes} below $2n$, the lcm could not be this
large: exponential growth of the lcm is a quantitative trace of
the abundance of \omterm{def:b1:arith:prime}{primes}.

\textbf{22.} $Z(n) = \sum_k\lfloor n/5^k\rfloor \approx \frac
n4$, so aim near $n = 4 \times 2026 = 8104$: $Z(8104) = 1620 +
324 + 64 + 12 + 2 = 2022$. Step up by multiples of $5$: $Z(8110)
= 2024$, $Z(8115) = 2025$, and
\[
Z(8120) = 1624 + 324 + 64 + 12 + 2 = 2026 .
\]
Since $Z$ is constant between multiples of $5$ and $Z(8119) =
Z(8115) = 2025$, the smallest $n$ with at least $2026$ trailing
zeros is $n = 8120$.

\textbf{23.} Base $7$: $50 = 49 + 1$, digits (low to high) $(1, 0,
1)$. Adding $50 + 50$: position $0$: $1 + 1 = 2 < 7$, no carry;
position $1$: $0 + 0 = 0$; position $2$: $1 + 1 = 2 < 7$, no
carry. Carry-free, so by Kummer $v_7\binom{100}{50} = 0$: $7 \nmid
\binom{100}{50}$. Legendre agrees: $v_7(100!) = \lfloor 100/7
\rfloor + \lfloor 100/49 \rfloor = 14 + 2 = 16$ and $v_7(50!) = 7
+ 1 = 8$, so $v_7\binom{100}{50} = 16 - 2\times8 = 0$.

\textbf{24.} (i) Unique factorization underlies the very
definition of $v_p$ and its additivity, hence Legendre's formula
and every \omterm{def:b1:arith:divides}{divisibility} conclusion
(\cref{prop:b1:arith:valuation}). (ii) Euclidean division produced
the truncation identity of question 2 and the split $m = p^km_1 +
m_0$ that isolates the carry (question 12). (iii) Counting: the
multiples-of-$m$ count (question 4), the digit-choice product
(question 18), and the row-sum bound $4^n \leq
(2n+1)\binom{2n}n$ (question 21) are all
\cref{ch:b1:counting}-style arguments.

\textbf{25.} Legendre converts ``what power of $p$ \omterm{def:b1:arith:divides}{divides}
$n!$'' into base-$p$ digit arithmetic; Kummer compresses the
answer for \omterm{def:b1:counting:objects}{binomial coefficients} into the carries of a single
addition --- \omterm{def:b1:arith:divides}{divisibility}, seemingly a global property of huge
numbers, is read off locally, digit by digit. Question 16 is the
paradigm: four carries, computed by hand, determine the exact
power of $3$ in a number with hundreds of digits. And question 21
shows the same circle of ideas brushing against deep waters: an
exponential lower bound for $\operatorname{lcm}(1, \dots, 2n)$ is
a first, fully elementary step toward the \omterm{def:b1:arith:prime}{prime number} theorem,
whose proof lies far beyond this volume. The entire toolkit ---
division, gcd, valuations --- is replayed for polynomials in
\cref{ch:b1:poly}, where the analogue of a digit expansion is
expansion in powers of $(X - a)$.
\end{solution}

\section*{Chapter \ref{ch:b1:structures} --- Algebraic Structures}

\begin{solution}{exo:b1:structures:1}
\emph{Stability:} $x * y = 1 \iff x + y - xy = 1 \iff (1-x)(1-y) = 0$,
impossible for $x, y \neq 1$. Indeed the key identity is
\[
1 - x * y = (1 - x)(1 - y):
\]
the map $\varphi(x) = 1 - x$ sends $(E, *)$ to $(\R^*, \times)$ with
$\varphi(x * y) = \varphi(x)\varphi(y)$ --- a bijective morphism. All
axioms now transport: associativity and commutativity follow from
those of $\times$; the identity is $\varphi^{-1}(1) = 0$ (check: $x *
0 = x$); the inverse of $x$ is $\varphi^{-1}\bigl((1-x)^{-1}\bigr) =
1 - \frac{1}{1-x} = \frac{x}{x - 1}$ (which is $\neq 1$). So $(E, *)$
is an abelian group.
\end{solution}

\begin{solution}{exo:b1:structures:2}
\begin{enumerate}
  \item Yes: product of positives is positive, identity $1$, inverse
        $\frac 1x$, associativity inherited from $\R^*$.
  \item No: not stable ($1 + 1 = 2 \notin \{-1,0,1\}$).
  \item Yes: the standard example.
  \item No: not stable (odd $+$ odd $=$ even), and no identity ($0$ is
        even).
\end{enumerate}
\end{solution}

\begin{solution}{exo:b1:structures:3}
Write $\mathrm{id}$, the transpositions $\tau_{12}, \tau_{13},
\tau_{23}$ (swapping the two named points), and the cycles $c =
(1\,2\,3)$ (i.e.\ $1 \mapsto 2 \mapsto 3 \mapsto 1$) and $c^2 =
(1\,3\,2)$. The table of $\sigma\rho$ (row $\sigma$, column $\rho$,
apply $\rho$ first):
\begin{center}
\renewcommand{\arraystretch}{1.2}
\begin{tabular}{c|cccccc}
$\sigma\backslash\rho$ & $\mathrm{id}$ & $c$ & $c^2$ & $\tau_{12}$ &
$\tau_{13}$ & $\tau_{23}$ \\
\hline
$\mathrm{id}$ & $\mathrm{id}$ & $c$ & $c^2$ & $\tau_{12}$ &
$\tau_{13}$ & $\tau_{23}$ \\
$c$ & $c$ & $c^2$ & $\mathrm{id}$ & $\tau_{13}$ & $\tau_{23}$ &
$\tau_{12}$ \\
$c^2$ & $c^2$ & $\mathrm{id}$ & $c$ & $\tau_{23}$ & $\tau_{12}$ &
$\tau_{13}$ \\
$\tau_{12}$ & $\tau_{12}$ & $\tau_{23}$ & $\tau_{13}$ & $\mathrm{id}$
& $c^2$ & $c$ \\
$\tau_{13}$ & $\tau_{13}$ & $\tau_{12}$ & $\tau_{23}$ & $c$ &
$\mathrm{id}$ & $c^2$ \\
$\tau_{23}$ & $\tau_{23}$ & $\tau_{13}$ & $\tau_{12}$ & $c^2$ & $c$ &
$\mathrm{id}$ \\
\end{tabular}
\end{center}
Non-commuting pair: $\tau_{12}\tau_{13} = c^2$ while
$\tau_{13}\tau_{12} = c$. (To check one entry:
$\tau_{12}\tau_{13}$ sends $1 \xmapsto{\tau_{13}} 3
\xmapsto{\tau_{12}} 3$, $3 \mapsto 1 \mapsto 2$, $2 \mapsto 2 \mapsto
1$: that is $1 \mapsto 3 \mapsto 2 \mapsto 1$, the cycle $c^2 =
(1\,3\,2)$.)
\end{solution}

\begin{solution}{exo:b1:structures:4}
$H$: $1 \in H$; for $z, w \in H$, $\abs{zw^{-1}} = \abs z / \abs w =
1$: the criterion applies. $\R_+^*$: same, $\abs{xy^{-1}}$ replaced by
positivity. Union: $\iu \in H$ and $2 \in \R_+^*$, but $2\iu$ has
modulus $2 \neq 1$ and is not a positive real: $2\iu \notin H \cup
\R_+^*$, so the union is not stable --- not a subgroup (as predicted
by \cref{exo:b1:structures:6}, neither subgroup contains the other).
\end{solution}

\begin{solution}{exo:b1:structures:5}
Morphism: $\eu^{\iu(\theta + \varphi)} =
\eu^{\iu\theta}\eu^{\iu\varphi}$ (\cref{thm:b1:complex:funceq}).
Kernel: $\eu^{\iu\theta} = 1 \iff \theta \in 2\pi\Z$, so $\ker f =
2\pi\Z \neq \{0\}$: not injective. Image: every unit complex number is
$\eu^{\iu\theta}$ for some $\theta$ (polar form), so
$\operatorname{im} f = \mathbb{U}$, the unit circle. The restriction
of $f$ to a half-open interval of length $2\pi$, such as
$\intco{0}{2\pi}$ or $\intoc{-\pi}{\pi}$, is injective (two angles
with the same image differ by a multiple of $2\pi$, and only one
representative of each class fits in the interval); no interval of
greater length works, since it contains two points at distance
$2\pi$.
\end{solution}

\begin{solution}{exo:b1:structures:6}
Intersection: $e \in H \cap K$, and $x, y \in H \cap K$ gives
$xy^{-1}$ in both $H$ and $K$. Union: if $H \subseteq K$ the union is
$K$, a subgroup (and symmetrically). Conversely, suppose neither
inclusion holds: pick $h \in H \setminus K$ and $k \in K \setminus
H$, and suppose $H \cup K$ were a subgroup; then $hk \in H \cup K$.
If $hk \in H$, then $k = h^{-1}(hk) \in H$: contradiction. If $hk \in
K$, then $h = (hk)k^{-1} \in K$: contradiction. So $H \cup K$ is not
a subgroup.
\end{solution}

\begin{solution}{exo:b1:structures:7}
Note first that $x^2 = e$ means $x^{-1} = x$ for every $x$. Then for
$x, y \in G$:
\[
xy = (xy)^{-1} = y^{-1} x^{-1} = yx ,
\]
using \cref{prop:b1:structures:rules}~(2). So $G$ is abelian.
\end{solution}

\begin{solution}{exo:b1:structures:8}
Units of $\Z/18\Z$: classes coprime to $18 = 2 \times 3^2$:
$\overline 1, \overline 5, \overline 7, \overline{11}, \overline{13},
\overline{17}$. Inverse of $\overline 5$: $5 \times 11 = 55 = 3\times
18 + 1$, so $\overline 5^{-1} = \overline{11}$.

$\overline 5 x = \overline 7$: multiply by $\overline{11}$: $x =
\overline{77} = \overline 5$ (since $77 = 4\times 18 + 5$). Unique
solution.

$\overline 6 x = \overline 3$: the equation $6x \equiv 3 \pmod{18}$
means $18 \mid 6x - 3$. But $6x - 3 = 3(2x - 1)$ is odd, while $18$
is even: an even number cannot divide an odd one. No solution.

$\overline 6 x = \overline{12}$: $6x \equiv 12 \pmod{18} \iff x
\equiv 2 \pmod 3$: solutions $x \in \{\overline 2, \overline 5,
\overline 8, \overline{11}, \overline{14}, \overline{17}\}$ --- six
of them.
\end{solution}

\begin{solution}{exo:b1:structures:9}
$\Z[\sqrt 2]$ contains $0$ and $1$, and is stable under subtraction
and product:
\[
(a + b\sqrt 2)(c + d\sqrt 2) = (ac + 2bd) + (ad + bc)\sqrt 2 ,
\]
so it is a subring of $\R$ (commutativity, associativity,
distributivity are inherited). Unit: $(1 + \sqrt 2)(-1 + \sqrt 2) =
2 - 1 = 1$, so $1 + \sqrt 2$ is invertible with inverse $\sqrt 2 - 1
\in \Z[\sqrt 2]$. Its powers $(1 + \sqrt 2)^n$ are strictly
increasing (the base is $> 1$), hence pairwise distinct, and each is
a unit ($\bigl((1+\sqrt2)^n\bigr)^{-1} = (\sqrt 2 - 1)^n$): the group
of units is infinite.
\end{solution}

\begin{solution}{exo:b1:structures:10}
$x + x = (x + x)^2 = x^2 + x^2 + x^2 + x^2 = 4x^2 = 4x$ --- so $2x =
4x$, giving $2x = 0$, i.e.\ $x + x = 0$ (each element is its own
additive inverse). Then
\[
x + y = (x+y)^2 = x^2 + xy + yx + y^2 = x + xy + yx + y ,
\]
so $xy + yx = 0$, i.e.\ $xy = -yx = yx$ (using $-z = z$). Hence $A$ is
commutative.

Example: on $\mathcal{P}(E)$, define $A + B = (A \cup B) \setminus (A
\cap B)$ (symmetric difference) and $A \times B = A \cap B$. One
checks: $(\mathcal{P}(E), +)$ is an abelian group with identity
$\emptyset$ and each set its own inverse; $\cap$ is associative,
commutative, with identity $E$; distributivity $A \cap (B + C) = (A
\cap B) + (A \cap C)$ holds (an element lies in the left side iff it
is in $A$ and in exactly one of $B, C$). And $A \cap A = A$: every
element is idempotent, as required.
\end{solution}

\begin{solution}{exo:b1:structures:11}
Write the hypothesis for $n$ and $n+1$:
\[
(xy)^{n+1} = x^{n+1} y^{n+1}
\quad\text{and}\quad
(xy)^{n+1} = (xy)(xy)^n = xy\,x^n y^n .
\]
Equating: $x^{n+1} y^{n+1} = x\,y\,x^n\,y^n$; cancel $x$ on the left
and $y^n$ on the right: $x^n y = y x^n$. The same computation one
degree higher ($n+1$ and $n+2$) gives $x^{n+1} y = y x^{n+1}$. Then
\[
y\,x^{n+1} = x^{n+1} y = x\,(x^n y) = x\,y\,x^n ,
\]
and cancelling $x^n$ on the right of $y x \cdot x^n = x y \cdot x^n$:
$yx = xy$. So $G$ is abelian.
\end{solution}

\section*{Chapter \ref{ch:b1:poly} --- Polynomials}

\begin{solution}{exo:b1:poly:1}
$X^5 - 1 = (X^2 + X + 1)(X^3 - X^2 + 1) + (-X - 2)$. Steps: subtract
$X^3 B$, then $-X^2 B$, then $B$; the remainder $-X - 2$ has degree
$1 < 2$. \emph{Check at $X = 1$:} $\;0 = 3 \times 1 + (-3)$.

$2X^4 + X^3 - X + 3 = (X^2 - 2)(2X^2 + X + 4) + (X + 11)$. \emph{Check
at $X = 0$:} $\;3 = (-2)(4) + 11$.
\end{solution}

\begin{solution}{exo:b1:poly:2}
$X^2 + X + 1 = (X - j)(X - j^2)$ with $j = \eu^{2\iu\pi/3}$, $j^3 =
1$. It divides $Q_n = X^{2n} + X^n + 1$ iff $j$ and $j^2$ are roots of
$Q_n$; since $Q_n$ has real coefficients, $Q_n(j^2) =
\conj{Q_n(j)}$, so the condition is just $Q_n(j) = 0$. Now $Q_n(j) =
j^{2n} + j^n + 1$ depends on $n$ mod $3$:
\begin{itemize}
  \item $n \equiv 0$: $Q_n(j) = 1 + 1 + 1 = 3 \neq 0$;
  \item $n \equiv 1$: $Q_n(j) = j^2 + j + 1 = 0$;
  \item $n \equiv 2$: $Q_n(j) = j^4 + j^2 + 1 = j + j^2 + 1 = 0$.
\end{itemize}
So $X^2 + X + 1 \mid X^{2n} + X^n + 1$ exactly when $3 \nmid n$.
\end{solution}

\begin{solution}{exo:b1:poly:3}
By \cref{prop:b1:poly:multiplicity}, $(X-1)^2 \mid P$ iff $P(1) =
P'(1) = 0$:
\[
P(1) = 2 + a + b = 0, \qquad P'(1) = 4 + 3a + 2b = 0 .
\]
Solving: $b = -a - 2$ and $4 + 3a - 2a - 4 = a = 0$, so $a = 0$, $b =
-2$: $P = X^4 - 2X^2 + 1 = (X^2 - 1)^2 = (X-1)^2 (X+1)^2$, which is
the real factorization.
\end{solution}

\begin{solution}{exo:b1:poly:4}
$X^3 - 1 = (X - 1)(X - j)(X - j^2)$ over $\C$ ($j = \eu^{2\iu\pi/3}$),
and $(X - 1)(X^2 + X + 1)$ over $\R$.

$X^4 + X^2 + 1 = (X^2 + X + 1)(X^2 - X + 1)$ over $\R$ (multiply out,
or note $X^4 + X^2 + 1 = (X^2+1)^2 - X^2$); over $\C$, each quadratic
splits: roots $j, j^2$ and $-j, -j^2$, i.e.\ $\eu^{\pm 2\iu\pi/3},
\eu^{\pm\iu\pi/3}$.

$X^6 - 1 = \prod_{k=0}^{5} (X - \eu^{\iu k\pi/3})$ over $\C$, and over
$\R$:
\[
X^6 - 1 = (X-1)(X+1)(X^2 + X + 1)(X^2 - X + 1),
\]
grouping the conjugate pairs $\eu^{\pm 2\iu\pi/3}$ and $\eu^{\pm
\iu\pi/3}$.
\end{solution}

\begin{solution}{exo:b1:poly:5}
\begin{enumerate}
  \item Let $p/q$ (lowest terms) be a root of the monic integer
        polynomial $X^3 + \dots + c_0$: clearing denominators in
        $P(p/q) = 0$ gives $p^3 = -q\,(\text{integer})$, so $q \mid
        p^3$; coprimality forces $q = \pm 1$: the root is an integer
        $p$, and $p \mid c_0$ (isolate $c_0$). Here the candidates
        divide $6$: testing, $P(1) = 0$, $P(2) = 0$, $P(3) = 0$. So
        $P = (X-1)(X-2)(X-3)$.
  \item Vieta: $s_1 = 6$, $s_2 = 11$, $s_3 = 6$. Sum of squares:
        $s_1^2 - 2s_2 = 36 - 22 = 14 = 1 + 4 + 9$, as expected. Sum of inverses:
        $\frac{s_2}{s_3} = \frac{11}{6} = 1 + \frac12 + \frac13$, as expected.
\end{enumerate}
\end{solution}

\begin{solution}{exo:b1:poly:6}
First division step of the Euclidean algorithm:
\[
(X + 1)(X^3 - X^2 + X - 1)
= X^4 - X^3 + X^2 - X + X^3 - X^2 + X - 1 = X^4 - 1 ,
\]
so the division of $X^4 - 1$ by $X^3 - X^2 + X - 1$ is exact
(quotient $X + 1$, remainder $0$), and the algorithm stops at once:
\[
\gcd(X^4 - 1,\; X^3 - X^2 + X - 1) = X^3 - X^2 + X - 1
\]
(already monic). The Bézout relation is the trivial one: $\gcd = 0
\cdot (X^4 - 1) + 1 \cdot (X^3 - X^2 + X - 1)$. Consistency check by
factoring: $X^3 - X^2 + X - 1 = (X - 1)(X^2 + 1)$, which is indeed
the product of the common irreducible factors of $X^4 - 1 =
(X-1)(X+1)(X^2+1)$.
\end{solution}

\begin{solution}{exo:b1:poly:7}
Since $P \geq 0$ on $\R$, its real roots have even multiplicity (at a
root of odd multiplicity, $P$ changes sign). Using
\cref{cor:b1:poly:factorization} and pairing, write
\[
P = c \prod_i (X - a_i)^{2k_i} \prod_j \bigl((X - z_j)(X - \conj
z_j)\bigr)^{n_j},
\]
with $c > 0$ (behavior at $+\infty$). Let
\[
S = \sqrt c\, \prod_i (X - a_i)^{k_i} \prod_j (X - z_j)^{n_j}
\in \C[X],
\]
so that $P = S\,\conj S$ where $\conj S$ has the conjugated
coefficients. Split $S = A + \iu B$ with $A, B \in \R[X]$: then
\[
P = (A + \iu B)(A - \iu B) = A^2 + B^2 .
\]
\end{solution}

\begin{solution}{exo:b1:poly:8}
Lagrange (\cref{thm:b1:poly:lagrange}) with nodes $0, 1, 2$:
\[
P = 1\cdot\frac{(X-1)(X-2)}{(0-1)(0-2)} + 3\cdot\frac{X(X-2)}{1\cdot(1-2)}
+ 2\cdot\frac{X(X-1)}{2\cdot 1}
= \frac{(X-1)(X-2)}{2} - 3X(X-2) + X(X-1).
\]
Expanding: $\frac{X^2 - 3X + 2}{2} - 3X^2 + 6X + X^2 - X =
-\frac{3}{2}X^2 + \frac{7}{2}X + 1$.

System: $P = aX^2 + bX + c$ with $c = 1$; $a + b + 1 = 3$; $4a + 2b +
1 = 2$. Subtracting twice the second from the third: $2a - 1 = -4$,
so $a = -\frac32$, $b = \frac72$. Same polynomial:
$P = -\frac32 X^2 + \frac72 X + 1$. (Check $P(2) = -6 + 7 + 1 = 2$.)
\end{solution}

\begin{solution}{exo:b1:poly:9}
$P = X^{2n+1} - 1$: $P' = (2n+1)X^{2n} \geq 0$, so the polynomial
function is increasing (strictly except at $0$), with limits
$\mp\infty$: it vanishes exactly once on $\R$ (at $x = 1$).

Let $E_n = \sum_{k=0}^{n} \frac{X^k}{k!}$. Then $E_n' = E_{n-1} = E_n
- \frac{X^n}{n!}$. A multiple root $a$ would satisfy $E_n(a) =
E_n'(a) = 0$ (\cref{prop:b1:poly:multiplicity}), hence
$\frac{a^n}{n!} = E_n(a) - E_n'(a) = 0$, so $a = 0$; but $E_n(0) = 1
\neq 0$. No multiple root.
\end{solution}

\begin{solution}{exo:b1:poly:10}
\begin{enumerate}
  \item Induction (two base cases hold). Using $\cos(n+1)\theta +
        \cos(n-1)\theta = 2\cos\theta\cos n\theta$:
        \[
        T_{n+1}(\cos\theta) = 2\cos\theta \cos n\theta -
        \cos(n-1)\theta = \cos(n+1)\theta .
        \]
  \item $T_n(\cos\theta) = 0$ iff $n\theta \equiv \frac\pi2 \pmod
        \pi$: the numbers
        \[
        x_k = \cos\Bigl(\frac{(2k+1)\pi}{2n}\Bigr),
        \qquad k = 0, 1, \dots, n-1,
        \]
        are $n$ distinct points of $\intoo{-1}{1}$ (the angles lie in
        $\intoo{0}{\pi}$ where $\cos$ is injective), all roots of
        $T_n$; since $\deg T_n = n$ (from the recurrence, with
        leading coefficient $2^{n-1}$ for $n \geq 1$, by induction),
        these are \emph{all} the roots, each simple.
  \item For $x = \cos\theta \in \intcc{-1}{1}$: $\abs{T_n(x)} =
        \abs{\cos n\theta} \leq 1$, with equality iff $n\theta \equiv
        0 \pmod\pi$, i.e.\ at the $n+1$ points $y_k =
        \cos\frac{k\pi}{n}$, $k = 0, \dots, n$, where $T_n(y_k) =
        (-1)^k$. (This equioscillation makes $2^{1-n}T_n$ the monic
        degree-$n$ polynomial of smallest sup-norm on
        $\intcc{-1}{1}$ --- proved in this chapter's weekend
        problem.)
\end{enumerate}
\end{solution}

\begin{solution}{exo:b1:poly:11}
Write $P = c\prod_i (X - a_i)^{m_i}$. The product rule (extended to
several factors) gives
\[
P' = c\sum_{i} m_i (X - a_i)^{m_i - 1} \prod_{k \neq i} (X -
a_k)^{m_k},
\]
and dividing by $P$: $\frac{P'}{P} = \sum_i \frac{m_i}{X - a_i}$
(as rational functions, i.e.\ away from the roots).

Let $w$ be a root of $P'$. If $w$ is one of the $a_i$, it lies in the
convex hull trivially. Otherwise, evaluating at $w$:
\[
0 = \sum_i \frac{m_i}{w - a_i}
= \sum_i m_i\, \frac{\conj w - \conj a_i}{\abs{w - a_i}^2} .
\]
Conjugating: $\sum_i \lambda_i (w - a_i) = 0$ where $\lambda_i =
\frac{m_i}{\abs{w - a_i}^2} > 0$. Hence
\[
w = \frac{\sum_i \lambda_i a_i}{\sum_i \lambda_i} :
\]
a convex combination (positive weights summing to $1$ after
normalization) of the roots $a_i$. So every root of $P'$ lies in the
convex hull of the roots of $P$.
\end{solution}

\begin{solution}{exo:b1:poly:12}
Sum the evaluations of $(1 + X)^n$ at the three cube roots of
unity:
\[
2^n + (1 + j)^n + (1 + j^2)^n
= \sum_{k=0}^n \binom nk\,\bigl(1 + j^k + j^{2k}\bigr)
= 3\sum_{k\,:\,3\mid k}\binom nk ,
\]
since $1 + j^k + j^{2k}$ is a geometric sum equal to $3$ when $3
\mid k$ and to $\frac{j^{3k} - 1}{j^k - 1} = 0$ otherwise. Now $1 +
j = \frac12 + \iu\frac{\sqrt3}2 = \eu^{\iu\pi/3}$ and $1 + j^2 =
\conj{1 + j} = \eu^{-\iu\pi/3}$, so $(1+j)^n + (1+j^2)^n =
2\cos\frac{n\pi}3$ and
\[
\sum_{k\geq0}\binom n{3k} = \frac{2^n + 2\cos\frac{n\pi}3}{3} .
\]
Checks: $n = 3$: $\frac{8 + 2\cos\pi}3 = 2 = \binom30 + \binom33$;
$n = 6$: $\frac{64 + 2}3 = 22 = 1 + 20 + 1$.
\end{solution}

\begin{solution}{pb:b1:poly:1}
\textbf{1.} $T_2 = 2X^2 - 1$; $T_3 = 2X(2X^2 - 1) - X = 4X^3 -
3X$; $T_4 = 2X\,T_3 - T_2 = 8X^4 - 8X^2 + 1$; $T_5 = 2X\,T_4 - T_3
= 16X^5 - 20X^3 + 5X$. The identity $T_3(\cos\theta) = \cos3\theta$
is exactly $\cos3\theta = 4\cos^3\theta - 3\cos\theta$ from
\cref{ex:b1:complex:cos3}.

\textbf{2.} True for $n = 1, 2$. If $T_{n-1}$, $T_n$ have degrees
$n-1$, $n$ and leading coefficients $2^{n-2}$, $2^{n-1}$, then
$2X\,T_n$ has degree $n+1$ and leading coefficient $2^n$, while
$T_{n-1}$ has lower degree: $T_{n+1}$ has degree $n + 1$, leading
coefficient $2^n$. Parity: if $T_{n-1}$ has the parity of $n - 1$
and $T_n$ that of $n$, then $2X\,T_n$ and $T_{n-1}$ both have the
parity of $n + 1$, hence so does $T_{n+1}$.

\textbf{3.} If $P(\cos\theta) = \cos n\theta$ for all $\theta$,
then $P$ and $T_n$ agree at every point of $\intcc{-1}1$ --- an
infinite set --- so $P - T_n$ has infinitely many roots and is the
zero polynomial (\cref{cor:b1:poly:nroots}).

\textbf{4.} For $x = \cos\theta$: $T_m(T_n(\cos\theta)) =
T_m(\cos n\theta) = \cos(mn\theta) = T_{mn}(\cos\theta)$, and
$2T_mT_n(\cos\theta) = 2\cos m\theta\cos n\theta = \cos(m+n)\theta
+ \cos\abs{m - n}\theta$. Both identities hold on $\intcc{-1}1$,
hence as polynomial identities by question 3's argument.

\textbf{5.} The $x_k$ are $n$ distinct simple roots and the
leading coefficient is $2^{n-1}$:
\[
T_n = 2^{n-1}\prod_{k=0}^{n-1}
\Bigl(X - \cos\frac{(2k+1)\pi}{2n}\Bigr) .
\]
Interlacing: the angles $0 < \frac{\pi}{2n} < \frac\pi n <
\frac{3\pi}{2n} < \frac{2\pi}n < \dots < \pi$ alternate between
the $y$-angles $\frac{k\pi}n$ and the $x$-angles
$\frac{(2k+1)\pi}{2n}$; since $\cos$ is strictly decreasing on
$\intcc0\pi$, the values interlace in the reverse order: $y_n <
x_{n-1} < y_{n-1} < \dots < x_0 < y_0$. Between two consecutive
extrema sits exactly one root, as a picture of $\cos n\theta$
suggests.

\textbf{6.} Induction with $2\cosh a\cosh b = \cosh(a + b) +
\cosh(a - b)$ (\cref{prop:b1:functions:hyprules}): $T_{n+1}(\cosh
t) = 2\cosh t\cosh nt - \cosh(n-1)t = \cosh(n+1)t$. For $x \geq
1$, write $x = \cosh t$ with $t \geq 0$; then $\eu^t = x +
\sqrt{x^2 - 1}$ and $\eu^{-t} = x - \sqrt{x^2 - 1}$, so
\[
T_n(x) = \cosh(nt)
= \frac{(x + \sqrt{x^2-1})^n + (x - \sqrt{x^2-1})^n}2 .
\]
For $x > 1$ the first term exceeds $\frac12(1)^n$ strictly and
grows geometrically: $T_n(x) > 1$.

\textbf{7.} De Moivre: $\cos n\theta = \Re\bigl((\cos\theta +
\iu\sin\theta)^n\bigr) = \sum_{2j \leq n}\binom n{2j}
\cos^{n-2j}\theta\,(\iu\sin\theta)^{2j}$, and $(\iu\sin\theta)^{2j}
= (-\sin^2\theta)^j = (\cos^2\theta - 1)^j$. Substituting $x =
\cos\theta$ and invoking question 3:
\[
T_n(x) = \sum_{0\leq 2j\leq n}\binom n{2j}x^{n-2j}(x^2 - 1)^j .
\]
For $n = 3$: $\binom30 x^3 + \binom32 x(x^2 - 1) = x^3 + 3x^3 -
3x = 4x^3 - 3x$, as in question 1.

\textbf{8.} $T_n(1) = \cos(n\cdot0) = 1$; $T_n(-1) = \cos(n\pi) =
(-1)^n$; $T_n(0) = \cos\frac{n\pi}2$, which is $0$ for odd $n$ and
$(-1)^{n/2}$ for even $n$.

\textbf{9.} Induction for $U_n(\cos\theta) =
\frac{\sin(n+1)\theta}{\sin\theta}$: true for $U_0 = 1$ and $U_1 =
2X$ ($\sin2\theta = 2\sin\theta\cos\theta$); the step is the
sum-to-product identity $\sin(n+2)\theta = 2\cos\theta\,
\sin(n+1)\theta - \sin n\theta$. Now differentiate
$T_n(\cos\theta) = \cos n\theta$ in $\theta$:
$-\sin\theta\,T_n'(\cos\theta) = -n\sin n\theta$, so for $\theta
\notin \pi\Z$:
\[
T_n'(\cos\theta) = n\,\frac{\sin n\theta}{\sin\theta}
= n\,U_{n-1}(\cos\theta) ,
\]
and the polynomials $T_n'$ and $nU_{n-1}$, agreeing on
$\intoo{-1}1$, are equal.

\textbf{10.} $\abs{\sin(n+1)\theta} = \abs{\sin n\theta\cos\theta
+ \cos n\theta\sin\theta} \leq \abs{\sin n\theta} +
\abs{\sin\theta}$, and induction gives $\abs{\sin n\theta} \leq
n\abs{\sin\theta}$. Hence $\abs{U_{n-1}} \leq n$ on $\intoo{-1}1$
and $\abs{T_n'} = n\abs{U_{n-1}} \leq n^2$ there; at $\pm1$ the
bound extends by taking limits (or directly: $U_{n-1}(1) = n$ from
the recurrence, $U_n(1) = n + 1$ by induction, and parity gives
$U_{n-1}(-1) = (-1)^{n-1}n$). Thus $T_n'(1) = n^2$ and $T_n'(-1)
= (-1)^{n-1}n^2$: the bound $n^2$ is attained at the endpoints.

\textbf{11.} Differentiate $\sin\theta\,T_n'(\cos\theta) = n\sin
n\theta$ (question 9) with respect to $\theta$:
\[
\cos\theta\,T_n'(\cos\theta) - \sin^2\theta\,T_n''(\cos\theta)
= n^2\cos n\theta = n^2\,T_n(\cos\theta) .
\]
With $x = \cos\theta$ and $\sin^2\theta = 1 - x^2$: $x\,T_n' - (1
- x^2)T_n'' = n^2T_n$ on $\intcc{-1}1$, hence everywhere:
$(1 - x^2)y'' - xy' + n^2y = 0$ for $y = T_n$. Check for $T_2 =
2x^2 - 1$: $(1 - x^2)(4) - x(4x) + 4(2x^2 - 1) = 4 - 4x^2 - 4x^2
+ 8x^2 - 4 = 0$.

\textbf{12.} $\widetilde T_n$ is monic (question 2) and
$\abs{\widetilde T_n} = 2^{1-n}\abs{T_n} \leq 2^{1-n}$ on
$\intcc{-1}1$, with $\widetilde T_n(y_k) = (-1)^k2^{1-n}$ at the
$n + 1$ points $y_k$ (\cref{exo:b1:poly:10}): the norm is exactly
$2^{1-n}$, attained with alternating signs.

\textbf{13.} $\widetilde T_n$ and $P$ are both monic of degree
$n$, so the leading terms cancel: $\deg D \leq n - 1$. At $y_k$:
$D(y_k) = (-1)^k2^{1-n} - P(y_k)$, and $\abs{P(y_k)} \leq \norm
P_\infty < 2^{1-n}$ forces the sign of $D(y_k)$ to be that of
$(-1)^k2^{1-n}$, strictly.

\textbf{14.} $D$ changes sign between $y_{k+1}$ and $y_k$ for each
$k = 0, \dots, n-1$: by the intermediate value property, $D$ has a
root in each of these $n$ pairwise disjoint open intervals --- $n$
distinct roots for a nonzero polynomial of degree $\leq n - 1$,
impossible. And $D = 0$ is impossible too (the norms differ).
Contradiction: no monic $P$ of degree $n$ has $\norm P_\infty <
2^{1-n}$, which is Chebyshev's theorem.

\textbf{15.} Now $\abs{P(y_k)} \leq 2^{1-n}$ only, so $(-1)^k
D(y_k) = 2^{1-n} - (-1)^kP(y_k) \geq 2^{1-n} - \abs{P(y_k)} \geq
0$. Suppose $D(y_k) = 0$ at an interior $y_k$ ($0 < k < n$): then
$P(y_k) = (-1)^k2^{1-n}$, so $\abs P$ attains its supremum
$2^{1-n}$ at the interior point $y_k$, whence $P'(y_k) = 0$
(interior extremum); and $T_n'(y_k) = nU_{n-1}(y_k) = 0$ since
$\sin(n\cdot\frac{k\pi}n) = 0$ --- so $\widetilde T_n'(y_k) = 0$
too, and $D'(y_k) = 0$: $y_k$ is a root of $D$ of multiplicity at
least $2$.

\textbf{16.} Count roots of $D$ with multiplicity. Let $z$ be the
number of interior points $y_k$ with $D(y_k) = 0$ (each a double
root, by question 15) and $e \in \{0, 1, 2\}$ the number of
endpoints ($y_0$ or $y_n$) with $D = 0$ (each a simple root at
least). A gap $(y_{k+1}, y_k)$ whose two endpoints both have $D
\neq 0$ carries strictly alternating signs, hence an interior
root. Each vanishing interior point spoils at most its two
adjacent gaps, each vanishing endpoint at most one gap: at least
$n - 2z - e$ gaps still contribute one root each, all distinct
from the $y$-roots. Total: at least $(n - 2z - e) + 2z + e = n$
roots with multiplicity, for a polynomial of degree $\leq n - 1$:
so $D = 0$ and $P = \widetilde T_n$. The minimizer is unique.

\textbf{17.} The affine map $t \mapsto x = \frac{a+b}2 +
\frac{b-a}2\,t$ is a bijection $\intcc{-1}1 \to \intcc ab$. If $P$
is monic of degree $n$, then $Q(t) = P(x(t))$ is a polynomial in
$t$ with leading coefficient $\bigl(\frac{b-a}2\bigr)^n$, and
$\sup_{\intcc ab}\abs P = \sup_{\intcc{-1}1}\abs Q$. The monic
polynomial $Q/\bigl(\frac{b-a}2\bigr)^n$ has sup-norm $\geq
2^{1-n}$ (questions 13--14), so
\[
\sup_{\intcc ab}\abs P \geq \Bigl(\frac{b-a}2\Bigr)^n 2^{1-n}
= 2\Bigl(\frac{b-a}4\Bigr)^n ,
\]
with equality exactly for $P(x) = \bigl(\frac{b-a}2\bigr)^n
\widetilde T_n\bigl(t(x)\bigr)$ (question 16).

\textbf{18.} $\widetilde T_3 = \frac{T_3}4 = X^3 - \frac34X$;
$\widetilde T_3{}' = 3X^2 - \frac34$ vanishes at $\pm\frac12$.
Values: $\widetilde T_3(-1) = -\frac14$, $\widetilde
T_3(-\tfrac12) = \frac14$, $\widetilde T_3(\tfrac12) = -\frac14$,
$\widetilde T_3(1) = \frac14$: four alternating extrema of
absolute value $\frac14$ --- so $\norm{\widetilde T_3}_\infty =
\frac14$, and by Chebyshev's theorem no monic cubic has smaller
sup-norm on $\intcc{-1}1$.

\textbf{19.} $\omega$ is monic of degree $n + 1$, so
$\norm\omega_\infty \geq 2^{-n}$ by Chebyshev's theorem (degree
$n+1$), with equality if and only if $\omega = \widetilde T_{n+1}
= 2^{-n}T_{n+1}$ (question 16), i.e.\ if and only if the nodes
are the $n + 1$ roots of $T_{n+1}$. With Chebyshev nodes the error
factor $\norm\omega_\infty$ is $2^{-n}$ --- the least possible.

\textbf{20.} $5 \times 36^\circ = 180^\circ$, so $T_5(c) =
\cos180^\circ = -1$: $16c^5 - 20c^3 + 5c + 1 = 0$. Testing $x =
-1$: $-16 + 20 - 5 + 1 = 0$, and expanding confirms
\[
16x^5 - 20x^3 + 5x + 1 = (x + 1)\bigl(4x^2 - 2x - 1\bigr)^2 .
\]
Since $c = \cos36^\circ \neq -1$, $c$ is a root of $4x^2 - 2x -
1$, whose roots are $\frac{1 \pm \sqrt5}4$; as $c > 0$,
\[
\cos36^\circ = \frac{1 + \sqrt5}4 .
\]
Consistency: $\cos72^\circ = T_2(c) = 2c^2 - 1 = 2\cdot\frac{3 +
\sqrt5}8 - 1 = \frac{\sqrt5 - 1}4$, the value found in
\cref{exo:b1:complex:8}.

\textbf{21.} $\sqrt{1.1^2 - 1} = \sqrt{0.21} \approx 0.458$, so
$x + \sqrt{x^2-1} \approx 1.558$ and $(1.558)^{10} \approx 84.5$,
while $(1.1 - 0.458)^{10} \approx 0.01$: $T_{10}(1.1) \approx
\frac{84.5 + 0.01}2 \approx 42$. A polynomial trapped in
$\intcc{-1}1$ on the interval has already grown past $40$ one
tenth beyond its edge: boundedness on a segment says nothing an
inch outside it.

\textbf{22.} In question 7's formula for $T_p$, the term $j = 0$
is $X^p$; every other term carries $\binom p{2j}$ with $0 < 2j <
p$ (note $2j \neq p$ since $p$ is odd), which is divisible by $p$
by the first step of the proof of \cref{thm:b1:arith:fermat}.
Hence every coefficient of $T_p - X^p$ is a multiple of $p$.
Checks: $T_3 - X^3 = 3X^3 - 3X = 3(X^3 - X)$; $T_5 - X^5 = 15X^5
- 20X^3 + 5X = 5(3X^5 - 4X^3 + X)$.

\textbf{23.} By question 17 with $\intcc ab = \intcc01$ and $n =
2$: minimal deviation $2\bigl(\frac14\bigr)^2 = \frac18$,
attained by $\bigl(\frac12\bigr)^2\widetilde T_2(2x - 1) =
\frac14\bigl((2x-1)^2 - \frac12\bigr) = x^2 - x + \frac18$. The
monic quadratic closest to zero on $\intcc01$ is $x^2 - x +
\frac18$, with sup-norm $\frac18$.

\textbf{24.} (i) Rigidity --- a polynomial with more roots than
its degree is zero --- powered the uniqueness principle
(question 3), the transfer of trigonometric identities to
polynomial identities (questions 4, 7, 9, 11), and both
root-counting arguments of the extremality proof (questions 14,
16). (ii) The trigonometry of \cref{ch:b1:complex} (de Moivre,
sum-to-product) and the hyperbolic functions of
\cref{ch:b1:functions} supplied every identity behind the family;
the substitution $x = \cos\theta$ is the bridge. (iii) The
divisibility $p \mid \binom p{2j}$ from \cref{ch:b1:arith} turned
the coefficient formula into the congruence of question 22.

\textbf{25.} Chebyshev's theorem converts an optimization over an
infinite-dimensional family (all monic polynomials) into finite
combinatorics: a competitor better than $\widetilde T_n$ would
differ from it by a low-degree polynomial forced to change sign
$n$ times --- one more root than its degree allows. The
equioscillation pattern is thus not a curiosity but the very
certificate of optimality, and the equality case sharpens
root-counting with multiplicities. The substitution $x =
\cos\theta$ deserves the last word: it transports the rigid,
discrete world of polynomials into the periodic world of
trigonometry, where roots and extrema of $T_n$ are simply the
regular grid of $\cos n\theta$. The two analysis facts borrowed
--- the intermediate value property (question 14; proved in
\cref{ch:b1:continuity}) and the vanishing derivative at an
interior extremum (question 15; proved in \cref{ch:b1:derivative})
--- are exactly the tools those later chapters will hand back,
closing the loop.
\end{solution}

\section*{Chapter \ref{ch:b1:fractions} --- Rational Fractions}

\begin{solution}{exo:b1:fractions:1}
$\dfrac{1}{X^2 - 1}$: simple \omterm{def:b1:fractions:field}{poles} $\pm 1$; cover-up:
$\dfrac{1/2}{X-1} - \dfrac{1/2}{X+1}$.

$\dfrac{X}{X^2 - 3X + 2} = \dfrac{X}{(X-1)(X-2)}$: cover-up gives
$\frac{1}{1-2} = -1$ at $1$ and $\frac{2}{2-1} = 2$ at $2$:
$\dfrac{-1}{X-1} + \dfrac{2}{X-2}$.

$\dfrac{X^2+1}{X(X-1)}$: the degree is $0$, so there is an integer
part: dividing, $X^2 + 1 = (X^2 - X) + (X + 1)$, so $F = 1 +
\frac{X+1}{X(X-1)}$. Cover-up on the remainder: $\frac{1}{-1} = -1$
at $0$, $\frac{2}{1} = 2$ at $1$:
\[
F = 1 - \frac{1}{X} + \frac{2}{X-1} .
\]
\end{solution}

\begin{solution}{exo:b1:fractions:2}
$\dfrac{1}{X(X^2+1)}$: shape $\frac aX + \frac{bX + c}{X^2 + 1}$.
Cover-up at $0$: $a = 1$. Limit of $xF$: $0 = a + b$, so $b = -1$.
Evaluation at $X = 1$: $\frac12 = 1 + \frac{c - 1}{2}$, so $c = 0$:
\[
\frac{1}{X(X^2+1)} = \frac 1X - \frac{X}{X^2+1} .
\]

$\dfrac{X^3}{X^2+X+1}$: division: $X^3 = (X^2+X+1)(X - 1) + 1$, so
\[
\frac{X^3}{X^2+X+1} = X - 1 + \frac{1}{X^2 + X + 1} ,
\]
already in real decomposed form (the quadratic has negative
discriminant).
\end{solution}

\begin{solution}{exo:b1:fractions:3}
$\dfrac{1}{X^2(X-1)}$: shape $\frac{a}{X^2} + \frac bX +
\frac{c}{X-1}$. Cover-up at the double \omterm{def:b1:fractions:field}{pole} $0$:
$a = \bigl[\frac{1}{X-1}\bigr]_{0} = -1$. Cover-up at $1$: $c = 1$.
Limit of $xF$: $0 = b + c$, so $b = -1$:
\[
\frac{1}{X^2(X-1)} = -\frac{1}{X^2} - \frac1X + \frac{1}{X-1} .
\]

$\dfrac{X+1}{(X-1)^3}$: with $Y = X - 1$, the numerator is $Y + 2$:
\[
\frac{Y + 2}{Y^3} = \frac{1}{Y^2} + \frac{2}{Y^3}
= \frac{1}{(X-1)^2} + \frac{2}{(X-1)^3} .
\]
\end{solution}

\begin{solution}{exo:b1:fractions:4}
The \omterm{def:b1:fractions:field}{poles} are the fourth \omterm{def:b1:complex:unity}{roots of unity} $1, \iu, -1, -\iu$, all
simple. \omterm{def:b1:fractions:field}{Simple-pole} formula with $B' = 4X^3$: coefficient at $a$ is
$\frac{1}{4a^3} = \frac{a}{4a^4} = \frac a4$ (using $a^4 = 1$). So,
over $\C$:
\[
\frac{1}{X^4 - 1}
= \frac{1/4}{X - 1} - \frac{1/4}{X + 1}
+ \frac{\iu/4}{X - \iu} - \frac{\iu/4}{X + \iu} .
\]
Grouping the conjugate pair (common denominator $X^2 + 1$):
$\frac{\iu}{4}\bigl(\frac{1}{X-\iu} - \frac{1}{X+\iu}\bigr) =
\frac{\iu}{4}\cdot\frac{2\iu}{X^2+1} = \frac{-1/2}{X^2+1}$. Over
$\R$:
\[
\frac{1}{X^4 - 1}
= \frac{1/4}{X-1} - \frac{1/4}{X+1} - \frac{1/2}{X^2 + 1} .
\]
\emph{Check at $X = 0$:} $-1 = -\frac14 - \frac14 - \frac12$.
\end{solution}

\begin{solution}{exo:b1:fractions:5}
$\dfrac{X^2}{(X^2+1)^2} = \dfrac{(X^2 + 1) - 1}{(X^2+1)^2} =
\dfrac{1}{X^2+1} - \dfrac{1}{(X^2+1)^2}$.

Hence a primitive:
\[
\int \frac{x^2\,\dd x}{(x^2+1)^2}
= \arctan x - \frac12\Bigl(\arctan x + \frac{x}{x^2+1}\Bigr) + C
= \frac12\arctan x - \frac{x}{2(x^2+1)} + C .
\]
\end{solution}

\begin{solution}{exo:b1:fractions:6}
The \omterm{def:b1:fractions:field}{poles} $0, -1, \dots, -n$ are simple. Cover-up at $-k$:
\[
c_k = \frac{n!}{\prod_{j \neq k} (-k + j)}
= \frac{n!}{\bigl(\prod_{j=0}^{k-1}(j - k)\bigr)
\bigl(\prod_{j=k+1}^{n}(j-k)\bigr)}
= \frac{n!}{(-1)^k k!\,(n-k)!} = (-1)^k \binom nk .
\]
So
\[
\frac{n!}{X(X+1)\cdots(X+n)}
= \sum_{k=0}^{n} \frac{(-1)^k \binom nk}{X + k} .
\]
(Sanity check for $n = 1$: $\frac{1}{X(X+1)} = \frac1X -
\frac{1}{X+1}$.)
\end{solution}

\begin{solution}{exo:b1:fractions:7}
$P = X^n - 1$ has the $n$ simple roots $\omega^k$
(\cref{thm:b1:complex:roots}), so the logarithmic-derivative identity
of \cref{exo:b1:poly:11} reads
\[
\sum_{k=0}^{n-1} \frac{1}{X - \omega^k}
= \frac{P'(X)}{P(X)} = \frac{n X^{n-1}}{X^n - 1} .
\]
At $X = 2$, $n = 4$: right side $= \frac{4 \times 8}{15} =
\frac{32}{15}$. Left side: $\frac{1}{2-1} + \frac{1}{2+1} +
\frac{1}{2 - \iu} + \frac{1}{2 + \iu} = 1 + \frac13 + \frac{4}{5} =
\frac{15 + 5 + 12}{15} = \frac{32}{15}$, using $\frac{1}{2-\iu} +
\frac{1}{2+\iu} = \frac{4}{5}$.
\end{solution}

\begin{solution}{exo:b1:fractions:8}
Cover-up: $\dfrac{1}{k(k+1)(k+2)} = \dfrac{1/2}{k} - \dfrac{1}{k+1} +
\dfrac{1/2}{k+2}$. Rewrite as a telescoping difference:
\[
\frac{1}{k(k+1)(k+2)}
= \frac12\Bigl(\frac{1}{k(k+1)} - \frac{1}{(k+1)(k+2)}\Bigr),
\]
(expand to check --- or subtract the two decompositions). Summing:
\[
S_n = \frac12\Bigl(\frac{1}{1 \times 2} - \frac{1}{(n+1)(n+2)}\Bigr)
= \frac14 - \frac{1}{2(n+1)(n+2)}
\xrightarrow[n \to \infty]{} \frac14 .
\]
\end{solution}

\begin{solution}{exo:b1:fractions:9}
Decompose, for $0 \leq m \leq n - 1$, the fraction $\frac{X^m}{P}$
(degree $m - n \leq -1$, \omterm{def:b1:fractions:field}{poles} simple): the \omterm{def:b1:fractions:field}{simple-pole} formula gives
\[
\frac{X^m}{P} = \sum_{i=1}^{n} \frac{x_i^m / P'(x_i)}{X - x_i} .
\]
Multiply by $X$ and let $X \to +\infty$: the left side tends to the
limit of $X^{m+1}/P$, which is $0$ if $m \leq n - 2$ and $1$ if $m =
n - 1$ ($P$ is \omterm{def:b1:poly:def}{monic} of degree $n$); the right side tends to
$\sum_i \frac{x_i^m}{P'(x_i)}$. Hence
\[
\sum_{i=1}^{n} \frac{x_i^{m}}{P'(x_i)} =
\begin{cases}
0 & \text{for } 0 \leq m \leq n-2,\\
1 & \text{for } m = n - 1,
\end{cases}
\]
which contains both announced identities ($m = 0$ requires $n \geq
2$). (Interpretation via \cref{thm:b1:poly:lagrange}: these sums are
the leading coefficients of the Lagrange interpolants of $X^m$, and
interpolating a \omterm{def:b1:poly:def}{polynomial} of degree $\leq n-1$ at $n$ points
reproduces it exactly.)
\end{solution}

\begin{solution}{exo:b1:fractions:10}
Shape $\frac aX + \frac b{X+1} + \frac c{(X+1)^2}$. Cover-up at
$0$: $a = 1$; cover-up at the double \omterm{def:b1:fractions:field}{pole}: $c = \bigl[\frac1X
\bigr]_{X=-1} = -1$; limit of $xF(x)$ at infinity: $0 = a + b$, so
$b = -1$:
\[
\frac{1}{X(X+1)^2} = \frac1X - \frac1{X+1} - \frac1{(X+1)^2} .
\]
Summing for $n = 1, \dots, N$: the first two bricks telescope to
$1 - \frac1{N+1}$, and the third contributes $-\sum_{k=2}^{N+1}
\frac1{k^2}$. Letting $N \to \infty$ and using $\sum_{k\geq1}
\frac1{k^2} = \frac{\pi^2}6$:
\[
\sum_{n=1}^{\infty}\frac{1}{n(n+1)^2}
= 1 - \Bigl(\frac{\pi^2}6 - 1\Bigr) = 2 - \frac{\pi^2}6
\approx 0.355 .
\]
\end{solution}

\begin{solution}{exo:b1:fractions:11}
In the variable $Y$: $\frac1{(Y+1)(Y+4)} = \frac{1/3}{Y+1} -
\frac{1/3}{Y+4}$ (cover-up at $-1$ and $-4$), so
\[
\frac{1}{(X^2+1)(X^2+4)}
= \frac13\,\frac1{X^2+1} - \frac13\,\frac1{X^2+4} .
\]
Likewise $\frac{Y}{(Y+1)(Y+4)} = \frac{-1/3}{Y+1} +
\frac{4/3}{Y+4}$, so
\[
\frac{X^2}{(X^2+1)(X^2+4)}
= -\frac13\,\frac1{X^2+1} + \frac43\,\frac1{X^2+4} .
\]
(Check at $X = 0$: $0 = \frac13(-1 + 1)$.) These are already the
real decompositions: the numerators over the irreducible quadratics
happen to be constants.
\end{solution}

\begin{solution}{exo:b1:fractions:12}
\begin{enumerate}
  \item On each interval avoiding the \omterm{def:b1:fractions:field}{poles}, $F'(x) = -\sum_i
        \frac{c_i}{(x - p_i)^2} < 0$: strictly decreasing. As $x
        \to \pm\infty$, every brick tends to $0$: $F \to 0$, from
        above at $+\infty$ (all bricks positive there) and from
        below at $-\infty$. As $x \to p_i^+$, the brick
        $\frac{c_i}{x - p_i}$ blows up to $+\infty$ and the others
        stay bounded: $F \to +\infty$; likewise $F \to -\infty$ as
        $x \to p_i^-$.
  \item Fix $\lambda > 0$. On $\intoo{-\infty}{p_1}$: $F$
        decreases from $0^-$ to $-\infty$, so $F < 0 < \lambda$:
        no solution. On each $\intoo{p_i}{p_{i+1}}$ ($1 \leq i
        \leq r-1$): $F$ decreases from $+\infty$ to $-\infty$,
        hence takes the value $\lambda$ exactly once (intermediate
        value property plus strict monotonicity). On
        $\intoo{p_r}{+\infty}$: $F$ decreases from $+\infty$ to
        $0^+$, again exactly one solution. Total: exactly $r$
        solutions, interlaced with the \omterm{def:b1:fractions:field}{poles}.
\end{enumerate}
\end{solution}

\begin{solution}{pb:b1:fractions:1}
\textbf{1.} Cover-up: $\frac1{X(X+1)} = \frac1X - \frac1{X+1}$.
The sum telescopes:
\[
\sum_{n=1}^{N}\Bigl(\frac1n - \frac1{n+1}\Bigr)
= 1 - \frac1{N+1} \longrightarrow 1 .
\]

\textbf{2.} $\frac1{X(X+2)} = \frac{1/2}X - \frac{1/2}{X+2}$.
Summing, the terms $\frac1n$ survive for $n = 1, 2$ and the terms
$-\frac1{n+2}$ survive for $n = N-1, N$:
\[
\sum_{n=1}^{N}\frac1{n(n+2)}
= \frac12\Bigl(1 + \frac12 - \frac1{N+1} - \frac1{N+2}\Bigr)
\longrightarrow \frac34 .
\]

\textbf{3.} If $F(X) = G(X) - G(X+1)$, then $\sum_{n=1}^N F(n) =
\sum_{n=1}^N\bigl(G(n) - G(n+1)\bigr) = G(1) - G(N+1)$: all
intermediate values cancel in pairs. For $F =
\frac1{X(X+1)(X+2)}$, the witness is $G(X) = \frac1{2X(X+1)}$:
\[
G(X) - G(X+1)
= \frac{(X+2) - X}{2X(X+1)(X+2)} = F(X) ,
\]
so $\sum_{n=1}^N F(n) = \frac14 - \frac1{2(N+1)(N+2)} \to
\frac14$, the value of \cref{exo:b1:fractions:8}.

\textbf{4.} Put the right-hand side over the common denominator
$X(X+1)\cdots(X+k)$:
\[
\frac1k\cdot\frac{(X + k) - X}{X(X+1)\cdots(X+k)}
= \frac{1}{X(X+1)\cdots(X+k)} ,
\]
which is the identity. So $F_k(X) = \frac1k\bigl(G_k(X) -
G_k(X+1)\bigr)$ with $G_k(X) = \frac1{X(X+1)\cdots(X+k-1)}$, and
question 3's mechanism gives
\[
\sum_{n=1}^{N}\frac1{n(n+1)\cdots(n+k)}
= \frac1k\Bigl(\frac1{k!} - G_k(N+1)\Bigr)
\longrightarrow \frac1{k\cdot k!} ,
\]
since $G_k(1) = \frac1{k!}$ and $G_k(N+1) \to 0$. For $k = 2$:
$\frac1{2\cdot2} = \frac14$, matching question 3.

\textbf{5.} \cref{exo:b1:fractions:6} gives
$\frac{n!}{X(X+1)\cdots(X+n)} = \sum_{j=0}^n
\frac{(-1)^j\binom nj}{X + j}$. Evaluate at $X = 1$: the left side
is $\frac{n!}{(n+1)!} = \frac1{n+1}$, the right side is
$\sum_j(-1)^j\binom nj\frac1{1+j}$: first identity. At $X = 2$:
the left side is $\frac{n!}{2\cdot3\cdots(n+2)} =
\frac{n!\cdot1}{(n+2)!} = \frac1{(n+1)(n+2)}$, the right side
$\sum_j(-1)^j\binom nj\frac1{2+j}$: second identity.

\textbf{6.} The \omterm{def:b1:fractions:field}{poles} $\omega^k$ are simple, and the \omterm{def:b1:fractions:field}{simple-pole}
formula of \cref{met:b1:fractions:compute} gives the coefficient
\[
\frac{1}{\bigl(nX^{n-1}\bigr)_{X = \omega^k}}
= \frac{1}{n\,\omega^{k(n-1)}}
= \frac{\omega^k}{n\,\omega^{kn}} = \frac{\omega^k}n ,
\]
using $\omega^{kn} = 1$. Hence $\frac1{X^n-1} = \frac1n\sum_k
\frac{\omega^k}{X - \omega^k}$.

\textbf{7.} For $n = 2$ ($\omega = -1$): $\frac12\bigl(
\frac1{X-1} - \frac1{X+1}\bigr) = \frac12\cdot\frac{2}{X^2-1} =
\frac1{X^2-1}$: correct. The coefficients sum to
$\frac1n\sum_k\omega^k = 0$ for $n \geq 2$
(\cref{prop:b1:complex:sumroots}). They must: $x\,F(x) \to
\sum_k c_k$ as $x \to \infty$ for any decomposition with simple
\omterm{def:b1:fractions:field}{poles}, while here $xF(x) = \frac{x}{x^n-1} \to 0$ since $n \geq
2$.

\textbf{8.} Cover-up for $\frac{X^{n-1}}{X^n - 1}$ at $\omega^k$:
$\frac{A(\omega^k)}{B'(\omega^k)} =
\frac{\omega^{k(n-1)}}{n\omega^{k(n-1)}} = \frac1n$, so
$\frac{X^{n-1}}{X^n-1} = \frac1n\sum_k\frac1{X - \omega^k}$ ---
the identity of \cref{exo:b1:fractions:7} again.

\textbf{9.} With $c = \omega^k$, $\conj c = \omega^{n-k}$ and
$c\conj c = 1$, $c + \conj c = 2\cos\theta_k$:
\[
\frac{c}{X - c} + \frac{\conj c}{X - \conj c}
= \frac{c(X - \conj c) + \conj c(X - c)}
{(X - c)(X - \conj c)}
= \frac{2\cos\theta_k\,X - 2}{X^2 - 2\cos\theta_k\,X + 1} .
\]
Grouping $k$ with $n - k$ in question 6: for odd $n$,
\[
\frac1{X^n - 1} = \frac1n\Biggl(\frac1{X-1} +
\sum_{k=1}^{(n-1)/2}
\frac{2\cos\theta_k X - 2}{X^2 - 2\cos\theta_k X + 1}\Biggr) ;
\]
for even $n$, the extra self-paired \omterm{def:b1:fractions:field}{pole} $\omega^{n/2} = -1$
contributes $\frac{-1}{X+1}$ inside the parenthesis and the pair
sum runs to $\frac n2 - 1$.

\textbf{10.} $n = 4$: $\theta_1 = \frac\pi2$, $\cos\theta_1 = 0$,
so the pair term is $\frac{-2}{X^2+1}$ and
\[
\frac1{X^4-1} = \frac14\Bigl(\frac1{X-1} - \frac1{X+1}
- \frac2{X^2+1}\Bigr) ,
\]
the decomposition of \cref{exo:b1:fractions:4}.

\textbf{11.} $P = \prod_{k=1}^{n-1}(X - \omega^k)$ (divide $X^n -
1$ by $X - 1$), so by \cref{exo:b1:poly:11}, $\frac{P'}P =
\sum_{k=1}^{n-1}\frac1{X-\omega^k}$. At $X = 1$: $P(1) = n$ and
$P'(1) = \sum_{j=1}^{n-1}j = \frac{n(n-1)}2$, whence
\[
\sum_{k=1}^{n-1}\frac1{1 - \omega^k}
= \frac{P'(1)}{P(1)} = \frac{n-1}2 .
\]

\textbf{12.} Half-angle: $1 - \eu^{\iu\theta} =
-2\iu\sin\frac\theta2\,\eu^{\iu\theta/2}$, so
\[
\frac1{1 - \eu^{\iu\theta}}
= \frac{\eu^{-\iu\theta/2}}{-2\iu\sin\frac\theta2}
= \frac{\iu\cos\frac\theta2 + \sin\frac\theta2}
{2\sin\frac\theta2}
= \frac12 + \frac\iu2\,
\frac{\cos\frac\theta2}{\sin\frac\theta2} .
\]
With $\theta = \theta_k = \frac{2k\pi}n$: $\frac1{1-\omega^k} =
\frac12 + \frac\iu2\cot\frac{k\pi}n$. Summing over $k = 1, \dots,
n-1$ and comparing with the real value $\frac{n-1}2$ of
question 11: the real parts already account for everything, so
$\sum_k\cot\frac{k\pi}n = 0$ --- as the symmetry $\cot\frac{(n-k)
\pi}n = -\cot\frac{k\pi}n$ also shows.

\textbf{13.} Differentiating $\frac{P'}P = \sum_k\frac1{X -
\omega^k}$:
\[
\frac{P''}P - \Bigl(\frac{P'}P\Bigr)^2
= -\sum_{k=1}^{n-1}\frac1{(X - \omega^k)^2} .
\]
At $X = 1$: $P''(1) = \sum_{j=2}^{n-1}j(j-1) = 2\binom n3 =
\frac{n(n-1)(n-2)}3$ (hockey-stick identity, or induction), so
\[
\sum_{k=1}^{n-1}\frac1{(1 - \omega^k)^2}
= \Bigl(\frac{n-1}2\Bigr)^2 - \frac{(n-1)(n-2)}3
= \frac{(n-1)\bigl(3(n-1) - 4(n-2)\bigr)}{12}
= \frac{(n-1)(5-n)}{12} .
\]

\textbf{14.} Squaring question 12's formula, with $c_k =
\cot\frac{k\pi}n$:
\[
\frac1{(1-\omega^k)^2}
= \Bigl(\frac12 + \frac\iu2 c_k\Bigr)^2
= \frac14 - \frac{c_k^2}4 + \frac\iu2\,c_k .
\]
Summing and using $\sum c_k = 0$ (question 12) and question 13:
$\frac{n-1}4 - \frac14\sum_k c_k^2 = \frac{(n-1)(5-n)}{12}$, so
\[
\sum_{k=1}^{n-1}\cot^2\frac{k\pi}n
= (n-1) - \frac{(n-1)(5-n)}3 = \frac{(n-1)(n-2)}3 .
\]
Then $\frac1{\sin^2t} = 1 + \cot^2t$ gives
$\sum_k\frac1{\sin^2\frac{k\pi}n} = (n-1) +
\frac{(n-1)(n-2)}3 = \frac{n^2-1}3$.

\textbf{15.} $n = 3$: $\cot^2\frac\pi3 + \cot^2\frac{2\pi}3 =
\frac13 + \frac13 = \frac23 = \frac{2\cdot1}3$; and
$\frac1{\sin^2} $ sums to $\frac43 + \frac43 = \frac83 =
\frac{9-1}3$. $n = 4$: $1 + 0 + 1 = 2 = \frac{3\cdot2}3$; and $2
+ 1 + 2 = 5 = \frac{16-1}3$. Both formulas check.

\textbf{16.} For $t \in \intoo0{\frac\pi2}$, the classical
comparison $\sin t < t < \tan t$ (area or convexity argument,
familiar from High School) yields, taking reciprocals, $\cot t <
\frac1t < \frac1{\sin t}$, all three positive there; squaring
preserves the order. With $t = \frac{k\pi}n$, $1 \leq k \leq m$,
$n = 2m+1$ (so $t < \frac\pi2$):
\[
\cot^2\frac{k\pi}n < \frac{n^2}{k^2\pi^2} <
\frac1{\sin^2\frac{k\pi}n} .
\]

\textbf{17.} By the symmetries $\cot^2\frac{(n-k)\pi}n =
\cot^2\frac{k\pi}n$ and likewise for $\sin^2$, question 14's sums
halve: $\sum_{k=1}^{m}\cot^2\frac{k\pi}n = \frac{(n-1)(n-2)}6 =
\frac{m(2m-1)}3$ and $\sum_{k=1}^m\frac1{\sin^2\frac{k\pi}n} =
\frac{n^2-1}6 = \frac{2m(m+1)}3$. Summing question 16 over $k =
1, \dots, m$ and multiplying by $\frac{\pi^2}{n^2}$:
\[
\frac{\pi^2}{(2m+1)^2}\cdot\frac{m(2m-1)}3
\;<\; \sum_{k=1}^{m}\frac1{k^2} \;<\;
\frac{\pi^2}{(2m+1)^2}\cdot\frac{2m(m+1)}3 .
\]
Both bounds tend to $\frac{\pi^2}6$ as $m \to \infty$ (the ratios
$\frac{m(2m-1)}{(2m+1)^2}$ and $\frac{2m(m+1)}{(2m+1)^2}$ both
tend to $\frac12$), so by the squeeze the increasing partial sums
converge and
\[
\sum_{k=1}^{\infty}\frac1{k^2} = \frac{\pi^2}6 .
\]

\textbf{18.} Square $\frac1{X^2-1} = \frac12\bigl(\frac1{X-1} -
\frac1{X+1}\bigr)$:
\[
\frac1{(X^2-1)^2} = \frac14\Bigl(\frac1{(X-1)^2} +
\frac1{(X+1)^2}\Bigr) - \frac12\cdot\frac1{(X-1)(X+1)} ,
\]
and re-decompose the cross term $\frac1{(X-1)(X+1)} =
\frac12\bigl(\frac1{X-1} - \frac1{X+1}\bigr)$ to get the stated
form. At $X = 0$: left side $1$; right side $\frac14(1 + 1) -
\frac14(-1 - 1) = \frac12 + \frac12 = 1$.

\textbf{19.} Sum question 18 over $n \geq 2$. With $S =
\sum_{k\geq1}\frac1{k^2} = \frac{\pi^2}6$: $\sum_{n\geq2}
\frac1{(n-1)^2} = S$; $\sum_{n\geq2}\frac1{(n+1)^2} = S - 1 -
\frac14$; and the gap-two telescope $\sum_{n\geq2}\bigl(
\frac1{n-1} - \frac1{n+1}\bigr) = 1 + \frac12 = \frac32$. Hence
\[
\sum_{n=2}^{\infty}\frac1{(n^2-1)^2}
= \frac14\Bigl(2S - \frac54\Bigr) - \frac14\cdot\frac32
= \frac S2 - \frac{11}{16}
= \frac{\pi^2}{12} - \frac{11}{16} \approx 0.135 .
\]
Numerically: $\frac19 + \frac1{64} + \frac1{225} + \frac1{576} +
\frac1{1225} + \dots \approx 0.1111 + 0.0156 + 0.0044 + 0.0017 +
0.0008 + \dots \approx 0.135$: consistent.

\textbf{20.} Differentiate question 8's identity
$\sum_k\frac1{X-\omega^k} = \frac{nX^{n-1}}{X^n-1}$:
\[
\sum_{k=0}^{n-1}\frac1{(X - \omega^k)^2}
= \frac{n^2X^{2n-2} - n(n-1)X^{n-2}(X^n - 1)}{(X^n - 1)^2} .
\]
Check at $n = 2$, $X = 2$: right side $\frac{4\cdot4 -
2\cdot1\cdot3}{9} = \frac{10}9$; left side $\frac1{(2-1)^2} +
\frac1{(2+1)^2} = 1 + \frac19 = \frac{10}9$.

\textbf{21.} $\frac1{\binom nk} = \frac{k!\,(n-k)!}{n!} =
\frac{k!}{(n-k+1)(n-k+2)\cdots n}$, a product of $k$ consecutive
integers downstairs. Substituting $j = n - k + 1$ (so $j$ runs
over $\N^*$ as $n$ runs from $k$):
\[
\sum_{n=k}^{\infty}\frac1{\binom nk}
= k!\sum_{j=1}^{\infty}\frac1{j(j+1)\cdots(j+k-1)}
= \frac{k!}{(k-1)\,(k-1)!} = \frac{k}{k-1} ,
\]
by question 4 applied with $k - 1$ in place of $k$ (valid since
$k - 1 \geq 1$).

\textbf{22.} For $k = 3$: $\frac1{\binom n3} =
\frac6{(n-2)(n-1)n}$, so by question 3's partial sum (shifted),
\[
\sum_{n=3}^{N}\frac1{\binom n3}
= 6\sum_{j=1}^{N-2}\frac1{j(j+1)(j+2)}
= 6\Bigl(\frac14 - \frac1{2(N-1)N}\Bigr)
= \frac32 - \frac3{(N-1)N}
\longrightarrow \frac32 ,
\]
the value $\frac k{k-1} = \frac32$ of question 21.

\textbf{23.} $n = 6$: pairs $k = 1$ ($\theta_1 = \frac\pi3$,
$2\cos\theta_1 = 1$) and $k = 2$ ($\theta_2 = \frac{2\pi}3$,
$2\cos\theta_2 = -1$), plus the real \omterm{def:b1:fractions:field}{poles} $\pm1$:
\[
\frac1{X^6-1} = \frac16\Bigl(\frac1{X-1} - \frac1{X+1}
+ \frac{X - 2}{X^2 - X + 1}
+ \frac{-X - 2}{X^2 + X + 1}\Bigr) .
\]
At $X = 0$: left side $-1$; right side $\frac16(-1 - 1 - 2 - 2)
= -1$: correct.

\textbf{24.} (i) Uniqueness legitimizes every identification of
coefficients --- cover-up, the conjugate-pair grouping of
question 9, and the differentiation tricks (questions 13, 20)
all rest on it. (ii) \omterm{def:b1:complex:unity}{Roots of unity} supplied the \omterm{def:b1:fractions:field}{poles} of $X^n -
1$, their symmetries ($k \leftrightarrow n-k$), and the
half-angle algebra of question 12
(\cref{met:b1:complex:trig}). (iii) The logarithmic derivative
$\frac{P'}P = \sum\frac{m_i}{X-a_i}$ converted information about
the roots of $P = 1 + X + \dots + X^{n-1}$ into the numerical
sums of questions 11 and 13 --- the hinge between Parts II and
III.

\textbf{25.} The decomposition is a purely algebraic identity,
true for every value of the variable at once; that is what makes
it an engine. Substituting integers and summing turned it into
telescopes (Part I); substituting \omterm{def:b1:complex:unity}{roots of unity} and grouping
conjugates turned it into trigonometric identities (Parts II and
III); and only at the very last step did analysis enter --- a
squeeze between two closed forms --- to deliver $\frac{\pi^2}6$,
a \omterm{def:b1:logic:statement}{statement} no finite substitution could reach. This division of
labor (algebra produces exact finite identities, analysis passes
to the limit) is the template for \cref{ch:b1:series}, where
telescoping and comparison become systematic, and for
\cref{ch:b1:integration}, where each brick acquires a primitive
and the same decompositions compute integrals instead of sums.
\end{solution}

\section*{Chapter \ref{ch:b1:reals} --- Real Numbers}

\begin{solution}{exo:b1:reals:1}
$A$: every element is $\leq 1$ and $1 \in A$: $\sup A = \max A = 1$.
Lower bounds: $0$ bounds below; for $\varepsilon > 0$, Archimedes
provides $n$ with $\frac 1n < \varepsilon$, so no positive number
bounds $A$ below: $\inf A = 0$, not attained (no min).

$B$: terms $0, -\frac12, \frac23, -\frac34, \frac45, \dots$ The even
terms $\frac{n}{n+1}$ ($n$ even) increase to $1$ without reaching it;
the odd terms $-\frac{n}{n+1}$ decrease to $-1$. So $\sup B = 1$ and
$\inf B = -1$, neither attained: no max, no min. (Bounds: $\abs{b}
< 1$ for all $b \in B$; and $\frac{n}{n+1} = 1 - \frac{1}{n+1} > 1 -
\varepsilon$ for $n$ large, similarly below.)

$C = \intoo{-\sqrt 3}{\sqrt 3}$: $\sup C = \sqrt 3$, $\inf C =
-\sqrt 3$, neither attained.
\end{solution}

\begin{solution}{exo:b1:reals:2}
Write $x = \lfloor x \rfloor + u$, $y = \lfloor y \rfloor + v$ with
$u, v \in \intco{0}{1}$. Then $x + y = \lfloor x \rfloor + \lfloor y
\rfloor + (u + v)$ with $u + v \in \intco{0}{2}$. If $u + v < 1$,
$\lfloor x + y\rfloor = \lfloor x\rfloor + \lfloor y \rfloor$; if
$1 \leq u + v < 2$, $\lfloor x+y \rfloor = \lfloor x \rfloor +
\lfloor y \rfloor + 1$. Both cases occur: $(x, y) = (0.2,\, 0.3)$
gives the left equality, $(0.7,\, 0.8)$ the right one.
\end{solution}

\begin{solution}{exo:b1:reals:3}
Let $k = \lfloor x \rfloor$, so $k \leq x < k + 1$. Multiplying by
$n$: $nk \leq nx < nk + n$, and taking floors (an increasing
operation on integers' side): $nk \leq \lfloor nx \rfloor \leq nk + n
- 1$. Dividing by $n$: $k \leq \frac{\lfloor nx \rfloor}{n} < k + 1$,
so the outer floor is $k$.
\end{solution}

\begin{solution}{exo:b1:reals:4}
Every element of $A$ is in $B$, so $\sup B$ bounds $A$ above:
therefore $\sup A \leq \sup B$ ($\sup A$ is the \emph{least} upper
bound). Symmetrically $\inf B \leq \inf A$. Finally $\inf A \leq
\sup A$ because $A$ is nonempty: any $a \in A$ sits between them.
\end{solution}

\begin{solution}{exo:b1:reals:5}
Let $s = \sup A$, $t = \sup B$. Every $a + b \leq s + t$: upper
bound. For $\varepsilon > 0$, choose $a > s - \frac\varepsilon2$ and
$b > t - \frac\varepsilon2$ (\cref{prop:b1:reals:epsilon}): then $a +
b > s + t - \varepsilon$. By the $\varepsilon$-characterization,
$\sup(A+B) = s + t$.

For $-A$: $m$ bounds $-A$ above $\iff$ $-m$ bounds $A$ below; the
least upper bound of $-A$ therefore corresponds to the greatest lower
bound of $A$: $\sup(-A) = -\inf A$.
\end{solution}

\begin{solution}{exo:b1:reals:6}
For every $x$: $f(x) + g(x) \leq \sup f + \sup g$; taking the sup of
the left side gives the inequality. Strict example: $E = \{0, 1\}$,
$f = \mathbf{1}_{\{0\}}$ (value $1$ at $0$, else $0$), $g =
\mathbf{1}_{\{1\}}$: $\sup(f + g) = 1 < 2 = \sup f + \sup g$.

No contradiction with \cref{exo:b1:reals:5}: there, $a \in A$ and $b
\in B$ vary \emph{independently}; here the same $x$ feeds both $f$
and $g$ --- the set $\{f(x) + g(x) : x \in E\}$ is smaller than the
set $\{f(x) + g(y) : x, y \in E\}$.
\end{solution}

\begin{solution}{exo:b1:reals:7}
$\sqrt 6$ is irrational: $6 = 2 \times 3$ is not a perfect square,
and $v_2(6q^2) = 1 + 2v_2(q)$ odd prevents $6q^2 = r^2$ (as in
\cref{exo:b1:arith:7}). Now suppose $x = \sqrt 2 + \sqrt 3 \in \Q$.
Then $x^2 = 5 + 2\sqrt 6 \in \Q$, so $\sqrt 6 = \frac{x^2 - 5}{2} \in
\Q$: contradiction. Hence $\sqrt 2 + \sqrt 3 \notin \Q$.
\end{solution}

\begin{solution}{exo:b1:reals:8}
Let $x < y$. Pick $n \in \N$ with $2^n > \frac{1}{y - x}$ (Archimedes:
$2^n \geq n + 1$ by an easy induction, so some power of $2$ exceeds
any real). Then, as in the proof of \cref{thm:b1:reals:density} with
$2^n$ in place of $n$: $m = \lfloor 2^n x \rfloor + 1$ satisfies $x <
\frac{m}{2^n} < y$. So $D$ is dense.
\end{solution}

\begin{solution}{exo:b1:reals:9}
\begin{enumerate}
  \item Suppose $\alpha > 0$. First, $\alpha \in G$. Suppose not: by
        the $\varepsilon$-characterization of the infimum with
        $\varepsilon = \alpha$, there is $g \in G$ with $\alpha < g <
        2\alpha$ (strict on the left since $\alpha \notin G$); then,
        with $\varepsilon = g - \alpha$, there is $h \in G$ with
        $\alpha < h < g$. Now $g - h \in G$ and $0 < g - h < g -
        \alpha < \alpha$: an element of $G \cap \intoo{0}{+\infty}$
        below its infimum, absurd. So $\alpha \in G$, and $\alpha\Z
        \subseteq G$ ($G$ is a group). Conversely, for $x \in G$, let
        $k = \lfloor x/\alpha \rfloor$: then $x - k\alpha \in G$ and
        $0 \leq x - k\alpha < \alpha$, and the definition of $\alpha$
        forces $x - k\alpha = 0$. Hence $G = \alpha\Z$.
  \item Suppose $\alpha = 0$, and let $x < y$. There is $g \in G$
        with $0 < g < y - x$. The multiple $kg$ with $k = \lfloor x/g
        \rfloor + 1$ satisfies $x < kg \leq x + g < y$, and $kg \in
        G$: density.
\end{enumerate}
$G = \Z + \sqrt 2\,\Z$ is a subgroup of $(\R, +)$. It is not of the
form $\alpha\Z$: otherwise $1 = p\alpha$ and $\sqrt 2 = q\alpha$
($p, q \in \Z$) would give $\sqrt 2 = \frac qp \in \Q$,
contradiction. By the dichotomy, $G$ is dense in $\R$.
\end{solution}

\begin{solution}{exo:b1:reals:10}
Let $s = \sup A > 0$, $t = \sup B > 0$. For $a \in A$, $b \in B$: $ab
\leq st$ (multiplying inequalities between \emph{positive} numbers).
For $0 < \varepsilon < \min(s, t)$: choose $a > s - \varepsilon$ and
$b > t - \varepsilon$; then
\[
ab > (s - \varepsilon)(t - \varepsilon)
= st - \varepsilon(s + t) + \varepsilon^2
> st - \varepsilon (s + t),
\]
and $\varepsilon(s+t)$ can be made arbitrarily small: by the
$\varepsilon$-characterization (in the form: no number $< st$ bounds
$AB$ above), $\sup AB = st$.

Positivity is essential: $A = B = \{-1, 0\}$ gives $AB = \{0, 1\}$,
$\sup AB = 1$, while $\sup A \cdot \sup B = 0 \times 0 = 0$.
\end{solution}

\begin{solution}{exo:b1:reals:11}
Write $s = \sup A$, $i = \inf A$. For $a, a' \in A$: $a \leq s$
and $a' \geq i$ give $a - a' \leq s - i$; by symmetry $\abs{a -
a'} \leq s - i$, so $s - i$ bounds the set of gaps above. For
$\varepsilon > 0$, choose $a > s - \frac\varepsilon2$ and $a' < i
+ \frac\varepsilon2$ (\cref{prop:b1:reals:epsilon} and its mirror
for the infimum): then $\abs{a - a'} \geq a - a' > s - i -
\varepsilon$. By the $\varepsilon$-characterization,
$\operatorname{diam} A = s - i$.

Every $a \in A$ satisfies $i \leq a \leq s$, so $A \subseteq
\intcc{i}{s}$, a closed interval of length $\operatorname{diam}
A$. If a closed interval $\intcc{u}{v}$ contains $A$, then $v$ is
an upper bound and $u$ a lower bound of $A$, so $u \leq i$ and $v
\geq s$: $\intcc{i}{s} \subseteq \intcc{u}{v}$. Hence
$\intcc{i}{s}$ is the smallest one.
\end{solution}

\begin{solution}{exo:b1:reals:12}
$0 \in E$, so $E \neq \emptyset$. If $x > \max(1, y)$ then $x^2 >
x > y$, so $E$ is bounded above by $\max(1, y)$: $s = \sup E$
exists (\cref{thm:b1:reals:sup}), and $s \geq \min(1, y) > 0$
because $\min(1, y) \in E$: indeed if $y \geq 1$ then $1^2 = 1
\leq y$, and if $y < 1$ then $y^2 < y$.

\emph{$s^2 < y$ is impossible.} Choose $0 < h < 1$ with $h <
\frac{y - s^2}{2s + 1}$. Then
\[
(s + h)^2 = s^2 + 2sh + h^2 \leq s^2 + (2s + 1)h < y ,
\]
so $s + h \in E$, contradicting that $s$ bounds $E$ above.

\emph{$s^2 > y$ is impossible.} Choose $0 < h < s$ with $h <
\frac{s^2 - y}{2s}$. Then $(s - h)^2 = s^2 - 2sh + h^2 > s^2 -
2sh > y$; every $x \in E$ satisfies $x^2 \leq y < (s - h)^2$,
hence $x < s - h$ (both are $\geq 0$): $s - h$ is an upper bound
of $E$ smaller than $s$, contradicting leastness.

Therefore $s^2 = y$. Uniqueness: if $0 < s < s'$ then $s^2 <
s'^2$, so two distinct positive roots cannot both square to $y$.
Monotonicity: if $0 < y < y'$, then $\sqrt y \neq \sqrt{y'}$, and
$\sqrt y > \sqrt{y'}$ would give $y > y'$ by squaring: so $\sqrt
y < \sqrt{y'}$.
\end{solution}

\begin{solution}{pb:b1:reals:1}
\textbf{1.} For $\frac 18$: $10 = 8 \cdot 1 + 2$, $20 = 8 \cdot 2
+ 4$, $40 = 8 \cdot 5 + 0$; digits $1, 2, 5$, remainder $0$, then
only zeros: $\frac 18 = 0.125$. For $\frac 17$: $10 = 7 \cdot 1 +
3$, $30 = 7 \cdot 4 + 2$, $20 = 7 \cdot 2 + 6$, $60 = 7 \cdot 8 +
4$, $40 = 7 \cdot 5 + 5$, $50 = 7 \cdot 7 + 1$: digits $1, 4, 2,
8, 5, 7$, remainders $3, 2, 6, 4, 5, 1$. The remainder has
returned to $r = 1$, so the six steps repeat verbatim forever:
$\frac 17 = 0.(142857)$, remainders cycling through $1, 3, 2, 6,
4, 5$.

\textbf{2.} $\frac 13$ in base $2$ ($r_0 = 1$): $2 = 3 \cdot 0 +
2$, $4 = 3 \cdot 1 + 1$, and $r = 1$ recurs: $\frac 13 =
(0.\overline{01})_2$. $\frac{5}{16}$ in base $2$: $10 = 16 \cdot
0 + 10$, $20 = 16 \cdot 1 + 4$, $8 = 16 \cdot 0 + 8$, $16 = 16
\cdot 1 + 0$: $\frac{5}{16} = (0.0101)_2$, terminating. $\frac 12$
in base $3$: $3 = 2 \cdot 1 + 1$, and $r = 1$ recurs at once:
$\frac 12 = (0.\overline{1})_3$.

\textbf{3.} The algorithm's remainder after $N$ steps is $r_N =
b^N p \bmod q$ (proved formally in question 14; here it is the
observation that each step multiplies the remainder by $b$ and
reduces mod $q$). All later digits are $0$ if and only if some
$r_N = 0$, i.e.\ $q \mid b^N p$; since $\gcd(p, q) = 1$, Gauss's
lemma gives $q \mid b^N$. If $q \mid b^N$, every prime factor of
$q$ divides $b^N$, hence divides $b$ (primality). Conversely, if
every prime of $q = p_1^{a_1} \cdots p_r^{a_r}$ divides $b$, then
with $A = \max_i a_i$ each $p_i^{a_i}$ divides $b^A$, and the
$p_i^{a_i}$ are pairwise coprime, so $q \mid b^A$. For $q = 20 =
2^2 \cdot 5$: both primes divide $10$ ($\frac{1}{20} = 0.05$),
but $2 \nmid 3$, so $\frac{1}{20}$ repeats forever in base $3$.

\textbf{4.} $1^2 = 1 < 2 < 4 = 2^2$ gives $s_0 = 1$. Then $1.4^2
= 1.96 < 2 < 2.25 = 1.5^2$: $\lfloor 10\sqrt 2 \rfloor = 14$,
$s_1 = 1.4$. Next $1.41^2 = 1.9881 < 2 < 2.0164 = 1.42^2$: $s_2 =
1.41$; $1.414^2 = 1.999396 < 2 < 2.002225 = 1.415^2$: $s_3 =
1.414$; $1.4142^2 = 1.99996164 < 2 < 2.00024449 = 1.4143^2$: $s_4
= 1.4142$. In each case the displayed inequalities say exactly
$s_n \leq \sqrt 2 < s_n + 10^{-n}$, which is the definition of
the floor of $10^n \sqrt 2$.

\textbf{5.} $A_0 = \lfloor x \rfloor = 0$ since $0 \leq x < 1$.
From $A_{n-1} \leq b^{n-1} x < A_{n-1} + 1$, multiply by $b$:
\[
b\,A_{n-1} \leq b^n x < b\,A_{n-1} + b .
\]
The integer $b\,A_{n-1}$ is $\leq b^n x$, so $b\,A_{n-1} \leq
A_n$; and $b^n x < b\,A_{n-1} + b$ with $b\,A_{n-1} + b$ an
integer forces $A_n \leq b\,A_{n-1} + b - 1$. Hence $0 \leq d_n =
A_n - b\,A_{n-1} \leq b - 1$: a digit.

\textbf{6.} Telescoping: $d_k b^{-k} = A_k b^{-k} - A_{k-1}
b^{-(k-1)}$, so
\[
\sum_{k=1}^n d_k b^{-k} = A_n b^{-n} - A_0 = s_n .
\]
Dividing $A_n \leq b^n x < A_n + 1$ by $b^n$ gives $s_n \leq x <
s_n + b^{-n}$.

\textbf{7.} Induction: $b^0 = 1 \geq 1$, and $b^{n+1} = b \cdot
b^n \geq 2(n + 1) \geq n + 2$. Monotonicity: $s_n - s_{n-1} = d_n
b^{-n} \geq 0$. Each $s_n \leq x$ (question 6): $x$ is an upper
bound of $\{s_n\}$. For $\varepsilon > 0$, the Archimedean
property provides $n$ with $n + 1 > \frac1\varepsilon$, hence
$b^{-n} < \varepsilon$, and then $s_n > x - b^{-n} > x -
\varepsilon$ by question 6. By \cref{prop:b1:reals:epsilon}, $x =
\sup_n s_n$.

\textbf{8.} Suppose $d_k = b - 1$ for all $k > N$. For $n > N$,
the finite geometric sum gives
\[
s_n = s_N + (b - 1)\sum_{k=N+1}^{n} b^{-k}
= s_N + b^{-N} - b^{-n} .
\]
So $x \geq s_n = s_N + b^{-N} - b^{-n}$ for every $n$; letting
the last term shrink below any $\varepsilon$ (question 7), $x
\geq s_N + b^{-N}$. But question 6 at rank $N$ says $x < s_N +
b^{-N}$: contradiction. The string $(d_n)$ is proper.

\textbf{9.} Bounded: $t_n \leq (b-1)\sum_{k=1}^n b^{-k} = 1 -
b^{-n} < 1$, and $(t_n)$ is nondecreasing, so $y = \sup t_n$
exists with $0 \leq y \leq 1$. Fix $n$. For the two-sided
estimate: $t_n \leq y$ is clear. By properness pick $m > n$ with
$e_m \leq b - 2$. For $p \geq m$:
\[
t_p - t_n = \sum_{k=n+1}^{p} e_k b^{-k}
\leq (b^{-n} - b^{-p}) - b^{-m} < b^{-n} - b^{-m},
\]
the middle sum losing at least $b^{-m}$ against the all-$(b-1)$
maximum; for $p \leq m$, $t_p \leq t_m \leq t_n + b^{-n} -
b^{-m}$ as well (monotonicity plus the case $p = m$). Hence every
$t_p \leq t_n + b^{-n} - b^{-m}$, so $y \leq t_n + b^{-n} -
b^{-m} < t_n + b^{-n}$. (With $n = 0$: $y < 1$, so $y \in
\intco{0}{1}$.) Now $b^n t_n = \sum_{k \leq n} e_k b^{n-k}$ is an
integer, and $b^n t_n \leq b^n y < b^n t_n + 1$: so $\lfloor b^n
y \rfloor = b^n t_n$. Finally the digits of $y$: $d_n(y) = b^n
t_n - b \cdot b^{n-1} t_{n-1} = b^n(t_n - t_{n-1}) = e_n$.

\textbf{10.} Question 9 says: (value of string) has digits (the
string); questions 5--8 say: (digits of $x$) form a proper string
whose truncations have supremum $x$ (question 7). So the two maps
compose to the identity in both orders: they are mutually inverse
bijections between $\intco{0}{1}$ and the proper strings. If two
proper strings had equal value, applying the digit map would make
them equal: uniqueness. This is the \emph{$b$-adic expansion
theorem}.

\textbf{11.} Let $e_n = b - 1$ for $n > M$, $M \geq 0$ minimal.
As in question 8, $t_n = t_M + b^{-M} - b^{-n}$ for $n \geq M$,
so the value is $\sup t_n = t_M + b^{-M}$. If $M = 0$ the value
is $0 + 1 = 1$: in base $10$, $0.999\dots = 1$ exactly --- not
approximately. If $M \geq 1$, minimality gives $e_M \leq b - 2$,
and the value is
\[
t_M + b^{-M} = \frac{b^M t_M + 1}{b^M} \in \intoo{0}{1},
\]
a $b$-adic fraction, whose \emph{proper} expansion is $e_1 \dots
e_{M-1}\,(e_M + 1)\,000\dots$ (a terminating string is proper,
and its value is the same number). Conversely a real with two
representations must have one improper (properness pins down the
representation, question 10), hence be of this form. And each
$m/b^N \in \intoo{0}{1}$, written with last nonzero digit $d_N$,
does have the improper twin $d_1 \dots d_{N-1}(d_N -
1)(b-1)(b-1)\dots$: exactly the $b$-adic fractions carry two
names, all other reals one.

\textbf{12.} Say the strings agree up to $m - 1$, with common
truncation $P = s_{m-1}$, and $d_m < e_m$. By question 9 (strict
upper estimate at rank $m$), $x < P + d_m b^{-m} + b^{-m} = P +
(d_m + 1)b^{-m} \leq P + e_m b^{-m} \leq y$, the last step
because $P + e_m b^{-m}$ is $y$'s truncation $t_m \leq y$. So
$d_m < e_m \implies x < y$; exchanging roles, $e_m < d_m \implies
y < x$; and since the strings differ at $m$, one of the two
holds. Both directions follow.

\textbf{13.} Let $z = bx - A_1 \in \intco{0}{1}$ (indeed $A_1
\leq bx < A_1 + 1$). For $n \geq 0$: $b^n z = b^{n+1} x - b^n
A_1$ with $b^n A_1 \in \Z$, so by $\lfloor u - K \rfloor =
\lfloor u \rfloor - K$ ($K$ integer),
\[
A_n(z) = A_{n+1}(x) - b^n A_1(x) .
\]
Hence $d_n(z) = A_n(z) - b\,A_{n-1}(z) = A_{n+1} - b^n A_1 -
b\,A_n + b^n A_1 = d_{n+1}(x)$. So the fractional part of $bx$
carries the shifted digits; iterating $m$ times, the fractional
part of $b^m x$ has digits $(d_{n+m})_{n \geq 1}$.

\textbf{14.} Euclidean division: $b^n p = q\,Q_n + r_n$ with $0
\leq r_n < q$. Divide by $q$: $b^n x = Q_n + \frac{r_n}{q}$ with
$0 \leq \frac{r_n}{q} < 1$, so $Q_n = \lfloor b^n x \rfloor =
A_n$, giving $A_n = \frac{b^n p - r_n}{q}$. For the recurrence:
$b^n p = b(q\,A_{n-1} + r_{n-1}) = q\,(b\,A_{n-1}) +
b\,r_{n-1}$, so $b^n p$ and $b\,r_{n-1}$ differ by a multiple of
$q$: $r_n = (b\,r_{n-1}) \bmod q$.

\textbf{15.} Divide $b\,r_{n-1}$ by $q$: $b\,r_{n-1} = q\,c +
r_n$ with $c = \lfloor b\,r_{n-1}/q \rfloor$. Substituting into
the display of question 14: $b^n p = q(b\,A_{n-1} + c) + r_n$,
and uniqueness of Euclidean division identifies $A_n = b\,A_{n-1}
+ c$, that is $d_n = c = \lfloor b\,r_{n-1}/q \rfloor$. Digit
$n$ depends only on $r_{n-1}$ --- the long division loop of Part
I, now certified.

\textbf{16.} The $q + 1$ remainders $r_0, \dots, r_q$ take values
in the $q$-element set $\intint{0}{q-1}$: by the pigeonhole
principle (\cref{cor:b1:counting:pigeonhole}) two coincide, say
$r_N = r_{N+T}$ with $0 \leq N < N + T \leq q$. Since $r_n$
determines $r_{n+1}$ (question 14), induction gives $r_{n+T} =
r_n$ for all $n \geq N$; since $r_{n-1}$ determines $d_n$
(question 15), $d_{n+T} = d_n$ for all $n \geq N + 1$. Every
rational's expansion is eventually periodic, with preperiod $\leq
q$ and period $\leq q$.

\textbf{17.} The digits of the fractional part of $b^T y$ are
$(d_{n+T}) = (d_n)$ (shift lemma, then pure periodicity): the
same proper string as $y$. By question 10 the values are equal:
$b^T y - \lfloor b^T y \rfloor = y$, so $(b^T - 1)\,y = \lfloor
b^T y \rfloor = A_T \in \N$ and
\[
y = \frac{A_T}{b^T - 1} ,
\]
rational with denominator dividing $b^T - 1$; the numerator $A_T$
is the integer whose base-$b$ digits are $d_1 \dots d_T$. Check:
$0.(142857) = \frac{142857}{999999}$, and $142857 \times 7 =
999999$, so this is $\frac 17$.

\textbf{18.} If $d_{n+T} = d_n$ for $n > N$, the fractional part
$z$ of $b^N x$ has digits $(d_{N+n})_{n\geq1}$ (shift lemma),
which are purely periodic; by question 17, $z \in \Q$. Then $b^N
x = A_N + z$ gives $x = (A_N + z)/b^N \in \Q$. With question 16:
$x$ rational $\iff$ expansion eventually periodic. The right side
mentions the base, the left does not: periodicity in one base is
equivalent to rationality, hence to periodicity in every base.

\textbf{19.} For $x = \frac 1q$, $r_n = b^n \bmod q$. If
$\gcd(b, q) = 1$, then $r_T = r_0 = 1$ if and only if $b^T \equiv
1 \pmod q$; such $T$ exists (pigeonhole gives $b^i \equiv b^j$,
$i < j$, and $b$ is invertible modulo $q$, so $b^{j-i} \equiv
1$), and the least one --- the multiplicative order --- makes the
remainders, hence the digits, purely periodic of period $T$.
No smaller period is possible: a period $T'$ would give $(b^{T'}
- 1)\frac1q \in \N$ (question 17), i.e.\ $q \mid b^{T'} - 1$.
For $q = 7$, $b = 10$: $10 \equiv 3$, $10^2 \equiv 2$, $10^3
\equiv 6$, $10^4 \equiv 4$, $10^5 \equiv 5$, $10^6 \equiv 1
\pmod 7$: order $6$, and indeed $\frac 17$ has period six.

\textbf{20.} The string has infinitely many $0$s ($x^*$'s digits
are mostly zero), so it is proper, and $x^*$ is well defined
(question 9). Suppose the digits eventually periodic with period
$T$ beyond $N$. Infinitely many digits equal $1$ (one per
triangular number), so some $1$ sits at a position $j > N$; then
periodicity puts a $1$ at every position $j + kT$: from $j$
onward, gaps between consecutive $1$s are at most $T$. But the
$1$s sit exactly at the triangular numbers, whose consecutive
gaps $\frac{(j+1)(j+2)}{2} - \frac{j(j+1)}{2} = j + 1$ exceed $T$
eventually: contradiction. Not eventually periodic, so by
question 18, $x^* \notin \Q$ --- irrationality read off the
rhythm of the digits alone.

\textbf{21.} Given $x < y$, question 7 provides $n$ with $b^{-n}
< y - x$; set $m = \lfloor b^n x \rfloor + 1$. Then $b^n x < m
\leq b^n x + 1 < b^n y$, so $x < \frac{m}{b^n} < y$: density,
for every base at once ($b = 2$ recovers
\cref{exo:b1:reals:8}). For $\frac pq \in \intoo{0}{1}$ in base
$b = q$: the first digit is $\lfloor q \cdot \frac pq \rfloor =
p$ and the fractional part of $q \cdot \frac pq = p$ is $0$: all
later digits vanish, a terminating expansion $\frac pq =
(0.p)_q$. Terminating depends on the base; periodicity ---
rationality --- does not (question 18).

\textbf{22.} Each $e_k \in \{1, 2\}$ is a digit of base $10$, and
the string never ends in all $9$s: proper. Its value $y$ lies in
$\intco{0}{1}$ and has digits exactly $(e_k)$ (question 9). Fix
$k$: the $k$-th digit of $y$ is $e_k$, chosen $\neq$ the $k$-th
digit of $x_k$, so the proper strings of $y$ and $x_k$ differ, so
$y \neq x_k$ (question 10: the coding is injective). Thus $y$ is
in no list: no map $\N^* \to \intco{0}{1}$ is surjective. The
reals, unlike the rationals, cannot be enumerated ---
uncountability, whose theory \cref{ch:b1:topology} develops.

\textbf{23.} If the expansions agree up to $n$, then $x$ and $y$
have the same truncation $s_n$, and question 6 puts both in
$\intco{s_n}{s_n + b^{-n}}$, an interval of length $b^{-n}$:
$\abs{x - y} < b^{-n}$. Converse: $x = 0.1$ and $y = 0.0999$
(terminating, hence proper) satisfy $\abs{x - y} = 10^{-4} <
10^{-3}$, yet their expansions differ at the very first digit.
The culprits are the $b$-adic fractions of question 11: near
them, a tiny move flips every displayed digit ($0.0999 \to
0.1000$), because they are precisely the reals where the
improper twin lurks.

\textbf{24.} $p = 1$, $q = 10$, $b = 2$, $r_0 = 1$: $2 = 10 \cdot
0 + 2$, $4 = 10 \cdot 0 + 4$, $8 = 10 \cdot 0 + 8$, $16 = 10
\cdot 1 + 6$, $12 = 10 \cdot 1 + 2$ --- and $r_5 = 2 = r_1$: the
remainders cycle $(2, 4, 8, 6)$ from index $1$. Digits: $d_1 =
0$, then the repeating block $d_2 d_3 d_4 d_5 = 0, 0, 1, 1$:
\[
\tfrac{1}{10} = (0.0\overline{0011})_2 ,
\]
preperiod $1$, period $4$. By question 3, a terminating base-$2$
expansion would need every prime of $10$ to divide $2$; the prime
$5$ refuses. So $0.1$ is \emph{not} representable by any finite
binary string --- a computer storing finitely many bits keeps
only a truncation, and the accumulated truncation errors are why
floating-point $0.1 + 0.2$ differs from $0.3$ in the last bits.

\textbf{25.} (i) Completeness produced the values: $x = \sup s_n$
and $y = \sup t_n$ (questions 7 and 9) --- over $\Q$ alone, the
proper string of $\sqrt 2$ would name nothing. (ii) The
Archimedean property made $b^{-n}$ eventually smaller than any
$\varepsilon$, forcing the truncations to close in on their
supremum (questions 7, 21). (iii) The uniqueness clause of the
floor identified $Q_n = A_n$ in question 14 and legitimized every
digit extraction $\lfloor u - K \rfloor = \lfloor u \rfloor - K$
(question 13). (iv) The pigeonhole principle, applied to finitely
many remainders, is the sole engine of periodicity (question 16).
Moral: proper strings code $\intco{0}{1}$ faithfully and turn
rationality into a visible rhythm; but addition of digit strings
requires carries propagating from infinitely far right, so no
finite-stage rule computes even the first digit of a sum ---
whereas the supremum interface of \cref{thm:b1:reals:sup} handles
all of analysis with one axiom. Digits are a magnificent
\emph{picture} of $\R$; the supremum is its \emph{engine}.
\end{solution}

\section*{Chapter \ref{ch:b1:seq} --- Sequences}

\begin{solution}{exo:b1:seq:1}
$\Bigl|\dfrac{2n+1}{n+3} - 2\Bigr| = \dfrac{5}{n+3}$. Given
$\varepsilon > 0$, take $N > \frac 5\varepsilon - 3$ (Archimedes):
for $n \geq N$, $\frac{5}{n+3} \leq \varepsilon$. Hence the limit is
$2$.

$((-1)^n)$: its \omterm{def:b1:seq:subsequence}{subsequences} $(u_{2n}) = (1)$ and $(u_{2n+1}) =
(-1)$ converge to different limits, so the sequence diverges
(\cref{prop:b1:seq:subsequences}). (Directly: any candidate $\ell$
fails for $\varepsilon = \frac12$, since consecutive terms are at
distance $2$.)
\end{solution}

\begin{solution}{exo:b1:seq:2}
Dividing by $n^2$: $\dfrac{1 - 3/n + 1/n^2}{2 + 5/n^2} \to
\dfrac12$.

$\sqrt{n+1} - \sqrt n = \dfrac{1}{\sqrt{n+1} + \sqrt n} \to 0$
(conjugate).

$\dfrac{2^n + n^3}{3^n - n^2} = \dfrac{(2/3)^n + n^3/3^n}{1 -
n^2/3^n} \to \dfrac{0 + 0}{1 - 0} = 0$, using $q^n \to 0$ for $\abs q
< 1$ and the polynomial-vs-geometric comparison
(\cref{prop:b1:functions:powerrules}).

$5^{1/n} = \eu^{\frac{\ln 5}{n}} \to \eu^0 = 1$.
\end{solution}

\begin{solution}{exo:b1:seq:3}
Bernoulli: $(1+h)^n \geq 1 + nh$ for $h \geq -1$, by induction ---
$(1+h)^{n+1} = (1+h)^n(1+h) \geq (1+nh)(1+h) = 1 + (n+1)h + nh^2
\geq 1 + (n+1)h$.

For $0 < \abs q < 1$: write $\frac{1}{\abs q} = 1 + h$, $h > 0$; then
$\abs{q}^n = \frac{1}{(1+h)^n} \leq \frac{1}{1 + nh} \to 0$, and the
squeeze gives $q^n \to 0$ (the case $q = 0$ is trivial). For $q = 1$:
constant sequence, limit $1$. For $q = -1$: diverges
(\cref{exo:b1:seq:1}). For $\abs q > 1$: $\abs q^n = (1 + h)^n \geq
1 + nh \to +\infty$, so $(q^n)$ is unbounded, hence divergent (to
$+\infty$ if $q > 1$; with alternating signs, no limit, if $q <
-1$).
\end{solution}

\begin{solution}{exo:b1:seq:4}
Fixed point: $\ell = \frac{\ell + 3}{2}$ gives $\ell = 3$. Then
\[
v_{n+1} = u_{n+1} - 3 = \frac{u_n + 3}{2} - 3 = \frac{u_n - 3}{2} =
\frac{v_n}{2}:
\]
$(v_n)$ is geometric with ratio $\frac12$, $v_0 = -3$. So $u_n = 3 -
\frac{3}{2^n} \to 3$.
\end{solution}

\begin{solution}{exo:b1:seq:5}
$H_{2n} - H_n = \sum_{k=n+1}^{2n} \frac 1k \geq n \cdot
\frac{1}{2n} = \frac12$ (each of the $n$ terms is $\geq
\frac{1}{2n}$). If $(H_n)$ were Cauchy, taking $\varepsilon =
\frac13$ would force $\abs{H_{2n} - H_n} \leq \frac13$ for large
$n$: contradiction. An increasing non-convergent sequence diverges to
$+\infty$ (\cref{thm:b1:seq:monotone}): $H_n \to +\infty$.
\end{solution}

\begin{solution}{exo:b1:seq:6}
Let $a = \lim u_{2n}$, $b = \lim u_{2n+1}$, $c = \lim u_{3n}$. The
sequence $(u_{6n})$ is a \omterm{def:b1:seq:subsequence}{subsequence} both of $(u_{2n})$ and of
$(u_{3n})$: its limit equals $a$ and $c$, so $a = c$. The sequence
$(u_{6n+3})$ is a \omterm{def:b1:seq:subsequence}{subsequence} of $(u_{2n+1})$ (odd indices) and of
$(u_{3n})$ (indices $6n + 3 = 3(2n+1)$): so $b = c$. Hence $a = b$,
and \cref{prop:b1:seq:subsequences} (evens and odds with equal
limits) gives the convergence of $(u_n)$.
\end{solution}

\begin{solution}{exo:b1:seq:7}
\emph{Stability and bounds:} $I = \intcc{0}{2}$ is stable: for $x \in
I$, $\sqrt{2 + x} \in \intcc{\sqrt 2}{2} \subseteq I$; and $u_0 = 0
\in I$.

\emph{Monotonicity:} $f(x) = \sqrt{2+x}$ is increasing and $u_1 =
\sqrt 2 > u_0$: by induction $(u_n)$ is increasing. Increasing and
bounded above by $2$: it converges
(\cref{thm:b1:seq:monotone}).

\emph{Limit:} $\ell = \sqrt{2 + \ell}$ with $\ell \geq 0$ gives
$\ell^2 - \ell - 2 = 0$, so $\ell = 2$.

\emph{Error bound:} multiplying by the conjugate,
\[
2 - u_{n+1} = 2 - \sqrt{2 + u_n}
= \frac{4 - (2 + u_n)}{2 + \sqrt{2+u_n}}
= \frac{2 - u_n}{2 + \sqrt{2 + u_n}}
\leq \frac{2 - u_n}{3},
\]
since $\sqrt{2 + u_n} \geq \sqrt 2 > 1$. By induction from $2 - u_0
= 2$: $\;0 \leq 2 - u_n \leq \frac{2}{3^n}$.
\end{solution}

\begin{solution}{exo:b1:seq:8}
Suppose $(u_n)$ does not converge to $\ell$: for some $\varepsilon_0
> 0$, infinitely many indices satisfy $\abs{u_n - \ell} >
\varepsilon_0$; they form a \omterm{def:b1:seq:subsequence}{subsequence} $(u_{\varphi(n)})$. This
\omterm{def:b1:seq:subsequence}{subsequence} is bounded, so by Bolzano--Weierstrass
(\cref{thm:b1:seq:bw}) it has a convergent \omterm{def:b1:seq:subsequence}{sub-subsequence}, whose
limit $\ell'$ satisfies $\abs{\ell' - \ell} \geq \varepsilon_0$
(pass the inequality to the limit, \cref{thm:b1:seq:order}). But a
\omterm{def:b1:seq:subsequence}{sub-subsequence} of $(u_n)$ is a convergent \omterm{def:b1:seq:subsequence}{subsequence} of $(u_n)$,
so by hypothesis $\ell' = \ell$: contradiction.
\end{solution}

\begin{solution}{exo:b1:seq:9}
Suppose $\eu = \frac pq$, $q \geq 1$. The strict inequalities $a_q <
\eu < a_q + \frac{1}{q\,q!}$ (strict since $(a_n)$ is strictly
increasing and $(b_n)$ strictly decreasing) multiplied by $q!$ give
\[
q!\,a_q \;<\; q!\,\frac pq \;<\; q!\,a_q + \frac 1q \leq q!\,a_q + 1.
\]
Now $N = q!\,a_q = \sum_{k=0}^{q} \frac{q!}{k!}$ is an integer (each
$\frac{q!}{k!}$ is a product of integers for $k \leq q$), and so is
$q!\,\frac pq = (q-1)!\,p$. The display thus places the integer
$(q-1)!\,p$ strictly between $N$ and $N + \frac 1q \leq N + 1$: an
integer strictly inside $\intoo{N}{N+1}$, which is impossible. Hence
$\eu \notin \Q$.
\end{solution}

\begin{solution}{exo:b1:seq:10}
\begin{enumerate}
  \item Let $\varepsilon > 0$ and $N$ with $\abs{u_k - \ell} \leq
        \frac{\varepsilon}{2}$ for $k > N$. For $n > N$:
        \[
        \abs{c_n - \ell}
        = \Bigl|\frac{\sum_{k=1}^{n}(u_k - \ell)}{n}\Bigr|
        \leq \frac{\sum_{k=1}^{N} \abs{u_k - \ell}}{n}
        + \frac{n - N}{n}\cdot\frac{\varepsilon}{2}
        \leq \frac{C}{n} + \frac{\varepsilon}{2},
        \]
        where $C = \sum_{k=1}^N \abs{u_k - \ell}$ is fixed. For $n$
        large, $\frac Cn \leq \frac\varepsilon2$: then $\abs{c_n -
        \ell} \leq \varepsilon$.
  \item $u_n = (-1)^n$: diverges, yet $c_n \to 0$ (partial sums
        bounded by $1$, divided by $n$).
  \item Apply (1) to the sequence $v_n = u_{n+1} - u_n \to \ell$: its
        Ces\`aro mean is $\frac{u_{n+1} - u_1}{n} \to \ell$
        (telescoping), and $\frac{u_{n+1}}{n} = \frac{u_{n+1} -
        u_1}{n} + \frac{u_1}{n} \to \ell$; renormalizing indices
        ($\frac{u_n}{n} = \frac{u_n}{n-1}\cdot\frac{n-1}{n}$) gives
        $\frac{u_n}{n} \to \ell$.
\end{enumerate}
\end{solution}

\begin{solution}{exo:b1:seq:11}
Let $L = \inf_{n \geq 1} \frac{u_n}{n} \geq 0$, and $\varepsilon >
0$. Choose $m$ with $\frac{u_m}{m} \leq L + \varepsilon$. Every $n$
writes $n = qm + r$, $0 \leq r < m$; subadditivity (iterated) gives
$u_n \leq q\,u_m + u_r$, so
\[
\frac{u_n}{n} \leq \frac{q m}{n}\cdot\frac{u_m}{m} + \frac{u_r}{n}
\leq \frac{u_m}{m} + \frac{\max(u_0, \dots, u_{m-1})}{n}
\leq L + \varepsilon + \frac{C_m}{n},
\]
using $qm \leq n$. For $n$ large, $\frac{C_m}{n} \leq \varepsilon$:
thus $L \leq \frac{u_n}{n} \leq L + 2\varepsilon$ for all large $n$,
which is the convergence to $L$.
\end{solution}

\begin{solution}{exo:b1:seq:12}
The \omterm{def:b1:structures:subgroup}{subgroup} $G = \Z + 2\pi\Z$ of $(\R, +)$ is dense: it is not
$\alpha\Z$, since $1 = p\alpha$, $2\pi = q\alpha$ would make $2\pi =
\frac qp$ rational --- and $\pi \notin \Q$ (admitted here; a proof is
sketched in \cref{ch:b1:integration}). By \cref{exo:b1:reals:9}, $G$
is dense in $\R$.

Let now $y \in \intcc{-1}{1}$ and $\theta = \arcsin y$. By density,
for every $\varepsilon > 0$ there are $n \in \Z$, $k \in \Z$ with
$\abs{(n + 2\pi k) - \theta} \leq \varepsilon$, i.e.\ $n$ is within
$\varepsilon$ of $\theta - 2\pi k$; then, $\sin$ being
$2\pi$-periodic and $1$-Lipschitz ($\abs{\sin a - \sin b} \leq
\abs{a - b}$, from the mean value inequality of
\cref{ch:b1:derivative}),
\[
\abs{\sin n - y} = \abs{\sin(n + 2\pi k) - \sin\theta}
\leq \abs{n + 2\pi k - \theta} \leq \varepsilon .
\]
One detail: $n$ ranges over $\Z$, but $\sin(-n) = -\sin n$ and $y$
was arbitrary in $\intcc{-1}{1}$, so nonnegative indices suffice
(replace $(n, y)$ by $(-n, -y)$ if needed). Hence $\{\sin n : n \in
\N\}$ is dense in $\intcc{-1}{1}$; a dense-in-a-segment sequence has
\omterm{def:b1:seq:subsequence}{subsequences} approaching distinct values, so it diverges.
\end{solution}

\begin{solution}{pb:b1:seq:1}
\textbf{1.} $\dfrac{n(n+1)/2}{n^2} = \dfrac{1 + 1/n}{2} \to
\dfrac12$, and $\dfrac{n(n+1)(2n+1)/6}{n^3} = \dfrac{(1 +
1/n)(2 + 1/n)}{6} \to \dfrac13$.

\textbf{2.} Upper: each of the $n$ terms is $\leq \sqrt n$, so
$T_n \leq n\sqrt n$. Lower: the terms with $k > \frac n2$ number
at least $\frac n2$, and each is $\geq \sqrt{n/2}$:
\[
T_n \geq \frac n2 \sqrt{\frac n2} = \frac{n^{3/2}}{2\sqrt 2} .
\]

\textbf{3.} For $k \geq N$, since $b_{k+1} - b_k > 0$:
$m(b_{k+1} - b_k) \leq a_{k+1} - a_k \leq M(b_{k+1} - b_k)$.
Summing for $k = N, \dots, n - 1$, both sides telescope:
\[
m\,(b_n - b_N) \leq a_n - a_N \leq M\,(b_n - b_N),
\]
and dividing by $b_n - b_N > 0$ gives the claim.

\textbf{4.} By definition of the limit there is $N$ with $\ell -
\varepsilon \leq \frac{a_{k+1} - a_k}{b_{k+1} - b_k} \leq \ell +
\varepsilon$ for all $k \geq N$; question 3 with $m = \ell -
\varepsilon$, $M = \ell + \varepsilon$ transfers the bracketing
to $\frac{a_n - a_N}{b_n - b_N}$.

\textbf{5.} Expanding the right side of the identity:
\[
\frac{a_N - \ell b_N}{b_n} + \frac{a_n - a_N}{b_n} -
\ell\,\frac{b_n - b_N}{b_n}
= \frac{a_n - \ell b_n}{b_n} = \frac{a_n}{b_n} - \ell .
\]
By question 4 the second factor of the product is bounded by
$\varepsilon$ in absolute value, and $0 < 1 - \frac{b_N}{b_n}
\leq 1$ for large $n$, so
\[
\Bigl|\frac{a_n}{b_n} - \ell\Bigr|
\leq \frac{\abs{a_N - \ell b_N}}{b_n} + \varepsilon
\leq 2\varepsilon
\]
as soon as $b_n \geq \frac{\abs{a_N - \ell b_N}}{\varepsilon}$,
which happens eventually since $b_n \to +\infty$. Hence
$\frac{a_n}{b_n} \to \ell$: the Ces\`aro--Stolz theorem.

\textbf{6.} Given $M$, choose $N$ with $\frac{a_{k+1} -
a_k}{b_{k+1} - b_k} \geq M$ for $k \geq N$; the lower half of
question 3 gives $a_n - a_N \geq M(b_n - b_N)$, so
\[
\frac{a_n}{b_n} \geq \frac{a_N}{b_n} +
M\Bigl(1 - \frac{b_N}{b_n}\Bigr) \longrightarrow M .
\]
Beyond some rank, $\frac{a_n}{b_n} \geq M - 1$; as $M$ was
arbitrary, $\frac{a_n}{b_n} \to +\infty$.

\textbf{7.} With $b_n = n$ and $a_n = u_1 + \dots + u_n$: the
increment quotient is $u_{n+1} \to \ell$, so the Ces\`aro mean
$\frac{a_n}{n}$ tends to $\ell$: part (1) of
\cref{exo:b1:seq:10}. With $a_n = u_n$: the increment quotient
is $u_{n+1} - u_n$, giving part (3). Converse: $a_n = (-1)^n$,
$b_n = n$ has $\frac{a_n}{b_n} \to 0$, yet $a_{n+1} - a_n = \pm
2$ alternates: the increment quotient has no limit.

\textbf{8.} Conjugation:
\[
\bigl((1+h)^{3/2} - 1\bigr)\bigl((1+h)^{3/2} + 1\bigr)
= (1+h)^3 - 1 = 3h + 3h^2 + h^3 .
\]
For $h = \frac1n$: $n\bigl((1 + \frac1n)^{3/2} - 1\bigr) =
\frac{3 + 3/n + 1/n^2}{(1 + 1/n)^{3/2} + 1}$, and $1 \leq (1 +
\frac1n)^{3/2} \leq (1 + \frac1n)^2 \to 1$ (squeeze), so the
denominator tends to $2$ and the expression to $\frac32$. Now
Stolz with $a_n = T_n$, $b_n = n^{3/2}$ (strictly increasing,
$\to +\infty$):
\[
\frac{\sqrt{n+1}}{(n+1)^{3/2} - n^{3/2}}
= \frac{\sqrt{n+1}}{\sqrt n} \cdot
\frac{1}{n\bigl((1 + \frac1n)^{3/2} - 1\bigr)}
\longrightarrow 1 \cdot \frac{2}{3},
\]
hence $T_n \sim \frac23\,n^{3/2}$. (The bracketing of question 2
had trapped the constant in $\intcc{0.35}{1}$; Stolz pins it.)

\textbf{9.} (G1) at $u = \ln(1+t)$: $1 + t = \eu^{\ln(1+t)} \geq
1 + \ln(1+t)$, so $\ln(1+t) \leq t$. (G1) at $u =
-\frac{t}{1+t}$: $\eu^{-t/(1+t)} \geq 1 - \frac{t}{1+t} =
\frac{1}{1+t} > 0$; taking $\ln$ (increasing):
$-\frac{t}{1+t} \geq -\ln(1+t)$, i.e.\ $\ln(1+t) \geq
\frac{t}{1+t}$.

\textbf{10.} $\ln$ is strictly increasing
(\cref{prop:b1:functions:expln}), and $\ln(2^k) = k\ln 2$ is
unbounded, so $\ln n \to +\infty$. Increments: with $t =
\frac1n$ in question 9,
\[
\frac 1{n+1} = \frac{1/n}{1 + 1/n} \leq
\ln\Bigl(1 + \frac1n\Bigr) \leq \frac1n
\quad\Longrightarrow\quad
\frac{n}{n+1} \leq \frac{1/(n+1)}{\ln(1 + 1/n)} \leq 1 ,
\]
so the increment quotient $\frac{H_{n+1} - H_n}{\ln(n+1) - \ln
n}$ tends to $1$; Stolz gives $H_n \sim \ln n$.

\textbf{11.} Set $x_n = \frac{u_{n+1}}{u_n} \to L$ and $t_n =
\frac{x_n}{L} - 1 \to 0$. Question 9: $\frac{t_n}{1 + t_n} \leq
\ln(1 + t_n) \leq t_n$, so $\ln x_n - \ln L = \ln(1 + t_n) \to
0$ by squeezing. Ces\`aro (question 7) applied to $(\ln x_k)$:
\[
\frac1n \sum_{k=0}^{n-1} \ln x_k = \frac{\ln u_n - \ln u_0}{n}
\longrightarrow \ln L ,
\]
so $\frac{\ln u_n}{n} \to \ln L$. With $h_n = \frac{\ln u_n}{n}
- \ln L \to 0$: $u_n^{1/n} = L\,\eu^{h_n}$, and (G1) squeezes
$1 + h_n \leq \eu^{h_n} \leq \frac{1}{1 - h_n}$ (for $h_n < 1$),
so $\eu^{h_n} \to 1$ and $u_n^{1/n} \to L$. Application: $u_n =
\binom{2n}{n}$ gives
\[
\frac{u_{n+1}}{u_n} = \frac{(2n+1)(2n+2)}{(n+1)^2}
= \frac{2(2n+1)}{n+1} \longrightarrow 4 ,
\qquad\text{so}\qquad
\binom{2n}{n}^{1/n} \to 4 .
\]

\textbf{12.} For $0 < x \leq 1$, (G2) gives $\sin x \geq x(1 -
\frac{x^2}{6}) \geq \frac{5x}{6} > 0$ and
\[
x - \sin x \geq \frac{x^3}{6} - \frac{x^5}{120}
= x^3\Bigl(\frac16 - \frac{x^2}{120}\Bigr)
\geq \frac{19}{120}\,x^3 > 0 :
\]
thus $0 < \sin x < x$ on $\intoc{0}{1}$. Always $u_1 = \sin u_0
\in \intcc{-1}{1}$. If $u_1 = 0$, then $u_n = 0$ for $n \geq 1$.
If $u_1 \in \intoc{0}{1}$: by induction $0 < u_{n+1} = \sin u_n
< u_n \leq 1$, so $(u_n)_{n\geq1}$ is strictly decreasing and
bounded below by $0$: it converges to some $\ell \in
\intco{0}{1}$ (\cref{thm:b1:seq:monotone}). The product-to-sum
formula and (G2) give $\abs{\sin a - \sin b} = 2\abs{\cos
\frac{a+b}{2}}\,\abs{\sin\frac{a-b}{2}} \leq \abs{a - b}$, so
$u_{n+1} = \sin u_n \to \sin \ell$: $\ell = \sin\ell$. If $\ell
> 0$ then $\sin\ell < \ell$: impossible. So $u_n \to 0$.

\textbf{13.} Dividing (G2) by $u_n > 0$:
\[
1 - \frac{u_n^2}{6} \leq \frac{\sin u_n}{u_n} \leq 1 -
\frac{u_n^2}{6} + \frac{u_n^4}{120} \leq 1 ,
\]
and $u_n \to 0$ squeezes $\frac{\sin u_n}{u_n} \to 1$. Dividing
$x - \sin x$ by $x^3$:
\[
\frac16 - \frac{u_n^2}{120} \leq \frac{u_n - \sin u_n}{u_n^3}
\leq \frac16 \longrightarrow \frac16 .
\]

\textbf{14.} Since $u_{n+1} = \sin u_n$:
\[
w_n = \frac{u_n^2 - \sin^2 u_n}{u_n^2 \sin^2 u_n}
= \frac{(u_n - \sin u_n)(u_n + \sin u_n)}{u_n^2 \sin^2 u_n}
= \frac{u_n - \sin u_n}{u_n^3}\cdot
\frac{u_n + \sin u_n}{u_n}\cdot
\Bigl(\frac{u_n}{\sin u_n}\Bigr)^{2}
\]
(check the powers of $u_n$: $3 + 1 + (-4)$ against the $u_n^2$
downstairs and $u_n^4$ upstairs). By question 13 the three
factors tend to $\frac16$, $2$, $1$: $w_n \to \frac13$.

\textbf{15.} $v_n = \frac{1}{u_n^2}$ has increments $v_{n+1} -
v_n = w_n \to \frac13$, so $\frac{v_n}{n} \to \frac13$ by
\cref{exo:b1:seq:10}~(3): $n u_n^2 \to 3$. Then
\[
\abs{\sqrt n\,u_n - \sqrt 3}
= \frac{\abs{n u_n^2 - 3}}{\sqrt n\,u_n + \sqrt 3}
\leq \frac{\abs{n u_n^2 - 3}}{\sqrt 3} \longrightarrow 0 :
\]
$\sqrt n\,u_n \to \sqrt 3$, i.e.\ $u_n \sim \sqrt{3/n}$.

\textbf{16.} Since $n u_n^2 \to 3$, eventually $2 \leq n u_n^2
\leq 4$, i.e.\ $\sqrt{2/n} \leq u_n \leq 2/\sqrt n$. If $u_n
\leq 10^{-2}$ with $n$ in that range, then $2/n \leq 10^{-4}$:
$n \geq 20\,000$; and $\sqrt{3/n} = 10^{-2}$ at $n = 30\,000$.
\omterm{ex:b1:seq:heron}{Heron's method} squares the error at each step --- the digit
count doubles --- because at its fixed point the relevant slope
is $< 1$ in \omterm{def:b1:complex:field}{modulus} (indeed the iteration is contracting).
Here $\sin' 0 = \cos 0 = 1$: the fixed point is neutral, no
geometric contraction exists, and the decay is governed by the
first nonlinear term $-\frac{x^3}{6}$, hence \omterm{def:b1:poly:def}{polynomial}. One
step of Heron gains more accuracy than ten thousand steps of
the sine.

\textbf{17.} For arbitrary $u_0$: $u_1 = \sin u_0 \in
\intcc{-1}{1}$. If $u_1 = 0$ the sequence vanishes from rank
$1$. If $u_1 > 0$, Part IV applies verbatim. If $u_1 < 0$, set
$v_n = -u_n$: oddness of $\sin$ gives $v_{n+1} = -\sin u_n =
\sin(-u_n) = \sin v_n$ with $v_1 \in \intoc{0}{1}$, so $v_n \sim
\sqrt{3/n}$, i.e.\ $u_n \sim -\sqrt{3/n}$. In all cases
$\abs{u_n} \sim \sqrt{3/n}$ (or the sequence is eventually $0$):
the fall is universal, only the sign remembers the start.

\textbf{18.} First $\frac{u_{n+1}}{u_n} = 1 - \frac{u_n -
u_{n+1}}{u_n^2}\,u_n \to 1 - a \cdot 0 = 1$. Then
\[
\frac{1}{u_{n+1}} - \frac{1}{u_n}
= \frac{u_n - u_{n+1}}{u_n u_{n+1}}
= \frac{u_n - u_{n+1}}{u_n^2}\cdot\frac{u_n}{u_{n+1}}
\longrightarrow a \cdot 1 = a ,
\]
and \cref{exo:b1:seq:10}~(3) gives $\frac{1}{n u_n} \to a$,
i.e.\ $n u_n \to \frac1a$.

\textbf{19.} $v_n = \frac{1}{u_n}$: $v_{n+1} = \frac{1 +
u_n}{u_n} = v_n + 1$, so $v_n = v_0 + n$ and
\[
u_n = \frac{u_0}{1 + n u_0} , \qquad
n u_n = \frac{n u_0}{1 + n u_0} \longrightarrow 1 .
\]
Lemma check: $u_n - u_{n+1} = \frac{u_n^2}{1 + u_n}$, so
$\frac{u_n - u_{n+1}}{u_n^2} = \frac{1}{1 + u_n} \to 1 = a$, and
question 18 predicts $n u_n \to 1$: exact agreement.

\textbf{20.} Positivity by induction ($\eu^{-u} > 0$);
decreasing since $\eu^{-u_n} < 1$ for $u_n > 0$; hence $u_n \to
\ell \geq 0$ (\cref{thm:b1:seq:monotone}). Continuity bridge:
with $h_n = \ell - u_n \to 0$, $\eu^{-u_n} = \eu^{-\ell}
\eu^{h_n} \to \eu^{-\ell}$ by the (G1) squeeze $1 + h_n \leq
\eu^{h_n} \leq \frac{1}{1 - h_n}$; so $\ell = \ell\,
\eu^{-\ell}$, and $\ell > 0$ would force $\eu^{-\ell} = 1$,
false: $\ell = 0$. For $t > 0$, (G1) gives $\eu^{-t} \geq 1 - t$
and $\eu^{-t} \leq \frac{1}{1 + t}$, so
\[
\frac{1}{1 + t} \leq \frac{1 - \eu^{-t}}{t} \leq 1 .
\]
With $t = u_n$: $\frac{u_n - u_{n+1}}{u_n^2} = \frac{1 -
\eu^{-u_n}}{u_n} \to 1$. Question 18 with $a = 1$: $n u_n \to
1$, so $u_n \sim \frac1n$.

\textbf{21.} As in question 18, $\frac{u_{n+1}}{u_n} = 1 -
\frac{u_n - u_{n+1}}{u_n^3}\,u_n^2 \to 1$. Then
\[
\frac{1}{u_{n+1}^2} - \frac{1}{u_n^2}
= \frac{(u_n - u_{n+1})(u_n + u_{n+1})}{u_n^2 u_{n+1}^2}
= \frac{u_n - u_{n+1}}{u_n^3}\cdot
\frac{u_n + u_{n+1}}{u_n}\cdot
\Bigl(\frac{u_n}{u_{n+1}}\Bigr)^2
\longrightarrow a \cdot 2 \cdot 1 = 2a ,
\]
and \cref{exo:b1:seq:10}~(3) gives $\frac{1}{n u_n^2} \to 2a$:
$n u_n^2 \to \frac{1}{2a}$. For the sine, $a = \frac16$
(question 13): $n u_n^2 \to 3$, exactly Part IV.

\textbf{22.} Increments: $a_{n+1} - a_n = 2^{-n} - 2^{-n-1} =
2^{-n-1} = b_{n+1} - b_n$, so the increment quotient is
constantly $1$. Yet $\frac{a_n}{b_n} = \frac{2 - 2^{-n}}{1 -
2^{-n}} \to 2 \neq 1$. The proof of question 5 breaks at the
boundary term: $\frac{a_N - \ell b_N}{b_n} \to 0$ needed $b_n
\to +\infty$; here (with $\ell = 1$) $a_N - b_N = 1$ and $b_n
\to 1$, so the term tends to $1$ --- precisely the residual gap
$2 - 1$.

\textbf{23.} Stolz with $A_n = \sum_{k=1}^n H_k$ and $B_n = n\ln
n$: $B_{n+1} - B_n = \ln(n+1) + n\ln(1 + \frac1n) > 0$ and $B_n
\to +\infty$. By question 9, $\frac{n}{n+1} \leq n\ln(1 +
\frac1n) \leq 1$, so $B_{n+1} - B_n = \ln(n+1) + \theta_n$ with
$\frac12 \leq \theta_n \leq 1$. Hence
\[
\frac{A_{n+1} - A_n}{B_{n+1} - B_n}
= \frac{H_{n+1}}{\ln(n+1)}\cdot
\frac{1}{1 + \theta_n/\ln(n+1)} \longrightarrow 1 \cdot 1 = 1
\]
(question 10 for the first factor; $\theta_n$ bounded and
$\ln(n+1) \to \infty$ for the second). Stolz concludes:
$\sum_{k=1}^n H_k \sim n\ln n$.

\textbf{24.} If $u_n \to \ell > 0$: as in question 11, $\ln u_n
\to \ln\ell$ (question 9 squeeze on $\ln\frac{u_n}{\ell}$), so
the Ces\`aro means $\frac1n\sum_{k=1}^n \ln u_k \to \ln\ell$,
and the exponential bridge of question 11 gives $(u_1\cdots
u_n)^{1/n} = \exp\bigl(\frac1n\sum\ln u_k\bigr) \to \ell$. If
$u_n \to +\infty$: for any $M$, eventually $u_n \geq \eu^M$, so
$\ln u_n \geq M$: $\ln u_n \to +\infty$; the $+\infty$ Ces\`aro
(question 6, $b_n = n$) gives $\frac1n\sum \ln u_k \to
+\infty$, and (G1) ($\eu^s \geq 1 + s$) sends the geometric mean
to $+\infty$. With $u_n = n$: $(n!)^{1/n} \to +\infty$.

\textbf{25.} (i) Completeness entered only through the monotone
limit theorem, to produce the limits in questions 12 and 20;
the Ces\`aro--Stolz theorem itself is pure
$\varepsilon$-management, valid over $\Q$. (ii) Stolz replaces
$\lim \frac{a_n}{b_n}$ by $\lim$ of the quotient of increments,
exactly as l'Hospital replaces $\lim\frac fg$ by
$\lim\frac{f'}{g'}$ --- the differential twin rests on the mean
value theorem of \cref{ch:b1:derivative}. (iii) Heuristic: if
$f(x) = x - a\,x^{p+1} + o(x^{p+1})$ at the neutral fixed point
$0$, then $\frac{1}{u_{n+1}^p} - \frac{1}{u_n^p} \to pa$ and
$u_n \sim (pan)^{-1/p}$: contact of order $p + 1$ yields decay
$n^{-1/p}$ --- the flatter the graph against the diagonal, the
slower the fall. (iv) The constant: (G2) supplies the cubic
coefficient $\frac16$; the factorization of question 14 doubles
it into the telescope increment $\frac13$; Ces\`aro turns
$\frac{1}{u_n^2}$ into $\frac n3$; inverting and taking roots
delivers $\sqrt{3/n}$.
\end{solution}

\section*{Chapter \ref{ch:b1:topology} --- Topology of the Real Line}

\begin{solution}{exo:b1:topology:1}
$\intoo{0}{1} \cup \intoo{2}{3}$: open (union of open sets), not
closed ($\frac 1n \to 0$ outside).

$\intco{0}{1}$: neither. Not open (no interval around $0$ inside);
not closed ($1 - \frac1n \to 1 \notin$ set).

$\{0\} \cup \intcc{1}{2}$: closed (finite union of closed sets), not
open (fails at $0$).

$\R \setminus \Z$: open ($\Z$ is closed), not closed: the sequence
$\bigl(\frac 1n\bigr)$ lies in it, but its limit $0$ belongs to $\Z$,
i.e.\ escapes the set.

$\Q \cap \intoo{0}{1}$: neither. Not open: every interval around a
rational contains irrationals. Not closed: it contains sequences
tending to the irrational $\frac{\sqrt 2}{2}$ (density).
\end{solution}

\begin{solution}{exo:b1:topology:2}
$A = \intoc{0}{1} \cup \{2\}$: $\mathring A = \intoo{0}{1}$,
$\overline A = \intcc{0}{1} \cup \{2\}$, $\partial A = \{0, 1, 2\}$.

$A = \R \setminus \Q$: $\mathring A = \emptyset$ (every interval
contains rationals), $\overline A = \R$ (density of the
irrationals), $\partial A = \R$.

$A = \bigl\{\frac{(-1)^n n}{n+1}\bigr\}$: the even terms tend to $1$,
the odd terms to $-1$, neither belongs to $A$. $\mathring A =
\emptyset$ (isolated points), $\overline A = A \cup \{-1, 1\}$,
$\partial A = \overline A$.
\end{solution}

\begin{solution}{exo:b1:topology:3}
\emph{Complements.} $F = \{a_1 < a_2 < \dots < a_k\}$: the complement
is the union of the open intervals $\intoo{-\infty}{a_1}$,
$\intoo{a_i}{a_{i+1}}$, $\intoo{a_k}{+\infty}$ --- open by
\cref{prop:b1:topology:openstable}.

\emph{Sequences.} Let $u_n \in F$, $u_n \to \ell$. With $\varepsilon
= \min\{\abs{a_i - a_j} : i \neq j\}/2 > 0$ (or any $\varepsilon$ if
$F$ is a singleton): beyond some rank, all terms are within
$\varepsilon$ of $\ell$, hence within $2\varepsilon$ of each other,
which forces them to be one single $a_i$ from that rank on; then
$\ell = a_i \in F$.
\end{solution}

\begin{solution}{exo:b1:topology:4}
$\subseteq$: $\overline A \cup \overline B$ is closed (finite union)
and contains $A \cup B$, so it contains the \emph{smallest} closed
superset $\overline{A \cup B}$. $\supseteq$: $A \subseteq A \cup B$
gives $\overline A \subseteq \overline{A \cup B}$ (closure is
monotone: adherent points of $A$ are adherent to the larger set), and
likewise for $B$.

Counterexample for intersections: $A = \intoo{0}{1}$, $B =
\intoo{1}{2}$: $\overline{A \cap B} = \overline\emptyset = \emptyset$
but $\overline A \cap \overline B = \{1\}$.
\end{solution}

\begin{solution}{exo:b1:topology:5}
Use the sequential characterization
(\cref{thm:b1:topology:seqclosed}). Let $S = \{u_n\} \cup \{\ell\}$
and $(v_k)$ a sequence in $S$ with $v_k \to m$; show $m \in S$. Two
cases. If some value $v \in S$ is taken by $(v_k)$ infinitely often,
then a constant subsequence gives $m = v \in S$. Otherwise every
value is taken finitely often; in particular, for each $n$, the term
$u_n$ appears finitely often, and $\ell$ too. Then for every $N$, the
indices $k$ with $v_k \in \{u_0, \dots, u_N, \ell\}$ are finitely
many: the remaining $v_k$ are terms $u_n$ with $n > N$. Given
$\varepsilon > 0$, choose $N$ with $\abs{u_n - \ell} \leq
\varepsilon$ for $n > N$: all but finitely many $v_k$ satisfy
$\abs{v_k - \ell} \leq \varepsilon$. Hence $v_k \to \ell$, so $m =
\ell \in S$.
\end{solution}

\begin{solution}{exo:b1:topology:6}
$U + A = \bigcup_{a \in A} (U + a)$, and each translate $U + a$ is
open (translating the interval certificates). A union of open sets is
open (\cref{prop:b1:topology:openstable}).

Closed sets: $\Z$ and $\sqrt 2\,\Z$ are closed (as $\alpha\Z$:
convergent sequences are eventually constant, cf.\
\cref{ex:b1:topology:closed}). Their sum $\Z + \sqrt 2\,\Z$ is dense
in $\R$ (\cref{exo:b1:reals:9}) but is not $\R$ (it is countable, or
simply: $\frac{\sqrt 2}{2} \notin \Z + \sqrt2\Z$, else $\sqrt 2$
would be rational --- writing $\frac{\sqrt2}{2} = m + n\sqrt 2$
forces $(2n - 1)\sqrt 2 = -2m$, so $\sqrt 2 \in \Q$ unless $n =
\frac12$, impossible). A dense proper subset is not closed: its
closure is $\R \neq$ itself.
\end{solution}

\begin{solution}{exo:b1:topology:7}
$\Z$: the neighborhood $\intoo{n - \frac12}{n + \frac12}$ of $n$
meets $\Z$ only at $n$.

$A = \{\frac1n : n \in \N^*\}$: around $\frac 1n$, the interval of
radius $\frac{1}{n} - \frac{1}{n+1} = \frac{1}{n(n+1)}$ (halved, say)
isolates it from its neighbors --- all points isolated. Yet $0 \in
\overline A \setminus A$: isolated points do not prevent adherent
outsiders.

If all points of $A$ are isolated: no point of $A$ is interior, since
an interior point has a whole interval of $A$-neighbors around it
(an interval is infinite), contradicting isolation. So $\mathring A
= \emptyset$.
\end{solution}

\begin{solution}{exo:b1:topology:8}
If $A \subseteq \intcc{-M}{M}$, the closed set $\intcc{-M}{M}$
contains $A$, hence contains $\overline A$ (smallest closed
superset): $\overline A$ is bounded.

Let $s = \sup A$ (finite). Since $A \subseteq \overline A$, $\sup
\overline A \geq s$. Conversely, $\overline A \subseteq
\intoc{-\infty}{s}$: the half-line is closed and contains $A$; so
every element of $\overline A$ is $\leq s$, giving $\sup\overline A
\leq s$. Equality.

$s \in \overline A$: by the $\varepsilon$-characterization
(\cref{prop:b1:reals:epsilon}), every interval $\intoo{s -
\varepsilon}{s + \varepsilon}$ contains an element of $A$: $s$ is
adherent.
\end{solution}

\begin{solution}{exo:b1:topology:9}
Suppose $A$ is open and closed, with $a \in A$ and $b \in \R
\setminus A$; without loss of generality $a < b$. The set $B = A \cap
\intcc{a}{b}$ is nonempty ($a$), bounded: let $s = \sup B$. By
\cref{exo:b1:topology:8}, $s \in \overline B \subseteq \overline A =
A$ ($A$ closed). Note $s \leq b$, and since $b \notin A$: $s < b$.
Now $A$ is open: some interval $\intoo{s - r}{s + r}$ lies in $A$,
and we may take $r < b - s$. Then $s + \frac{r}{2}$ belongs to $A
\cap \intcc{a}{b} = B$ and exceeds $s$ --- contradicting $s = \sup
B$. Hence no such pair $(a, b)$ exists: one of $A$, $\R \setminus A$
is empty.
\end{solution}

\begin{solution}{exo:b1:topology:10}
$I_x$ is a union of open intervals all containing $x$: it is open,
and it is an interval, being convex --- if $u < z < v$ with $u, v \in
I_x$, then $u$ and $v$ lie in open subintervals $J_u \ni x$, $J_v \ni
x$ of $U$, and $J_u \cup J_v$ is an interval (both contain $x$)
inside $U$ containing $z$; so $z \in I_x$
(\cref{prop:b1:reals:intervals}).

If $I_x \cap I_y \neq \emptyset$: $I_x \cup I_y$ is then an open
interval (convex: two overlapping intervals) contained in $U$
containing $x$ and $y$, so $I_x \cup I_y \subseteq I_x$ and
$\subseteq I_y$ by maximality of each: $I_x = I_y$.

So $U$ is the disjoint union of the distinct sets $I_x$ (each $x \in
U$ lies in its own $I_x$). Countability: each nonempty open interval
$I$ of the family contains a rational $q_I$
(\cref{thm:b1:reals:density}), and distinct disjoint intervals get
distinct rationals: the family injects into $\Q$, which is countable
(it is indexed by pairs of integers). Hence at most countably many
intervals.
\end{solution}

\section*{Chapter \ref{ch:b1:continuity} --- Limits and Continuity}

\begin{solution}{exo:b1:continuity:1}
Take $u_n = \frac{1}{2\pi n + \pi/2}$ and $v_n = \frac{1}{2\pi n}$:
both tend to $0^+$, yet $\sin\frac{1}{u_n} = 1$ and
$\sin\frac{1}{v_n} = 0$. Two sequences, two different limits of
images: by \cref{thm:b1:continuity:seqchar}, no limit at $0^+$.

$x \sin\frac1x$: squeezed by $\abs{x\sin\frac1x} \leq \abs x \to 0$,
so the limit at $0$ exists and equals $0$.
\end{solution}

\begin{solution}{exo:b1:continuity:2}
$f = \lfloor\cdot\rfloor$ is continuous on $\R \setminus \Z$
(locally constant) and discontinuous at each $n \in \Z$: left limit
$n - 1$, value $n$.

$g(x) = x - \lfloor x\rfloor$: same discontinuity points (the
identity is continuous, so $g$ inherits the jumps of $f$); at $n \in
\Z$, left limit $1 \neq 0 = g(n)$.

$h(x) = \lfloor x\rfloor + (x - \lfloor x\rfloor)^2$: on $\intco{n}{n
+ 1}$, $h(x) = n + (x - n)^2$, continuous there; at $x = n$, the left
limit is $(n-1) + 1 = n = h(n)$: the jumps cancel. $h$ is continuous
on $\R$ (and strictly increasing).
\end{solution}

\begin{solution}{exo:b1:continuity:3}
$P(x) = x^5 - 3x + 1$: $P(-2) = -32 + 6 + 1 = -25 < 0$; $P(0) = 1 >
0$; $P(1) = -1 < 0$; $P(2) = 27 > 0$. Three sign changes on the
disjoint segments $\intcc{-2}{0}$, $\intcc{0}{1}$, $\intcc{1}{2}$:
by \cref{thm:b1:continuity:ivt}, at least three roots. (Being of
degree $5$, $P$ has at most five; a variation study would show
exactly three.)
\end{solution}

\begin{solution}{exo:b1:continuity:4}
Let $P = a_{2m+1}X^{2m+1} + \dots$ with $a_{2m+1} > 0$ (else replace
$P$ by $-P$). Factoring the dominant term, $P(x) = a_{2m+1}
x^{2m+1}\bigl(1 + o(1)\bigr)$ as $x \to \pm\infty$: so $P(x) \to
+\infty$ at $+\infty$ and $-\infty$ at $-\infty$. Pick $a$ with
$P(a) < 0$ and $b$ with $P(b) > 0$: the intermediate value theorem on
$\intcc{a}{b}$ provides a root.
\end{solution}

\begin{solution}{exo:b1:continuity:5}
Let $g(x) = f(x) - x$, continuous on $\intcc{0}{1}$. Since $f$ maps
into $\intcc{0}{1}$: $g(0) = f(0) \geq 0$ and $g(1) = f(1) - 1 \leq
0$. By \cref{thm:b1:continuity:ivt}, $g(c) = 0$ for some $c$: a
fixed point.

Necessity of the hypotheses: on $\intcc{0}{1}$, the discontinuous map
$f(x) = 1$ for $x \leq \frac12$, $f(x) = 0$ for $x > \frac12$ has no
fixed point; on the interval $\intoo{0}{1}$ (not a segment), $f(x) =
\frac x2$ is continuous into $\intoo{0}{1}$ with no fixed point
(the candidate $0$ is missing); on $\R$, $f(x) = x + 1$.
\end{solution}

\begin{solution}{exo:b1:continuity:6}
$g(x) = \cos x - x$ is continuous, $g(0) = 1 > 0$, $g(1) = \cos 1 -
1 < 0$: a solution exists in $\intoo{0}{1}$
(\cref{thm:b1:continuity:ivt}). Uniqueness: $g$ is strictly
decreasing on $\R$ --- for $x \leq 0$, $g(x) \geq 1 - x > 0$ has no
zero anyway; and $g'(x) = -\sin x - 1 \leq 0$ with equality only at
isolated points ($x \equiv -\frac\pi2 \bmod 2\pi$), so $g$ is
strictly decreasing (\cref{ch:b1:derivative}; alternatively: on
$\intcc{0}{1}$, $\cos$ is strictly decreasing and $-x$ too, so $g$
is). A strictly monotone function vanishes at most once.
\end{solution}

\begin{solution}{exo:b1:continuity:7}
Fix $M = f(0) + 1$. There is $A > 0$ with $f(x) \geq M$ for $\abs x
\geq A$ (definition of the two infinite limits; take the larger
threshold). On the segment $\intcc{-A}{A}$, the extreme value theorem
(\cref{thm:b1:continuity:evt}) provides $c$ with $f(c) =
\inf_{\intcc{-A}{A}} f \leq f(0)$. For $\abs x \geq A$: $f(x) \geq M
> f(0) \geq f(c)$. So $f(c)$ is the global minimum.
\end{solution}

\begin{solution}{exo:b1:continuity:8}
On the segment $\intcc{0}{T}$, $f$ is bounded and attains its bounds
(\cref{thm:b1:continuity:evt}); by periodicity, these are the bounds
on all of $\R$, still attained.

Let $g(x) = f(x + \frac T2) - f(x)$, continuous. Then
\[
g(0) + g\bigl(\tfrac T2\bigr)
= \bigl(f(\tfrac T2) - f(0)\bigr) + \bigl(f(T) - f(\tfrac T2)\bigr)
= f(T) - f(0) = 0 :
\]
$g(0)$ and $g(\frac T2)$ have opposite signs (or one vanishes), so
the intermediate value theorem on $\intcc{0}{T/2}$ gives $c$ with
$g(c) = 0$, i.e.\ $f(c + \frac T2) = f(c)$.
\end{solution}

\begin{solution}{exo:b1:continuity:9}
First the inequality: for $0 \leq y \leq x$,
\[
\bigl(\sqrt y + \sqrt{x - y}\bigr)^2 = x + 2\sqrt{y(x-y)} \geq x,
\]
so $\sqrt x \leq \sqrt y + \sqrt{x - y}$, i.e.\ $\sqrt x - \sqrt y
\leq \sqrt{x - y}$. Hence $\abs{\sqrt x - \sqrt y} \leq
\sqrt{\abs{x - y}}$ for all $x, y \geq 0$.

Uniform continuity: given $\varepsilon > 0$, take $\delta =
\varepsilon^2$; then $\abs{x - y} \leq \delta$ implies $\abs{\sqrt x
- \sqrt y} \leq \sqrt\delta = \varepsilon$.

Not Lipschitz near $0$: $\frac{\sqrt x - \sqrt 0}{x - 0} =
\frac{1}{\sqrt x} \to +\infty$ as $x \to 0^+$, so no constant $k$
can dominate all difference quotients.
\end{solution}

\begin{solution}{exo:b1:continuity:10}
Suppose $f$ is injective, continuous, and not strictly monotonic.
Then there are $a < b < c$ in $I$ with $f(b)$ not between $f(a)$ and
$f(c)$ --- indeed, if for \emph{all} triples the middle value were
between the outer ones, $f$ would be monotonic (compare any two
pairs; a short case check). Say $f(b) > \max(f(a), f(c))$ (the other
case is symmetric, replace $f$ by $-f$). Choose $v$ with $\max(f(a),
f(c)) < v < f(b)$. By the intermediate value theorem applied on
$\intcc{a}{b}$ and on $\intcc{b}{c}$, there are $u_1 \in
\intoo{a}{b}$ and $u_2 \in \intoo{b}{c}$ with $f(u_1) = v = f(u_2)$:
two distinct points with equal images, contradicting injectivity.
\end{solution}

\begin{solution}{exo:b1:continuity:11}
$f(0) = f(0+0) = 2f(0)$ gives $f(0) = 0$; $f(-x) = -f(x)$ from $0 =
f(x - x)$. Set $\alpha = f(1)$. Induction: $f(n) = n\alpha$ for $n
\in \N$, then for $n \in \Z$ by oddness. For $q \in \N^*$:
$q\,f(\frac pq) = f(p) = p\alpha$ (add $\frac pq$ to itself $q$
times), so $f(\frac pq) = \alpha\frac pq$: $f = \alpha\,
\mathrm{id}$ on $\Q$.

Let now $x \in \R$ and $(r_n)$ a sequence of rationals with $r_n \to
x$ (density, \cref{thm:b1:reals:density}, applied in nested
intervals; or $r_n = \frac{\lfloor nx\rfloor}{n}$). Continuity:
$f(x) = \lim f(r_n) = \lim \alpha r_n = \alpha x$.
\end{solution}

\begin{solution}{exo:b1:continuity:12}
Let $\varepsilon > 0$. By the limit at $+\infty$, there is $A$ with
$\abs{f(x) - \ell} \leq \frac\varepsilon2$ for $x \geq A$; hence for
$x, y \geq A$: $\abs{f(x) - f(y)} \leq \varepsilon$ (no closeness
needed).

On the segment $\intcc{0}{A + 1}$, Heine's theorem
(\cref{thm:b1:continuity:heine}) gives $\delta_0 > 0$ for this
$\varepsilon$; set $\delta = \min(\delta_0, 1)$.

Now take any $x, y \geq 0$ with $\abs{x - y} \leq \delta$, say $x
\leq y$. If $y \leq A + 1$: both lie in the segment, and the Heine
$\delta_0$ applies. Otherwise $y > A + 1$, and then $x \geq y - 1 >
A$: both lie in $\intco{A}{+\infty}$, where the limit argument
applies. In both cases $\abs{f(x) - f(y)} \leq \varepsilon$: uniform
continuity.
\end{solution}

\begin{solution}{pb:b1:continuity:1}
\textbf{1.} For $n \in \N$: $f(nx) = n f(x)$ by induction
($f((n+1)x) = f(nx) + f(x)$). Also $f(0) = 2f(0)$ gives $f(0) =
0$, and $0 = f(x - x) = f(x) + f(-x)$ gives oddness, so $f(nx) =
nf(x)$ for $n \in \Z$. For $r = \frac pq$: $q\,f\bigl(\tfrac
pq x\bigr) = f(px) = p\,f(x)$, so $f(rx) = r f(x)$: $f$ is
$\Q$-linear.

\textbf{2.} Additivity gives, for all $x$ and $h$:
\[
f(x + h) - f(x) = f(h) = f(x_0 + h) - f(x_0) .
\]
As $h \to 0$, the right side tends to $0$ by continuity at
$x_0$; hence $f(x + h) \to f(x)$: continuity at every $x$. Then
\cref{exo:b1:continuity:11} yields $f(x) = f(1)\,x$.

\textbf{3.} Two continuous additive functions are of the form
$cx$ and $c'x$ (question 2); if they agree on the dense subgroup
$\Z + \sqrt2\,\Z$ (\cref{exo:b1:reals:9}), then $c\,g = c'g$ for
some $g \neq 0$ in it: $c = c'$, the functions are equal. Without
continuity: $\Q$-linearity only ties values at
$\Q$-combinations, and $1, \sqrt2$ are $\Q$-independent ($\sqrt2
\notin \Q$), so $f(1) = 0$, $f(\sqrt2) = 1$ is consistent and
forces, on the subgroup,
\[
f(m + n\sqrt2) = m\,f(1) + n\,f(\sqrt2) = n .
\]

\textbf{4.} For $t \in \intcc{0}{\ell}$: $a + t \in
\intcc{a}{b}$, so $f(t) = f(a + t) - f(a) \leq M - f(a) =: M'$.

\textbf{5.} For $t \in \intcc{0}{\ell}$, also $\ell - t \in
\intcc{0}{\ell}$, and additivity gives $f(t) = f(\ell) - f(\ell
- t) \geq f(\ell) - M'$. Hence $\abs{f} \leq C$ on
$\intcc{0}{\ell}$ with $C = \max\bigl(\abs{M'}, \abs{f(\ell) -
M'}\bigr)$.

\textbf{6.} For $t \in \intcc{0}{\ell/n}$: $nt \in
\intcc{0}{\ell}$ and $f(t) = \frac{f(nt)}{n}$ (question 1), so
$\abs{f(t)} \leq \frac Cn$. Given $\varepsilon > 0$, choose $n >
\frac C\varepsilon$: for $0 \leq h \leq \frac\ell n$,
$\abs{f(h)} \leq \varepsilon$, and for negative $h$ use $f(h) =
-f(-h)$. So $f(h) \to 0 = f(0)$ as $h \to 0$: continuity at $0$,
hence everywhere (question 2), hence $f(x) = f(1)x$.

\textbf{7.} If $f$ is nondecreasing on $\intcc{a}{b}$, then
$f(a) \leq f(x) \leq f(b)$ there: bounded, so linear by question
6, $f(x) = cx$; and $c(b - a) = f(b) - f(a) \geq 0$ forces $c
\geq 0$.

\textbf{8.} (a)$\Rightarrow$(b)$\Rightarrow$(c): trivial.
(c)$\Rightarrow$(a): question 2. (a)$\Rightarrow$(d): a linear
function is monotone everywhere. (d)$\Rightarrow$(e): a monotone
function on $\intcc{a}{b}$ is bounded there by its endpoint
values. (e)$\Rightarrow$(a): questions 4--6. The five
statements are equivalent --- the regularity ladder.

\textbf{9.} Question 4 used only the upper bound $M$; question 5
\emph{derived} the lower bound from the upper one via the
reflection $f(t) = f(\ell) - f(\ell - t)$; question 6 then ran on
$\abs f \leq C$. So: additive and bounded \emph{above} on one
nondegenerate interval already implies linear. For a bound
below, apply this to $-f$ (additive, bounded above). Strongest
rung (e): a one-sided bound on one interval suffices.

\textbf{10.} If $\frac{f(x)}{x}$ were one constant $c$ for all
$x \neq 0$, $f$ would be linear; so there are nonzero $u, v$
with $\frac{f(u)}u \neq \frac{f(v)}v$, i.e.\ $u f(v) - v f(u)
\neq 0$: the determinant of the vectors $(u, f(u))$, $(v,
f(v))$ is nonzero, and they span $\R^2$.

\textbf{11.} For $r, s \in \Q$: $f(ru + sv) = r f(u) + s f(v)$
(question 1 twice plus additivity), so the graph contains
\[
\bigl(ru + sv,\; r f(u) + s f(v)\bigr)
= r\,(u, f(u)) + s\,(v, f(v)) .
\]
Given $(x_0, y_0)$ and $\varepsilon > 0$: the $2 \times 2$
system $a(u, f(u)) + b(v, f(v)) = (x_0, y_0)$ has a (unique)
real solution $(a, b)$ since the determinant is nonzero. Pick
rationals $r_n \to a$, $s_n \to b$: then $r_n(u, f(u)) +
s_n(v, f(v)) \to (x_0, y_0)$ coordinatewise, and each of these
points lies on the graph: the graph is dense in $\R^2$.

\textbf{12.} Let $I$ be a nondegenerate interval, $x_0$ its
midpoint, $M$ arbitrary: density provides a graph point within
$\min\bigl(\frac{\abs I}{2}, 1\bigr)$ of $(x_0, M + 1)$, i.e.\
$x \in I$ with $f(x) > M$: unbounded on $I$, hence (question 8)
discontinuous at every point and monotone on no interval; and
for any target $y_0$, graph points near $(x_0, y_0)$ give $f(x)$
arbitrarily close to $y_0$ with $x \in I$: $f(I)$ is dense in
$\R$. This is question 8 read backwards: since all rungs are
equivalent, a nonlinear additive function must fail \emph{all}
of them, everywhere.

\textbf{13.} Well defined: $m + n\sqrt2 = m' + n'\sqrt2$ forces
$(n - n')\sqrt2 = m' - m \in \Z$, so $n = n'$ (else $\sqrt2 \in
\Q$) and $m = m'$. Additivity on $G$ is then clear coordinate by
coordinate. Unboundedness near $0^+$: fix $\varepsilon \in
\intoo{0}{1}$ and $N \in \N$. For each fixed $n$ with $\abs n
\leq N$, the condition $m + n\sqrt2 \in \intoo{0}{1}$ pins $m$
inside an interval of length $1$: at most one integer $m$ per
$n$, so at most $2N + 1$ elements of $G \cap \intoo{0}{1}$ have
$\abs{f} \leq N$. But $G \cap \intoo{0}{\varepsilon}$ is
infinite ($G$ is dense, \cref{exo:b1:reals:9}); hence it
contains some $g$ with $\abs{f(g)} > N$: $f$ is unbounded on
every right neighborhood of $0$. Extending $f$ to an additive
function on $\R$ means choosing values coherently on a family of
reals that is $\Q$-linearly independent and spans $\R$ over $\Q$
--- a Hamel basis; producing one requires the axiom of choice,
and no explicit formula can do it: constructively we own the
monster only on $G$.

\textbf{14.} $f(x) = f(\frac x2 + \frac x2) = f(\frac x2)^2 \geq
0$. If $f(x_0) = 0$, then $f(x) = f(x - x_0)f(x_0) = 0$ for all
$x$: excluded. So $f > 0$ and $g = \ln \circ f$ is continuous
(\cref{prop:b1:continuity:algebra}) with $g(x + y) = g(x) +
g(y)$: by \cref{exo:b1:continuity:11}, $g(x) = cx$, so $f(x) =
\eu^{cx}$. Conversely each $\eu^{cx}$ works: the continuous
morphisms $(\R, +) \to (\R^*, \times)$ are exactly the
exponentials.

\textbf{15.} $g(u) = f(\eu^u)$ is continuous and $g(u + v) =
f(\eu^u \eu^v) = g(u) + g(v)$: $g(u) = cu$, and every $x > 0$
writes $x = \eu^u$ with $u = \ln x$: $f(x) = c\ln x$.

\textbf{16.} $h = \ln \circ f$ is continuous on
$\intoo{0}{+\infty}$ with $h(xy) = h(x) + h(y)$: by question 15,
$h(x) = c\ln x$, so $f(x) = \eu^{c\ln x} = x^c$.

\textbf{17.} $x = y = 0$: $f(0) = 2f(0) + f(0)^2$, so $f(0)(1 +
f(0)) = 0$. If $f(0) = -1$: setting $y = 0$, $f(x) = f(x) +
f(0) + f(x)f(0) = -1$ for all $x$: the constant $f \equiv -1$
(which indeed satisfies the equation). Otherwise $f(0) = 0$; $h
= 1 + f$ is continuous, $h(0) = 1$, and
\[
h(x + y) = 1 + f(x) + f(y) + f(x)f(y) = h(x)\,h(y) :
\]
by question 14, $h(x) = \eu^{cx}$, i.e.\ $f(x) = \eu^{cx} - 1$
(the case $c = 0$ giving $f \equiv 0$). Complete list: $f
\equiv -1$ and $f(x) = \eu^{cx} - 1$, $c \in \R$.

\textbf{18.} $x = y = 0$: $2f(0) = 4f(0)$, so $f(0) = 0$. $x =
0$: $f(y) + f(-y) = 2f(y)$, so $f$ is even. $y = x$: $f(2x) =
4f(x)$. Induction using $(x, y) \to (nx, x)$:
\[
f((n{+}1)x) = 2f(nx) + 2f(x) - f((n{-}1)x)
= (2n^2 + 2 - (n-1)^2) f(x) = (n+1)^2 f(x) .
\]
Then $f(x) = f\bigl(q\cdot\frac xq\bigr) = q^2 f\bigl(\frac
xq\bigr)$ gives $f\bigl(\frac pq x\bigr) = \frac{p^2}{q^2}f(x)$:
$f(r) = r^2 f(1)$ on $\Q$ (evenness handles signs). The two
continuous functions $f$ and $x \mapsto f(1)x^2$ agree on the
dense set $\Q$, hence everywhere (question 24): $f(x) =
f(1)\,x^2$; every $cx^2$ satisfies the equation.

\textbf{19.} $g = f - f(0)$ is continuous, $g(0) = 0$, and
satisfies Jensen's equation (the constants cancel). Taking $y =
0$: $g\bigl(\frac x2\bigr) = \frac{g(x)}{2}$. Then for all $x,
y$:
\[
\frac{g(x + y)}{2} = g\Bigl(\frac{x+y}{2}\Bigr)
= \frac{g(x) + g(y)}{2} ,
\]
so $g$ is additive and continuous: $g(x) = cx$ (question 2), and
$f(x) = cx + d$ with $d = f(0)$. All affine functions satisfy
Jensen: the list is complete.

\textbf{20.} Induction on $n$. For $n = 0$: $\lambda \in \{0,
1\}$, trivial. Assume the inequality for all weights
$\frac{k}{2^n}$. A weight $\lambda = \frac{k}{2^{n+1}}$ with $k$
even reduces to level $n$; for $k = 2j + 1$, $\lambda$ is the
midpoint of $\lambda_1 = \frac{j}{2^n}$ and $\lambda_2 =
\frac{j+1}{2^n}$. With $z_i = \lambda_i x + (1 - \lambda_i)y$:
$\lambda x + (1-\lambda)y = \frac{z_1 + z_2}{2}$, so
\[
f\bigl(\lambda x + (1{-}\lambda)y\bigr)
\leq \frac{f(z_1) + f(z_2)}{2}
\leq \frac{(\lambda_1 + \lambda_2)f(x) + (2 - \lambda_1 -
\lambda_2)f(y)}{2}
= \lambda f(x) + (1 - \lambda)f(y) .
\]

\textbf{21.} Fix $x, y$. The maps $\lambda \mapsto f(\lambda x +
(1 - \lambda)y)$ and $\lambda \mapsto \lambda f(x) + (1 -
\lambda)f(y)$ are continuous on $\intcc{0}{1}$ (composition and
algebra, \cref{prop:b1:continuity:algebra}). The inequality
holds on the dyadic weights, which are dense in $\intcc{0}{1}$
(\cref{exo:b1:reals:8}); for arbitrary $\lambda$ take dyadics
$\lambda_n \to \lambda$ and pass to the limit
(\cref{thm:b1:continuity:seqchar} and \cref{thm:b1:seq:order}):
the convexity inequality holds for every $\lambda \in
\intcc{0}{1}$ --- midpoint convexity plus continuity equals
convexity (the notion of \cref{ch:b1:derivative}).

\textbf{22.} A nonlinear additive $f$ satisfies
$f\bigl(\frac{x+y}{2}\bigr) = \frac{f(x) + f(y)}{2}$ exactly
(question 1 with $r = \frac12$, then additivity): it is
midpoint-convex, even midpoint-affine. If it satisfied the full
convexity inequality on some interval $\intcc{x}{y}$, then for
$\lambda \in \intcc{0}{1}$: $f(\lambda x + (1-\lambda)y) \leq
\max(f(x), f(y))$ --- bounded above on a nondegenerate interval,
hence linear by question 9: contradiction. So continuity in
question 21 is not a luxury: without it, midpoint convexity
controls only the countable dyadic skeleton, and the continuum
in between runs wild.

\textbf{23.} $p(x) = \frac{x^2}{2}$ satisfies $p(x+y) = p(x) +
p(y) + xy$. If $f$ is any continuous solution, $g = f - p$ is
continuous and additive, so $g(x) = cx$:
\[
f(x) = \frac{x^2}{2} + cx , \qquad c \in \R ,
\]
and each of these is a solution: the list is complete.

\textbf{24.} General principle: if $u, v$ are continuous and
agree on a dense $D \subseteq \R$, then for $x \in \R$ choose
$d_n \in D$ with $d_n \to x$
(\cref{prop:b1:topology:closureprops}); $u(x) = \lim u(d_n) =
\lim v(d_n) = v(x)$. Now let $f$ be continuous and additive on
the dense subgroup $G$. Fix $x, y \in \R$ and take $g_n \to x$,
$g'_n \to y$ with $g_n, g'_n \in G$; then $g_n + g'_n \to x + y$
and, by sequential continuity at $x + y$, $x$ and $y$:
\[
f(x + y) = \lim f(g_n + g'_n) = \lim\bigl(f(g_n) +
f(g'_n)\bigr) = f(x) + f(y) :
\]
$f$ is additive on all of $\R$ (hence linear, by question 2).

\textbf{25.} (i) For additive $f$: linear $\iff$ continuous
$\iff$ continuous at one point $\iff$ monotone on some interval
$\iff$ bounded (even one-sidedly) on some interval. (ii) Density
of the graph in the plane shows the failure is not a local
defect but a global explosion: above every subinterval the
values smear over all of $\R$, so every regularity property
fails everywhere at once. (iii) The catch: $cx$, $\eu^{cx}$,
$c\ln x$, $x^c$, $cx^2$, and $cx + d$ --- six characterizations,
one method: transport the equation to Cauchy's, prove the
$\Q$-skeleton by induction, upgrade to $\R$ by density plus
continuity. (iv) Density of $\Q$ (or the dyadics) carried the
upgrades in \cref{exo:b1:continuity:11} and in question 21; it
carried nothing in question 13, because without continuity
values do not propagate from a dense set to its closure ---
density transfers information only along continuity.
\end{solution}

\section*{Chapter \ref{ch:b1:derivative} --- Differentiation}

\begin{solution}{exo:b1:derivative:1}
$x^x = \eu^{x\ln x}$ on $\intoo{0}{+\infty}$: derivative $(\ln x +
1)\,x^x$.

$\ln(x + \sqrt{x^2+1})$ on $\R$ (the argument is always $> 0$):
derivative $\frac{1}{\sqrt{x^2+1}}$ (computed in
\cref{prop:b1:functions:invhyp} --- it is
$\operatorname{arsinh}$).

$\arctan\frac1x$ on $\R^*$: derivative $\frac{-1/x^2}{1 + 1/x^2} =
\frac{-1}{1 + x^2}$ (consistent with
\cref{prop:b1:functions:arcidentities}~(2): the function is
$\pm\frac\pi2 - \arctan x$ on each half-line).

$\sqrt{1 + \eu^{2x}}$ on $\R$: derivative
$\frac{\eu^{2x}}{\sqrt{1 + \eu^{2x}}}$.
\end{solution}

\begin{solution}{exo:b1:derivative:2}
At $0$: $\bigl|\frac{f(h) - 0}{h}\bigr| = \abs{h \sin\frac1h} \leq
\abs h \to 0$, so $f'(0) = 0$. For $x \neq 0$, the usual rules give
$f'(x) = 2x\sin\frac1x - \cos\frac1x$. Along $x_n = \frac{1}{2\pi
n}$: $f'(x_n) = 0 - 1 \to -1$; along $y_n = \frac{1}{(2n+1)\pi}$:
$f'(y_n) = 0 + 1 \to 1$. Two sequences tending to $0$ with different
limits of $f'$: no limit (\cref{thm:b1:continuity:seqchar}), so $f'$
is not continuous at $0$ and $f$ is differentiable without being
$C^1$.
\end{solution}

\begin{solution}{exo:b1:derivative:3}
$\ln(1+x) < x$ for $x > 0$: concavity tangent inequality at $0$
(strict away from the contact point since $\ln$ is strictly concave;
or apply the mean value theorem: $\ln(1+x) = \frac{x}{1+c}$ for some
$c \in \intoo{0}{x}$, and $\frac{x}{1+c} < x$). The same mean value
identity gives the lower bound: $\frac{x}{1+c} > \frac{x}{1+x}$.

Consequence: with $x/n$ in place of $x$,
\[
\frac{x/n}{1 + x/n} < \ln\Bigl(1 + \frac xn\Bigr) < \frac xn
\quad\implies\quad
\frac{x}{1 + x/n} < n \ln\Bigl(1 + \frac xn\Bigr) < x .
\]
The left member tends to $x$: by the squeeze, $n\ln(1 + \frac xn)
\to x$, and by continuity of $\exp$, $\bigl(1 + \frac xn\bigr)^n =
\eu^{n\ln(1 + x/n)} \to \eu^x$.
\end{solution}

\begin{solution}{exo:b1:derivative:4}
Let $x_1 < x_2 < \dots < x_k$ be distinct roots of $P$. On each
$\intcc{x_i}{x_{i+1}}$, Rolle (\cref{thm:b1:derivative:rolle})
produces $c_i \in \intoo{x_i}{x_{i+1}}$ with $P'(c_i) = 0$: that is
$k - 1$ roots of $P'$, distinct because the open intervals are
disjoint --- and interlaced by construction.

If $P$ (degree $n$) has all roots real, write them with
multiplicities $m_1 + \dots + m_k = n$. Each root of multiplicity
$m_i \geq 2$ is a root of $P'$ of multiplicity $m_i - 1$
(\cref{prop:b1:poly:multiplicity}), contributing $\sum (m_i - 1) = n
- k$; Rolle contributes $k - 1$ more, all distinct from these. Total
$\geq n - 1 = \deg P'$: all roots of $P'$ are real.
\end{solution}

\begin{solution}{exo:b1:derivative:5}
Fix $\varepsilon > 0$ and $A$ with $\abs{f'(t) - \ell} \leq
\varepsilon$ for $t \geq A$. For $x > A$, the mean value theorem on
$\intcc{A}{x}$ gives $c \in \intoo{A}{x}$ with
\[
f(x) = f(A) + f'(c)(x - A),
\qquad\text{so}\qquad
\Bigl|\frac{f(x)}{x} - \ell\Bigr|
\leq \frac{\abs{f(A)} + \abs\ell A}{x} + \abs{f'(c) - \ell}
\cdot\frac{x - A}{x} \leq \frac{C_A}{x} + \varepsilon .
\]
For $x$ large, $\frac{C_A}{x} \leq \varepsilon$: hence
$\frac{f(x)}{x} \to \ell$.

Yes: $f(x+1) - f(x) = f'(c_x)$ with $c_x \in \intoo{x}{x+1}$ (mean
value theorem on $\intcc{x}{x+1}$), and $c_x \to +\infty$, so
$f(x+1) - f(x) \to \ell$.
\end{solution}

\begin{solution}{exo:b1:derivative:6}
Induction on $n$. For $n = 1$: Rolle. If the claim holds for $n - 1$:
$f$ vanishes at $n+1$ points, so by Rolle applied on the $n$ gaps,
$f'$ vanishes at $n$ distinct points; the induction hypothesis
applied to $f'$ ($n-1$ times differentiable, $n$ zeros) makes
$(f')^{(n-1)} = f^{(n)}$ vanish somewhere.

Application: if $P$ of degree $\leq n$ vanishes at $n+1$ points, then
$P^{(n)}$, a constant equal to $n!$ times the leading coefficient,
vanishes: the leading coefficient is $0$, and one concludes by
downward induction (or directly: all coefficients vanish).
\end{solution}

\begin{solution}{exo:b1:derivative:7}
By the mean value theorem on $\intcc{a}{x_0}$ and on
$\intcc{x_0}{b}$:
\[
f'(c_1) = \frac{f(x_0) - f(a)}{x_0 - a} = \frac{f(x_0)}{x_0 - a} > 0,
\qquad
f'(c_2) = \frac{f(b) - f(x_0)}{b - x_0} = \frac{-f(x_0)}{b - x_0} < 0,
\]
with $c_1 < x_0 < c_2$. Then the mean value theorem applied to $f'$
on $\intcc{c_1}{c_2}$ gives $c$ with
\[
f''(c) = \frac{f'(c_2) - f'(c_1)}{c_2 - c_1} < 0 . \qedhere
\]
\end{solution}

\begin{solution}{exo:b1:derivative:8}
$f'(x) = \frac{1 - \ln x}{x^2}$: $f$ increases on $\intoc{0}{\eu}$,
decreases on $\intco{\eu}{+\infty}$, with maximum $f(\eu) =
\frac1\eu$; limits $-\infty$ at $0^+$ and $0$ at $+\infty$ (growth
comparison).

For $\eu \leq a < b$: $f$ strictly decreasing there gives
$\frac{\ln a}{a} > \frac{\ln b}{b}$, i.e.\ $b \ln a > a \ln b$,
i.e.\ $a^b > b^a$.

With $a = \eu < b = \pi$: $\eu^\pi > \pi^\eu$.

Small pairs: $(2, 3)$: $2^3 = 8 < 9 = 3^2$ --- reversed! The reason:
$2 < \eu$, and on $\intoo{0}{\eu}$ the function $f$ is
\emph{increasing}, so the comparison flips when both numbers sit
below $\eu$, and is unpredictable across $\eu$ ($f(2) = f(4)$
explains the tie $2^4 = 4^2 = 16$).
\end{solution}

\begin{solution}{exo:b1:derivative:9}
\emph{Jensen for $\ln$, by induction on $n$.} Claim: for positive
$x_i$ and weights $\lambda_i > 0$ with $\sum \lambda_i = 1$:
$\ln\bigl(\sum \lambda_i x_i\bigr) \geq \sum \lambda_i \ln x_i$. For
$n = 2$ this is concavity. Step: with $\Lambda = \lambda_1 + \dots +
\lambda_{n-1} = 1 - \lambda_n$ and $y = \sum_{i<n}
\frac{\lambda_i}{\Lambda} x_i$,
\[
\ln\Bigl(\sum_{i \leq n} \lambda_i x_i\Bigr)
= \ln\bigl(\Lambda y + \lambda_n x_n\bigr)
\geq \Lambda \ln y + \lambda_n \ln x_n
\geq \Lambda \sum_{i<n} \frac{\lambda_i}{\Lambda}\ln x_i
+ \lambda_n \ln x_n,
\]
using concavity ($n = 2$) then the induction hypothesis.

With $\lambda_i = \frac 1n$ and $x_i = a_i$: $\ln\frac{\sum a_i}{n}
\geq \frac 1n \sum \ln a_i = \ln\sqrt[n]{a_1\cdots a_n}$;
exponentiate. Equality: $\ln$ is \emph{strictly} concave ($\ln'' < 0$),
so equality at each step forces the averaged points to coincide ---
i.e.\ all $a_i$ equal; and if all are equal, equality is clear.
\end{solution}

\begin{solution}{exo:b1:derivative:10}
Let $g(x) = f(x) - vx$: differentiable, with $g'(a) = f'(a) - v < 0$
and $g'(b) = f'(b) - v > 0$. By the extreme value theorem, $g$
attains its minimum on $\intcc{a}{b}$ at some $c$. It is not at $a$:
since $g'(a) < 0$, points just right of $a$ have $g < g(a)$. It is
not at $b$: since $g'(b) > 0$, points just left of $b$ have $g <
g(b)$. So $c$ is interior, and \cref{prop:b1:derivative:fermat}
gives $g'(c) = 0$, i.e.\ $f'(c) = v$.
\end{solution}

\begin{solution}{exo:b1:derivative:11}
\emph{Uniqueness:} two fixed points $\ell \neq \ell'$ would give
$\abs{\ell - \ell'} = \abs{f(\ell) - f(\ell')} \leq k\abs{\ell -
\ell'} < \abs{\ell - \ell'}$, absurd.

\emph{Existence:} $g(x) = f(x) - x$ satisfies, by the mean value
inequality, $f(x) \leq f(0) + k\abs x$; so for $x \geq
\frac{\abs{f(0)}}{1 - k}$, $g(x) \leq f(0) + kx - x \leq 0$, and
symmetrically $g(-x) \geq 0$ for $x$ large. The intermediate value
theorem gives a zero $\ell$ of $g$: a fixed point.

\emph{Convergence:} the mean value inequality again:
\[
\abs{u_{n+1} - \ell} = \abs{f(u_n) - f(\ell)} \leq k\abs{u_n - \ell},
\]
so by induction $\abs{u_n - \ell} \leq k^n \abs{u_0 - \ell} \to 0$.
\end{solution}

\begin{solution}{exo:b1:derivative:12}
\begin{enumerate}
  \item If $g(b) = g(a)$, Rolle would give an interior zero of
        $g'$: excluded. Set $\lambda = \frac{f(b) - f(a)}{g(b) -
        g(a)}$ and $h = f - \lambda g$: $h$ is continuous on
        $\intcc{a}{b}$, differentiable inside, and $h(b) - h(a)
        = f(b) - f(a) - \lambda(g(b) - g(a)) = 0$. Rolle
        provides $c$ with $h'(c) = 0$, i.e.\ $f'(c) =
        \lambda\,g'(c)$; divide by $g'(c) \neq 0$.
  \item For $x > a$ close to $a$, part (1) on $\intcc{a}{x}$
        (where $g' \neq 0$) gives $g(x) \neq 0$ and $c_x \in
        \intoo{a}{x}$ with
        \[
        \frac{f(x)}{g(x)} = \frac{f(x) - f(a)}{g(x) - g(a)}
        = \frac{f'(c_x)}{g'(c_x)} .
        \]
        As $x \to a^+$, $c_x \to a^+$ (squeeze), so the right
        side tends to $\ell$: $\frac{f}{g} \to \ell$.
  \item $\frac{f(x)}{g(x)} = x\sin\frac1x \to 0$, while
        $\frac{f'(x)}{g'(x)} = 2x\sin\frac1x - \cos\frac1x$ has
        no limit at $0$ (\cref{exo:b1:derivative:2}):
        l'Hospital's rule transfers information only from
        $\frac{f'}{g'}$ to $\frac fg$, never back.
\end{enumerate}
\end{solution}

\begin{solution}{pb:b1:derivative:1}
\textbf{1.} $\bigl|\frac ab - \frac pq\bigr| = \frac{\abs{aq -
bp}}{bq}$, and $aq - bp$ is a nonzero integer when the fractions
differ: the distance is $\geq \frac{1}{bq}$. So no rational
other than $x$ itself enters the punctured interval of radius
$\frac{1}{bq}$ around $x = \frac ab$.

\textbf{2.} If $\bigl|\sqrt2 - \frac pq\bigr| \geq 1 \geq
\frac{1}{4q^2}$, done. Otherwise $\frac pq \in \intoo{\sqrt2 -
1}{\sqrt2 + 1}$, so $0 < \sqrt2 + \frac pq < 2\sqrt2 + 1 < 4$.
Since $\sqrt 2 \notin \Q$, $p^2 - 2q^2$ is a nonzero integer, and
\[
\Bigl|\sqrt2 - \frac pq\Bigr|
= \frac{\abs{2q^2 - p^2}}{q^2\,\bigl(\sqrt2 + \frac pq\bigr)}
\geq \frac{1}{4q^2} .
\]

\textbf{3.} $(p + 2q)^2 - 2(p + q)^2 = -(p^2 - 2q^2)$: the value
$\pm1$ propagates. From $(1,1)$:
\[
(3,2),\ (7,5),\ (17,12),\ (41,29),\ (99,70),
\]
with $p^2 - 2q^2$ alternating $-1, +1, \dots$
For these, $\frac pq \geq 1$, so $\sqrt2 + \frac pq > 2$ and
\[
\Bigl|\sqrt2 - \frac pq\Bigr| =
\frac{1}{q^2(\sqrt2 + p/q)} < \frac{1}{2q^2} ,
\]
with $q \to \infty$: infinitely many order-$2$ approximations.
With question 2, the exponent $2$ is exact for $\sqrt 2$.

\textbf{4.} The $N + 1$ numbers $kx - \lfloor kx\rfloor$ ($0
\leq k \leq N$) lie in the $N$ boxes $\intco{\frac
jN}{\frac{j+1}{N}}$: two share a box
(\cref{cor:b1:counting:pigeonhole}), say for $i < j$. With $q =
j - i \leq N$ and $p = \lfloor jx\rfloor - \lfloor ix\rfloor$:
$\abs{qx - p} < \frac1N$, so $\bigl|x - \frac pq\bigr| <
\frac{1}{Nq} \leq \frac{1}{q^2}$. Letting $N \to \infty$: since
$x$ is irrational, each fixed fraction has positive distance to
$x$, while $\frac{1}{Nq} \leq \frac 1N \to 0$ forces new
fractions to appear: infinitely many distinct $\frac pq$ with
$\bigl|x - \frac pq\bigr| < \frac{1}{q^2}$.

\textbf{5.} Start from any nonzero integer $P_0$ with $P_0(x) =
0$. If $P_0$ has a rational root $\frac ab$, the factor theorem
(\cref{thm:b1:poly:factor}) writes $P_0 = \bigl(X - \frac
ab\bigr)Q$ with $Q \in \Q[X]$;
since $x \neq \frac ab$ ($x$ irrational), $Q(x) = 0$, and
clearing denominators gives a nonzero \emph{integer} polynomial
of smaller degree vanishing at $x$. The degree drops at each
step, so the process stops: we reach $P \in \Z[X]$, $P(x) = 0$,
with no rational root, of some degree $d$. If $d \leq 1$, $P =
uX + v$ would make $x = -\frac vu$ rational: so $d \geq 2$.

\textbf{6.} $q^d\,P\bigl(\frac pq\bigr) = a_d p^d + a_{d-1}
p^{d-1} q + \dots + a_0 q^d$ is an integer, and it is nonzero
because $P$ has no rational root: $\bigl|P\bigl(\frac
pq\bigr)\bigr| \geq q^{-d}$.

\textbf{7.} Note $M > 0$: $P'$ is a nonzero polynomial ($d \geq
2$), so it cannot vanish identically on
$\intcc{x-1}{x+1}$. If $\bigl|x - \frac pq\bigr| > 1$, then it
exceeds $\frac{C}{q^d}$ trivially. Otherwise $\frac pq \in
\intcc{x-1}{x+1}$ and the mean value theorem
(\cref{thm:b1:derivative:mvt}) gives $c$ between $x$ and
$\frac pq$ with
\[
\Bigl|P\Bigl(\frac pq\Bigr)\Bigr|
= \Bigl|P\Bigl(\frac pq\Bigr) - P(x)\Bigr|
= \abs{P'(c)}\,\Bigl|x - \frac pq\Bigr|
\leq M\,\Bigl|x - \frac pq\Bigr| ,
\]
so with question 6: $\bigl|x - \frac pq\bigr| \geq
\frac{1}{Mq^d} \geq \frac{C}{q^d}$.

\textbf{8.} Suppose $x = \frac ab$ is Liouville. Pick $n$ with
$2^{n-1} > b$ and the corresponding $\frac pq$, $q \geq 2$:
\[
0 < \Bigl|x - \frac pq\Bigr| < \frac{1}{q^n}
= \frac{1}{q^{n-1}\,q} \leq \frac{1}{2^{n-1} q} <
\frac{1}{bq} ,
\]
contradicting question 1. So Liouville numbers are irrational.

\textbf{9.} If a Liouville $x$ were algebraic: it is irrational
(question 8), so questions 5--7 provide $d \geq 2$ and $C > 0$
with $\bigl|x - \frac pq\bigr| \geq \frac{C}{q^d}$ always. For
each $n$, the Liouville approximant gives $\frac{C}{q^d} <
q^{-n}$, i.e.\ $C < q^{d-n} \leq 2^{d-n}$ (as $q \geq 2$). For
$n$ large, $2^{d-n} < C$: contradiction. Liouville numbers are
transcendental.

\textbf{10.} Ones at positions $1, 2, 6, 24$; all other digits
among the first $25$ vanish:
\[
L = 0.1100010000\,0000000000\,00010\dots
\]

\textbf{11.} For $m > k$, the positions $n!$ with $n > k$ are
distinct integers $\geq (k+1)!$, so the finite geometric sum
gives
\[
t_m - t_k = \sum_{n=k+1}^{m} 10^{-n!}
\leq \sum_{j = (k+1)!}^{m!} 10^{-j}
< 10^{-(k+1)!}\,\frac{1}{1 - \frac1{10}}
= \frac{10}{9}\,10^{-(k+1)!} ;
\]
taking the supremum over $m$: $L - t_k \leq
\frac{10}{9}10^{-(k+1)!} < 2\cdot10^{-(k+1)!}$. Lower bound: $L
\geq t_{k+1} = t_k + 10^{-(k+1)!}$.

\textbf{12.} $p_k = 10^{k!}\,t_k \in \N$, $q_k = 10^{k!}$, and
$(k+1)! = (k+1)\,k!$ gives $10^{-(k+1)!} = q_k^{-(k+1)}$:
question 11 reads
\[
0 < L - \frac{p_k}{q_k} < \frac{2}{q_k^{\,k+1}} .
\]
Given $n$: for $k \geq n$, $2\,q_k^{-(k+1)} \leq q_k^{-n}$
(indeed $q_k^{\,k+1-n} \geq q_k \geq 10 > 2$), and $q_k \geq 2$:
the definition of question 8 is met. $L$ is a Liouville number.

\textbf{13.} By question 9, $L$ is transcendental --- the first
number in history proved transcendental (Liouville, 1844).
Digit cross-check: the string has infinitely many ones with
consecutive gaps $(k+1)! - k! = k\cdot k! \to \infty$, so it is
not eventually periodic, and $L \notin \Q$ by the periodicity
criterion of \cref{pb:b1:reals:1} (question 18) --- consistent.

\textbf{14.} With digits $d_n \in \intint{1}{9}$ at the
factorial positions: the tail bound of question 11 scales by at
most $9$: $0 < L' - t'_k \leq 9\cdot\frac{10}{9}\,10^{-(k+1)!} =
10\,q_k^{-(k+1)}$ (positivity because the digit at position
$(k+1)!$ is nonzero). For $k \geq n$: $10\,q_k^{-(k+1)} \leq
q_k^{-n}$ since $q_k^{\,k+1-n} \geq 10$: again a Liouville
number, hence transcendental. These values are pairwise
distinct for distinct digit choices (the strings are proper ---
zeros abound --- and proper strings determine their value,
\cref{pb:b1:reals:1}, question 10). Given any list $k \mapsto
x_k$ of them, choose the $k$-th factorial digit in
$\intint{1}{9}$ different from that of $x_k$: a number of the
same shape missing from the list. Uncountably many explicit
transcendentals.

\textbf{15.} First a lemma: \emph{if $\bigl|x - \frac pq\bigr|
\geq \frac{C}{q^s}$ for all $\frac pq \neq x$, then $x$ is not
approximable to any order $\mu > s$.} Indeed infinitely many
$\frac pq \neq x$ with $\bigl|x - \frac pq\bigr| <
\frac{c}{q^\mu}$ would force $\frac{C}{q^s} < \frac{c}{q^\mu}$,
i.e.\ $q^{\mu - s} < \frac cC$: the $q$ are bounded, and
boundedly many fractions lie within distance $1$ of $x$ ---
finitely many candidates, not infinitely many. Now the
hierarchy: rationals are approximable to order $1$ ($\frac
pq$ with $p = \lfloor qx\rfloor + 1$ gives error $\leq \frac1q <
\frac2q$) and to no order $\mu > 1$ (question 1 gives the
hypothesis of the lemma with $s = 1$, $C = \frac1b$); $\sqrt2$:
order $2$ (question 3) and no more (question 2 and the lemma);
every irrational: at least $2$ (question 4); algebraic of degree
$d$: at most $d$ (question 7 and the lemma); Liouville numbers:
every order (question 12's display, with $c = 2$).

\textbf{16.} With $r = \frac ab$: $\frac{p_k}{q_k} + \frac ab =
\frac{b p_k + a q_k}{b q_k} =: \frac{P_k}{Q_k}$, $Q_k = b q_k
\geq 2$, and
\[
\Bigl|(L + r) - \frac{P_k}{Q_k}\Bigr| = L - \frac{p_k}{q_k}
< 2\,q_k^{-(k+1)} = 2\,b^{\,k+1} Q_k^{-(k+1)} .
\]
Given $n$: for large $k$, $Q_k^{\,k+1-n} \geq Q_k = b\,10^{k!}
\geq 2\,b^{\,k+1}$ (the factorial crushes the power), so the
error is $< Q_k^{-n}$: $L + r$ is Liouville. Since $\Q$ is dense
and each $L + r$ is transcendental, transcendental numbers are
dense in $\R$.

\textbf{17.} For $h \geq 1$ there are finitely many $P \in
\Z[X]$ with $\deg P + \sum_i \abs{a_i} \leq h$ (degree $\leq h$
and each coefficient in $\intint{-h}{h}$: at most
$(2h+1)^{h+1}$). Every nonzero integer polynomial has such a
height, and has at most $\deg P$ real roots: the algebraic
numbers form a countable union (over $h$) of finite sets, hence
can be listed as a single sequence. If the transcendentals could
also be listed, interleaving the two lists would list $\R$,
contradicting \cref{pb:b1:reals:1} (question 22). So
transcendental numbers form an uncountable set. Comparison:
Cantor proves \emph{most} reals are transcendental yet exhibits
none; Liouville exhibits one, with effective constants (Part V)
--- existence by abundance versus existence by construction.

\textbf{18.} From question 2, the lemma hypothesis holds with $s
= 2$, $C = \frac14$. If $\sqrt2$ were Liouville, then for $n =
3$: $\frac{1}{4q^2} < q^{-3}$ forces $q < 4$, so $q \in \{2,
3\}$; only finitely many $\frac pq$ with these $q$ lie within
$1$ of $\sqrt2$, each at some positive distance
$\geq \varepsilon_0$ ($\sqrt2$ irrational); choosing $n$ with
$2^{-n} < \varepsilon_0$ leaves no admissible $\frac pq$ at all:
contradiction. The same argument with $\frac{C}{q^d}$ shows no
algebraic number is Liouville --- question 9 in effective
clothing.

\textbf{19.} $99^2 - 2\cdot70^2 = 9801 - 9800 = 1$. Hence
\[
\sqrt2 - \frac{99}{70} =
\frac{2 - (99/70)^2}{\sqrt2 + 99/70}
= \frac{-1}{4900\,\bigl(\sqrt2 + \tfrac{99}{70}\bigr)} ,
\qquad
\Bigl|\sqrt2 - \frac{99}{70}\Bigr|
= \frac{1}{4900 \times 2.8284\dots} \approx 7.2\cdot10^{-5} :
\]
$\frac{99}{70} = 1.414285\dots$ against $\sqrt2 =
1.414213\dots$ --- five correct digits.

\textbf{20.} Rational-root test for $P = X^3 - 2$: candidates
$\pm1, \pm2, \pm\frac12$, none a root. So $d = 3$ and Part II
applies to $x = 2^{1/3} = 1.2599\dots$ On $\intcc{x - 1}{x + 1}
\subseteq \intcc{0.25}{2.26}$: $\abs{P'(t)} = 3t^2 \leq 3\,(1 +
2^{1/3})^2 < 3\times(2.26)^2 = 15.32 < 16$, so $M < 16$ and $C
\geq \frac{1}{16}$:
\[
\Bigl|2^{1/3} - \frac pq\Bigr| \geq \frac{1}{16\,q^3}
\qquad\text{for all rationals.}
\]

\textbf{21.} If $\bigl|2^{1/3} - \frac pq\bigr| < 10^{-6}$, then
$\frac{1}{16 q^3} < 10^{-6}$, i.e.\ $q^3 > \frac{10^6}{16} =
62\,500$; since $39^3 = 59\,319 < 62\,500 \leq 64\,000 = 40^3$:
$q \geq 40$.

\textbf{22.} Run Part III in base $2$: $B = \sup_k \sum_{n\leq
k} 2^{-n!}$, $q_k = 2^{k!}$, and the geometric tail (ratio
$\frac12$) gives $0 < B - \frac{p_k}{q_k} \leq 2\cdot2^{-(k+1)!}
= 2\,q_k^{-(k+1)} \leq q_k^{-n}$ for $k \geq n$. So $B$ is
Liouville, hence transcendental: nothing in the argument is
decimal.

\textbf{23.} With ones at positions $3^k$: $q_k = 10^{3^k}$ and
the tail bound gives $0 < x^\dagger - \frac{p_k}{q_k} <
2\cdot10^{-3^{k+1}} = 2\,q_k^{-3}$ (as $3^{k+1} = 3\cdot3^k$):
infinitely many approximations of order $3$. By the lemma of
question 15: order $3 > 1$ rules out rationality, and order $3 >
2$ rules out being a quadratic irrational (whose Liouville
inequality has $s = d = 2$). But an algebraic number of degree
$\geq 3$ is only repelled at order $d \geq 3$: Liouville's
method cannot separate $x^\dagger$ from the cubics. The gap is
closed by Roth's theorem --- every algebraic irrational has
approximation order exactly $2$ --- a twentieth-century result
far beyond this volume; granted it, $x^\dagger$ too is
transcendental.

\textbf{24.} There are at most $(2H+1)^{d+1}$ tuples $(a_0,
\dots, a_d)$ with entries in $\intint{-H}{H}$, and each nonzero
polynomial among them has at most $d$ real roots: at most
$d\,(2H+1)^{d+1}$ algebraic numbers arise --- the finiteness
that let question 17 enumerate them all.

\textbf{25.} (i) The single analytic ingredient is the mean
value theorem, in question 7, converting the vanishing $P(x) =
0$ into the Lipschitz repulsion $\abs{P(p/q)} \leq M\abs{x -
p/q}$. (ii) The tension: integrality pushes $\abs{P(p/q)}$ up to
$q^{-d}$, smoothness pulls it down to $M\abs{x - p/q}$ --- a
rational too close to $x$ would be crushed between the two.
(iii) Liouville constructs one transcendental with effective
constants; Cantor shows almost all reals are transcendental
without naming one: construction versus cardinality. (iv) The
weekend problem of \cref{ch:b1:integration} proves the
irrationality of $\pi$ by the same squeeze --- an integral that
would be a positive integer yet is trapped in $\intoo{0}{1}$ ---
with integration replacing differentiation as the analytic
half.
\end{solution}

\section*{Chapter \ref{ch:b1:integration} --- Integration on a Segment}

\begin{solution}{exo:b1:integration:1}
$\dfrac{1}{x^2-4} = \dfrac{1/4}{x - 2} - \dfrac{1/4}{x+2}$
(cover-up), so
\[
\int_0^1 \frac{\dd x}{x^2 - 4}
= \frac14\Bigl[\ln\abs{x-2} - \ln\abs{x+2}\Bigr]_0^1
= \frac14\Bigl(\ln\frac{1}{3} - \ln 1\Bigr) = -\frac{\ln 3}{4}.
\]

By parts ($u' = 1$, $v = \ln t$): $\int_1^{\eu} \ln t\,\dd t =
\bigl[t\ln t\bigr]_1^{\eu} - \int_1^{\eu} \dd t = \eu - (\eu - 1) =
1$.

By parts ($u' = \sin t$, $v = t$): $\int_0^\pi t\sin t\,\dd t =
\bigl[-t\cos t\bigr]_0^\pi + \int_0^\pi \cos t\,\dd t = \pi + 0 =
\pi$.
\end{solution}

\begin{solution}{exo:b1:integration:2}
Substitution $u = t^2 + 1$, $\dd u = 2t\,\dd t$:
\[
\int_0^1 \frac{t\,\dd t}{(t^2+1)^2}
= \frac12 \int_1^2 \frac{\dd u}{u^2}
= \frac12\Bigl[-\frac1u\Bigr]_1^2 = \frac14 .
\]
With $u = \sin t$: $\int_0^{\pi/2} \cos^3 t\,\dd t = \int_0^{\pi/2}
(1 - \sin^2 t)\cos t\,\dd t = \bigl[\sin t - \frac{\sin^3
t}{3}\bigr]_0^{\pi/2} = 1 - \frac13 = \frac23$.
\end{solution}

\begin{solution}{exo:b1:integration:3}
$u_n = \dfrac1n \sum_{k=1}^{n} \dfrac{1}{1 + (k/n)^2}$: a Riemann sum
of $x \mapsto \frac{1}{1+x^2}$ on $\intcc{0}{1}$, so $u_n \to
\int_0^1 \frac{\dd x}{1+x^2} = \arctan 1 = \dfrac\pi4$.

$\ln v_n = \dfrac1n \sum_{k=1}^{n} \ln\dfrac{n+k}{n} = \dfrac 1n
\sum_{k=1}^n \ln\Bigl(1 + \dfrac kn\Bigr) \to \int_0^1 \ln(1+x)\,\dd
x = \bigl[(1+x)\ln(1+x) - x\bigr]_0^1 = 2\ln 2 - 1$. Hence $v_n \to
\eu^{2\ln 2 - 1} = \dfrac 4\eu$. (Check of the identification:
$\frac{(2n)!}{n!\,n^n} = \prod_{k=1}^{n} \frac{n+k}{n}$.)
\end{solution}

\begin{solution}{exo:b1:integration:4}
The limit is $0$. Let $M = \sup \abs f$ and $\varepsilon \in
\intoo{0}{1}$. Cut at $1 - \varepsilon$:
\[
\Bigl| \int_0^1 x^n f \Bigr|
\leq \int_0^{1 - \varepsilon} x^n \abs f + \int_{1-\varepsilon}^1
x^n \abs f
\leq M\,(1-\varepsilon)^n + M\varepsilon .
\]
Since $(1 - \varepsilon)^n \to 0$ (\cref{exo:b1:seq:3}), the limsup
of the left side is $\leq M\varepsilon$ for every $\varepsilon$:
the integral tends to $0$.
\end{solution}

\begin{solution}{exo:b1:integration:5}
$Q(\lambda) = \int_a^b (f + \lambda g)^2 = \int f^2 + 2\lambda \int
fg + \lambda^2 \int g^2 \geq 0$ for all $\lambda$. If $\int g^2 =
0$, then $g = 0$ (strict positivity,
\cref{thm:b1:integration:props}~(4)) and the inequality is $0 \leq
0$. Otherwise $Q$ is a genuine quadratic, everywhere $\geq 0$: its
discriminant is $\leq 0$, i.e.\ $\bigl(\int fg\bigr)^2 \leq \int
f^2 \int g^2$.

Equality iff the discriminant vanishes iff $Q(\lambda_0) = 0$ for
some $\lambda_0$, i.e.\ $\int (f + \lambda_0 g)^2 = 0$, i.e.\ (strict
positivity again) $f = -\lambda_0 g$: equality holds exactly when
$f$ and $g$ are proportional.
\end{solution}

\begin{solution}{exo:b1:integration:6}
Let $\Phi(a) = \int_a^{a+T} f$. By the fundamental theorem
(\cref{thm:b1:integration:ftc}), $\Phi$ is differentiable with
$\Phi'(a) = f(a + T) - f(a) = 0$: constant.

For $x > 0$, write $x = nT + r$, $0 \leq r < T$ ($n = \lfloor x/T
\rfloor$). Chasles:
\[
\int_0^x f = n \int_0^T f + \int_{nT}^{nT + r} f,
\qquad
\Bigl| \int_{nT}^{nT+r} f \Bigr| \leq T \sup_{\intcc{0}{T}} \abs f
= C .
\]
Then $\frac 1x \int_0^x f = \frac{nT}{x}\cdot\frac 1T \int_0^T f +
O\bigl(\frac 1x\bigr)$, and $\frac{nT}{x} \to 1$: the limit is
$\frac1T \int_0^T f$.
\end{solution}

\begin{solution}{exo:b1:integration:7}
Let $g(x) = f(x) - x$: continuous, with
\[
\int_0^1 g = \int_0^1 f - \frac12 = 0 .
\]
If $g$ never vanished, the intermediate value theorem would force a
constant sign (a continuous function on an interval taking both signs
vanishes); say $g > 0$. Then, by strict positivity
(\cref{thm:b1:integration:props}~(4) applied to $g > 0$, giving
$\int g > 0$): contradiction with $\int g = 0$. So $g(c) = 0$ for
some $c$: $f(c) = c$.
\end{solution}

\begin{solution}{exo:b1:integration:8}
\begin{enumerate}
  \item By parts with $u' = \sin t$, $v = \sin^{n-1} t$:
        \[
        W_n = \bigl[-\cos t\sin^{n-1}t\bigr]_0^{\pi/2}
        + (n-1)\int_0^{\pi/2} \cos^2 t\,\sin^{n-2} t\,\dd t
        = (n-1)(W_{n-2} - W_n),
        \]
        so $nW_n = (n-1)W_{n-2}$. From $W_0 = \frac\pi2$, $W_1 = 1$:
        \[
        W_{2p} = \frac{(2p-1)(2p-3)\cdots 1}{(2p)(2p-2)\cdots 2}\,
        \frac{\pi}{2} = \frac{(2p)!}{4^p (p!)^2}\,\frac\pi2,
        \qquad
        W_{2p+1} = \frac{(2p)(2p-2)\cdots 2}{(2p+1)(2p-1)\cdots 3}
        = \frac{4^p (p!)^2}{(2p+1)!} .
        \]
  \item On $\intoo{0}{\frac\pi2}$, $0 < \sin t < 1$, so $\sin^{n+1} <
        \sin^n$ and $(W_n)$ is (strictly) decreasing, positive.
        Sandwiching with the recurrence:
        \[
        \frac{n}{n+1} = \frac{W_{n+1}}{W_{n-1}}
        \leq \frac{W_{n+1}}{W_n} \leq 1
        \quad\implies\quad \frac{W_{n+1}}{W_n} \to 1 .
        \]
        Invariant: $a_n = (n+1)W_{n+1}W_n$ satisfies $a_n = a_{n-1}$
        by the recurrence $(n+1)W_{n+1} = nW_{n-1}$, so $a_n = a_0 =
        1 \cdot W_1 W_0 = \frac\pi2$. Then
        \[
        n W_n^2 \sim (n+1) W_{n+1} W_n = \frac\pi2
        \quad\implies\quad
        W_n \sim \sqrt{\frac{\pi}{2n}} .
        \]
\end{enumerate}
\end{solution}

\begin{solution}{exo:b1:integration:9}
\begin{enumerate}
  \item On $\intoo{0}{\pi}$: $x > 0$, $a - bx = b(\frac ab - x) =
        b(\pi - x) > 0$ and $\sin x > 0$, so the integrand is $> 0$
        and $I_n > 0$ (strict positivity). Bound: on
        $\intcc{0}{\pi}$, $x \leq \pi$ and $a - bx \leq a$, so $P
        \leq \frac{\pi^n a^n}{n!}$ and $I_n \leq
        \pi\,\frac{(\pi a)^n}{n!}$, which tends to $0$ (the
        factorial beats the geometric term: it is the general term
        of the convergent exponential series, cf.\
        \cref{ex:b1:seq:e}); in particular $I_n < 1$ for large $n$.
  \item Expand $x^n(a - bx)^n = \sum_{j=0}^{n} \binom nj a^{n-j}
        (-b)^j x^{n+j}$: so $P = \frac{1}{n!}\sum_j c_j x^{n+j}$
        with integer $c_j$. Then $P^{(k)}(0) = 0$ for $k < n$
        (valuation) and, for $n \leq k \leq 2n$, $P^{(k)}(0) =
        \frac{k!}{n!} c_{k-n}$, an integer since $n! \mid k!$.
        Moreover $P(\pi - x) = P(x)$ (substitute: $\pi - x$ swaps
        the factors, using $a - b(\pi - x) = bx$), so $P^{(k)}(\pi)
        = \pm P^{(k)}(0)$: integers as well.
  \item With $Q = P - P'' + P^{(4)} - \dots$ (finite: $P$ has degree
        $2n$): $Q + Q'' = P$, and
        \[
        \bigl(Q'\sin x - Q\cos x\bigr)' = (Q + Q'')\sin x = P\sin x .
        \]
        Hence $I_n = \bigl[Q'\sin x - Q\cos x\bigr]_0^{\pi} = Q(\pi)
        + Q(0)$, a sum of values $P^{(2k)}$ at $0$ and $\pi$: an
        integer by (2).
  \item For $n$ large, $I_n$ is an integer with $0 < I_n < 1$:
        impossible. The assumption $\pi = \frac ab$ fails: $\pi$ is
        irrational.
\end{enumerate}
\end{solution}

\begin{solution}{exo:b1:integration:10}
$C^1$ case: by parts,
\[
\int_a^b f(t)\sin\lambda t\,\dd t
= \Bigl[-f(t)\frac{\cos\lambda t}{\lambda}\Bigr]_a^b
+ \frac{1}{\lambda}\int_a^b f'(t)\cos\lambda t\,\dd t,
\]
bounded in absolute value by $\frac{2\sup\abs f + (b - a)\sup\abs{f'}}
{\lambda} \to 0$.

Continuous case: let $\varepsilon > 0$ and pick a step function
$\varphi$ with $\abs{f - \varphi} \leq \varepsilon$
(\cref{thm:b1:integration:approx} provides $\varphi \leq f \leq
\psi$ with gap $\leq\varepsilon$; take $\varphi$). Then
\[
\Bigl|\int_a^b f\sin\lambda t\Bigr|
\leq \int_a^b \abs{f - \varphi}
+ \Bigl|\int_a^b \varphi \sin\lambda t\Bigr|
\leq (b-a)\varepsilon
+ \sum_i \abs{c_i}\,\Bigl|\int_{x_{i-1}}^{x_i} \sin\lambda t\,\dd
t\Bigr| ,
\]
and each $\bigl|\int \sin \lambda t\bigr| = \bigl|\frac{\cos\lambda
x_{i-1} - \cos\lambda x_i}{\lambda}\bigr| \leq \frac{2}{\lambda}$:
the second term tends to $0$. Hence the limsup is $\leq
(b-a)\varepsilon$ for every $\varepsilon$: the limit is $0$.
\end{solution}

\begin{solution}{exo:b1:integration:11}
Let $m = \min f$ and $M = \max f$, attained by the extreme value
theorem. Since $g \geq 0$: $m\,g \leq fg \leq M\,g$, so by
monotonicity
\[
m \int_a^b g \;\leq\; \int_a^b fg \;\leq\; M \int_a^b g .
\]
If $\int_a^b g = 0$: strict positivity
(\cref{thm:b1:integration:props}~(4)) forces $g \equiv 0$, both
sides vanish, and any $c$ works. Otherwise $t = \frac{\int
fg}{\int g}$ lies in $\intcc{m}{M} = f(\intcc{a}{b})$
(\cref{thm:b1:continuity:evt,thm:b1:continuity:ivt}), so $t =
f(c)$ for some $c$.

Sign matters: on $\intcc{-1}{1}$ with $f(t) = g(t) = t$: $\int
fg = \int_{-1}^1 t^2 = \frac23$, while $f(c)\int_{-1}^1 t\,\dd t
= 0$ for every $c$.
\end{solution}

\begin{solution}{exo:b1:integration:12}
If $f \equiv 0$ the claim is vacuous (every point is a zero). So
assume $f \not\equiv 0$ and suppose it has at most $n$ distinct
zeros in $\intoo{a}{b}$; let $z_1 < \dots < z_m$ ($m \leq n$) be
those zeros where $f$ \emph{changes sign} (possibly none). Set
$P(t) = \prod_{i=1}^{m}(t - z_i)$ (empty product $= 1$), of
degree $m \leq n$. On each subinterval cut by the $z_i$, both
$f$ and $P$ have constant sign, and both flip sign when crossing
some $z_i$: the product $fP$ has one constant sign on all of
$\intoo{a}{b}$. Being continuous, not identically zero, of
constant sign, it has $\bigl|\int_a^b fP\bigr| > 0$ (strict
positivity applied to $\abs{fP}$). But $\int f P$ is a linear
combination of the moments $\int f\,t^k$, $k \leq n$, all zero:
contradiction. Hence $f$ has at least $n + 1$ distinct zeros in
$\intoo{a}{b}$.
\end{solution}

\begin{solution}{pb:b1:integration:1}
\textbf{1.} Let $N = \lceil 2c \rceil$. For $n \geq N$:
$\frac{c^{n+1}/(n+1)!}{c^n/n!} = \frac{c}{n+1} \leq \frac12$, so
$0 < \frac{c^n}{n!} \leq \frac{c^N}{N!}\,2^{-(n - N)} \to 0$:
squeeze.

\textbf{2.} Fix $k$; induction on $m$. For $m = 0$: $\int_0^1
x^k = \frac{1}{k+1} = \frac{k!\,0!}{(k+1)!}$. Step, by parts
($u = (1-x)^m$, $v' = x^k$):
\[
\int_0^1 x^k (1-x)^m \dd x
= \Bigl[\frac{x^{k+1}}{k+1}(1-x)^m\Bigr]_0^1
+ \frac{m}{k+1}\int_0^1 x^{k+1}(1-x)^{m-1}\dd x
= \frac{m}{k+1}\cdot\frac{(k+1)!\,(m-1)!}{(k+m+1)!} ,
\]
which is $\frac{k!\,m!}{(k+m+1)!}$.

\textbf{3.} $k = m = n$: $\int_0^1 (x(1-x))^n =
\frac{(n!)^2}{(2n+1)!} = \frac{1}{(2n+1)\binom{2n}{n}}$. Since
$x(1-x) \leq \frac14$ on $\intcc{0}{1}$, the integral is $\leq
4^{-n}$, whence $\binom{2n}{n} \geq \frac{4^n}{2n+1}$ ---
consistent with $\binom{2n}{n}^{1/n} \to 4$
(\cref{pb:b1:seq:1}).

\textbf{4.} The integrand is continuous, $\geq 0$, and positive
on $\intoo{0}{1}$, hence not identically zero: its integral is
$> 0$ (\cref{thm:b1:integration:props}~(4)). Upper bound:
$(x(1-x))^n \leq 4^{-n}$ and $g \leq \sup\abs g$, then
monotonicity.

\textbf{5.} $A_0 = \eu - 1$; $A_1 = [x\eu^x]_0^1 - \int_0^1
\eu^x = \eu - (\eu - 1) = 1$. By parts: $A_n = [x^n \eu^x]_0^1 -
n\int_0^1 x^{n-1}\eu^x = \eu - n A_{n-1}$. Bounds: the integrand
is positive, so $A_n > 0$; and $\eu^x \leq \eu$ gives $A_n \leq
\eu\int_0^1 x^n = \frac{\eu}{n+1}$.

\textbf{6.} $A_0 = -1 + 1\cdot\eu$. If $A_{n-1} = \alpha_{n-1} +
\beta_{n-1}\eu$ with integer entries, then
\[
A_n = \eu - n\alpha_{n-1} - n\beta_{n-1}\eu
= \underbrace{(-n\,\alpha_{n-1})}_{\alpha_n}
+ \underbrace{(1 - n\,\beta_{n-1})}_{\beta_n}\,\eu ,
\]
both integers.

\textbf{7.} If $\eu = \frac pq$: $q A_n = q\alpha_n + p\beta_n
\in \Z$, and $0 < qA_n \leq \frac{q\eu}{n+1} < 1$ for $n$ large:
an integer strictly between $0$ and $1$ --- impossible. So $\eu
\notin \Q$. In \cref{exo:b1:seq:9} the trapped integer was
$q!\,\frac pq - q!\,a_q$; here it is $qA_n$: same trap,
integral fuel.

\textbf{8.} $n = 0$: $R_0 = \int_0^1 \eu^t = \eu - 1$, so $\eu =
1 + R_0$. By parts ($u = \eu^t$, $v = -\frac{(1-t)^{n+1}}{n+1}$):
\[
R_n = \frac{1}{n!}\Bigl(\Bigl[-\frac{(1-t)^{n+1}}{n+1}
\eu^t\Bigr]_0^1 + \frac{1}{n+1}\int_0^1 (1-t)^{n+1}\eu^t\Bigr)
= \frac{1}{(n+1)!} + R_{n+1} ,
\]
so the formula propagates from $n$ to $n + 1$. Bounds: $1 \leq
\eu^t \leq \eu$ on $\intcc{0}{1}$ and $\int_0^1 (1-t)^n =
\frac{1}{n+1}$ give $\frac{1}{(n+1)!} \leq R_n \leq
\frac{\eu}{(n+1)!}$.

\textbf{9.} $\sum_{k=0}^{9} \frac{1}{k!} = \frac{986410}{362880}
= 2.71828152\dots$, and
\[
\frac{1}{10!} = 2.76\cdot10^{-7} \leq R_9 \leq
\frac{2.75}{10!} = 7.58\cdot10^{-7} ,
\]
so $2.7182818 \leq \eu \leq 2.7182823$: with proof, $\eu =
2.718282$ to six decimals (true value $2.7182818\dots$).

\textbf{10.} $f(1 - x) = \frac{(1-x)^n x^n}{n!} = f(x)$. On
$\intoo{0}{1}$: $0 < x(1-x) \leq \frac14$, so $0 < f \leq
\frac{1}{4^n\,n!}$.

\textbf{11.} $x^n(1-x)^n = \sum_{j=0}^{n} (-1)^j\binom
nj\,x^{n+j}$, so $f = \frac{1}{n!}\sum_j c_j\,x^{n+j}$ with
$c_j \in \Z$. Hence $f^{(k)}(0) = 0$ for $k < n$ or $k > 2n$,
and for $n \leq k \leq 2n$: $f^{(k)}(0) = \frac{k!}{n!}\,
c_{k-n}$, an integer because $n! \mid k!$. The symmetry gives
$f^{(k)}(1) = (-1)^k f^{(k)}(0) \in \Z$.

\textbf{12.} $b^n \pi^{2n-2k} = b^n\bigl(\frac
ab\bigr)^{\,n-k} = a^{\,n-k}\,b^{\,k} \in \Z$, so $G(0) =
\sum_k (-1)^k a^{n-k} b^k f^{(2k)}(0)$ and likewise $G(1)$ are
integers by question 11.

\textbf{13.} In $\pi^2 G + G''$, the term $k$ of $\pi^2 G$
carries $\pi^{2n-2k+2}f^{(2k)}$ and the term $j = k - 1$ of
$G''$ carries $(-1)^{k-1}\pi^{2n-2k+2}f^{(2k)}$: everything
cancels except $k = 0$ in the first sum and $j = n$ in the
second, i.e.
\[
G'' + \pi^2 G = b^n\bigl(\pi^{2n+2} f + (-1)^n f^{(2n+2)}\bigr)
= b^n \pi^{2n+2} f = \pi^2 a^n f
\]
($f$ has degree $2n$, so $f^{(2n+2)} = 0$; and $b^n\pi^{2n} =
a^n$). Then
\[
\bigl(G'\sin\pi x - \pi G\cos\pi x\bigr)'
= (G'' + \pi^2 G)\sin \pi x = \pi^2 a^n f\sin\pi x .
\]

\textbf{14.} Integrating over $\intcc{0}{1}$:
\[
\pi^2 a^n \int_0^1 f\sin(\pi x)\,\dd x
= \bigl[G'\sin\pi x - \pi G\cos\pi x\bigr]_0^1
= \pi\bigl(G(1) + G(0)\bigr) ,
\]
so $\pi a^n \int_0^1 f\sin\pi x = G(0) + G(1) \in \Z$. On
$\intoo{0}{1}$, $f > 0$ and $\sin\pi x > 0$: the left side is
positive, so $G(0) + G(1) \geq 1$; and $\sin \leq 1$ with
question 10 bounds it by $\frac{\pi a^n}{4^n\,n!}$.

\textbf{15.} By question 1 (with $c = \frac a4$),
$\frac{\pi a^n}{4^n n!} \to 0$: for large $n$ it is $< 1$,
contradicting $G(0) + G(1) \geq 1$. So no fraction $\frac ab$
equals $\pi^2$: Legendre's theorem. If $\pi$ were rational,
$\pi^2$ would be too: so $\pi \notin \Q$ --- and the implication
only runs this way ($\sqrt2$ is irrational with rational
square), which is why $\pi^2 \notin \Q$ is strictly stronger
than \cref{exo:b1:integration:9}.

\textbf{16.} Integrality: questions 11--12 (endpoint
derivatives); smallness: questions 10 and 1; bridge: questions
13--14 (the telescoped double integration by parts). Without
$\frac{1}{n!}$, the endpoint data remain integers (even more
easily), but the bound becomes $\frac{\pi a^n}{4^n}$, which
tends to $0$ only when $a < 4$ --- and every candidate has $a =
b\,\pi^2 > 9$. The factorial is exactly what outruns the
geometric growth $a^n$: no factorial, no theorem.

\textbf{17.} For $a \leq 10$, the integer $G(0) + G(1)$ is
positive and at most $\pi\,(10/4)^n/n!$. At $n = 7$:
$2.5^7 = 610.35\dots$, so the bound is $\frac{\pi \times
610.35}{5040} \approx 0.38 < 1$ (at $n = 6$ it is still
$1.07$): the contradiction lands at the seventh stage,
explicitly.

\textbf{18.} Two integrations by parts:
\[
\int_0^1 x(1-x)\sin\pi x\,\dd x
= \frac1\pi\int_0^1 (1 - 2x)\cos\pi x\,\dd x
= \frac{2}{\pi^2}\int_0^1 \sin\pi x\,\dd x
= \frac{2}{\pi^2}\cdot\frac{2}{\pi} = \frac{4}{\pi^3}
\]
(the bracket terms vanish: $x(1-x)$ at $0, 1$, and $\sin\pi x$
at $0, 1$). Symbolically, the $n = 1$ telescope (no assumption
on $\pi$) reads $\pi^3\int_0^1 f_1\sin\pi x = -(f_1''(0) +
f_1''(1))$ with $f_1 = x(1 - x)$, $f_1'' = -2$: right side $4$
--- the two computations agree.

\textbf{19.} $\eu^x$ solves $y' = y$ and $\sin\pi x$ solves
$y'' = -\pi^2 y$: linear equations with constant coefficients,
so integration by parts cycles the kernel back to itself and
keeps all boundary data inside $\Z + \Z\eu$ (resp.\ integer
polynomials in $\pi^2$). That closure property is what the
machine needs. Refined with kernels adapted to several points at
once, the same mechanism yields Hermite's theorem ($\eu$
transcendental, 1873) and Lindemann's ($\pi$ transcendental,
1882) --- beyond this volume.

\textbf{20.} Polynomial division (or multiply back and check):
\[
x^4(1-x)^4 = (x^6 - 4x^5 + 5x^4 - 4x^2 + 4)(1 + x^2) - 4 .
\]
Integrating the displayed identity divided by $1 + x^2$:
\[
\int_0^1 \frac{x^4(1-x)^4}{1+x^2}\dd x
= \Bigl(\frac17 - \frac46 + 1 - \frac43 + 4\Bigr)
- 4\arctan 1
= \frac{22}{7} - \pi ,
\]
using $\frac17 - \frac23 + 1 - \frac43 + 4 = 3 + \frac17$ and
$\arctan 1 = \frac\pi4$.

\textbf{21.} The integrand is continuous, positive on
$\intoo{0}{1}$: the integral is $> 0$, so $\pi < \frac{22}{7}$.
Moreover $\frac12 \leq \frac{1}{1+x^2} \leq 1$ on
$\intcc{0}{1}$ and $\int_0^1 (x(1-x))^4 = \frac{(4!)^2}{9!} =
\frac{1}{630}$ (question 3):
\[
\frac{1}{1260} \leq \frac{22}{7} - \pi \leq \frac{1}{630}
\quad\Longrightarrow\quad
3.14126 < \frac{22}{7} - \frac{1}{630} \leq \pi \leq
\frac{22}{7} - \frac{1}{1260} < 3.14207 .
\]

\textbf{22.} Modulo $x^2 + 1$: $x^2 \equiv -1$, so $x^{4m} =
(x^2)^{2m} \equiv 1$ and $(1 - x)^2 = 1 - 2x + x^2 \equiv -2x$,
hence $(1-x)^{4m} \equiv (-2x)^{2m} = 4^m (x^2)^m \equiv
(-4)^m$: the remainder is the constant $(-4)^m$, and the
quotient $Q_m$ has integer coefficients (division by a monic
integer polynomial). Dividing the identity by $1 + x^2$ and
integrating:
\[
J_m := \int_0^1 \frac{(x(1-x))^{4m}}{1+x^2}\dd x
= s_m + (-4)^m\,\frac{\pi}{4},
\qquad s_m = \int_0^1 Q_m \in \Q .
\]
Solving for $\pi$: with $r_m = (-1)^{m+1}\,4^{\,1-m} s_m \in
\Q$, $\;\abs{\pi - r_m} = 4^{\,1-m} J_m \leq
4^{\,1-m}\cdot4^{-4m} = 4^{\,1-5m}$. For $m = 1$: $s_1 =
\frac{22}{7}$, $r_1 = \frac{22}{7}$, bound $4^{-4} =
\frac{1}{256}$ --- questions 20--21 again.

\textbf{23.} $\bigl|\pi - \frac{22}{7}\bigr| = \frac{22}{7} -
\pi \approx 1.26\cdot10^{-3}$, sixteen times better than the
order-$2$ benchmark $\frac{1}{7^2} \approx 2.0\cdot10^{-2}$
guaranteed by Dirichlet (\cref{pb:b1:derivative:1}, question 4);
and $\bigl|\pi - \frac{355}{113}\bigr| \approx 2.7\cdot10^{-7}$
beats $\frac{1}{113^2} \approx 7.8\cdot10^{-5}$ by a factor
$\approx 300$. No contradiction with anything proved: Liouville
inequalities \emph{lower-bound} approximation errors only for
algebraic numbers, and no such bound for $\pi$ is available at
this level --- $\pi$ is free to be approximated spectacularly
well.

\textbf{24.} Lemma: suppose $x = \frac pq$ and $0 < \abs{a_n +
b_n x} \to 0$. Then $\abs{a_n + b_n x} = \frac{\abs{q a_n + p
b_n}}{q}$, with $q a_n + p b_n$ a nonzero integer (nonzero
because the absolute value is $> 0$): so $\abs{a_n + b_n x}
\geq \frac1q$ for every $n$, contradicting the convergence to
$0$. Instances: question 7 ($x = \eu$, $a_n = \alpha_n$, $b_n =
\beta_n$); \cref{exo:b1:seq:9} ($x = \eu$ again, with $a_q =
-q!\sum_{k \leq q}\frac{1}{k!}$, $b_q = q!$); and
\cref{pb:b1:derivative:1}, question 1, is its geometric form.
In Part III and in \cref{exo:b1:integration:9} the trap runs
\emph{inside} the contradiction: assuming rationality converts
an expression into an integer, which the analysis then squeezes
into $\intoo{0}{1}$ --- the same principle, transposed.

\textbf{25.} (i) Integrality lives in the endpoint calculus of
polynomials (questions 6, 11--12), smallness in the bounds
$\sup$-monotonicity gives (questions 4, 10), the bridge in
integration by parts (questions 8, 13--14) --- all three are
theorems of this chapter. (ii) The integral supplies what the
mean value theorem could not: an \emph{exact} identity between
the analytic object and the arithmetic data (equality, not just
an inequality with an unknown $c$), which is why the machine
reaches $\pi^2$ while \cref{pb:b1:derivative:1} reached only
approximation exponents. (iii) Extracted: $\eu \notin \Q$,
$\pi^2 \notin \Q$ (hence $\pi \notin \Q$), $\eu = 2.718282$
certified, $\frac{22}{7} - \frac{1}{630} \leq \pi \leq
\frac{22}{7} - \frac{1}{1260}$, and $\binom{2n}{n} \geq
\frac{4^n}{2n+1}$. (iv) Frontier: Hermite and Lindemann push
the same machine to transcendence; and Ap\'ery (1979) ran the
integer trap on $\zeta(3)$ --- the machine is still producing
twentieth-century mathematics.
\end{solution}

\section*{Chapter \ref{ch:b1:taylor} --- Taylor Formulas and Asymptotic
Expansions}

\begin{solution}{exo:b1:taylor:1}
\[
\eu^{2x} = 1 + 2x + 2x^2 + \frac{4x^3}{3} + o(x^3);
\qquad
\ln(1 - x) = -x - \frac{x^2}{2} - \frac{x^3}{3} - \frac{x^4}{4} +
o(x^4);
\]
\[
\sqrt{1+x} = 1 + \frac x2 - \frac{x^2}{8} + \frac{x^3}{16} + o(x^3);
\qquad
\frac{1}{1 + x^2} = 1 - x^2 + x^4 - x^6 + o(x^6),
\]
the last by substituting $-x^2$ into the geometric expansion.
\end{solution}

\begin{solution}{exo:b1:taylor:2}
$\eu^x - 1 - x = \frac{x^2}{2} + o(x^2)$: limit $\dfrac12$.

$\cos x = 1 - \frac{x^2}{2} + \frac{x^4}{24} + o(x^4)$ and
$\sqrt{1 - x^2} = 1 - \frac{x^2}{2} - \frac{x^4}{8} + o(x^4)$:
difference $\frac{x^4}{24} + \frac{x^4}{8} = \frac{x^4}{6} + o(x^4)$:
limit $\dfrac16$.

$\ln(1+x) - \sin x = \bigl(x - \frac{x^2}{2}\bigr) - x + o(x^2) =
-\frac{x^2}{2} + o(x^2)$: limit $-\dfrac12$.
\end{solution}

\begin{solution}{exo:b1:taylor:3}
$\frac{1}{1+t^2} = 1 - t^2 + t^4 + o(t^5)$; integrating from $0$ to
$x$ (\cref{met:b1:taylor:compute}~(5)):
\[
\arctan x = x - \frac{x^3}{3} + \frac{x^5}{5} + o(x^5)\ \ (\text{even
}o(x^6)\text{, by oddness}).
\]
$(1 - t^2)^{-1/2} = 1 + \frac{t^2}{2} + \frac38 t^4 + o(t^5)$ (binomial
expansion with $\alpha = -\frac12$, $x = -t^2$:
$\binom{-1/2}{2} = \frac{(-\frac12)(-\frac32)}{2} = \frac38$);
integrating:
\[
\arcsin x = x + \frac{x^3}{6} + \frac{3x^5}{40} + o(x^5) .
\]
\end{solution}

\begin{solution}{exo:b1:taylor:4}
Taylor--Lagrange (\cref{thm:b1:taylor:lagrange}) for $\exp$ at $a =
0$, $x = 1$: the $(n+1)$-st derivative is $\eu^t \leq \eu < 3$ on
$\intcc{0}{1}$, so
\[
\Bigl|\eu - \sum_{k=0}^{n} \frac{1}{k!}\Bigr| \leq
\frac{3}{(n+1)!} .
\]
For $6$ exact decimals, want $\frac{3}{(n+1)!} < 5\times 10^{-7}$,
i.e.\ $(n+1)! > 6\times 10^{6}$: since $10! = 3\,628\,800$ and $11! =
39\,916\,800$, $n + 1 = 11$, i.e.\ $n = 10$ suffices.
\end{solution}

\begin{solution}{exo:b1:taylor:5}
$n\ln\bigl(1 + \frac1n\bigr) = n\Bigl(\frac1n - \frac{1}{2n^2} +
\frac{1}{3n^3} + o\bigl(\frac{1}{n^3}\bigr)\Bigr) = 1 - \frac{1}{2n}
+ \frac{1}{3n^2} + o\bigl(\frac{1}{n^2}\bigr)$. Exponentiating, with
$u = -\frac{1}{2n} + \frac{1}{3n^2}$ and $\eu^u = 1 + u + \frac{u^2}2
+ o(u^2)$:
\[
\Bigl(1 + \frac1n\Bigr)^n = \eu\cdot \eu^{u}
= \eu\Bigl(1 - \frac{1}{2n} + \frac{1}{3n^2} + \frac{1}{8n^2}
+ o\Bigl(\frac{1}{n^2}\Bigr)\Bigr)
= \eu\Bigl(1 - \frac{1}{2n} + \frac{11}{24n^2} +
o\Bigl(\frac{1}{n^2}\Bigr)\Bigr).
\]
Limit $\eu$; the error is $\sim \dfrac{\eu}{2n}$: slow (one digit per
tenfold increase of $n$).
\end{solution}

\begin{solution}{exo:b1:taylor:6}
$f(x) = x^2 - x^4 = x^2(1 + o(1))$: first term $x^2$, $p = 2$ even,
coefficient $> 0$: local minimum at $0$ (not global: $f(2) = -12$).

$g(x) = x^3 + x^5$: first term $x^3$, odd $p$: no extremum; $g$
crosses its (horizontal) tangent: inflection at $0$.

$h = \exp$ at $a = 1$: $h(x) = \eu + \eu(x-1) + \frac{\eu}{2}(x-1)^2
+ o((x-1)^2)$; the difference with the tangent is
$\frac{\eu}{2}(x-1)^2 + o(\cdot) > 0$ locally. Globally: $\eu^x -
\eu x \geq 0$ for all $x$ by convexity
(\cref{thm:b1:derivative:convexchar}~(3)): the graph lies above every
tangent, with equality only at the contact point.
\end{solution}

\begin{solution}{exo:b1:taylor:7}
For $x \to +\infty$:
\[
f(x) = x\Bigl(1 + \frac1x\Bigr)^{1/3}
= x\Bigl(1 + \frac{1}{3x} - \frac{1}{9x^2} +
o\Bigl(\frac{1}{x^2}\Bigr)\Bigr)
= x + \frac13 - \frac{1}{9x} + o\Bigl(\frac1x\Bigr):
\]
asymptote $y = x + \frac13$, curve below it near $+\infty$. As $x
\to -\infty$, the same computation is valid (the cube root is
defined for all reals, and $\frac1x \to 0$): same asymptote $y = x +
\frac13$, but now $-\frac{1}{9x} > 0$: curve \emph{above} the line.
\end{solution}

\begin{solution}{exo:b1:taylor:8}
$u_n = \sqrt n\bigl(\sqrt{1 + \tfrac1n} - 1\bigr) = \sqrt
n\bigl(\frac{1}{2n} + o(\frac1n)\bigr) \sim \dfrac{1}{2\sqrt n}$.

$v_n = \ln\bigl(1 + \frac1n\bigr) \sim \dfrac 1n$.

$w_n$: with $h = \frac1n \to 0$, $\sin h - \tan h = \bigl(h -
\frac{h^3}{6}\bigr) - \bigl(h + \frac{h^3}{3}\bigr) + o(h^3) =
-\frac{h^3}{2} + o(h^3)$, so $w_n \sim -\dfrac{1}{2n^3}$.
\end{solution}

\begin{solution}{exo:b1:taylor:9}
Taylor--Lagrange at order $1$ around $x$, on both sides:
\[
f(x + h) = f(x) + h f'(x) + R_+,\quad
f(x - h) = f(x) - h f'(x) + R_-,
\qquad \abs{R_\pm} \leq \frac{h^2}{2} \sup \abs{f''} .
\]
Subtracting: $f(x+h) - f(x-h) = 2h f'(x) + (R_+ - R_-)$, so
\[
\abs{f'(x)} \leq \frac{\abs{f(x+h) - f(x-h)}}{2h} +
\frac{h}{2}\sup_{\intcc{x-h}{x+h}}\abs{f''} .
\]
With global bounds: $\abs{f'(x)} \leq \frac{M_0}{h} + \frac{M_2
h}{2}$ for every $h > 0$. The right side is minimized at $h =
\sqrt{2M_0/M_2}$ (derivative zero), with value $\sqrt{2M_0M_2}$
--- hence $\abs{f'} \leq \sqrt{2M_0M_2}$. (If $M_2 = 0$, let $h \to
\infty$: $f' = 0$, consistent.)
\end{solution}

\begin{solution}{exo:b1:taylor:10}
On $\intoo{0}{\pi}$: $0 < \sin u < u$, so $(u_n)$ is strictly
decreasing, positive, hence convergent; the limit is a fixed point
of $\sin$ in $\intcc{0}{\pi}$, and $\sin \ell = \ell$ forces $\ell =
0$ (since $\sin x < x$ for $x > 0$).
\begin{enumerate}
  \item Using $\sin u = u - \frac{u^3}{6} + o(u^3)$ as $u \to 0$:
        \[
        v_{n+1} - v_n = \frac{1}{\sin^2 u_n} - \frac{1}{u_n^2}
        = \frac{1}{u_n^2}\Bigl(\Bigl(1 - \frac{u_n^2}{6} +
        o(u_n^2)\Bigr)^{\!-2} - 1\Bigr)
        = \frac{1}{u_n^2}\Bigl(\frac{u_n^2}{3} + o(u_n^2)\Bigr)
        \longrightarrow \frac13 .
        \]
  \item By \cref{exo:b1:seq:10}~(3) (Cesàro for differences),
        $\frac{v_n}{n} \to \frac13$, i.e.\ $v_n \sim \frac n3$, i.e.\
        $u_n^2 \sim \frac 3n$: since $u_n > 0$,
        \[
        u_n \sim \sqrt{\frac{3}{n}} .
        \]
\end{enumerate}
\end{solution}

\section*{Chapter \ref{ch:b1:series} --- Numerical Series}

\begin{solution}{exo:b1:series:1}
$\dfrac{1}{n(n+1)} = \dfrac1n - \dfrac{1}{n+1}$: telescoping, $S_N =
1 - \frac{1}{N+1} \to 1$. Convergent, sum $1$.

$\ln\bigl(1 - \frac{1}{n^2}\bigr) = \ln\frac{(n-1)(n+1)}{n^2} =
\ln\frac{n-1}{n} - \ln\frac{n}{n+1}$: telescoping again, $S_N =
\ln\frac12 - \ln\frac{N}{N+1} \to -\ln 2$. Convergent, sum $-\ln 2$.

$\dfrac{3^n + 4^n}{5^n} = \bigl(\frac35\bigr)^n +
\bigl(\frac45\bigr)^n$: two convergent \omterm{ex:b1:series:geometric}{geometric series}, sum
$\frac{1}{1 - 3/5} + \frac{1}{1 - 4/5} = \frac52 + 5 = \frac{15}{2}$.
\end{solution}

\begin{solution}{exo:b1:series:2}
\omterm{thm:b1:series:ratio}{Ratio test} throughout.

$\frac{u_{n+1}}{u_n} = \frac{(n+1)^2}{2n^2} \to \frac12 < 1$:
convergent.

$\frac{u_{n+1}}{u_n} = \frac{(n+1)!\,n^n}{n!\,(n+1)^{n+1}} =
\bigl(\frac{n}{n+1}\bigr)^n = \bigl(1 + \frac1n\bigr)^{-n} \to
\frac1\eu < 1$: convergent.

With the factor $2^n$: ratio $\to \frac2\eu < 1$: convergent.

With $3^n$: ratio $\to \frac3\eu > 1$: divergent (terms tend to
$+\infty$).
\end{solution}

\begin{solution}{exo:b1:series:3}
All nonnegative terms; use equivalents
(\cref{thm:b1:series:positive}).

$\sin\frac{1}{n^2} \sim \frac{1}{n^2}$: convergent (Riemann $\alpha
= 2$).

$1 - \cos\frac1n \sim \frac{1}{2n^2}$: convergent.

$\frac{1}{\sqrt{n(n+1)}} \sim \frac1n$: divergent.

$\frac{\ln n}{n^2} = \frac{1}{n^{3/2}}\cdot\frac{\ln n}{n^{1/2}}$
and $\frac{\ln n}{\sqrt n} \to 0$
(\cref{prop:b1:functions:powerrules}): so $\frac{\ln n}{n^2} \leq
\frac{1}{n^{3/2}}$ for large $n$: convergent.
\end{solution}

\begin{solution}{exo:b1:series:4}
For $n \geq 2$: $\frac{1}{n^2} \leq \frac{1}{n(n-1)} =
\frac{1}{n-1} - \frac1n$. Hence
\[
\sum_{n=1}^{N} \frac{1}{n^2} \leq 1 + \sum_{n=2}^{N}
\Bigl(\frac{1}{n-1} - \frac1n\Bigr) = 1 + 1 - \frac1N < 2 :
\]
partial sums increasing and bounded by $2$: convergence
(\cref{thm:b1:series:positive}), sum $\leq 2$. (The exact value
$\frac{\pi^2}{6}$ is a second-year celebration.)
\end{solution}

\begin{solution}{exo:b1:series:5}
$\sum \frac{(-1)^n}{\sqrt n}$: alternating with $\frac{1}{\sqrt n}
\downarrow 0$: convergent (\cref{thm:b1:series:alternating}); not
absolutely ($\alpha = \frac12 \leq 1$).

$\sum \frac{(-1)^n}{n + (-1)^n}$: the sequence $\frac{1}{n +
(-1)^n}$ is \emph{not} decreasing ($\frac{1}{n+1}$ then
$\frac{1}{n}$ alternate badly), so the test does not apply directly.
Expand:
\[
\frac{(-1)^n}{n + (-1)^n}
= \frac{(-1)^n}{n}\cdot\frac{1}{1 + \frac{(-1)^n}{n}}
= \frac{(-1)^n}{n} - \frac{1}{n^2} +
O\Bigl(\frac{1}{n^3}\Bigr):
\]
the first \omterm{def:b1:series:def}{series} converges (alternating), $\sum \frac{1}{n^2}$
converges, the $O(n^{-3})$ converges absolutely: the sum of three
convergent \omterm{def:b1:series:def}{series} converges.

$\sin\bigl(\pi\sqrt{n^2+1}\bigr)$: write $\sqrt{n^2 + 1} = n +
\frac{1}{2n} + \varepsilon_n$ with $\varepsilon_n = O(n^{-3})$;
then, by $\pi$-periodicity of $\sin$ up to sign,
\[
\sin\bigl(\pi\sqrt{n^2+1}\bigr)
= (-1)^n \sin\Bigl(\frac{\pi}{2n} + \pi\varepsilon_n\Bigr) .
\]
Set $\theta_n = \frac{\pi}{2n} + \pi\varepsilon_n$ and $a_n =
\sin\theta_n$. For large $n$, $\theta_n \in \intoo{0}{\frac\pi2}$
and
\[
\theta_n - \theta_{n+1} = \frac{\pi}{2n(n+1)} +
\pi(\varepsilon_n - \varepsilon_{n+1})
= \frac{\pi}{2n^2} + O\Bigl(\frac{1}{n^3}\Bigr) > 0
\]
eventually, so $(\theta_n)$ decreases to $0$; since $\sin$ is
increasing on $\intcc{0}{\frac\pi2}$, $(a_n)$ decreases to $0$ as
well. The alternating test applies: convergent --- not absolutely,
since $a_n \sim \frac{\pi}{2n}$.
\end{solution}

\begin{solution}{exo:b1:series:6}
$f(t) = \frac{1}{t(\ln t)^\alpha}$ is positive, \omterm{def:b1:continuity:continuous}{continuous},
decreasing on $\intco{2}{+\infty}$. Substituting $u = \ln t$:
\[
\int_2^x \frac{\dd t}{t(\ln t)^\alpha}
= \int_{\ln 2}^{\ln x} \frac{\dd u}{u^\alpha},
\]
bounded as $x \to \infty$ iff $\alpha > 1$
(\cref{thm:b1:series:riemann}'s computation). By \omterm{thm:b1:integration:def}{integral}
comparison: convergence iff $\alpha > 1$. (These \emph{Bertrand-type}
\omterm{def:b1:series:def}{series} show how fine the \omterm{def:b1:topology:closure}{boundary} of convergence is: $n\ln n$
diverges, $n(\ln n)^{1.01}$ converges.)
\end{solution}

\begin{solution}{exo:b1:series:7}
By the tangent-line bounds of \cref{exo:b1:derivative:3} rewritten
via expansions: for $x = \frac1n \in \intoc{0}{1}$, Taylor--Lagrange
for $\ln(1+x)$ at order $1$ gives $\ln(1 + x) = x - \frac{x^2}{2(1 +
c)^2}$ for some $c \in \intoo{0}{x}$, so
\[
0 \leq u_n = \frac1n - \ln\Bigl(1 + \frac1n\Bigr) \leq
\frac{1}{2n^2} .
\]
Comparison with the \omterm{thm:b1:series:riemann}{Riemann series}: $\sum u_n$ converges. Its partial
sum telescopes the logarithms:
\[
\sum_{n=1}^{N} u_n = H_N - \ln(N+1)
\]
(since $\sum_{n\leq N} \ln\frac{n+1}{n} = \ln(N+1)$). So $H_N -
\ln(N+1)$ converges; adding $\ln\frac{N+1}{N} \to 0$, the sequence
$H_N - \ln N$ converges. Its limit is $\gamma \approx 0.5772$.
\end{solution}

\begin{solution}{exo:b1:series:8}
$\dfrac{1}{n(n+2)} = \dfrac{1/2}{n} - \dfrac{1/2}{n+2}$: the partial
sum telescopes with a lag of $2$,
\[
S_N = \frac12\Bigl(1 + \frac12 - \frac{1}{N+1} -
\frac{1}{N+2}\Bigr) \longrightarrow \frac34 .
\]

$\sum \frac{n}{2^n}$: for $\abs x < 1$, differentiating the finite
geometric sum and passing to the limit (all \omterm{def:b1:series:def}{series} here converge
absolutely, \omterm{thm:b1:series:ratio}{ratio test}): from $\sum_{n\geq0} x^n = \frac{1}{1-x}$,
one gets by direct computation with partial sums
\[
\sum_{n=1}^{N} n x^{n-1}
= \frac{1 - (N+1)x^N + N x^{N+1}}{(1 - x)^2}
\xrightarrow[N\to\infty]{} \frac{1}{(1-x)^2}
\quad (\abs x < 1),
\]
(the \omterm{def:b1:topology:closure}{boundary} terms $N x^N \to 0$). At $x = \frac12$:
$\sum_{n\geq1} n\bigl(\frac12\bigr)^{n-1} = 4$, so $\sum_{n\geq0}
\frac{n}{2^n} = \frac12 \times 4 = 2$.
\end{solution}

\begin{solution}{exo:b1:series:9}
Group the terms of $\sum u_n$ in packets between powers of $2$.
Upper packets: for $2^k \leq n < 2^{k+1}$ there are $2^k$ terms,
each $\leq u_{2^k}$:
\[
\sum_{n=1}^{2^{K+1}-1} u_n
= \sum_{k=0}^{K} \sum_{n=2^k}^{2^{k+1}-1} u_n
\leq \sum_{k=0}^{K} 2^k u_{2^k} .
\]
Lower packets: each term of the same packet is $\geq u_{2^{k+1}}$,
so $\sum_{n=2^k}^{2^{k+1}-1} u_n \geq 2^k u_{2^{k+1}} = \frac12
\cdot 2^{k+1} u_{2^{k+1}}$, whence
\[
\sum_{n=1}^{2^{K+1}-1} u_n \geq \frac12 \sum_{k=1}^{K+1} 2^{k}
u_{2^{k}} .
\]
Both partial-sum comparisons go both ways (nonnegative terms,
\cref{thm:b1:series:positive}): the two \omterm{def:b1:series:def}{series} have the same nature.

Riemann: $u_n = n^{-\alpha}$ gives $2^k u_{2^k} =
2^{k(1-\alpha)}$, a \omterm{ex:b1:series:geometric}{geometric series}, convergent iff $2^{1 - \alpha}
< 1$ iff $\alpha > 1$. Bertrand
(\cref{exo:b1:series:6}): $u_n = \frac{1}{n(\ln n)^\alpha}$ gives
$2^k u_{2^k} = \frac{1}{(k\ln 2)^\alpha}$, a \omterm{thm:b1:series:riemann}{Riemann series} in $k$:
convergent iff $\alpha > 1$.
\end{solution}

\begin{solution}{exo:b1:series:10}
Finite geometric identity with ratio $-t^2$:
\[
\frac{1}{1 + t^2} = \sum_{k=0}^{n-1} (-1)^k t^{2k}
+ \frac{(-1)^n t^{2n}}{1 + t^2} .
\]
Integrate over $\intcc{0}{1}$ (the left side integrates to $\arctan
1 = \frac\pi4$, \cref{prop:b1:functions:arcderiv}):
\[
\frac{\pi}{4} = \sum_{k=0}^{n-1} \frac{(-1)^k}{2k+1}
+ (-1)^n \int_0^1 \frac{t^{2n}}{1+t^2}\,\dd t,
\qquad
0 \leq \int_0^1 \frac{t^{2n}}{1+t^2}\,\dd t \leq \int_0^1 t^{2n}\dd
t = \frac{1}{2n+1} .
\]
Letting $n \to \infty$ proves the formula, and the displayed bound
on the \omterm{thm:b1:integration:def}{integral} is exactly the remainder bound: after summing up to
$N$ (i.e.\ $n = N + 1$ terms), $\abs{R_N} \leq \frac{1}{2N + 3}$.
\end{solution}

\begin{solution}{exo:b1:series:11}
$n^{1/n} = \eu^{\frac{\ln n}{n}} \to \eu^0 = 1$
(\cref{prop:b1:functions:powerrules}). Hence
\[
\frac{1}{n^{1 + 1/n}} = \frac{1}{n}\,\eu^{-\frac{\ln n}{n}}
\sim \frac{1}{n} ,
\]
and the equivalents test (\cref{thm:b1:series:positive})
compares with the divergent harmonic \omterm{def:b1:series:def}{series}: \emph{divergent},
although every exponent $1 + \frac1n$ exceeds $1$. The Riemann
criterion concerns a \emph{fixed} exponent $\alpha$; an exponent
sliding down to $1$ can lose all its margin, as here.
\end{solution}

\begin{solution}{exo:b1:series:12}
Let $\varepsilon > 0$. By the Cauchy criterion for the
convergent \omterm{def:b1:series:def}{series} (\cref{thm:b1:seq:complete} applied to the
partial sums), there is $N$ with $\sum_{k=n+1}^{2n} u_k \leq
\varepsilon$ for $n \geq N$. By monotonicity each of these $n$
terms is $\geq u_{2n}$:
\[
n\,u_{2n} \leq \sum_{k=n+1}^{2n} u_k \leq \varepsilon
\quad\Longrightarrow\quad 2n\,u_{2n} \leq 2\varepsilon ,
\]
and for odd indices $(2n+1)\,u_{2n+1} \leq (2n+1)\,u_{2n} \leq
2\bigl(2n\,u_{2n}\bigr) \leq 4\varepsilon$ for $n \geq N$: in
both parities, $n u_n \to 0$.

Converse false: $u_n = \frac{1}{n\ln n}$ has $n u_n =
\frac{1}{\ln n} \to 0$, yet the \omterm{def:b1:series:def}{series} diverges
(\cref{exo:b1:series:6}, $\alpha = 1$). Monotonicity necessary:
let $u_n = \frac1n$ when $n$ is a perfect square and $u_n =
2^{-n}$ otherwise: the \omterm{def:b1:series:def}{series} converges (the square terms sum
like $\sum \frac{1}{k^2}$, the rest geometrically), but $n u_n =
1$ along the squares.
\end{solution}

\begin{solution}{pb:b1:series:1}
\textbf{1.} $a_{n+1} - a_n = \frac{1}{n+1} - \ln\frac{n+1}{n}
\leq 0$ because $\ln(1 + \frac1n) \geq \frac{1/n}{1 + 1/n} =
\frac{1}{n+1}$; and $b_{n+1} - b_n = \frac{1}{n+1} -
\ln\frac{n+2}{n+1} \geq 0$ because $\ln(1 + \frac{1}{n+1}) \leq
\frac{1}{n+1}$. Their gap $a_n - b_n = \ln(1 + \frac1n) \to 0$:
adjacent (\cref{thm:b1:seq:adjacent}), with common limit $\lim
a_n = \gamma$ (\cref{exo:b1:series:7}), and $b_n \leq \gamma
\leq a_n$.

\textbf{2.} $b_{10} = 2.928968 - \ln 11 = 0.5311$ and $a_{10} =
2.928968 - \ln 10 = 0.6264$: so $\gamma \in
\intcc{0.5311}{0.6264}$. The gap is $\ln 1.1 \approx 0.095$ and
shrinks like $\frac1n$: four decimals ($\text{gap} \leq
10^{-4}$) would need $n \approx 10^4$ --- the brackets are
correct but slow.

\textbf{3.} Telescoping $a_n - a_{m+1} = \sum_{k=n}^{m} w_k$ and
letting $m \to \infty$: $a_n - \gamma = \sum_{k\geq n} w_k$
(limit of partial sums). Moreover
\[
w_k = \int_k^{k+1} \frac{\dd t}{t} - \frac{1}{k+1}
= \int_k^{k+1} \Bigl(\frac1t - \frac{1}{k+1}\Bigr)\dd t
= \int_k^{k+1} \frac{k + 1 - t}{t\,(k+1)}\,\dd t .
\]
On $\intcc{k}{k+1}$: $\frac{1}{(k+1)^2} \leq \frac{1}{t(k+1)}
\leq \frac{1}{k(k+1)}$, and $\int_k^{k+1}(k + 1 - t)\dd t =
\frac12$: hence $\frac{1}{2(k+1)^2} \leq w_k \leq
\frac{1}{2k(k+1)}$.

\textbf{4.} Upper: $\sum_{k \geq n} \frac{1}{2k(k+1)} =
\frac12\sum_{k\geq n}\bigl(\frac1k - \frac{1}{k+1}\bigr) =
\frac{1}{2n}$ (telescoping). Lower: $\frac{1}{2(k+1)^2} \geq
\frac{1}{2(k+1)(k+2)}$, whose sum telescopes to
$\frac{1}{2(n+1)}$. With question 3:
\[
\frac{1}{2(n+1)} \leq H_n - \ln n - \gamma \leq \frac{1}{2n} .
\]

\textbf{5.} Subtract $\frac{1}{2n}$: $\gamma_n - \gamma \in
\intcc{\frac{1}{2(n+1)} - \frac{1}{2n}}{0} =
\intcc{-\frac{1}{2n(n+1)}}{0}$: the corrected estimate is exact
to $O\bigl(\frac{1}{n^2}\bigr)$, and always from below.

\textbf{6.} $\gamma_{100} = 5.1873775 - \ln 100 - 0.005 =
0.5772073$, with $0 \leq \gamma - \gamma_{100} \leq
\frac{1}{20200} < 5\cdot10^{-5}$: hence $0.577207 \leq \gamma
\leq 0.577257$, i.e.\ $\gamma = 0.5772 \pm 5\cdot10^{-5}$ (true
value $0.5772156\dots$) --- four certified decimals from a
hundred terms, against ten thousand for question 2.

\textbf{7.} First dividend:
\[
H_{2m} - H_m = \Bigl(\ln 2m + \gamma + \frac{1}{4m}\Bigr)
- \Bigl(\ln m + \gamma + \frac{1}{2m}\Bigr) +
O\Bigl(\frac{1}{m^2}\Bigr)
= \ln 2 - \frac{1}{4m} + O\Bigl(\frac{1}{m^2}\Bigr).
\]
Second: the even reciprocals up to $2m$ sum to $\frac12 H_m$, so
$\sum_{j=1}^{m}\frac{1}{2j-1} = H_{2m} - \frac12 H_m = \frac12
\ln m + \ln 2 + \frac\gamma2 + o(1)$, and $\sum_{j=1}^m
\frac{1}{2j} = \frac12\ln m + \frac\gamma2 + o(1)$.

\textbf{8.} Splitting off the even terms twice:
$\sum_{k=1}^{2m}\frac{(-1)^{k-1}}{k} = H_{2m} - 2\cdot\frac12
H_m = H_{2m} - H_m$. By question 7 this equals $\ln 2 -
\frac{1}{4m} + O(m^{-2})$: the sum is $\ln 2$
(\cref{ex:b1:series:ln2} again) \emph{with} its speed.

\textbf{9.} For $N = 2m$: $S - S_N = \frac{1}{4m} + O(m^{-2}) =
\frac{1}{2N} + O(N^{-2})$. For $N = 2m + 1$: $S_{N} = S_{2m} +
\frac{1}{2m+1}$, so
\[
S - S_N = \Bigl(\frac{1}{4m} - \frac{1}{2m+1}\Bigr) +
O\Bigl(\frac{1}{m^2}\Bigr)
= -\frac{1}{4m} + O\Bigl(\frac{1}{m^2}\Bigr)
= -\frac{1}{2N} + O\Bigl(\frac{1}{N^2}\Bigr).
\]
In both cases $S - S_N \sim \frac{(-1)^N}{2N}$: half the
worst-case bound $a_{N+1}$, with a known, alternating sign.

\textbf{10.} Averaging kills the oscillating leading term:
\[
S - \tilde S_N = \frac{(S - S_N) + (S - S_{N+1})}{2}
= \frac{(-1)^N}{2}\Bigl(\frac{1}{2N} - \frac{1}{2(N+1)}\Bigr)
+ O\Bigl(\frac{1}{N^2}\Bigr) = O\Bigl(\frac{1}{N^2}\Bigr).
\]
Numerically: $S_{10} = 0.645635$, $S_{11} = 0.736544$, $\tilde
S_{10} = 0.691089$, and $\ln 2 = 0.693147$: the error drops from
$4.8\cdot10^{-2}$ to $2.1\cdot10^{-3}$ --- one addition, twenty
times better.

\textbf{11.} The alternating error changes sign at every step,
so consecutive partial sums straddle the limit and their mean
cancels the first-order term; the bracket error $H_n - \ln n -
\gamma \approx \frac{1}{2n}$ has constant sign, so no averaging
along $n$ can cancel it. Averaging the two \emph{brackets} does
help: $\frac{a_n + b_n}{2} = H_n - \ln\sqrt{n(n+1)}$, and since
$\ln\sqrt{n(n+1)} = \ln\bigl(n + \frac12\bigr) + O(n^{-2})$,
\[
H_n - \ln\Bigl(n + \frac12\Bigr)
= \Bigl(H_n - \ln n - \frac{1}{2n}\Bigr) + \frac{1}{8n^2} +
O\Bigl(\frac{1}{n^3}\Bigr) = \gamma + O\Bigl(\frac{1}{n^2}\Bigr)
\]
(question 5 and $\ln(1 + \frac{1}{2n}) = \frac{1}{2n} -
\frac{1}{8n^2} + O(n^{-3})$). Check at $n = 10$: $H_{10} -
\ln 10.5 = 0.57759$, already within $4\cdot10^{-4}$ of
$\gamma$.

\textbf{12.} $\sum_{j\leq m} \frac{1}{2j-1} \geq \sum_{j \leq m}
\frac{1}{2j} = \frac12 H_m \to +\infty$: both the positive and
the negative part of the alternating harmonic \omterm{def:b1:series:def}{series} diverge.

\textbf{13.} Write $u_n^\pm = \frac{\abs{u_n} \pm u_n}{2} \geq
0$, so $u_n = u_n^+ - u_n^-$ and $\abs{u_n} = u_n^+ + u_n^-$. If
$\sum u_n^+$ converged, then $\sum u_n^- = \sum (u_n^+ - u_n)$
would converge (difference of convergent \omterm{def:b1:series:def}{series}), hence $\sum
\abs{u_n}$ too: contradiction with conditional convergence. By
symmetry both $\sum u_n^\pm$ diverge (to $+\infty$): an infinite
reservoir of positive and of negative mass.

\textbf{14.} Let $S = \sum u_n$, $\varepsilon > 0$, and $N$ with
$\sum_{n > N} \abs{u_n} \leq \varepsilon$ (Cauchy criterion for
$\sum\abs{u_n}$). Let $M_0$ be large enough that $\sigma(\{0,
\dots, M_0\}) \supseteq \{0, \dots, N\}$. For $M \geq M_0$, the
difference $\sum_{m \leq M} u_{\sigma(m)} - \sum_{n \leq N} u_n$
is a finite sum of \emph{distinct} terms $u_n$ with $n > N$,
hence of absolute value $\leq \varepsilon$; and $\abs{S -
\sum_{n\leq N} u_n} \leq \varepsilon$ as well. So the rearranged
partial sums are within $2\varepsilon$ of $S$ eventually:
$\sum u_{\sigma(n)} = S$. \omterm{thm:b1:series:absolute}{Absolute convergence} is
rearrangement-proof.

\textbf{15.} Each phase of the greedy procedure ends after
finitely many terms, because the remaining positive
(respectively negative) terms alone have divergent partial sums
(question 12): the running sum must eventually cross $t$. The
procedure therefore alternates infinitely many finite phases,
consuming the positive terms in order and the negative terms in
order: every term is used exactly once --- a rearrangement.
After the first crossing, between two consecutive crossings the
partial sums move monotonically toward $t$, and at a crossing
they overshoot by at most the term just added; since the terms
used at the $j$-th crossing have index at least $j$ in their
class, these overshoots tend to $0$. Hence the partial sums
converge to $t$: every real number is the sum of some
rearrangement.

\textbf{16.} The positive slots receive $\frac{1}{2j-1}$ for
$j = 1, 2, \dots$ in order, the negative slots $\frac{1}{2j}$ in
order: every term of the alternating harmonic \omterm{def:b1:series:def}{series} appears
exactly once. For $(p, q) = (1, 2)$, the blocks are
$\bigl(1, -\frac12, -\frac14\bigr)$, $\bigl(\frac13, -\frac16,
-\frac18\bigr)$, $\bigl(\frac15, -\frac1{10},
-\frac1{12}\bigr)$, \dots\ --- the displayed \omterm{def:b1:series:def}{series}.

\textbf{17.} Since $\frac{1}{4k-2} = \frac12\cdot\frac{1}{2k-1}$:
\[
\frac{1}{2k-1} - \frac{1}{4k-2} - \frac{1}{4k}
= \frac12\,\frac{1}{2k-1} - \frac12\,\frac{1}{2k}
= \frac12\Bigl(\frac{1}{2k-1} - \frac{1}{2k}\Bigr).
\]
Summing over $k = 1, \dots, K$: $T_{3K} = \frac12
\sum_{k=1}^{K}\bigl(\frac{1}{2k-1} - \frac{1}{2k}\bigr) =
\frac12 S_{2K}$: at every third partial sum, the rearranged
\omterm{def:b1:series:def}{series} is \emph{exactly} half the original one.

\textbf{18.} After $K$ complete blocks, the rearranged partial
sum is $\sum_{j=1}^{pK}\frac{1}{2j-1} -
\sum_{j=1}^{qK}\frac{1}{2j}$, and question 7 evaluates it:
\[
\Bigl(\frac{\ln(pK)}{2} + \ln 2 + \frac\gamma2\Bigr)
- \Bigl(\frac{\ln(qK)}{2} + \frac\gamma2\Bigr) + o(1)
= \ln 2 + \frac12\ln\frac pq + o(1) :
\]
the $\gamma$s cancel, the $\ln K$s cancel, the ratio
$\frac pq$ survives.

\textbf{19.} A partial sum inside block $K + 1$ differs from the
$K$-block sum by at most $p + q$ terms, each of absolute value
$\leq \frac{1}{2qK}$-ish, hence by $O\bigl(\frac1K\bigr) \to 0$:
the full sequence of partial sums has the same limit $\ln 2 +
\frac12\ln\frac pq$. For $(1, 2)$: $\ln 2 + \frac12\ln\frac12 =
\frac{\ln 2}{2} = 0.34657\dots$, and indeed $T_9 = 0.30833$
creeps toward it: by question 17, $T_{3K} = \frac12 S_{2K}$
converges with exactly half the alternating-harmonic error.
Same terms, half the sum.

\textbf{20.} $(1,1)$: $\ln 2 + \frac12\ln 1 = \ln 2$ --- the
original order, consistency. $(2,1)$: $\frac32\ln 2 \approx
1.0397$. The $(p,q)$ menu reaches exactly the countable \omterm{def:b1:topology:dense}{dense}
family $\ln 2 + \frac12\ln r$, $r \in \Q_{>0}$; Riemann's greedy
recipe (question 15) reaches \emph{every} real. Structure buys
formulas; greed buys totality.

\textbf{21.} The partial sums telescope: $\sum_{k=1}^{N}
\bigl(\frac1k - \ln\frac{k+1}{k}\bigr) = H_N - \ln(N+1) = b_N
\to \gamma$: the \omterm{def:b1:series:def}{series} of \cref{exo:b1:series:7} sums exactly
to \omterm{pb:b1:series:1}{Euler's constant}.

\textbf{22.} Greedy for $t = 1$: the first positive term brings
the sum exactly to $1$, not beyond it, so a second positive is
taken to cross: $1, \frac13$ (sum $1.3333 > 1$), then
$-\frac12$ ($0.8333$), $\frac15$ ($1.0333$), $-\frac14$
($0.7833$), $\frac17, \frac19$ ($1.0373$), $-\frac16$
($0.8706$), $\frac1{11}, \frac1{13}$ ($1.0384$), $-\frac18$
($0.9134$), $\frac1{15}$ ($0.9801$), \dots\ --- the sums breathe
around $1$ with ever smaller amplitude, two positives now
needed per cycle since the negatives are larger.

\textbf{23.} Build blocks: at stage $j$, append enough unused
positive terms to raise the partial sum by at least $1$
(possible: the remaining positive terms have divergent sums,
question 12), then append the single negative term $-\frac{1}
{2j}$. Every positive term is eventually used (each stage uses
at least one), every negative one too (one per stage): a
rearrangement. Each stage changes the sum by $\geq 1 -
\frac{1}{2j} \geq \frac12$: the partial sums exceed
$\frac{j}{2}$ after stage $j$ and the increments within a stage
are positive except the last, bounded by $\frac{1}{2j} \to 0$:
divergence to $+\infty$.

\textbf{24.} By the law, $H_N \geq 20 \iff \ln N \geq 20 -
\gamma - \frac{1}{2N} + O(N^{-2})$: the threshold $N^*$
satisfies $\ln N^* = 20 - \gamma + o(1)$, i.e.\ $N^* =
\eu^{20 - \gamma}(1 + o(1)) \approx \eu^{19.4228} \approx
2.72\cdot10^{8}$ --- inside the crude window
$\intcc{1.8\cdot10^8}{4.9\cdot10^8}$ of
\cref{ex:b1:series:harmonicstack}, and pinned by $\gamma$.

\textbf{25.} (i) In $H_n = \ln n + \gamma + \frac{1}{2n} +
O(n^{-2})$: the $\ln n$ is the \omterm{thm:b1:integration:def}{integral}, $\gamma$ the price of
replacing a sum by an \omterm{thm:b1:integration:def}{integral} (a genuinely new constant of
analysis), and $\frac{1}{2n}$ the first correction --- the
trapezoid's shadow. (ii) Conditional convergence leans on
cancellation between two infinite reservoirs (question 13), so
reordering re-weights the reservoirs; \omterm{thm:b1:series:absolute}{absolute convergence} has
finite total mass, and question 14's tail estimate is
order-blind. (iii) The $(p,q)$ sums were \emph{computed}: the
$\gamma$-law turned each rearranged partial sum into $\ln 2 +
\frac12\ln\frac pq + o(1)$, with $\gamma$ itself canceling ---
an asymptotic bookkeeping exercise, not an abstract argument.
(iv) Next: products and unconditional summability for power
\omterm{def:b1:series:def}{series} in the Year 2 volume; and the Year 3 volume's weekend
problem on Stirling's formula, where sum-versus-integral
bookkeeping, pushed one order further, produces $\sqrt{2\pi}$
itself.
\end{solution}

\section*{Chapter \ref{ch:b1:vspaces} --- Vector Spaces}

\begin{solution}{exo:b1:vspaces:1}
\begin{enumerate}
  \item Yes: contains $0$, and the defining equation is linear
        (stable under $x + \lambda y$).
  \item No: does not contain $(0,0,0)$.
  \item No: $(1, 0)$ and $(0, -1)$ belong ($xy = 0$), their sum $(1,
        -1)$ does not ($xy = -1 < 0$).
  \item Yes: $0$ vanishes at $1$; $(P + \lambda Q)(1) = P(1) +
        \lambda Q(1) = 0$.
  \item Yes: the zero function is bounded; if $\abs f \leq M$ and
        $\abs g \leq M'$, then $\abs{f + \lambda g} \leq M +
        \abs\lambda M'$.
\end{enumerate}
\end{solution}

\begin{solution}{exo:b1:vspaces:2}
$(1,2,1) = a(1,0,1) + b(1,1,0)$ requires $a + b = 1$, $b = 2$, $a =
1$: incompatible ($a + b = 3 \neq 1$): not in the \omterm{def:b1:vspaces:span}{span}. $(2,1,1) =
a(1,0,1) + b(1,1,0)$: $b = 1$, $a = 1$, $a + b = 2$: consistent, so
$(2,1,1) = (1,0,1) + (1,1,0)$, in the \omterm{def:b1:vspaces:span}{span}.

Equation: $(x, y, z) = (a + b, b, a)$ means $x = y + z$: the \omterm{def:b1:vspaces:span}{span} is
the plane $\{x - y - z = 0\}$.
\end{solution}

\begin{solution}{exo:b1:vspaces:3}
First family: $\lambda(1,1,0) + \mu(1,0,1) + \nu(0,1,1) = 0$ gives
$\lambda + \mu = 0$, $\lambda + \nu = 0$, $\mu + \nu = 0$: adding,
$2(\lambda + \mu + \nu) = 0$, and subtracting each equation,
$\lambda = \mu = \nu = 0$: \omterm{def:b1:vspaces:free}{free}.

Second: $(2,4,6) = 2(1,2,3)$: linked.

Third: four vectors in $\R^3$ --- necessarily linked once dimension
is available (\cref{ch:b1:findim}); directly: $(0,1,1) = -(1,0,0) +
0\cdot(1,1,0) + (1,1,1)$, indeed $(-1,0,0) + (1,1,1) = (0,1,1)$: a
nontrivial relation.
\end{solution}

\begin{solution}{exo:b1:vspaces:4}
\omterm{def:b1:poly:def}{Polynomials} $1, (X-1), (X-1)^2, (X-1)^3$ have distinct degrees $0,
1, 2, 3$: \omterm{def:b1:vspaces:free}{free} (\cref{prop:b1:vspaces:freecriteria}~(1)); four \omterm{def:b1:vspaces:free}{free}
vectors generating (every $P \in \R_3[X]$ expands in powers of $X -
1$, e.g.\ by Taylor for \omterm{def:b1:poly:def}{polynomials}, cf.\ the proof of
\cref{prop:b1:poly:multiplicity}): a \omterm{def:b1:vspaces:free}{basis}. For $X^3$, Taylor at
$1$: $P = X^3$, $P(1) = 1$, $P'(1) = 3$, $P''(1) = 6$, $P'''(1) =
6$:
\[
X^3 = 1 + 3(X - 1) + 3(X-1)^2 + (X-1)^3 ,
\]
\omterm{prop:b1:vspaces:coordinates}{coordinates} $(1, 3, 3, 1)$ (Pascal's row, as expected from $X^3 =
((X-1)+1)^3$).
\end{solution}

\begin{solution}{exo:b1:vspaces:5}
$F \cap G$: the conditions $x = y = z$ and $x = t = 0$ together give
$x = 0$, hence $y = z = 0$, and $t = 0$: the intersection is
$\{0\}$. Sum: given $(x,y,z,t)$, seek $(a,a,a,b) \in F$ and
$(0,c,d,0) \in G$ summing to it: $a = x$, $b = t$, $c = y - x$, $d =
z - x$: always possible. So $\R^4 = F \oplus G$, and
\[
(1,2,3,4) = (1,1,1,4) + (0,1,2,0) .
\]
\end{solution}

\begin{solution}{exo:b1:vspaces:6}
An element of $G$ is $\lambda u$; if it converges, then (as
$\lambda u_n = \lambda(-1)^n$ has the two subsequential limits $\pm
\lambda$) necessarily $\lambda = 0$: $F \cap G = \{0\}$.

$F + G$ is not everything: it consists of sequences of the form
$c_n + \lambda(-1)^n$ with $(c_n)$ convergent. The sequence $v_n =
n$ is not of this form ($v_n - \lambda(-1)^n$ is unbounded, never
convergent). So $F \oplus G \subsetneq$ (space of all sequences).
\end{solution}

\begin{solution}{exo:b1:vspaces:7}
($\supseteq$) Both $F \cap G$ and $F \cap H$ lie in $F$, and their
sum lies in $G + (F \cap H)$: inclusion follows since the left side
is a \omterm{def:b1:vspaces:subspace}{subspace} containing both pieces --- concretely, an element $g +
h$ with $g \in F\cap G$, $h \in F \cap H$ is in $F$ (sum of two
elements of $F$) and in $G + (F \cap H)$.

($\subseteq$) Let $x \in F$ with $x = g + h$, $g \in G$, $h \in F
\cap H$. Then $g = x - h \in F$ (difference of elements of $F$), so
$g \in F \cap G$, and $x = g + h \in (F \cap G) + (F \cap H)$.

Counterexample to full distributivity in $\R^2$: $F =
\operatorname{Vect}(1,1)$, $G = \operatorname{Vect}(1,0)$, $H =
\operatorname{Vect}(0,1)$. Then $G + H = \R^2$, so $F \cap (G+H) =
F$, while $F \cap G = F \cap H = \{0\}$: the right side is $\{0\}
\neq F$.
\end{solution}

\begin{solution}{exo:b1:vspaces:8}
\begin{enumerate}
  \item Suppose $\sum_{i} \lambda_i \eu^{a_i x} = 0$ for all $x$.
        Multiply by $\eu^{-a_p x}$: $\lambda_p + \sum_{i < p}
        \lambda_i \eu^{(a_i - a_p)x} \to \lambda_p$ as $x \to
        +\infty$ (each exponent $a_i - a_p < 0$). The left side is
        identically $0$, so $\lambda_p = 0$; iterate downward.
  \item Let $a\cos x + b \sin x + c \cos 2x + d\sin 2x = 0$ for all
        $x$. Evaluate at $x = 0$: $a + c = 0$; at $x = \pi$: $-a + c
        = 0$; hence $a = c = 0$, and the relation reduces to $b\sin
        x + d \sin 2x = 0$. Evaluate at $x = \frac\pi2$: $b = 0$;
        then at $x = \frac\pi4$: $d = 0$.
  \item Suppose $\sum \lambda_i \abs{x - a_i} = 0$ for all $x$. The
        function $\sum_{i \neq j} \lambda_i\abs{x - a_i}$ is
        \omterm{def:b1:derivative:def}{differentiable} at $a_j$ (each term is, away from its own
        corner), so $-\lambda_j \abs{x - a_j}$, their difference,
        must be \omterm{def:b1:derivative:def}{differentiable} at $a_j$ too --- which forces
        $\lambda_j = 0$ ($\abs{\,\cdot\,}$ has a corner). This holds
        for every $j$.
\end{enumerate}
\end{solution}

\begin{solution}{exo:b1:vspaces:9}
Let $h \in H$. Since $h \in H \subseteq F + H = F + G$, write $h = f
+ g$ with $f \in F$, $g \in G$. Then $f = h - g \in H$ (both terms
in $H$, using $G \subseteq H$), so $f \in F \cap H = F \cap G
\subseteq G$, and $h = f + g \in G$. Hence $H \subseteq G$, and with
the hypothesis $G \subseteq H$: equality.

Counterexample without $G \subseteq H$: in $\R^2$, take $F =
\operatorname{Vect}(1,0)$, $G = \operatorname{Vect}(0,1)$, $H =
\operatorname{Vect}(1,1)$: then $F + G = F + H = \R^2$ and $F \cap G
= F \cap H = \{0\}$, yet $G \neq H$.
\end{solution}

\begin{solution}{exo:b1:vspaces:10}
Both \omterm{def:b1:logic:sets}{sets} are \omterm{def:b1:vspaces:subspace}{subspaces} (the defining conditions are linear and
hold for $0$). Decomposition: for $P \in \R[X]$,
\[
P(X) = \underbrace{\frac{P(X) + P(-X)}{2}}_{\in\,\mathcal P}
+ \underbrace{\frac{P(X) - P(-X)}{2}}_{\in\,\mathcal I},
\]
and a \omterm{def:b1:poly:def}{polynomial} both even and odd satisfies $P = -P$, so $P = 0$:
the sum is direct and equals $\R[X]$.

Now let $P = \sum_k a_k X^k$ be even. Then $P(X) - P(-X) = 2
\sum_{k \text{ odd}} a_k X^k$ is the zero \omterm{def:b1:poly:def}{polynomial}, so every
odd-degree coefficient vanishes (\cref{def:b1:poly:def}): $P \in
\operatorname{Vect}(1, X^2, X^4, \dots)$, i.e.\ $P = Q(X^2)$ for a
\omterm{def:b1:poly:def}{polynomial} $Q$. Conversely every \omterm{def:b1:poly:def}{polynomial} in $X^2$ is even.
\end{solution}

\begin{solution}{exo:b1:vspaces:11}
Suppose $a(x_1 + x_2) + b(x_2 + x_3) + c(x_3 + x_1) = 0$.
Regrouping on the \omterm{def:b1:vspaces:free}{free family} $(x_1, x_2, x_3)$:
\[
(a + c)\,x_1 + (a + b)\,x_2 + (b + c)\,x_3 = 0
\implies a + c = a + b = b + c = 0 .
\]
Subtracting the first two equations gives $c = b$; the third then
gives $2b = 0$, so $b = c = 0$, then $a = 0$: the family is \omterm{def:b1:vspaces:free}{free}.

For four vectors the analogous family is \emph{always} linked:
\[
(x_1 + x_2) - (x_2 + x_3) + (x_3 + x_4) - (x_4 + x_1) = 0
\]
is a nontrivial null combination (coefficients $1, -1, 1, -1$),
whatever $(x_1, x_2, x_3, x_4)$ is. The parity of the cycle length
decides.
\end{solution}

\begin{solution}{exo:b1:vspaces:12}
\begin{enumerate}
  \item If $F_1 \subseteq F_2$ or $F_2 \subseteq F_1$, the union
        is one of the two, hence proper. Otherwise pick $x \in F_1
        \setminus F_2$ and $y \in F_2 \setminus F_1$, and consider
        $x + y$. If $x + y \in F_1$, then $y = (x + y) - x \in
        F_1$: contradiction. If $x + y \in F_2$, then $x \in F_2$:
        contradiction. So $x + y \notin F_1 \cup F_2$, and $E \neq
        F_1 \cup F_2$.
  \item Suppose for contradiction that $E = F_1 \cup \dots \cup
        F_k$, with $k$ chosen \emph{minimal} among all such
        coverings. Minimality forbids $F_1 \subseteq F_2 \cup
        \dots \cup F_k$ (else drop $F_1$), so there is $x \in F_1$
        with $x \notin F_i$ for all $i \geq 2$. Since $F_1$ is
        proper, pick $y \notin F_1$. For each scalar $t$, the
        vector $y + t x$ lies in some $F_i$. It never lies in
        $F_1$: otherwise $y = (y + tx) - tx \in F_1$ (as $x \in
        F_1$). The \omterm{def:b1:structures:field}{field} is infinite, so choose $k$ distinct
        scalars $t_1, \dots, t_k$: the $k$ vectors $y + t_j x$
        fall into the $k - 1$ \omterm{def:b1:vspaces:subspace}{subspaces} $F_2, \dots, F_k$, and two
        of them, say $y + t x$ and $y + t' x$ with $t \neq t'$,
        lie in the same $F_i$ ($i \geq 2$). Then their difference
        $(t - t')x \in F_i$, so $x \in F_i$: contradiction. Hence
        no finite covering by proper \omterm{def:b1:vspaces:subspace}{subspaces} exists.
\end{enumerate}
\end{solution}

\begin{solution}{pb:b1:vspaces:1}
\textbf{1.} $L_i$ is a product of $n$ linear factors divided by a
nonzero constant (the $x_i$ are distinct), so $\deg L_i = n$.
Evaluating at $x_j$ with $j \neq i$: the factor $X - x_j$ of the
numerator vanishes, so $L_i(x_j) = 0$. At $x_i$, numerator and
denominator coincide: $L_i(x_i) = 1$.

\textbf{2.} Suppose $\sum_i \lambda_i L_i = 0$. Evaluate at $x_j$:
all terms die except $\lambda_j L_j(x_j) = \lambda_j$, so
$\lambda_j = 0$ for every $j$: the family is \omterm{def:b1:vspaces:free}{free}.

\textbf{3.} Let $D = P - \sum_i P(x_i) L_i$. Then $\deg D \leq n$
and, by question 1, $D(x_j) = P(x_j) - P(x_j) = 0$ for the $n + 1$
distinct points $x_0, \dots, x_n$. A nonzero \omterm{def:b1:poly:def}{polynomial} of degree
$\leq n$ has at most $n$ roots (\cref{cor:b1:poly:nroots}), so $D
= 0$. Thus every $P \in \R_n[X]$ is a combination of the $L_i$:
the family is generating, and with question 2, a \omterm{def:b1:vspaces:free}{basis}.

\textbf{4.} Given $y_0, \dots, y_n$, the \omterm{def:b1:poly:def}{polynomial} $P = \sum_i
y_i L_i$ has degree $\leq n$ and interpolates. Uniqueness: an
interpolating $P$ has, by question 3, \omterm{prop:b1:vspaces:coordinates}{coordinates} $(P(x_0), \dots,
P(x_n)) = (y_0, \dots, y_n)$ in the \omterm{def:b1:vspaces:free}{basis} $(L_i)$, and \omterm{prop:b1:vspaces:coordinates}{coordinates}
in a \omterm{def:b1:vspaces:free}{basis} are unique (\cref{prop:b1:vspaces:coordinates}). The
\omterm{prop:b1:vspaces:coordinates}{coordinates} of $P$ in the Lagrange \omterm{def:b1:vspaces:free}{basis} are its \emph{values at
the nodes} --- this is the whole point of the \omterm{def:b1:vspaces:free}{basis}.

\textbf{5.} Apply question 3 to $P = X^k$ ($0 \leq k \leq n$):
\[
X^k = \sum_{i=0}^{n} x_i^{k} L_i ,
\]
and $k = 0$ gives $\sum_i L_i = 1$.

\textbf{6.} $\deg N_k = k$ exactly: the family $(N_0, \dots, N_n)$
is a staircase of degrees in $\R_n[X]$, hence a \omterm{def:b1:vspaces:free}{basis} by
\cref{ex:b1:vspaces:staircase} (freeness from
\cref{prop:b1:vspaces:freecriteria}~(1), generation by finite
descent on the degree).

\textbf{7.} From the recurrence,
\[
f[x_0, x_1] = \frac{f(x_1) - f(x_0)}{x_1 - x_0},
\qquad
f[x_0, x_1, x_2] = \frac{f[x_1, x_2] - f[x_0, x_1]}{x_2 - x_0}.
\]
For $f(x) = x^2$:
\[
f[x_0, x_1] = \frac{x_1^2 - x_0^2}{x_1 - x_0} = x_0 + x_1,
\]
and then
\[
f[x_0, x_1, x_2]
= \frac{(x_1 + x_2) - (x_0 + x_1)}{x_2 - x_0}
= \frac{x_2 - x_0}{x_2 - x_0} = 1 .
\]

\textbf{8.} $\deg S \leq n$ since $Q, R$ have degree $\leq n - 1$.
At $x_0$: $S(x_0) = \frac{-(x_0 - x_n) R(x_0)}{x_n - x_0} =
R(x_0) = f(x_0)$. At $x_n$: $S(x_n) = \frac{(x_n - x_0)
Q(x_n)}{x_n - x_0} = Q(x_n) = f(x_n)$. At an \omterm{def:b1:topology:closure}{interior} node $x_i$
($1 \leq i \leq n-1$), both $Q$ and $R$ take the value $f(x_i)$,
so
\[
S(x_i) = \frac{(x_i - x_0) - (x_i - x_n)}{x_n - x_0}\, f(x_i)
= f(x_i) .
\]

\textbf{9.} Induction on the number of nodes. One node: the
interpolant is the constant $f(x_0) = f[x_0]$. Assume the claim
for $k$ nodes and let $S$ interpolate at $x_0, \dots, x_k$; by
uniqueness (question 4), $S$ is given by Aitken's lemma from $R$
(nodes $x_0, \dots, x_{k-1}$) and $Q$ (nodes $x_1, \dots, x_k$).
The coefficient of $X^{k}$ in $S$ is
\[
\frac{[X^{k-1}]\,Q - [X^{k-1}]\,R}{x_k - x_0}
= \frac{f[x_1, \dots, x_k] - f[x_0, \dots, x_{k-1}]}{x_k - x_0}
= f[x_0, \dots, x_k]
\]
by the induction hypothesis and the defining recurrence.

\textbf{10.} Let $P_k$ interpolate $f$ at $x_0, \dots, x_k$. The
difference $P_k - P_{k-1}$ has degree $\leq k$ and vanishes at
$x_0, \dots, x_{k-1}$, so by the factor theorem applied $k$ times
(\cref{thm:b1:poly:factor}) it equals $c\,N_k$ for a constant $c$;
comparing coefficients of $X^{k}$ and using question 9, $c =
f[x_0, \dots, x_k]$. Telescoping from $P_0 = f(x_0) N_0$ gives
Newton's formula. For the closed form, write $P_k = \sum_{i \leq
k} f(x_i) L_i$ (Lagrange, on the nodes $x_0, \dots, x_k$) and read
the coefficient of $X^{k}$: each $L_i$ contributes
$\frac{1}{\prod_{j \neq i}(x_i - x_j)}$, whence
\[
f[x_0, \dots, x_k] = \sum_{i=0}^{k}
\frac{f(x_i)}{\prod_{j \neq i,\, j \leq k}(x_i - x_j)} .
\]
The right side is invariant under any \omterm{def:b1:counting:objects}{permutation} of the nodes, so
the divided difference does not depend on their ordering.

\textbf{11.} If $P = a X^m + (\text{lower degrees})$, the binomial
theorem gives
\[
\Delta P = a\bigl((X+1)^m - X^m\bigr) + \dots
= a\,m\,X^{m-1} + (\text{lower degrees}),
\]
because $(X+1)^m - X^m = m X^{m-1} + \dots$ and the lower-degree
part of $P$ contributes degree $\leq m - 2$ after $\Delta$ (or
degree $\leq m-2$ terms). So $\deg \Delta P = m - 1$ with leading
coefficient $m a$. A constant $c$ gives $\Delta c = c - c = 0$.

\textbf{12.} $\deg B_k = k$: staircase, hence a \omterm{def:b1:vspaces:free}{basis} of
$\R_n[X]$ (\cref{ex:b1:vspaces:staircase}). For $\Delta B_k$
($k \geq 1$), factor the common product:
\begin{align*}
k!\,\Delta B_k
&= (X+1)X\cdots(X-k+2) - X(X-1)\cdots(X-k+1) \\
&= X(X-1)\cdots(X-k+2)\,\bigl[(X+1) - (X-k+1)\bigr] \\
&= k\,X(X-1)\cdots(X-k+2),
\end{align*}
so $\Delta B_k = \frac{X(X-1)\cdots(X-k+2)}{(k-1)!} = B_{k-1}$.

\textbf{13.} Write $P = \sum_{k=0}^{n} c_k B_k$ (\omterm{def:b1:vspaces:free}{basis}, question
12). Apply $\Delta^{j}$: by question 12, $\Delta^{j} P = \sum_{k
\geq j} c_k B_{k-j}$. Evaluate at $0$: $B_0(0) = 1$ and $B_m(0) =
0$ for $m \geq 1$ (the factor $X$ vanishes), so
$\bigl(\Delta^{j}P\bigr)(0) = c_j$. This is the forward-difference
formula.

\textbf{14.} Induction on $k$. For $k = 0$ the identity reads
$P(0) = P(0)$. Assume it for $k$ and apply it to $\Delta P$:
\[
\bigl(\Delta^{k+1} P\bigr)(0)
= \sum_{j=0}^{k} (-1)^{k-j}\binom kj \bigl(P(j+1) - P(j)\bigr).
\]
Collect the coefficient of $P(i)$: it is $(-1)^{k-i+1}\binom
k{i-1}\cdot(-1)^{0}$ from the first sum (shifted) and
$-(-1)^{k-i}\binom ki$ from the second --- together
\[
(-1)^{k+1-i}\Bigl(\binom k{i-1} + \binom ki\Bigr)
= (-1)^{k+1-i}\binom{k+1}i
\]
by Pascal's rule, which is the identity at rank $k + 1$.

\textbf{15.} Iterating question 11 from degree $n$, leading
coefficient $a_n$: after one $\Delta$, degree $n-1$ and leading
coefficient $n a_n$; after two, $n(n-1)a_n$; after $n$ steps,
degree $0$ and value $n(n-1)\cdots 1\, a_n = n!\,a_n$, a constant.
One more $\Delta$ kills it: $\Delta^{n+1}P = 0$.

\textbf{16.} If $m \geq k$: $B_k(m) = \binom mk \in \N$. If $0
\leq m < k$: one factor of $m(m-1)\cdots(m-k+1)$ is zero, so
$B_k(m) = 0$. If $m = -q$ with $q \geq 1$:
\[
B_k(-q) = \frac{(-q)(-q-1)\cdots(-q-k+1)}{k!}
= (-1)^k\,\frac{q(q+1)\cdots(q+k-1)}{k!}
= (-1)^k \binom{q+k-1}{k},
\]
an integer. So every $B_k$ \omterm{def:b1:logic:map}{maps} $\Z$ into $\Z$.

\textbf{17.} ($\Leftarrow$) If $P = \sum_k c_k B_k$ with $c_k \in
\Z$, then for $m \in \Z$, $P(m) = \sum_k c_k B_k(m) \in \Z$ by
question 16. ($\Rightarrow$) If $P$ is integer-valued, its
\omterm{prop:b1:vspaces:coordinates}{coordinates} are $c_k = \bigl(\Delta^k P\bigr)(0) = \sum_{j=0}^k
(-1)^{k-j}\binom kj P(j)$ (questions 13 and 14), an integer
combination of the integers $P(0), \dots, P(k)$. This is
P\'olya's characterization of \omterm{pb:b1:vspaces:1}{integer-valued polynomials}.

\textbf{18.} Set $Q(X) = P(X + a)$, a \omterm{def:b1:poly:def}{polynomial} of degree $\leq
n$ with $Q(0), Q(1), \dots, Q(n) \in \Z$. Its \omterm{prop:b1:vspaces:coordinates}{coordinates} in
$(B_k)_{k \leq n}$ are $c_k = \sum_{j \leq k}(-1)^{k-j}\binom kj
Q(j) \in \Z$ (question 14 uses only the values at $0, \dots, k
\leq n$). By question 17 ($\Leftarrow$), $Q$ is integer-valued on
all of $\Z$, hence so is $P(X) = Q(X - a)$.

\textbf{19.} A product of $k$ consecutive integers is $m(m-1)
\cdots(m-k+1) = k!\,B_k(m)$ for some $m \in \Z$, and $B_k(m) \in
\Z$ by question 16: the product is divisible by $k!$.

\textbf{20.} Values of $P = \frac{X(X+1)(2X+1)}{6}$ at $0,1,2,3$:
$0, 1, 5, 14$. Difference table: $\Delta$ row $1, 4, 9$;
$\Delta^2$ row $3, 5$; $\Delta^3$ row $2$. Hence, by question 13,
\[
P = 0\cdot B_0 + 1\cdot B_1 + 3\,B_2 + 2\,B_3 ,
\]
with integer \omterm{prop:b1:vspaces:coordinates}{coordinates}: $P$ is integer-valued (question 17),
while its monomial coefficients $\frac13, \frac12, \frac16$ are
not integers. Direct computation:
\begin{align*}
\Delta P &= \frac{(X+1)(X+2)(2X+3) - X(X+1)(2X+1)}{6} \\
&= \frac{(X+1)\bigl[(X+2)(2X+3) - X(2X+1)\bigr]}{6}
= \frac{(X+1)(6X+6)}{6} = (X+1)^2 .
\end{align*}
Telescoping $P(m) = \sum_{j=0}^{m-1}\Delta P(j) = \sum_{j=1}^{m}
j^2$ (with $P(0) = 0$): the sum-of-squares formula.

\textbf{21.} The values $2^i$ at $i = 0, \dots, n$ have difference
table constantly $1$ on the left edge: $\Delta^k$ of the sequence
$(2^i)$ is $(2^i)$ again (since $2^{i+1} - 2^i = 2^i$), so
$\bigl(\Delta^k P\bigr)(0) = 2^0 = 1$ for all $k \leq n$, and $P =
B_0 + B_1 + \dots + B_n$ by question 13. Then
\[
P(n+1) = \sum_{k=0}^{n}\binom{n+1}{k}
= 2^{n+1} - \binom{n+1}{n+1} = 2^{n+1} - 1 \neq 2^{n+1}:
\]
the pattern breaks at the first uncontrolled point.

\textbf{22.} By question 12, $B_k = \Delta B_{k+1}$, so
\[
\sum_{j=0}^{m-1} B_k(j)
= \sum_{j=0}^{m-1}\bigl(B_{k+1}(j+1) - B_{k+1}(j)\bigr)
= B_{k+1}(m) - B_{k+1}(0) = B_{k+1}(m).
\]
For $j < k$ the terms $B_k(j)$ vanish, so the sum really starts at
$j = k$: $\sum_{j=k}^{m-1}\binom jk = \binom m{k+1}$, the
hockey-stick identity.

\textbf{23.} Difference tables (or direct expansion) give
\[
X^2 = B_1 + 2 B_2, \qquad X^3 = B_1 + 6 B_2 + 6 B_3
\]
(check: $B_1 + 2B_2 = X + X(X-1) = X^2$; at $X = 1, 2, 3$ the
second gives $1, 8, 27$). Question 22 then yields
\[
\sum_{j=0}^{m-1} j^2 = B_2(m) + 2B_3(m)
= \binom m2 + 2\binom m3 = \frac{m(m-1)(2m-1)}{6},
\]
\[
\sum_{j=0}^{m-1} j^3 = B_2(m) + 6B_3(m) + 6B_4(m)
= \binom m2 + 6\binom m3 + 6\binom m4 .
\]
Expanding the last expression: $\binom m2 + 6\binom m3 + 6\binom
m4 = \frac{m(m-1)}{2}\bigl[1 + 2(m-2) +
\frac{(m-2)(m-3)}{2}\bigr] = \frac{m^2(m-1)^2}{4} = \binom m2^2$.
Replacing $m$ by $m + 1$: $1^3 + \dots + m^3 =
\bigl(\frac{m(m+1)}2\bigr)^2 = (1 + \dots + m)^2$, Nicomachus'
identity.

\textbf{24.} Nodes $0, 1, 2$, \omterm{def:b1:poly:def}{polynomial} $X^2$. Monomial \omterm{def:b1:vspaces:free}{basis}
$(1, X, X^2)$: \omterm{prop:b1:vspaces:coordinates}{coordinates} $(0, 0, 1)$. Lagrange \omterm{def:b1:vspaces:free}{basis}:
\omterm{prop:b1:vspaces:coordinates}{coordinates} are the values $(0, 1, 4)$ (question 4). Newton \omterm{def:b1:vspaces:free}{basis}
$(1, X, X(X-1))$: \omterm{pb:b1:vspaces:1}{divided differences} $f[0] = 0$, $f[0,1] = 1$,
$f[0,1,2] = \frac{3 - 1}{2} = 1$ (question 9), so \omterm{prop:b1:vspaces:coordinates}{coordinates}
$(0, 1, 1)$ --- indeed $X + X(X-1) = X^2$. Three bases, three
coordinate vectors, one \omterm{def:b1:poly:def}{polynomial}.

\textbf{25.} (i) Question 4 is automatic because $(L_i)$ is a
\emph{\omterm{def:b1:vspaces:free}{basis}}: existence and uniqueness of interpolation are
exactly existence and uniqueness of \omterm{prop:b1:vspaces:coordinates}{coordinates}. (ii) The Newton
\omterm{def:b1:vspaces:free}{basis} is a staircase, so \omterm{prop:b1:vspaces:coordinates}{coordinates} are computed by successive
divisions --- each new node adds one term without disturbing the
previous ones --- whereas the Lagrange \omterm{prop:b1:vspaces:coordinates}{coordinates} of $P$ are the
values $P(x_i)$, available with no computation at all. (iii) Both
bases are \omterm{def:b1:vspaces:free}{free} by the same two criteria of
\cref{prop:b1:vspaces:freecriteria}: distinct degrees for Newton,
evaluation at the nodes for Lagrange. (iv) P\'olya's theorem says
that ``$P(\Z) \subseteq \Z$'', a property of values, is equivalent
to integrality of the \omterm{prop:b1:vspaces:coordinates}{coordinates} in the \omterm{def:b1:vspaces:free}{basis} $(B_k)$ --- the
arithmetic of a \omterm{def:b1:poly:def}{polynomial} becomes visible only in the \omterm{def:b1:vspaces:free}{basis}
adapted to the question.
\end{solution}

\section*{Chapter \ref{ch:b1:findim} --- Finite Dimension}

\begin{solution}{exo:b1:findim:1}
\begin{enumerate}
  \item $z = -x - y$: $F = \{(x, y, -x-y)\} =
        \operatorname{Vect}\bigl((1,0,-1), (0,1,-1)\bigr)$; the two
        vectors are free (coordinates): $\dim F = 2$.
  \item $G = \{(x, x, 2t, t)\} =
        \operatorname{Vect}\bigl((1,1,0,0), (0,0,2,1)\bigr)$: free,
        $\dim G = 2$.
  \item $P(1) = P'(1) = 0$ means $(X-1)^2 \mid P$
        (\cref{prop:b1:poly:multiplicity}): $P = (X-1)^2(aX + b)$.
        Basis $\bigl((X-1)^2, X(X-1)^2\bigr)$, dimension $2$.
\end{enumerate}
\end{solution}

\begin{solution}{exo:b1:findim:2}
$(3,5,7) = (1,1,1) + 2(1,2,3)$ and $(0,1,2) = (1,2,3) - (1,1,1)$:
both are combinations of the first two, which are free (not
proportional). Rank $2$; basis of the span:
$\bigl((1,1,1), (1,2,3)\bigr)$.
\end{solution}

\begin{solution}{exo:b1:findim:3}
Try appending $e_1 = (1,0,0,0)$ and $e_3 = (0,0,1,0)$. The family
$\bigl((1,1,0,0), (0,0,1,1), e_1, e_3\bigr)$ is free: a null
combination $\alpha(1,1,0,0) + \beta(0,0,1,1) + \gamma e_1 + \delta
e_3 = 0$ reads $(\alpha + \gamma, \alpha, \beta + \delta, \beta) =
0$, so $\alpha = \beta = 0$, then $\gamma = \delta = 0$. Four free
vectors in dimension $4$: a basis
(\cref{prop:b1:findim:twoofthree}).
\end{solution}

\begin{solution}{exo:b1:findim:4}
Degrees $0, 1, 2, 3$ pairwise distinct: free
(\cref{prop:b1:vspaces:freecriteria}), four vectors in dimension
$4$: basis. For $X^3$: expand downward,
\[
X(X-1)(X-2) = X^3 - 3X^2 + 2X,
\qquad
X(X-1) = X^2 - X,
\]
so $X^3 - X(X-1)(X-2) = 3X^2 - 2X$; and $3X^2 - 2X = 3(X^2 - X) + X
= 3\,X(X-1) + X$. Hence
\[
X^3 = X(X-1)(X-2) + 3\,X(X-1) + 1\cdot X + 0\cdot 1 :
\]
coordinates $(0, 1, 3, 1)$ on $\bigl(1, X, X(X-1),
X(X-1)(X-2)\bigr)$. (These are Stirling numbers in disguise.)
\end{solution}

\begin{solution}{exo:b1:findim:5}
Grassmann: $\dim(F \cap G) = 4 + 5 - \dim(F + G)$, and $F + G$ is a
subspace of $E$ containing $G$: $5 \leq \dim(F+G) \leq 7$. Hence
$\dim(F \cap G) \in \{2, 3, 4\}$. Realizations in $\R^7$ with
canonical basis $(e_1, \dots, e_7)$, taking $G =
\operatorname{Vect}(e_1, \dots, e_5)$:
\begin{itemize}
  \item $F = \operatorname{Vect}(e_1, e_2, e_6, e_7)$: $F + G = \R^7$,
        intersection $\operatorname{Vect}(e_1, e_2)$, dimension $2$;
  \item $F = \operatorname{Vect}(e_1, e_2, e_3, e_6)$: intersection
        of dimension $3$;
  \item $F = \operatorname{Vect}(e_1, e_2, e_3, e_4) \subseteq G$:
        dimension $4$.
\end{itemize}
\end{solution}

\begin{solution}{exo:b1:findim:6}
$H$ is a subspace (\cref{exo:b1:vspaces:1}~(4) verbatim). By the
factor theorem (\cref{thm:b1:poly:factor}), $P \in H \iff P =
(X-1)Q$ with $\deg Q \leq n - 1$: the map $Q \mapsto (X-1)Q$ is a
linear bijection from $\R_{n-1}[X]$ onto $H$, so a basis of $H$ is
\[
\bigl((X-1),\ (X-1)X,\ (X-1)X^2,\ \dots,\ (X-1)X^{n-1}\bigr),
\qquad \dim H = n .
\]
A supplementary line: $\operatorname{Vect}(1)$ (constants). Indeed
$H \cap \operatorname{Vect}(1) = \{0\}$ (a nonzero constant does not
vanish at $1$) and dimensions add up to $n + 1$: by Grassmann, $H
\oplus \operatorname{Vect}(1) = \R_n[X]$.
\end{solution}

\begin{solution}{exo:b1:findim:7}
Deletion: removing $u_p$, the span can only shrink; and it shrinks
by at most one dimension, since adding $u_p$ back to a basis of the
smaller span gives a generating family of the larger with at most
one extra vector. Symmetrically, appending a vector $v$: the new
span contains the old with at most one extra generator, so its
dimension is $r$ (if $v$ was already in the span) or $r + 1$
(otherwise, by \cref{prop:b1:vspaces:freecriteria}~(2) a basis
extends). Both statements together give the ``rank is
$1$-Lipschitz'' conclusion.
\end{solution}

\begin{solution}{exo:b1:findim:8}
Along a strictly increasing chain, dimensions strictly increase
($F_i \subseteq F_{i+1}$, $F_i \neq F_{i+1}$ and
\cref{thm:b1:findim:subspaces}: equality of dimensions would force
equality of spaces). So $\dim F_1 < \dim F_2 < \dots < \dim F_k$ is
a strictly increasing sequence of integers in $\intint{0}{n}$: at
most $n + 1$ values, $k \leq n + 1$. Maximal chain in $\R^n$:
\[
\{0\} \subsetneq \operatorname{Vect}(e_1) \subsetneq
\operatorname{Vect}(e_1, e_2) \subsetneq \dots \subsetneq \R^n ,
\]
of length exactly $n + 1$.
\end{solution}

\begin{solution}{exo:b1:findim:9}
$F + G$ strictly contains $F$ (since $G \not\subseteq F$), so
$\dim(F + G) = n$ and Grassmann gives $\dim(F \cap G) = (n-1) +
(n-1) - n = n - 2$.

General claim, by induction on $k$: the intersection $I_k$ of $k$
hyperplanes has $\dim I_k \geq n - k$. True for $k = 1$. Step: $I_{k+1}
= I_k \cap H_{k+1}$, and Grassmann inside $E$:
\[
\dim(I_k \cap H_{k+1}) = \dim I_k + (n - 1) - \dim(I_k + H_{k+1})
\geq \dim I_k + (n-1) - n \geq n - k - 1 . \qedhere
\]
\end{solution}

\section*{Chapter \ref{ch:b1:linmaps} --- Linear Maps}

\begin{solution}{exo:b1:linmaps:1}
(1) Linear: coordinates are linear expressions. (2) Not linear:
$u(2(1,1)) = 4 \neq 2 = 2u(1,1)$. (3) Linear: differentiation and
multiplication by $X$ are, and sums of linear maps are. (4) Linear:
evaluation respects pointwise operations.
\end{solution}

\begin{solution}{exo:b1:linmaps:2}
Kernel: solve $x + y - z = 0$, $2x + y + z = 0$, $3x + 2y = 0$.
From the third, $y = -\frac{3x}{2}$; the first gives $z = x + y =
-\frac x2$; check in the second: $2x - \frac{3x}{2} - \frac x2 = 0$:
satisfied. So $\ker u = \operatorname{Vect}\bigl((2, -3, -1)\bigr)$
(taking $x = 2$), dimension $1$.

Rank--nullity: $\operatorname{rk} u = 3 - 1 = 2$. Image: spanned by
the images of the canonical basis, $u(e_1) = (1,2,3)$, $u(e_2) =
(1,1,2)$, $u(e_3) = (-1,1,0)$; the first two are free, and the rank
is $2$: basis $\bigl((1,2,3), (1,1,2)\bigr)$.

Not injective ($\ker \neq \{0\}$), not surjective (rank $2 < 3$):
consistent with \cref{cor:b1:linmaps:samedim}.
\end{solution}

\begin{solution}{exo:b1:linmaps:3}
\emph{Kernel:} $P = P'$ forces $\deg P = \deg P'$ unless $P = 0$;
but $\deg P' < \deg P$ for $P \neq 0$: so $\ker u = \{0\}$, and $u$,
an injective endomorphism of the finite-dimensional $\R_n[X]$, is an
isomorphism (\cref{cor:b1:linmaps:samedim}).

\emph{Inverse:} let $v(P) = P + P' + P'' + \dots + P^{(n)}$ (a
finite sum on $\R_n[X]$). Then
\[
v\bigl(u(P)\bigr) = \sum_{k=0}^{n} (P - P')^{(k)}
= \sum_{k=0}^{n} P^{(k)} - \sum_{k=0}^{n} P^{(k+1)}
= P - P^{(n+1)} = P ,
\]
telescoping, since $P^{(n+1)} = 0$. So $v = u^{-1}$.
\end{solution}

\begin{solution}{exo:b1:linmaps:4}
$\operatorname{im}(v \circ u) = v(\operatorname{im} u) \subseteq
\operatorname{im} v$: rank $\leq \operatorname{rk} v$. And $v$
restricted to $\operatorname{im} u$ has image
$\operatorname{im}(vu)$ with rank--nullity inside
$\operatorname{im} u$: $\operatorname{rk}(vu) \leq
\dim\operatorname{im} u = \operatorname{rk} u$.
\end{solution}

\begin{solution}{exo:b1:linmaps:5}
$u^2 = 0$ means $u(u(x)) = 0$ for all $x$: every $u(x)$ lies in
$\ker u$, i.e.\ $\operatorname{im} u \subseteq \ker u$. Then
rank--nullity:
\[
\dim E = \dim\ker u + \operatorname{rk} u \geq 2\operatorname{rk} u .
\]
Equality example in $\R^2$: $u(x, y) = (y, 0)$: $u^2 = 0$,
$\operatorname{rk} u = 1 = \frac{\dim E}{2}$.
\end{solution}

\begin{solution}{exo:b1:linmaps:6}
$(pq)^2 = pqpq = ppqq = pq$ (commutation): a projection
(\cref{thm:b1:linmaps:projchar}).

Image: $\operatorname{im}(pq) \subseteq \operatorname{im} p$, and
$= \operatorname{im}(qp) \subseteq \operatorname{im} q$: contained
in the intersection. Conversely, if $x \in \operatorname{im} p \cap
\operatorname{im} q$, then $p(x) = x$ and $q(x) = x$ (fixed points
characterize the image of a projection), so $pq(x) = x$: $x \in
\operatorname{im}(pq)$.

Kernel: $\ker p \subseteq \ker(qp) = \ker(pq)$ and likewise $\ker q
\subseteq \ker(pq)$: the sum is contained. Conversely, let $pq(x) =
0$, and write
\[
x = \underbrace{q(x)}_{\in\, \ker p} +
\underbrace{(x - q(x))}_{\in\, \ker q} :
\]
the first term satisfies $p(q(x)) = 0$, so it lies in $\ker p$; the
second is in $\ker q$ since $q(x - q(x)) = q(x) - q^2(x) = 0$. Hence
$x \in \ker p + \ker q$.
\end{solution}

\begin{solution}{exo:b1:linmaps:7}
Note first the general inclusions $\ker u \subseteq \ker u^2$ and
$\operatorname{im} u^2 \subseteq \operatorname{im} u$, and by
rank--nullity, (2) $\iff$ (3) (equal kernels $\iff$ equal ranks
$\iff$ equal images, given the inclusions).

(1 $\Rightarrow$ 2): let $u^2(x) = 0$; then $u(x) \in \ker u \cap
\operatorname{im} u = \{0\}$, so $x \in \ker u$.

(2 $\Rightarrow$ 1): by Grassmann and rank--nullity, $\dim(\ker u +
\operatorname{im} u) = \dim\ker u + \operatorname{rk} u -
\dim(\ker u \cap \operatorname{im} u) = \dim E - \dim(\ker u \cap
\operatorname{im} u)$: the sum is $E$ iff the intersection is
$\{0\}$. Let $y \in \ker u \cap \operatorname{im} u$: $y = u(x)$ and
$u(y) = 0$, so $u^2(x) = 0$, so (by (2)) $u(x) = 0$: $y = 0$. Hence
$E = \ker u \oplus \operatorname{im} u$.
\end{solution}

\begin{solution}{exo:b1:linmaps:8}
If $\varphi = 0$: then $\ker\varphi = E \subseteq \ker\psi$ forces
$\psi = 0 = 0\cdot\varphi$. If $\varphi \neq 0$: $\ker\varphi$ is a
hyperplane; pick $a \notin \ker\varphi$. Set $\lambda =
\frac{\psi(a)}{\varphi(a)}$. The form $\psi - \lambda\varphi$
vanishes on $\ker\varphi$ (both do, by the inclusion) and at $a$:
it vanishes on $\ker\varphi \oplus Ka = E$. So $\psi =
\lambda\varphi$.
\end{solution}

\begin{solution}{exo:b1:linmaps:9}
Suppose $\lambda_0 x + \lambda_1 u(x) + \dots + \lambda_{n-1}
u^{n-1}(x) = 0$. Apply $u^{n-1}$: all terms with a factor $u^{\geq
n}$ die, leaving $\lambda_0 u^{n-1}(x) = 0$, so $\lambda_0 = 0$.
Apply $u^{n-2}$ to the remaining relation: $\lambda_1 u^{n-1}(x) =
0$, so $\lambda_1 = 0$; and so on. The family is free; being of size
$n = \dim E$, it is a basis (\cref{prop:b1:findim:twoofthree}). (In
this basis, $u$ acts as a shift --- the model of maximal
nilpotence.)
\end{solution}

\begin{solution}{exo:b1:linmaps:10}
\begin{enumerate}
  \item $f \circ f = -\mathrm{id}$ is bijective, so $f$ is
        (\cref{prop:b1:logic:comp} adapted: $f$ has the two-sided
        inverse $-f$). If $f(x) = \lambda x$ with $x \neq 0$:
        applying $f$, $-x = \lambda^2 x$, so $\lambda^2 = -1$:
        impossible in $\R$.
  \item Build the family greedily. Take $x_1 \neq 0$: $(x_1,
        f(x_1))$ is free by (1). If $\operatorname{Vect}$ of the
        current family $\bigl(x_1, f(x_1), \dots, x_k,
        f(x_k)\bigr)$, call it $V_k$ --- a subspace stable under $f$
        (each generator maps to another generator or its negative:
        $f(f(x_i)) = -x_i$) --- is not all of $E$, choose $x_{k+1}
        \notin V_k$. \emph{Claim: the enlarged family is free.}
        Suppose $\alpha x_{k+1} + \beta f(x_{k+1}) + v = 0$ with $v
        \in V_k$ and $(\alpha, \beta) \neq (0,0)$. Apply $f$:
        $\alpha f(x_{k+1}) - \beta x_{k+1} + f(v) = 0$ with $f(v)
        \in V_k$. Eliminate $f(x_{k+1})$ between the two relations
        (multiply the first by $\alpha$, the second by $-\beta$, and
        add):
        \[
        (\alpha^2 + \beta^2)\, x_{k+1} \in V_k ,
        \]
        and $\alpha^2 + \beta^2 \neq 0$ forces $x_{k+1} \in V_k$:
        contradiction. So the construction continues, adding vectors
        \emph{two at a time}, until $V_k = E$: the final family is a
        basis of even size, and $n$ is even.
\end{enumerate}
\end{solution}

\section*{Chapter \ref{ch:b1:matrices} --- Matrices}

\begin{solution}{exo:b1:matrices:1}
\[
AB = \begin{pmatrix} 2 & 1\\ 1 & 0\end{pmatrix},
\quad
BA = \begin{pmatrix} 0 & 1\\ 1 & 2\end{pmatrix},
\quad
A^2 - B^2 = \begin{pmatrix} 1 & 4\\ 0 & 1\end{pmatrix} - I
= \begin{pmatrix} 0 & 4\\ 0 & 0\end{pmatrix},
\]
\[
(A+B)(A-B) = A^2 - AB + BA - B^2
= \begin{pmatrix} 0 & 4\\ 0 & 0\end{pmatrix} +
\begin{pmatrix} -2 & 0\\ 0 & 2 \end{pmatrix}
= \begin{pmatrix} -2 & 4\\ 0 & 2\end{pmatrix}.
\]
They differ by $BA - AB \neq 0$: the identity $(a+b)(a-b) = a^2 -
b^2$ requires commutativity, which fails here.
\end{solution}

\begin{solution}{exo:b1:matrices:2}
$M$: $L_2 \leftarrow L_2 - 2L_1$ kills the second row; $L_3
\leftarrow L_3 - L_1$ gives $(0, -1, -2)$. Two pivots:
$\operatorname{rk} M = 2$.

$N$: $L_3 \leftarrow L_3 - L_1$ gives $(0,1,1,1) = L_2$; then $L_3
\leftarrow L_3 - L_2 = 0$. Two pivots: $\operatorname{rk} N = 2$.
\end{solution}

\begin{solution}{exo:b1:matrices:3}
Reducing $(A \mid I_3)$: $L_2 \leftarrow L_2 - 2L_1$, $L_3
\leftarrow L_3 - L_1$:
\[
\begin{pmatrix}
1 & 0 & 1 & 1 & 0 & 0\\
0 & 1 & -1 & -2 & 1 & 0\\
0 & 1 & 0 & -1 & 0 & 1
\end{pmatrix}
\xrightarrow{L_3 \leftarrow L_3 - L_2}
\begin{pmatrix}
1 & 0 & 1 & 1 & 0 & 0\\
0 & 1 & -1 & -2 & 1 & 0\\
0 & 0 & 1 & 1 & -1 & 1
\end{pmatrix},
\]
then $L_1 \leftarrow L_1 - L_3$, $L_2 \leftarrow L_2 + L_3$:
\[
A^{-1} = \begin{pmatrix}
0 & 1 & -1\\
-1 & 0 & 1\\
1 & -1 & 1
\end{pmatrix}.
\]
\emph{Check:} first row of $A$ times first column of $A^{-1}$: $1
\cdot 0 + 0\cdot(-1) + 1\cdot 1 = 1$; times second column: $1 - 0 -
1 = 0$; times third: $-1 + 0 + 1 = 0$.
\end{solution}

\begin{solution}{exo:b1:matrices:4}
$u(1) = 1$, $u(X) = X + 1$, $u(X^2) = X^2 + 2X + 1$: columns of
coordinates in $(1, X, X^2)$ give
\[
M = \begin{pmatrix}
1 & 1 & 1\\
0 & 1 & 2\\
0 & 0 & 1
\end{pmatrix}.
\]
$u$ is invertible because it has the obvious inverse $P \mapsto P(X
- 1)$ (composition of substitutions). Its matrix is obtained the
same way from $u^{-1}(X^k) = (X-1)^k$:
\[
M^{-1} = \begin{pmatrix}
1 & -1 & 1\\
0 & 1 & -2\\
0 & 0 & 1
\end{pmatrix}.
\]
\end{solution}

\begin{solution}{exo:b1:matrices:5}
$N = \begin{pmatrix} 0 & 1\\ 0 & 0\end{pmatrix}$, $N^2 = 0$. Since
$2I$ and $N$ commute, the binomial expansion truncates after two
terms:
\[
A^k = (2I + N)^k = 2^k I + k\,2^{k-1} N
= \begin{pmatrix} 2^k & k\,2^{k-1}\\ 0 & 2^k\end{pmatrix}.
\]
(Check $k = 2$: $A^2 = \begin{pmatrix}4 & 4\\ 0 & 4\end{pmatrix}$,
correct by direct product.)
\end{solution}

\begin{solution}{exo:b1:matrices:6}
Traces: $\operatorname{tr}(AB - BA) = \operatorname{tr}(AB) -
\operatorname{tr}(BA) = 0$ (\cref{def:b1:matrices:transpose}),
while $\operatorname{tr}(I_n) = n \neq 0$ in $\R$ or $\C$. No
solution. (On infinite-dimensional spaces the identity \emph{is}
realizable --- differentiation and multiplication by $x$ satisfy it
--- precisely because no trace exists there.)
\end{solution}

\begin{solution}{exo:b1:matrices:7}
Telescoping product, all powers of $A$ commuting:
\[
(I - A)(I + A + \dots + A^{m-1}) = I - A^m = I ,
\]
and \cref{prop:b1:matrices:ring} upgrades the one-sided inverse. For
the application: the given matrix is $I + N$ with
\[
N = \begin{pmatrix} 0 & 2 & 3\\ 0 & 0 & 2\\ 0&0&0 \end{pmatrix},
\quad
N^2 = \begin{pmatrix} 0&0&4\\ 0&0&0\\ 0&0&0\end{pmatrix},
\quad N^3 = 0 ,
\]
so, replacing $A$ by $-N$ in the formula:
\[
(I + N)^{-1} = I - N + N^2 =
\begin{pmatrix}
1 & -2 & 1\\
0 & 1 & -2\\
0 & 0 & 1
\end{pmatrix}.
\]
\end{solution}

\begin{solution}{exo:b1:matrices:8}
$A^2 = A$: the endomorphism $a$ is a projection
(\cref{thm:b1:linmaps:projchar}), $E = \operatorname{im} a \oplus
\ker a$ with $\dim\operatorname{im} a = r = \operatorname{rk} A$. In
a basis adapted to this decomposition ($r$ vectors of the image,
then a basis of the kernel), the matrix of $a$ is
$\begin{pmatrix} I_r & 0\\ 0 & 0\end{pmatrix}$, of trace $r$. The
trace is invariant under change of basis:
$\operatorname{tr}(P^{-1}MP) = \operatorname{tr}(MPP^{-1}) =
\operatorname{tr} M$ by the cyclic identity. Hence
$\operatorname{tr} A = r = \operatorname{rk} A$.
\end{solution}

\begin{solution}{exo:b1:matrices:9}
$J^2 = nJ$ (each entry of $J^2$ sums $n$ ones). Seek $M^{-1} =
\alpha I + \beta J$:
\[
(aI + bJ)(\alpha I + \beta J)
= a\alpha\, I + (a\beta + b\alpha + nb\beta)\, J .
\]
This equals $I$ iff $a\alpha = 1$ and $a\beta + b\alpha + nb\beta =
0$, i.e.\ $\alpha = \frac1a$ and $\beta(a + nb) = -\frac ba$. If $a
\neq 0$ and $a + nb \neq 0$:
\[
M^{-1} = \frac 1a I - \frac{b}{a(a + nb)}\, J .
\]
Conversely, if $a = 0$: $M = bJ$ has rank $\leq 1 < n$ (for $n \geq
2$): not invertible ($n = 1$ is the scalar case). If $a + nb = 0$:
the vector $v = (1, \dots, 1)^{\mathsf T}$ satisfies $Mv = (a +
nb)v = 0$ with $v \neq 0$: not invertible. So $M \in GL_n \iff a
\neq 0$ and $a + nb \neq 0$.
\end{solution}

\begin{solution}{exo:b1:matrices:10}
\emph{Sum:} $\operatorname{im}(A + B) \subseteq \operatorname{im} A
+ \operatorname{im} B$ (each $(A+B)x = Ax + Bx$), and Grassmann
bounds the dimension of a sum by the sum of dimensions.

\emph{Sylvester:} let $a$ be the map of $A$ restricted to $V =
\operatorname{im} B$ (dimension $\operatorname{rk} B$). Its image is
$\operatorname{im}(AB)$ ($a(Bx) = ABx$), and rank--nullity in $V$:
\[
\operatorname{rk} B = \dim\ker(a_{|V}) + \operatorname{rk}(AB) .
\]
Now $\ker(a_{|V}) \subseteq \ker A$, of dimension $n -
\operatorname{rk} A$: so
\[
\operatorname{rk}(AB) \geq \operatorname{rk} B - (n -
\operatorname{rk} A) = \operatorname{rk} A + \operatorname{rk} B -
n . \qedhere
\]
\end{solution}

\section*{Chapter \ref{ch:b1:det} --- Determinants and Linear Systems}

\begin{solution}{exo:b1:det:1}
$3 \times 2 - 1 \times 5 = 1$.

Second: $L_2 \leftarrow L_2 - L_1$, $L_3 \leftarrow L_3 - L_2$ (on
the original rows) give rows $(1,2,3), (3,3,3), (3,3,3)$: two equal
rows, determinant $0$. (Sarrus confirms: $45 + 84 + 96 - 105 - 48 -
72 = 0$.)

Third: it is Vandermonde with $x = 1, 2, 3$
(\cref{ex:b1:det:vandermonde}): $(2-1)(3-1)(3-2) = 2$.
\end{solution}

\begin{solution}{exo:b1:det:2}
The determinant equals (add all columns to the first, factor)
$(\lambda + 2)$ times
\[
\begin{vmatrix}
1 & 1 & \lambda\\ 1 & \lambda & 1\\ 1 & 1 & 1
\end{vmatrix}
= -(\lambda - 1)^2
\]
(clear with $L_1 \leftarrow L_1 - L_3$, $L_2 \leftarrow L_2 - L_3$
and expand), giving $\det = -(\lambda+2)(\lambda-1)^2$. \omterm{def:b1:vspaces:free}{Basis}
$\iff \det \neq 0 \iff \lambda \notin \{1, -2\}$.
\end{solution}

\begin{solution}{exo:b1:det:3}
First system: $\det = -7$;
$x = \frac{1}{-7}\begin{vmatrix} 5 & 1\\ 4 & -2\end{vmatrix}
= \frac{-14}{-7} = 2$,
$y = \frac{1}{-7}\begin{vmatrix} 2 & 5\\ 3 & 4\end{vmatrix}
= \frac{-7}{-7} = 1$. Check: $2(2) + 1 = 5$; $3(2) - 2 = 4$.

Second system: after $L_2 - L_1$ and $L_3 - 2L_1$, the rows become
$(1,1,1)$, $(0,-2,0)$, $(0,-1,-3)$, so
\[
\det A = \begin{vmatrix} 1&1&1\\ 1&-1&1\\ 2&1&-1\end{vmatrix}
= 1 \times \begin{vmatrix} -2 & 0\\ -1 & -3\end{vmatrix} = 6 .
\]
Cramer, replacing columns by $(6,2,1)^{\mathsf T}$:
\[
x = \frac{6}{6} = 1, \qquad
y = \frac{12}{6} = 2, \qquad
z = \frac{18}{6} = 3
\]
(the numerators computed the same way). Check: $1 + 2 + 3 = 6$; $1 -
2 + 3 = 2$; $2 + 2 - 3 = 1$.
\end{solution}

\begin{solution}{exo:b1:det:4}
Reduce the augmented matrix: $L_2 \leftarrow L_2 - 2L_1$, $L_3
\leftarrow L_3 - L_1$:
\[
\begin{pmatrix}
1 & 2 & -1 & 1 & 1\\
0 & 0 & 3 & -3 & 3\\
0 & 0 & 3 & -3 & 3
\end{pmatrix}
\to
\begin{pmatrix}
1 & 2 & -1 & 1 & 1\\
0 & 0 & 1 & -1 & 1\\
0 & 0 & 0 & 0 & 0
\end{pmatrix}.
\]
Pivot unknowns $x, z$; \omterm{def:b1:vspaces:free}{free} unknowns $y, t$. Back-substitution: $z =
1 + t$, $x = 1 - 2y + z - t = 2 - 2y$. Solution \omterm{def:b1:logic:sets}{set}:
\[
\{(2 - 2y,\; y,\; 1 + t,\; t) : y, t \in \R\}
= (2, 0, 1, 0) + \operatorname{Vect}\bigl((-2,1,0,0),\,
(0,0,1,1)\bigr),
\]
an affine plane (dimension $2 = 4 - \operatorname{rk} 2$) of $\R^4$.
\end{solution}

\begin{solution}{exo:b1:det:5}
Induction; $n = 1$ is empty product $= 1$. For the step, perform
$C_k \leftarrow C_k - x_1 C_{k-1}$ for $k = n, n-1, \dots, 2$ (in
this order, so each operation uses a not-yet-modified column). The
first row becomes $(1, 0, \dots, 0)$; in row $i \geq 2$, the $k$-th
entry becomes $x_i^{k-1} - x_1 x_i^{k-2} = x_i^{k-2}(x_i - x_1)$.
Expanding along the first row and factoring $(x_i - x_1)$ out of
each row $i$:
\[
V(x_1, \dots, x_n)
= \prod_{i=2}^{n} (x_i - x_1)\cdot V(x_2, \dots, x_n),
\]
and the induction hypothesis completes the product $\prod_{i<j}(x_j
- x_i)$.
\end{solution}

\begin{solution}{exo:b1:det:6}
Expanding $D_n$ along the first row: $D_n = 2 D_{n-1} -
1\cdot\begin{vmatrix} 1 & \ast\\ 0 & D_{n-2}\text{-block}
\end{vmatrix}$; the second determinant, expanded along its first
column, is $D_{n-2}$. So $D_n = 2D_{n-1} - D_{n-2}$, i.e.\ $D_n -
D_{n-1} = D_{n-1} - D_{n-2}$: the differences are constant, equal to
$D_2 - D_1 = 1$. Hence $D_n = D_1 + (n - 1) = n + 1$.
(Check: $D_2 = 3$, and the $3\times3$ case is
\cref{ex:b1:det:cofactor}: $D_3 = 4$.)
\end{solution}

\begin{solution}{exo:b1:det:7}
$C_1 \leftarrow C_1 + C_2 + C_3$ makes the first column constant
$(m+2)$; factor it:
\[
\det = (m+2)\begin{vmatrix}
1 & 1 & m\\ 1 & m & 1\\ 1 & 1 & 1
\end{vmatrix}
\overset{L_1 - L_3,\ L_2 - L_3}{=}
(m+2)\begin{vmatrix}
0 & 0 & m-1\\ 0 & m-1 & 0\\ 1 & 1 & 1
\end{vmatrix}
= (m+2)\cdot\bigl(-(m-1)^2\bigr)
\]
(expand along the first column: the single entry $1$ carries sign
$+$, and the remaining $2 \times 2$ determinant is $0 \cdot 0 -
(m-1)(m-1) = -(m-1)^2$).

Discussion. $m \notin \{1, -2\}$: Cramer system; by the symmetry of
the equations, $x = y = z$, and each equation gives $(m + 2)x = 1$:
unique solution $\bigl(\frac{1}{m+2}, \frac{1}{m+2},
\frac{1}{m+2}\bigr)$. $m = 1$: the three equations all read $x + y
+ z = 1$: solutions form the affine plane $x + y + z = 1$. $m = -2$:
summing the three equations gives $0 = 3$: empty solution \omterm{def:b1:logic:sets}{set}.
\end{solution}

\begin{solution}{exo:b1:det:8}
($\Rightarrow$) If $A^{-1}$ has integer entries: $\det A \cdot \det
A^{-1} = 1$ with both determinants integers (sums of products of
entries): two integers with product $1$ are both $\pm1$.

($\Leftarrow$) The cofactor formula $A^{-1} = \frac{1}{\det
A}\operatorname{Com}(A)^{\mathsf T}$ (checked for $n \leq 3$ by
direct expansion, admitted in general) has
$\operatorname{Com}(A)$ with integer entries (each cofactor is an
integer determinant); dividing by $\det A = \pm 1$ keeps integers.
\end{solution}

\begin{solution}{exo:b1:det:9}
Add all columns to the first: each entry of the new first column is
$a + nb$; factor it out, so that the first column is all ones. Then
the \omterm{met:b1:matrices:gauss}{row operations} $L_i \leftarrow L_i - L_1$ ($i \geq 2$) clear
every entry below the top-left $1$ and leave $a$ on the diagonal and
$0$ elsewhere in those rows: the matrix is upper triangular with
diagonal $(1, a, \dots, a)$. Hence
\[
\det(aI + bJ) = (a + nb)\, a^{\,n-1} ,
\]
nonzero iff $a \neq 0$ and $a + nb \neq 0$: the condition of
\cref{exo:b1:matrices:9}.
\end{solution}

\begin{solution}{exo:b1:det:10}
Block operations (each a composition of the corresponding $n$ scalar
operations, allowed by \cref{thm:b1:det:props}~(1)):
\[
\begin{vmatrix} A & B\\ B & A\end{vmatrix}
\overset{C_1 \leftarrow C_1 + C_2}{=}
\begin{vmatrix} A + B & B\\ A + B & A\end{vmatrix}
\overset{L_2 \leftarrow L_2 - L_1}{=}
\begin{vmatrix} A + B & B\\ 0 & A - B\end{vmatrix}
= \det(A+B)\,\det(A-B),
\]
using the block-triangular rule for the last step.
\end{solution}

\begin{solution}{exo:b1:det:11}
$C_1 \leftarrow C_1 + C_2 + C_3$ makes the first column constant
$(a + b + c)$; factor it out, then $L_2 \leftarrow L_2 - L_1$,
$L_3 \leftarrow L_3 - L_1$:
\[
\Delta = (a+b+c)\begin{vmatrix}
1 & b & c\\
0 & a - b & b - c\\
0 & c - b & a - c
\end{vmatrix}
= (a+b+c)\bigl[(a-b)(a-c) + (b-c)^2\bigr],
\]
and expanding, $(a-b)(a-c) + (b-c)^2 = a^2 + b^2 + c^2 - ab - bc
- ca$. Over $\C$, with $j^3 = 1$ and $1 + j + j^2 = 0$:
\begin{align*}
(a + jb + j^2c)(a + j^2b + jc)
&= a^2 + b^2 + c^2 + (j + j^2)(ab + bc + ca)\\
&= a^2 + b^2 + c^2 - ab - bc - ca ,
\end{align*}
whence the complete factorization. (Structurally: the column
$(1, j, j^2)^{\mathsf T}$ satisfies $M\,(1, j, j^2)^{\mathsf T}
= (a + jb + j^2c)(1, j, j^2)^{\mathsf T}$, and similarly for
$j^2$ and $1$: the three factors are the three ``eigenvalues''
of the circulant, a story systematized in the Year~2 volume.)
\end{solution}

\begin{solution}{exo:b1:det:12}
Write $r = \operatorname{rk} A$.

\emph{An invertible $r \times r$ submatrix exists.} Choose $r$
\omterm{def:b1:vspaces:free}{free} columns of $A$ and let $B \in \mathcal{M}_{n,r}$ be the
matrix they form: $\operatorname{rk} B = r$. Since row rank
equals column rank (\cref{thm:b1:matrices:rank}), $B$ has $r$
\omterm{def:b1:vspaces:free}{free} rows; keeping those rows yields an $r \times r$ submatrix
of $A$ of rank $r$, i.e.\ invertible.

\emph{No larger one does.} Let $S$ be an invertible $s \times s$
submatrix, taken from columns $j_1, \dots, j_s$ and rows $i_1,
\dots, i_s$ of $A$. If a combination $\sum_k \lambda_k
C_{j_k} = 0$ of the corresponding \emph{full} columns vanishes,
then reading only the rows $i_1, \dots, i_s$ gives $\sum_k
\lambda_k S_k = 0$ on the columns of $S$, so all $\lambda_k = 0$
($S$ invertible): the columns $C_{j_1}, \dots, C_{j_s}$ of $A$
are \omterm{def:b1:vspaces:free}{free}, and $s \leq \operatorname{rk} A = r$.

Hence $\operatorname{rk} A$ is exactly the largest size of an
invertible submatrix.
\end{solution}

\begin{solution}{pb:b1:det:1}
\textbf{1.} $V(1,2,3,4) = (2-1)(3-1)(4-1)(3-2)(4-2)(4-3) = 1
\cdot 2\cdot 3\cdot 1\cdot 2\cdot 1 = 12$. Interpolation at
distinct nodes asks for the coefficients of $P$ solving $W c =
y$ with $W = (x_i^{\,j-1})$, and $\det W = V \neq 0$: a Cramer
system.

\textbf{2.} Work through the columns from left to right. $C_1$
is the constant column $P_0(x_i) = 1$ ($P_0$ \omterm{def:b1:poly:def}{monic} of degree
$0$). Assume columns $1, \dots, j-1$ have already been reduced
to the pure powers $1, x_i, \dots, x_i^{\,j-2}$. Since $P_{j-1}
= X^{j-1} + \sum_{k < j-1}\alpha_k X^k$, subtracting from $C_j$
the combination $\sum_k \alpha_k\,(\text{column of } x_i^k)$
--- an operation that does not change the determinant --- leaves
the pure power column $x_i^{\,j-1}$. After the last column the
matrix is the Vandermonde matrix: $\det = V(x_1, \dots, x_n)$.

\textbf{3.} The \omterm{def:b1:poly:def}{polynomials} $(j-1)!\,B_{j-1}$ are \omterm{def:b1:poly:def}{monic} of
degree $j - 1$, so question 2 gives
\[
\det\bigl(B_{j-1}(m_i)\bigr)
= \frac{V(m_1, \dots, m_n)}{0!\,1!\cdots(n-1)!} .
\]
The left side is the determinant of a matrix with \emph{integer}
entries ($B_k$ is integer-valued on $\Z$: questions 16--17 of
the weekend problem \cref{pb:b1:vspaces:1}), hence an integer;
and it is positive since $V(m_1, \dots, m_n) > 0$ for $m_1 <
\dots < m_n$. So the superfactorial \omterm{def:b1:arith:divides}{divides} the product of all
pairwise differences.

\textbf{4.} Factor $x_i$ out of each row $i$:
$\det(x_i^{\,j})_{j = 1..n} = x_1\cdots x_n\,
\det(x_i^{\,j-1}) = x_1\cdots x_n\,V$.

\textbf{5.} $(W^{\mathsf T}W)_{ij} = \sum_k x_k^{\,i-1}
x_k^{\,j-1} = p_{i+j-2}$: $S = W^{\mathsf T}W$. Hence $\det S =
\det(W^{\mathsf T})\det W = V^2$
(\cref{thm:b1:det:props}~(2),(4)). For real $x_i$: $\det S = V^2
\geq 0$, and $S$ is invertible iff $V \neq 0$ iff the $x_i$ are
pairwise distinct --- a sign-definite test computable from the
power sums alone.

\textbf{6.} The interpolation conditions $\sum_{k}
c_k\,x_i^{\,k} = y_i$ form the system $Wc = y$; $\det W = V \neq
0$ gives existence and uniqueness at once. This is the third
proof in the book: explicit formula in
\cref{thm:b1:poly:lagrange}, kernel argument in
\cref{ex:b1:linmaps:interpolation}, Cramer here.

\textbf{7.} Cramer: $c_{n-1} = \det W'/\det W$ where $W'$ is $W$
with its last column replaced by $y$. Expanding $\det W'$ along
that column:
\[
\det W' = \sum_{i=1}^n (-1)^{i+n} y_i\,V(x_1, \dots, \widehat{x_i},
\dots, x_n) .
\]
Now $V = V(\setminus i)\cdot\prod_{j<i}(x_i - x_j)\prod_{j>i}
(x_j - x_i)$, and converting the second product costs
$(-1)^{n-i}$:
\[
(-1)^{i+n}\,\frac{V(\setminus i)}{V}
= \frac{(-1)^{i+n}(-1)^{n-i}}{\prod_{j\neq i}(x_i - x_j)}
= \frac{1}{\prod_{j\neq i}(x_i - x_j)} ,
\]
whence $c_{n-1} = \sum_i y_i/\prod_{j \neq i}(x_i - x_j)$ ---
the divided-difference formula again.

\textbf{8.} $L_3 \leftarrow L_3 - L_1$ gives the rows $(1, x_1,
x_1^2)$, $(0, 1, 2x_1)$, $(0,\ x_2 - x_1,\ (x_2-x_1)(x_2+x_1))$;
expanding along the first column and factoring $(x_2 - x_1)$:
\[
(x_2 - x_1)\begin{vmatrix} 1 & 2x_1\\ 1 & x_2 + x_1
\end{vmatrix} = (x_2 - x_1)(x_2 - x_1) = (x_2 - x_1)^2 .
\]
Nonzero for $x_1 \neq x_2$: the \omterm{def:b1:det:system}{linear system} expressing
$P(x_1) = u$, $P'(x_1) = v$, $P(x_2) = w$ on the coefficients of
$P \in \R_2[X]$ is Cramer --- Hermite interpolation with a
doubled node is well posed.

\textbf{9.} $P = a + bX + cX^2$ with $a = P(0) = 1$, $b = P'(0)
= 0$, $a + b + c = P(1) = 2$: $c = 1$, so $P = 1 + X^2$, unique.
Consistency: here $x_1 = 0$, $x_2 = 1$ and the determinant of
question 8 is $(1 - 0)^2 = 1 \neq 0$.

\textbf{10.} Direct computation:
\[
C_2 = \frac{1}{(a_1+b_1)(a_2+b_2)} - \frac{1}{(a_1+b_2)(a_2+b_1)}
= \frac{(a_1+b_2)(a_2+b_1) - (a_1+b_1)(a_2+b_2)}
{\prod_{i,j}(a_i+b_j)} ,
\]
and the numerator expands to $a_1b_1 + a_2b_2 - a_1b_2 - a_2b_1
= (a_2 - a_1)(b_2 - b_1)$: differences over sums.

\textbf{11.} For $i < n$, the new entry of row $i$ is
\[
\frac{1}{a_i + b_j} - \frac{1}{a_n + b_j}
= \frac{a_n - a_i}{(a_i + b_j)(a_n + b_j)} .
\]
Factor $(a_n - a_i)$ out of each row $i < n$, then
$\frac1{a_n + b_j}$ out of each column $j$: what remains has
entries $\frac1{a_i + b_j}$ in rows $i < n$ and constant $1$ in
row $n$ --- the matrix $M$, with the announced prefactor.

\textbf{12.} On $M$, for $j < n$ the operation $C_j \leftarrow
C_j - C_n$ turns row $n$ into $(0, \dots, 0, 1)$ and, in row $i
< n$,
\[
\frac{1}{a_i + b_j} - \frac{1}{a_i + b_n}
= \frac{b_n - b_j}{(a_i + b_j)(a_i + b_n)} .
\]
Factor $(b_n - b_j)$ from each column $j < n$ and $\frac1{a_i +
b_n}$ from each row $i < n$, then expand along the last row
(sign $(-1)^{n+n} = +1$): the remaining determinant is
$C_{n-1}$. Collecting the factors of questions 11--12:
\[
C_n = \frac{\prod_{i<n}(a_n - a_i)\,\prod_{j<n}(b_n - b_j)}
{\prod_{j}(a_n + b_j)\,\prod_{i<n}(a_i + b_n)}\;C_{n-1},
\]
and the induction (base $C_1 = \frac1{a_1+b_1}$) assembles
exactly Cauchy's double alternant: the factors $(a_j - a_i)(b_j
- b_i)$ for all pairs, over all sums $(a_i + b_j)$.

\textbf{13.} The formula vanishes iff some $a_j = a_i$ or $b_j =
b_i$: the Cauchy matrix is invertible iff both families are
pairwise distinct. Hilbert: $a_i = i$, $b_j = j - 1$. For $n =
2$: numerator $(2-1)(1-0) = 1$, denominator $1\cdot2\cdot2\cdot3
= 12$: $\det H_2 = \frac1{12}$. For $n = 3$: numerator
$\bigl[(1)(2)(1)\bigr]^2 = 4$, denominator $(1\cdot2\cdot3)
(2\cdot3\cdot4)(3\cdot4\cdot5) = 6\cdot24\cdot60 = 8640$:
$\det H_3 = \frac{4}{8640} = \frac1{2160}$. Inverse for $n = 2$:
\[
H_2^{-1} = 12\begin{pmatrix} \frac13 & -\frac12\\[2pt]
-\frac12 & 1\end{pmatrix}
= \begin{pmatrix} 4 & -6\\ -6 & 12 \end{pmatrix},
\]
all integers (a phenomenon true for every $n$).

\textbf{14.} The matrix of the system is the Cauchy matrix,
invertible by question 13 when the $b_j$ (and the $a_i$) are
pairwise distinct: unique solution. Interpretation: a rational
function $R = \sum_j \frac{c_j}{X + b_j}$ with simple \omterm{def:b1:fractions:field}{poles} is
determined by $n$ of its values $R(a_1), \dots, R(a_n)$, and
conversely any such data sheet is realized exactly once ---
the sampling counterpart of the partial-fraction existence and
uniqueness theorem (\cref{thm:b1:fractions:complex}).

\textbf{15.} Vieta for $X^3 + pX + q$: $\lambda_1 + \lambda_2 +
\lambda_3 = 0$, $\sum_{i<j}\lambda_i\lambda_j = p$, so $p_1 = 0$
and $p_2 = p_1^2 - 2p = -2p$. Each root satisfies $\lambda^3 =
-p\lambda - q$; summing: $p_3 = -p\,p_1 - 3q = -3q$. Multiplying
by $\lambda$ and summing: $p_4 = -p\,p_2 - q\,p_1 = 2p^2$.

\textbf{16.} By question 5 (the identity $S = W^{\mathsf T}W$
and $\det S = V^2$ are valid over $\C$),
\[
V^2 = \begin{vmatrix}
3 & 0 & -2p\\
0 & -2p & -3q\\
-2p & -3q & 2p^2
\end{vmatrix}
= 3\bigl(-4p^3 - 9q^2\bigr) + (-2p)\bigl(0 - 4p^2\bigr)
= -4p^3 - 27q^2 ,
\]
expanding along the first row.

\textbf{17.} A repeated root means two equal $\lambda_i$, i.e.\
$V = 0$, i.e.\ $\operatorname{disc} = -4p^3 - 27q^2 = 0$. For
$X^3 - 3X + 2$: $4(-3)^3 + 27\cdot4 = -108 + 108 = 0$, matching
the double root $1$ of $(X-1)^2(X+2)$.

\textbf{18.} The nonreal roots of a real cubic come in conjugate
pairs, so exactly two cases occur when $\operatorname{disc} \neq
0$. Three distinct real roots: $V$ is real and nonzero, so
$\operatorname{disc} = V^2 > 0$. One real root $\lambda_1$ and
$\lambda_3 = \conj{\lambda_2} \notin \R$: then
\[
(\lambda_2 - \lambda_1)(\lambda_3 - \lambda_1) =
\abs{\lambda_2 - \lambda_1}^2 > 0,
\qquad
\lambda_3 - \lambda_2 = -2\iu\,\operatorname{Im}\lambda_2 \neq 0,
\]
so $V$ is a nonzero purely imaginary number and
$\operatorname{disc} = V^2 < 0$. The two signs characterize the
two cases.

\textbf{19.} Take $b_i = a_i$ in the double alternant: the
numerator is $\prod_{i<j}(a_j - a_i)^2 > 0$ and the denominator
$\prod_{i,j}(a_i + a_j) > 0$ (all entries positive): the
determinant is positive. (In later language: the kernel
$\frac1{x+y}$ is positive definite.)

\textbf{20.} By questions 2--3, the determinant equals
$V(2,4,7)/(0!\,1!\,2!) = \frac{(4-2)(7-2)(7-4)}{2} =
\frac{30}{2} = 15$. Directly, the matrix is
\[
\begin{pmatrix}
1 & 2 & 1\\
1 & 4 & 6\\
1 & 7 & 21
\end{pmatrix},
\qquad
\det = (84 - 42) - 2(21 - 6) + (7 - 4) = 42 - 30 + 3 = 15 .
\]

\textbf{21.} Suppose $\sum_i c_i\,(\lambda_i^{\,k})_{k} = 0$ as
a sequence. Reading off $k = 0, 1, \dots, n-1$ gives $W^{\mathsf
T}c = 0$ with $W = (\lambda_i^{\,j-1})$ invertible ($\det = V
\neq 0$, distinct $\lambda_i$): $c = 0$. The geometric sequences
are \omterm{def:b1:vspaces:free}{free}.

\textbf{22.} $a = b = (1, 2, 3)$: numerator $\bigl[(2-1)(3-1)
(3-2)\bigr]^2 = 4$; denominator $\prod_{i,j}(i + j) =
(2\cdot3\cdot4)(3\cdot4\cdot5)(4\cdot5\cdot6) = 24\cdot60\cdot120
= 172800$. Hence $\det\bigl(\frac1{i+j}\bigr) = \frac{4}{172800}
= \frac1{43200}$.

\textbf{23.} If $x_i = x_j$, the swap of the two variables fixes
the point but must change the sign of $F$: $F = -F$, so $F = 0$
there. \omterm{def:b1:arith:divides}{Divisibility}: view $F$ as a \omterm{def:b1:poly:def}{polynomial} in the single
variable $x_n$ with coefficients in the other variables; it
vanishes at the $n - 1$ ``values'' $x_1, \dots, x_{n-1}$, so
repeated factoring (\cref{thm:b1:poly:factor}) gives $F =
\prod_{i<n}(x_n - x_i)\cdot G$ with $G$ \omterm{def:b1:poly:def}{polynomial}. The
prefactor is invariant under swaps of two indices $i, j < n$, so
$G$ is alternating in $x_1, \dots, x_{n-1}$, and induction
completes: $\prod_{i<j}(x_j - x_i)$ \omterm{def:b1:arith:divides}{divides} $F$.

\textbf{24.} $D = \det(x_i^{\,j-1})$ is a \omterm{def:b1:poly:def}{polynomial} in the
$x_i$; swapping two variables swaps two rows, so $D$ is
alternating, and by question 23, $D = c\,\prod_{i<j}(x_j - x_i)$
for some \omterm{def:b1:poly:def}{polynomial} $c$. Total degrees: $D$ has degree $\leq 0 +
1 + \dots + (n-1) = \binom n2$, the product has degree exactly
$\binom n2$: $c$ is a constant. The monomial $x_2\,x_3^2\cdots
x_n^{\,n-1}$ has coefficient $1$ in $D$ (diagonal product) and
$1$ in the product (choose the larger-index variable in every
factor): $c = 1$, and Vandermonde's formula falls out with no
induction.

\textbf{25.} (i) Alternation --- the axiom ``two equal columns
kill the determinant'' --- is the engine: it produced every
factor $(x_j - x_i)$, $(a_j - a_i)$, $(b_j - b_i)$ in the
problem. (ii) The identity $\det S = V^2$ replaces the
individually complex, unreachable roots by their power sums,
which are real \omterm{def:b1:poly:def}{polynomials} in the coefficients, so distinctness
becomes the sign of a computable real number. (iii) The
\omterm{ex:b1:det:vandermonde}{Vandermonde determinant} governs \omterm{def:b1:poly:def}{polynomial} interpolation, the
\omterm{pb:b1:det:1}{Cauchy determinant} governs partial fractions and sampled
rational functions (with the \omterm{pb:b1:det:1}{Hilbert matrix} as its most famous
special case). (iv) Any alternating \omterm{def:b1:poly:def}{polynomial} is divisible by
$\prod_{i<j}(x_j - x_i)$, and a degree count then pins such a
\omterm{def:b1:poly:def}{polynomial} down up to a constant --- which is why this product
keeps reappearing wherever a determinant vanishes on
coincidences. The theorem of Part III is \emph{Cauchy's double
alternant}.
\end{solution}

\section*{Chapter \ref{ch:b1:euclid} --- Euclidean Spaces}

\begin{solution}{exo:b1:euclid:1}
$\langle u, v\rangle = 2 - 4 + 2 = 0$; $\norm u = \norm v = 3$. The
vectors are orthogonal: the angle is $\frac\pi2$. Cauchy--Schwarz:
$\abs 0 \leq 9$, comfortably.
\end{solution}

\begin{solution}{exo:b1:euclid:2}
Expand both squared norms with the identity of
\cref{thm:b1:euclid:cs} and add: the cross terms cancel.

Sup norm: take $x = (1, 0)$, $y = (0, 1)$. Then $\norm{x + y}_\infty
= \norm{x - y}_\infty = 1$ and the identity would require $1 + 1 =
2(1) + 2(1) = 4$: false. A norm violating the parallelogram
identity comes from no inner product.
\end{solution}

\begin{solution}{exo:b1:euclid:3}
$e_1 = \frac{1}{\sqrt 2}(1,1,0)$.

$w_2 = (1,0,1) - \langle (1,0,1), e_1\rangle e_1 = (1,0,1) -
\frac{1}{2}(1,1,0) = \bigl(\tfrac12, -\tfrac12, 1\bigr)$;
$\norm{w_2} = \sqrt{\tfrac32}$, so $e_2 =
\frac{1}{\sqrt 6}(1, -1, 2)$.

$w_3 = (0,1,1) - \langle\cdot, e_1\rangle e_1 - \langle\cdot,
e_2\rangle e_2 = (0,1,1) - \frac12 (1,1,0) - \frac16(1,-1,2) =
\bigl(-\tfrac23, \tfrac23, \tfrac23\bigr)$; normalizing, $e_3 =
\frac{1}{\sqrt 3}(-1, 1, 1)$.

(Check: pairwise products vanish; each has norm $1$.)
\end{solution}

\begin{solution}{exo:b1:euclid:4}
$F^\perp = \{(x,y,z) : x + y + z = 0\}$, basis $\bigl((1,-1,0),
(1,0,-1)\bigr)$. With $e = \frac{1}{\sqrt3}(1,1,1)$: $p_F(x) =
\langle x, e\rangle e = \frac{x_1 + x_2 + x_3}{3}\,(1,1,1)$, so
\[
\operatorname{Mat}(p_F) = \frac13
\begin{pmatrix} 1&1&1\\ 1&1&1\\ 1&1&1\end{pmatrix}.
\]
For $x = (1,2,3)$: $p_F(x) = (2,2,2)$ and $d(x, F) = \norm{(1,2,3)
- (2,2,2)} = \norm{(-1,0,1)} = \sqrt 2$.
\end{solution}

\begin{solution}{exo:b1:euclid:5}
Let $p = \alpha(1,0,1) + \beta(0,1,1) = (\alpha, \beta, \alpha +
\beta)$. Orthogonality of $x - p$ to the generators:
\[
\langle x - p, (1,0,1)\rangle = (1 - \alpha) + (4 - \alpha - \beta)
= 0,
\qquad
\langle x - p, (0,1,1)\rangle = (1 - \beta) + (4 - \alpha - \beta)
= 0,
\]
i.e.\ $2\alpha + \beta = 5$ and $\alpha + 2\beta = 5$: $\alpha =
\beta = \frac53$. So $p_F(x) = \bigl(\frac53, \frac53,
\frac{10}{3}\bigr)$ and
\[
d(x, F) = \norm{x - p} =
\norm{\bigl(-\tfrac23, -\tfrac23, \tfrac23\bigr)} =
\frac{2}{\sqrt 3} .
\]
Gram--Schmidt is unnecessary because the defining property of the
projection --- $x - p \perp F$ --- is itself a linear system
(``normal equations'') on the coefficients in \emph{any} generating
family.
\end{solution}

\begin{solution}{exo:b1:euclid:6}
Cauchy--Schwarz with $g = 1$:
\[
\Bigl(\int_0^1 f \cdot 1\Bigr)^2 \leq \int_0^1 f^2 \cdot \int_0^1
1^2 = \int_0^1 f^2 ,
\]
equality iff $f$ is proportional to $1$, i.e.\ constant. Discrete
version: in $\R^n$ with $x = (a_1, \dots, a_n)$, $y = (1, \dots,
1)$:
$\bigl(\sum a_i\bigr)^2 \leq \norm x^2\, \norm y^2 = n \sum a_i^2$,
equality iff all $a_i$ equal.
\end{solution}

\begin{solution}{exo:b1:euclid:7}
By \cref{thm:b1:euclid:projection}, $E = F \oplus F^\perp$, so
dimensions add: $\dim F^\perp = \dim E - \dim F$. Inclusion $F
\subseteq (F^\perp)^\perp$ is immediate (vectors of $F$ are
orthogonal to everything in $F^\perp$). Dimensions:
$\dim(F^\perp)^\perp = \dim E - \dim F^\perp = \dim F$; an
inclusion with equal (finite) dimensions is an equality
(\cref{thm:b1:findim:subspaces}).
\end{solution}

\begin{solution}{exo:b1:euclid:8}
$A$: columns unit, determinant $+1$: a rotation, with $\cos\theta =
\frac{1}{\sqrt2}$, $\sin\theta = \frac{1}{\sqrt2}$: $\theta =
\frac\pi4$. Hence $A^8 = R_{8\pi/4} = R_{2\pi} = I$.

$B$: determinant $\frac{1}{25}(-9 - 16) = -1$: a reflection
$S_\theta$ with $\cos\theta = \frac35$, $\sin\theta = \frac45$; its
axis makes the angle $\frac\theta2$ with the $x$-axis, i.e.\ the
line directed by $\bigl(\cos\frac\theta2, \sin\frac\theta2\bigr)$
--- concretely the axis is
$\operatorname{Vect}\bigl((2, 1)\bigr)$, since $B(2,1)^{\mathsf T} =
\frac15(6+4, 8-3)^{\mathsf T} = (2,1)^{\mathsf T}$. As a reflection,
$B^2 = I$.
\end{solution}

\begin{solution}{exo:b1:euclid:9}
The quantity is $\norm{f - (a + bX)}^2$ in $C(\intcc{0}{1})$ with
the integral inner product: minimal exactly at the orthogonal
projection $a + bX = p_F(f) = X - \frac16$
(\cref{ex:b1:euclid:bestapprox}). The minimum is
\[
\norm{f - p_F(f)}^2 = \int_0^1 \Bigl(x^2 - x + \frac16\Bigr)^{\!2}
\dd x .
\]
Expand: $\int_0^1 (x^2 - x + \frac16)^2 = \int_0^1 \bigl(x^4 - 2x^3
+ \frac{4}{3}x^2 - \frac13 x + \frac{1}{36}\bigr)\dd x = \frac15 -
\frac12 + \frac49 - \frac16 + \frac{1}{36} = \frac{1}{180}$.
So the minimum equals $\dfrac{1}{180}$.
\end{solution}

\begin{solution}{exo:b1:euclid:10}
Let $y \in \ker(u - \mathrm{id})$ (i.e.\ $u(y) = y$) and $x \in E$.
Then, using preservation of the inner product ($\langle u(a),
u(b)\rangle = \langle a, b\rangle$):
\[
\langle x - u(x),\, y\rangle
= \langle x, y\rangle - \langle u(x), y\rangle
= \langle x, y\rangle - \langle u(x), u(y)\rangle
= \langle x, y\rangle - \langle x, y\rangle = 0 :
\]
every vector of $\operatorname{im}(u - \mathrm{id})$ is orthogonal
to every fixed vector.

Therefore $\operatorname{im}(u - \mathrm{id}) \subseteq \ker(u -
\mathrm{id})^{\perp}$, and by rank--nullity plus
\cref{exo:b1:euclid:7}, both have dimension $\dim E - \dim\ker(u -
\mathrm{id})$: they are equal. Then $E = \ker(u - \mathrm{id})
\oplus \ker(u - \mathrm{id})^\perp = \ker(u - \mathrm{id}) \oplus
\operatorname{im}(u - \mathrm{id})$.
\end{solution}

\section*{Chapter \ref{ch:b1:curves} --- Plane Curves}

\begin{solution}{exo:b1:curves:1}
$M(-t) = (t^2, -t^3)$: reflection in the $x$-axis; study $t \geq 0$.
Both $x' = 2t$ and $y' = 3t^2$ are $\geq 0$: the branch moves
right and up, from $(0,0)$ to infinity. At $t = 0$ the velocity
vanishes; expansions $x = t^2$, $y = t^3$ show the tangent is the
$x$-axis ($y/x = t \to 0$) with $y$ changing sign while $x \geq 0$:
a \emph{cusp} pointing left. The curve is the semi-cubical parabola
$y^2 = x^3$.
\end{solution}

\begin{solution}{exo:b1:curves:2}
$M(t + 2\pi) = M(t) + (2\pi, 0)$: the curve repeats, translated by
one wheel circumference; study one period. $x'(t) = 1 - \cos t \geq
0$, $y'(t) = \sin t$: the arch rises on $\intoo{0}{\pi}$, falls on
$\intoo{\pi}{2\pi}$, culminating at $(\pi, 2)$. Singular points
where $x' = y' = 0$: $t \in 2\pi\Z$, on the ground. Near $t = 0$:
\[
x(t) = \frac{t^3}{6} + o(t^3),
\qquad
y(t) = \frac{t^2}{2} + o(t^2):
\]
$x$ changes sign, $y \geq 0$, and $\frac{x}{y^{3/2}}$ bounded: the
tangent is \emph{vertical} (direction $(0,1)$), a cusp where the
tracked point momentarily has zero speed --- the physical signature
of rolling without slipping.
\end{solution}

\begin{solution}{exo:b1:curves:3}
At $t = \frac\pi4$: $M = \bigl(\frac{\sqrt2}{4},
\frac{\sqrt2}{4}\bigr)$, tangent direction $(-\cos t, \sin t) =
\frac{1}{\sqrt2}(-1, 1)$ (\cref{ex:b1:curves:astroid}). Tangent
line: $y - \frac{\sqrt2}{4} = -(x - \frac{\sqrt2}{4})$, i.e.\ $x +
y = \frac{\sqrt2}{2}$. Intercepts: $\bigl(\frac{\sqrt2}{2},
0\bigr)$ and $\bigl(0, \frac{\sqrt2}{2}\bigr)$; the segment between
them has length $\sqrt{\frac12 + \frac12} = 1$.

(General point $t$: the tangent at $(\cos^3 t, \sin^3 t)$ cuts the
axes at $(\cos t, 0)$ and $(0, \sin t)$ --- check that the line
through these points has direction $(-\cos t, \sin t)$ and passes
through $M(t)$ --- and the cut segment has length $\sqrt{\cos^2 t +
\sin^2 t} = 1$.)
\end{solution}

\begin{solution}{exo:b1:curves:4}
$r = \cos 2\theta$: period $\pi$ in $\theta$, symmetric in both
axes; $r$ vanishes at $\theta = \pm\frac\pi4$ (tangents at the
origin along the diagonals) and is negative for $\theta \in
\intoo{\frac\pi4}{\frac{3\pi}{4}}$, where the points plot on the
opposite ray --- producing four petals along the axes directions
$\theta = 0, \frac\pi2, \pi, \frac{3\pi}{2}$, each of maximal
radius $1$.

$r\cos\theta = 1$ is the equation $x = 1$: the polar curve $r =
\frac{1}{\cos\theta}$ is the vertical line $x = 1$ (swept once for
$\theta \in \intoo{-\frac\pi2}{\frac\pi2}$).
\end{solution}

\begin{solution}{exo:b1:curves:5}
$r(\theta) = 1 + 2\cos\theta$ vanishes for $\cos\theta = -\frac12$:
$\theta = \pm\frac{2\pi}{3}$. Symmetry in the $x$-axis; study
$\theta \in \intcc{0}{\pi}$. For $\theta \in
\intco{0}{\frac{2\pi}{3}}$: $r > 0$, decreasing from $3$ to $0$:
outer arc, entering the origin tangent to the ray $\theta =
\frac{2\pi}{3}$. For $\theta \in \intoc{\frac{2\pi}{3}}{\pi}$: $r <
0$: the points $r\vec u(\theta)$ sit on the opposite ray (angle
$\theta - \pi \in \intoc{-\frac\pi3}{0}$), with distance $\abs r$
growing from $0$ to $1$: this draws a small \emph{inner loop}
through the origin. At $\theta = \pi$, $r = -1$ and the point is
$-\vec u(\pi) = (1, 0)$: the loop closes on the positive $x$-axis.
Sketch: a big heart-like outer curve
with maximum reach $3$ at $\theta = 0$, plus a loop inside it
through $O$ and $(1,0)$.
\end{solution}

\begin{solution}{exo:b1:curves:6}
Complete squares:
\[
9(x^2 - 4x) + 25(y^2 - 2y) = 164
\iff 9(x-2)^2 + 25(y-1)^2 = 164 + 36 + 25 = 225 .
\]
Dividing by $225$: $\frac{(x-2)^2}{25} + \frac{(y-1)^2}{9} = 1$: an
ellipse of center $(2, 1)$, semi-axes $a = 5$ (horizontal), $b =
3$; $c = \sqrt{25 - 9} = 4$: foci $(2 \pm 4,\, 1) = (-2, 1)$ and
$(6, 1)$; eccentricity $e = \frac45$.
\end{solution}

\begin{solution}{exo:b1:curves:7}
Focus at the origin, directrix $D: x = d$ with $d = \frac pe > 0$.
For a point $M = (r\cos\theta, r\sin\theta)$ with $r > 0$:
\[
MF = r,
\qquad
d(M, D) = \abs{d - r\cos\theta},
\]
and the conic condition $MF = e\,d(M, D)$, in the regime where $M$
is on the focus side of the directrix ($r\cos\theta < d$), reads $r
= e(d - r\cos\theta)$, i.e.
\[
r(1 + e\cos\theta) = ed = p,
\qquad
r = \frac{p}{1 + e\cos\theta} .
\]
For $e < 1$ the denominator never vanishes: all $\theta$ allowed
(ellipse). For $e = 1$: $\theta \neq \pi$ (parabola, open towards
the directrix's far side). For $e > 1$: need $1 + e\cos\theta > 0$,
i.e.\ $\theta \in \intoo{-\theta_0}{\theta_0}$ with $\theta_0 =
\arccos\bigl(-\frac1e\bigr)$: one \emph{branch} of the hyperbola
(the other branch corresponds to the sign choice $r < 0$, or to the
second focus).
\end{solution}

\begin{solution}{exo:b1:curves:8}
\emph{Gardener:} attach a string of length $2a$ to two stakes $F,
F'$ and keep it taut with the tracing point $M$: the constraint is
exactly $MF + MF' = 2a$, so the traced curve is the ellipse.

\emph{Reflection property:} let $t \mapsto M(t)$ be a regular
parametrization and $u(t) = \frac{M(t) - F}{\norm{M(t) - F}}$,
$u'(t)$ the analogous unit vector to $F'$. Differentiating
$\norm{M - F} + \norm{M - F'} = 2a$, using $\frac{\dd}{\dd
t}\norm{M - F} = \bigl\langle \frac{M - F}{\norm{M-F}},\,
M'\bigr\rangle$ (chain rule on $\sqrt{\langle\cdot,\cdot\rangle}$):
\[
\bigl\langle u + u',\, M'\bigr\rangle = 0 .
\]
So the tangent direction $M'$ is orthogonal to the bisector
direction $u + u'$ of the two focal rays: the tangent makes equal
angles with $MF$ and $MF'$. (A light ray from one focus reflects
off the ellipse to the other focus.)
\end{solution}

\begin{solution}{exo:b1:curves:9}
\begin{enumerate}
  \item Domain: $\cos 2\theta \geq 0$: $\theta \in
        \intcc{-\frac\pi4}{\frac\pi4} \cup
        \intcc{\frac{3\pi}{4}}{\frac{5\pi}{4}}$. Symmetries: $\theta
        \mapsto -\theta$ ($x$-axis) and $\theta \mapsto \pi -
        \theta$ ($y$-axis): study $\intcc{0}{\frac\pi4}$; $r$
        decreases from $1$ to $0$. At $\theta = \pm\frac\pi4$: $r =
        0$, tangents at the origin along the diagonals. The curve is
        the $\infty$ symbol: two symmetric loops meeting at $O$,
        reaching $(\pm 1, 0)$.
  \item With $F, F' = (\pm c, 0)$, $c = \frac{1}{\sqrt 2}$, and $M
        = (r\cos\theta, r\sin\theta)$:
        \[
        MF^2\, MF'^2
        = \bigl((r\cos\theta - c)^2 + r^2\sin^2\theta\bigr)
        \bigl((r\cos\theta + c)^2 + r^2\sin^2\theta\bigr)
        = (r^2 + c^2)^2 - (2rc\cos\theta)^2 ,
        \]
        by the identity $(A - B)(A + B) = A^2 - B^2$ with $A = r^2
        + c^2$, $B = 2rc\cos\theta$. With $c^2 = \frac12$:
        \[
        MF^2 MF'^2 = \Bigl(r^2 + \frac12\Bigr)^2 - 2r^2\cos^2\theta
        = r^4 + r^2\bigl(1 - 2\cos^2\theta\bigr) + \frac14
        = r^4 - r^2\cos 2\theta + \frac14 .
        \]
        On the lemniscate, $r^2 = \cos 2\theta$: the first two terms
        cancel, leaving $MF^2MF'^2 = \frac14$, i.e.\ $MF \cdot MF' =
        \frac12$. Conversely, the computation read backwards shows
        the locus equation $MF\,MF' = \frac12$ is $r^2 = \cos
        2\theta$ (for $r \neq 0$; and $O$ satisfies both).
\end{enumerate}
\end{solution}

\begin{solution}{exo:b1:curves:10}
$r = 1 + \cos\theta$, $r' = -\sin\theta$:
\[
r^2 + r'^2 = 1 + 2\cos\theta + \cos^2\theta + \sin^2\theta
= 2 + 2\cos\theta = 4\cos^2\frac\theta2 ,
\]
so $\sqrt{r^2 + r'^2} = 2\abs{\cos\frac\theta2}$ and, by the
symmetry in the $x$-axis,
\[
L = \int_0^{2\pi} 2\Bigl|\cos\frac\theta2\Bigr|\,\dd\theta
= 2\int_0^{\pi} 2\cos\frac\theta2\,\dd\theta
= 2\Bigl[4\sin\frac\theta2\Bigr]_0^{\pi} = 8 .
\]
Another algebraic, $\pi$-free perimeter.
\end{solution}

\begin{solution}{exo:b1:curves:11}
The point $(a\cos t, b\sin t)$ satisfies the equation:
$\cos^2 t + \sin^2 t = 1$. The line's normal vector is
$\bigl(\frac{\cos t}a, \frac{\sin t}b\bigr)$, and its product
with the velocity $(-a\sin t, b\cos t)$ is $-\sin t\cos t +
\sin t\cos t = 0$: the line passes through the point with the
tangent direction --- it is the tangent. For $(x_0, y_0)$ on the
ellipse, write $\cos t = \frac{x_0}a$, $\sin t = \frac{y_0}b$ and
substitute:
\[
\frac{x\,x_0}{a^2} + \frac{y\,y_0}{b^2} = 1 ,
\]
obtained from the ellipse equation by ``splitting'' $x^2 \mapsto
x\,x_0$ and $y^2 \mapsto y\,y_0$.
\end{solution}

\begin{solution}{exo:b1:curves:12}
\begin{enumerate}
  \item $r' = k\eu^{k\theta}$, so $\tan V = \frac{r}{r'} =
        \frac1k$: constant. The spiral cuts every ray from the
        origin at the same angle $V = \arctan\frac1k$.
  \item In complex notation the spiral is $\{\eu^{k\theta}
        \eu^{\iu\theta} : \theta \in \R\}$. Rotating by $c$
        multiplies by $\eu^{\iu c}$:
        \[
        \eu^{k\theta}\eu^{\iu(\theta + c)}
        = \eu^{-kc}\,\eu^{k(\theta + c)}\eu^{\iu(\theta+c)} ,
        \]
        and as $\theta + c$ runs over $\R$ this describes
        $\eu^{-kc}$ times the spiral: rotation $=$ scaling. No
        other smooth curve but lines and circles has this
        property.
  \item $\sqrt{r^2 + r'^2} = \sqrt{1 + k^2}\,\eu^{k\theta}$, so
        on $\intcc{A}{\theta_0}$ the length is
        $\frac{\sqrt{1+k^2}}{k}\bigl(\eu^{k\theta_0} -
        \eu^{kA}\bigr)$, and as $A \to -\infty$:
        \[
        L = \frac{\sqrt{1 + k^2}}{k}\,\eu^{k\theta_0} < \infty :
        \]
        infinitely many turns around the origin, finite total
        length (the turns shrink geometrically).
\end{enumerate}
\end{solution}

\begin{solution}{pb:b1:curves:1}
\textbf{1.} Rolling without slipping means the contact point has
traveled a distance equal to the arc of wheel unrolled: after
turning by $t$, the center is at $(Rt, R)$. The marked point sits
on the rim at angle $t$ \emph{behind} the downward vertical
(the wheel turns clockwise while advancing):
\[
M(t) = \Omega(t) + R(-\sin t, -\cos t)
= \bigl(R(t - \sin t),\ R(1 - \cos t)\bigr),
\]
which is correct at $t = 0$ ($M = (0,0)$) and at $t = \pi$ ($M =
(\pi R, 2R)$, the top).

\textbf{2.} $f'(t) = R(1 - \cos t,\ \sin t)$ and
\[
\norm{f'}^2 = R^2\bigl((1-\cos t)^2 + \sin^2 t\bigr)
= 2R^2(1 - \cos t) = 4R^2\sin^2\frac t2 ,
\]
so $\norm{f'} = 2R\abs{\sin\frac t2}$. Singular points at $t \in
2\pi\Z$: the cusps on the ground (\cref{exo:b1:curves:2}); the
top of the arch is $t = \pi$, where the speed $2R$ is maximal.

\textbf{3.} Half-angle factorizations:
\[
f'(t) = 2R\sin\frac t2\,\Bigl(\sin\frac t2,\ \cos\frac t2\Bigr),
\quad
M - C = 2R\sin\frac t2\,\Bigl(-\cos\frac t2,\ \sin\frac t2\Bigr),
\]
\[
T - M = \bigl(R\sin t,\ R(1 + \cos t)\bigr)
= 2R\cos\frac t2\,\Bigl(\sin\frac t2,\ \cos\frac t2\Bigr).
\]
For $0 < t < 2\pi$, $\sin\frac t2 \neq 0$: $M - C$ is orthogonal
to the tangent direction $\bigl(\sin\frac t2, \cos\frac
t2\bigr)$ (their product is $-\sin\frac t2\cos\frac t2 +
\sin\frac t2\cos\frac t2 = 0$), and $T - M$ is parallel to it.
The chord to the top of the wheel \emph{is} the tangent; the
chord to the contact point is the normal.

\textbf{4.} At each instant the wheel pivots about its contact
point (that point has zero velocity: rolling without slipping),
so every rigid point of the wheel moves orthogonally to the line
joining it to $C$, with speed (angular velocity $1$) equal to
that distance. Check: $\norm{M - C} = 2R\abs{\sin\frac t2} =
\norm{f'(t)}$.

\textbf{5.} $2R\,y(t) = 2R\cdot 2R\sin^2\frac t2 =
4R^2\sin^2\frac t2 = \norm{f'(t)}^2$. The speed at height $y$ is
$\sqrt{2R\,y}$ --- formally the law $v = \sqrt{2gh}$ of free
fall, with the ground playing the ceiling; Part IV turns this
observation into clockwork.

\textbf{6.} By \cref{def:b1:curves:arclength} and question 2
($\sin\frac t2 \geq 0$ on $\intcc{0}{2\pi}$):
\[
L = \int_0^{2\pi} 2R\sin\frac t2\,\dd t
= 2R\Bigl[-2\cos\frac t2\Bigr]_0^{2\pi} = 2R(2 + 2) = 8R .
\]
Wren's theorem: four diameters exactly.

\textbf{7.} $s(t) = \int_0^t 2R\sin\frac u2\,\dd u =
4R\bigl(1 - \cos\frac t2\bigr)$; $s(2\pi) = 4R(1+1) = 8R$,
consistent.

\textbf{8.} $\sigma(t) = \abs{s(t) - s(\pi)} = \abs{4R(1 -
\cos\frac t2) - 4R} = 4R\abs{\cos\frac t2}$, so $\sigma^2 =
16R^2\cos^2\frac t2$; and $2R - y = R(1 + \cos t) =
2R\cos^2\frac t2$, whence $8R(2R - y) = 16R^2\cos^2\frac t2 =
\sigma^2$.

\textbf{9.} At the cusp $t = 0$ (or $2\pi$): $\sigma = 4R$, half
of $8R$: the apex halves the arch. In both computations the
speed is $\abs{\text{trigonometric polynomial in } t/2}$, whose
antiderivative is again trigonometric: the length is a
difference of \emph{values} of cosines --- rational numbers
times $R$ --- with no arc of circle to measure, hence no $\pi$.
The circle itself has constant speed, so its length integral
produces the full interval length $2\pi$; the cycloid's speed
vanishes at the ends and integrates algebraically.

\textbf{10.} The arch is swept with $x$ increasing from $0$ to
$2\pi R$ ($x' = R(1 - \cos t) \geq 0$, vanishing only at
isolated points). Substituting $x = x(t)$ in the area integral
$\int_0^{2\pi R} y\,\dd x$ (\cref{thm:b1:integration:parts})
gives $A = \int_0^{2\pi} y(t)\,x'(t)\,\dd t$.

\textbf{11.}
\[
A = \int_0^{2\pi} R(1 - \cos t)\cdot R(1 - \cos t)\,\dd t
= R^2\int_0^{2\pi}\bigl(1 - 2\cos t + \cos^2 t\bigr)\dd t
= R^2\Bigl(2\pi - 0 + \pi\Bigr) = 3\pi R^2 ,
\]
using $\int_0^{2\pi}\cos^2 = \pi$. Exactly three wheel areas:
Galileo's balance said ``about $3$''; Roberval's computation
says ``exactly''.

\textbf{12.} With $x = \cos^3 t$, $y = \sin^3 t$: $y\,x' =
-3\sin^4 t\cos^2 t$. Linearize:
\[
\sin^4 t\cos^2 t = \frac{1 - \cos 2t}{2}\cdot\frac{\sin^2
2t}{4} = \frac{\sin^2 2t}{8} - \frac{\sin^2 2t\cos 2t}{8} ,
\]
and over $\intcc{0}{2\pi}$: $\int\sin^2 2t = \pi$, $\int \sin^2
2t\cos 2t = \bigl[\frac{\sin^3 2t}{6}\bigr] = 0$. So $\oint
y\,\dd x = -3\cdot\frac\pi8$, and the enclosed area is
$\frac{3\pi}8$ (the sign records the counterclockwise
orientation).

\textbf{13.} For $(\cos t, \sin t)$ with $t$ from $\pi$ to $0$,
$x$ increases from $-1$ to $1$ and
\[
\int_\pi^0 \sin t\cdot(-\sin t)\,\dd t
= \int_0^\pi \sin^2 t\,\dd t = \frac\pi2 :
\]
the formula returns the area of the upper half-disc, as it must.

\textbf{14.} $x' = R(1 + \cos t) = 2R\cos^2\frac t2$, $y' =
R\sin t = 2R\sin\frac t2\cos\frac t2$, so $\norm{f'} =
2R\cos\frac t2$ (nonnegative on $\intcc{-\pi}{\pi}$). Arc from
the bottom: $s(t) = \int_0^t 2R\cos\frac u2\,\dd u =
4R\sin\frac t2$. Then
\[
y = R(1 - \cos t) = 2R\sin^2\frac t2
= 2R\Bigl(\frac{s}{4R}\Bigr)^{\!2} = \frac{s^2}{8R} .
\]

\textbf{15.} Energy conservation with $v = \frac{\dd s}{\dd\tau}$
and $y = \frac{s^2}{8R}$, $y_0 = \frac{s_0^2}{8R}$:
\[
\Bigl(\frac{\dd s}{\dd\tau}\Bigr)^{\!2} = 2g(y_0 - y)
= \frac{2g}{8R}\bigl(s_0^2 - s^2\bigr)
= \frac{g}{4R}\bigl(s_0^2 - s^2\bigr).
\]

\textbf{16.} Differentiating in $\tau$: $2s's'' =
-\frac{g}{4R}\,2ss'$, so wherever $s' \neq 0$ (hence everywhere
by continuity), $s'' = -\frac{g}{4R}s$: the harmonic oscillator.
By \cref{thm:b1:diffeq:homogeneous2}, $s(\tau) =
A\cos\omega\tau + B\sin\omega\tau$ with $\omega =
\sqrt{g/(4R)}$; the initial conditions $s(0) = s_0$, $s'(0) = 0$
give $s(\tau) = s_0\cos\omega\tau$.

\textbf{17.} The bead reaches the bottom when $s = 0$, i.e.\ at
$\omega\tau = \frac\pi2$:
\[
T_{\downarrow} = \frac{\pi}{2\omega}
= \frac\pi2\sqrt{\frac{4R}{g}} = \pi\sqrt{\frac Rg}\,,
\]
independent of $s_0$: released from anywhere in the bowl, all
beads arrive at the same instant --- the tautochrone.

\textbf{18.} Substituting $s = s_0\cos\omega\tau$: the left side
is $s_0^2\omega^2\sin^2\omega\tau$ and the right side is
$\frac{g}{4R}s_0^2(1 - \cos^2\omega\tau) =
s_0^2\omega^2\sin^2\omega\tau$: exact equality, so no spurious
solution was introduced. For a circular arc, $y$ is not
proportional to $s^2$ ($y = \ell(1 - \cos\frac s\ell) =
\frac{s^2}{2\ell} - \frac{s^4}{24\ell^3} + \dots$): the
restoring term is only \emph{approximately} linear, so the
period of a circular pendulum drifts with amplitude, while the
cycloid's is rigorously constant.

\textbf{19.} $\pi\sqrt{R/g} = 1$ gives $R = \frac{g}{\pi^2}
\approx \frac{9.81}{9.87} \approx 0.994$ m --- almost exactly
one meter. The full oscillation takes $\frac{2\pi}\omega =
2\pi\sqrt{4R/g}$, which is precisely the small-angle formula
$2\pi\sqrt{\ell/g}$ for a pendulum of length $\ell = 4R$:
Huygens suspended his pendulum from cycloidal cheeks of exactly
that proportion.

\textbf{20.} At $t = \frac\pi2$, $R = 1$: $M = \bigl(\frac\pi2 -
1,\ 1\bigr)$, $T = \bigl(\frac\pi2,\ 2\bigr)$, so $T - M = (1,
1)$. And $f'(\frac\pi2) = (1 - 0,\ 1) = (1, 1)$: the chord to
the top is the velocity, as promised.

\textbf{21.} For the trochoid, $x'(t) = R - d\cos t$ and $y'(t)
= d\sin t$. If $d < R$: $x' \geq R - d > 0$, the point always
advances; the velocity never vanishes (its first component is
positive): no singular point, a regular wave. If $d > R$: $x'(0)
= R - d < 0 < R + d = x'(\pi)$, so the point moves backwards
near the contact instants and forwards elsewhere: the curve
crosses itself in loops. A point on the flange of a railway
wheel (below the rail head, $d > R$) travels backwards at every
turn.

\textbf{22.} The chord from the cusp $(\pi R, 2R)$ to the origin
has length $L = R\sqrt{\pi^2 + 4}$ and slope angle $\alpha$ with
$\sin\alpha = \frac{2R}{L}$. Sliding from rest with constant
acceleration $g\sin\alpha$: $L = \frac12 g\sin\alpha\,T^2$, so
\[
T = \sqrt{\frac{2L}{g\sin\alpha}} = \sqrt{\frac{2L^2}{2Rg}}
= \frac{L}{\sqrt{Rg}} = \sqrt{\pi^2 + 4}\,\sqrt{\frac Rg}
\approx 3.72\sqrt{\frac Rg}\,,
\]
against $\pi\sqrt{R/g} \approx 3.14\sqrt{R/g}$ for the cycloid:
the curved path is faster. It is in fact the fastest possible
--- the brachistochrone --- a theorem of the calculus of
variations, beyond this volume.

\textbf{23.} $\dfrac{\dd y}{\dd x} = \dfrac{y'}{x'} =
\dfrac{\sin t}{1 - \cos t} = \cot\dfrac t2$ (half-angle
formulas). Then
\[
\frac{\dd^2y}{\dd x^2}
= \frac{\dd}{\dd t}\Bigl(\cot\frac t2\Bigr)\cdot\frac{1}{x'(t)}
= -\frac{1}{2\sin^2\frac t2}\cdot\frac{1}{2R\sin^2\frac t2}
= -\frac{1}{4R\sin^4\frac t2} < 0 :
\]
concave throughout the arch; as $t \to 0^+$ or $2\pi^-$ the
slope $\cot\frac t2 \to \pm\infty$: vertical tangents at the
cusps.

\textbf{24.} The chord direction is $T - M \parallel
\bigl(\sin\frac t2, \cos\frac t2\bigr)$, which tends to $(0, 1)$
as $t \to 0^{+}$: the tangent at the cusp is vertical ---
recovered from pure geometry, with no Taylor expansion.

\textbf{25.} (i) Wren's theorem, $L = 8R$; Roberval's area, $A =
3\pi R^2$ (Galileo's conjectured ratio $3$); Huygens'
tautochrone, $T_\downarrow = \pi\sqrt{R/g}$. (ii) Every
computation ran on the half-angle factorizations $f' =
2R\sin\frac t2(\sin\frac t2, \cos\frac t2)$ and their bowl
analogue --- one identity powering tangent, length and clock
alike. (iii) The intrinsic relation $y = s^2/8R$ converted the
energy equation into $s'' = -\frac{g}{4R}s$, a linear equation
with constant coefficients whose solutions are exactly
isochronous. (iv) Part I used the differential calculus of
\cref{ch:b1:derivative}, Part II--III the integral of
\cref{ch:b1:integration}, Part IV the differential equations of
\cref{ch:b1:diffeq} --- the cycloid is this book's curriculum
rolled into one curve.
\end{solution}

\section*{Chapter \ref{ch:b1:multivar} --- Functions of Two Variables}

\begin{solution}{exo:b1:multivar:1}
$\dfrac{\partial f}{\partial x} = 2xy + y\,\eu^{xy}$,
$\dfrac{\partial f}{\partial y} = x^2 + x\,\eu^{xy}$.

$\dfrac{\partial g}{\partial x} = \dfrac{2x}{x^2+y^2}$,
$\dfrac{\partial g}{\partial y} = \dfrac{2y}{x^2+y^2}$.

$\dfrac{\partial h}{\partial x} = \dfrac{-y/x^2}{1 + y^2/x^2} =
\dfrac{-y}{x^2+y^2}$,
$\dfrac{\partial h}{\partial y} = \dfrac{x}{x^2+y^2}$.
\end{solution}

\begin{solution}{exo:b1:multivar:2}
In polar \omterm{prop:b1:vspaces:coordinates}{coordinates} ($\rho \to 0$):

$\abs{f} = \dfrac{\rho^3\abs{\cos^2\theta\sin\theta}}{\rho^2} \leq
\rho \to 0$: \omterm{def:b1:continuity:continuous}{continuous}.

$g = \cos\theta\sin\theta$: independent of $\rho$, taking different
values along different rays (cf.\
\cref{ex:b1:multivar:trap}): no limit, not \omterm{def:b1:continuity:continuous}{continuous}.

$\abs h \leq \dfrac{\rho^3(\abs{\cos^3\theta} +
\abs{\sin^3\theta})}{\rho^2} \leq 2\rho \to 0$: \omterm{def:b1:continuity:continuous}{continuous}.
\end{solution}

\begin{solution}{exo:b1:multivar:3}
Along $y = mx$: $f(x, mx) = \dfrac{m x^3}{x^4 + m^2 x^2} =
\dfrac{mx}{x^2 + m^2} \to 0$ (for $m \neq 0$; along $y = 0$ and the
$y$-axis, $f = 0$). So every straight-line limit is $0$. But on the
parabola $y = x^2$:
\[
f(x, x^2) = \frac{x^4}{x^4 + x^4} = \frac12 :
\]
the sequence $\bigl(\frac1n, \frac{1}{n^2}\bigr) \to (0,0)$ has
$f \to \frac12 \neq 0$. Not \omterm{def:b1:continuity:continuous}{continuous}: lines do not suffice to
test two-variable limits.
\end{solution}

\begin{solution}{exo:b1:multivar:4}
$\frac{\partial f}{\partial x} = 3x^2y^2 + y\cos(xy)$; then
\[
\frac{\partial^2 f}{\partial y\,\partial x}
= 6x^2 y + \cos(xy) - xy\sin(xy) .
\]
$\frac{\partial f}{\partial y} = 2x^3 y + x\cos(xy)$; then
\[
\frac{\partial^2 f}{\partial x\,\partial y}
= 6x^2 y + \cos(xy) - xy\sin(xy) :
\]
equal, as Schwarz promises.
\end{solution}

\begin{solution}{exo:b1:multivar:5}
By \cref{thm:b1:multivar:chain} with $(x, y) = (\cos t, \sin t)$:
\[
g'(t) = -\sin t\,\frac{\partial f}{\partial x}(\cos t, \sin t)
+ \cos t\,\frac{\partial f}{\partial y}(\cos t, \sin t)
= \bigl\langle \nabla f,\ (-\sin t, \cos t)\bigr\rangle .
\]
At an extremum of $g$, $g'(t) = 0$: $\nabla f$ is \omterm{def:b1:euclid:orthogonal}{orthogonal} to the
tangent vector $(-\sin t, \cos t)$ of the circle, hence parallel to
the radius vector $(\cos t, \sin t)$ (the plane \omterm{def:b1:euclid:orthogonal}{orthogonal
complement} of a unit vector is the line it \omterm{def:b1:vspaces:span}{spans}, taken
perpendicular). This is the simplest instance of a Lagrange
multiplier.
\end{solution}

\begin{solution}{exo:b1:multivar:6}
$f$: $\nabla f = (2x + y - 3,\; x + 2y) = 0$: $y = -\frac x2$ and
$2x - \frac x2 = 3$: $x = 2$, $y = -1$. Second order: $r = 2$, $s =
1$, $t = 2$: $rt - s^2 = 3 > 0$, $r > 0$: local (indeed global ---
quadratic) minimum at $(2, -1)$, value $f(2,-1) = -3$.

$g$: $\nabla g = (2x,\; -2y + 4) = 0$: point $(0, 2)$; $r = 2$, $s
= 0$, $t = -2$: $rt - s^2 = -4 < 0$: saddle.

$h$: $\nabla h = \bigl(4x^3 - 4(x - y),\; 4y^3 + 4(x-y)\bigr) = 0$.
Adding: $x^3 + y^3 = 0$, so $y = -x$; substituting: $4x^3 - 8x = 0$:
$x \in \{0, \pm\sqrt2\}$. \omterm{thm:b1:multivar:critical}{Critical points}: $(0,0)$,
$(\sqrt2, -\sqrt2)$, $(-\sqrt2, \sqrt2)$.
At $(\pm\sqrt2, \mp\sqrt2)$: $r = 12x^2 - 4 = 20$, $s = 4$, $t =
20$: $rt - s^2 > 0$, $r > 0$: local minima (value $h = 4 + 4 - 16 =
-8$). At $(0,0)$: $r = t = -4$, $s = 4$: $rt - s^2 = 0$: silent.
Examine: $h(x, x) = 2x^4 > 0$ and $h(x, -x) = 2x^4 - 8x^2 < 0$ for
small $x \neq 0$: both signs in every \omterm{def:b1:topology:open}{neighborhood} --- no extremum
at the origin.
\end{solution}

\begin{solution}{exo:b1:multivar:7}
$x^2 - y^2$: second partials $2$ and $-2$: sum $0$. $xy$: both pure
second partials vanish. $\eu^x\cos y$: $\frac{\partial^2}{\partial
x^2} = \eu^x\cos y$, $\frac{\partial^2}{\partial y^2} =
-\eu^x\cos y$: sum $0$. $\ln(x^2+y^2)$: from
\cref{exo:b1:multivar:1},
\[
\frac{\partial^2}{\partial x^2}\ln(x^2+y^2)
= \frac{2(x^2+y^2) - 4x^2}{(x^2+y^2)^2}
= \frac{2(y^2 - x^2)}{(x^2+y^2)^2},
\]
and the $y$-version is its opposite: sum $0$.
\end{solution}

\begin{solution}{exo:b1:multivar:8}
\emph{\omterm{def:b1:linmaps:projection}{Projection}:} the plane $P$ has normal $n = (1, 2, -1)$ and
passes through $A = (4, 0, 0)$. The closest point to the origin is
$O$ projected: $p = O + \frac{\langle A - O, n\rangle}{\norm
n^2}\,n = \frac{4}{6}(1,2,-1) = \bigl(\frac23, \frac43,
-\frac23\bigr)$, at distance $\frac{4}{\sqrt 6}$. (Formula: distance
from origin to plane $\langle X, n \rangle = c$ is
$\frac{\abs c}{\norm n}$ with $c = 4$.)

\emph{Minimization:} $f(x,y) = x^2 + y^2 + (x + 2y - 4)^2$. Writing
$w = x + 2y - 4$ for the last factor:
\[
\nabla f = \bigl(2x + 2w,\; 2y + 4w\bigr) = 0
\iff x = -w \text{ and } y = -2w .
\]
Substituting into the definition of $w$: $w = -w - 4w - 4$, so $w =
-\frac23$, giving $x = \frac23$, $y = \frac43$, and $z = w =
-\frac23$. Same point as by \omterm{def:b1:linmaps:projection}{projection}, at distance
$\sqrt{\frac49 + \frac{16}{9} + \frac49} = \frac{2\sqrt 6}{3} =
\frac{4}{\sqrt6}$; it is a minimum since $f \to \infty$ at infinity
(a positive-definite quadratic plus \omterm{def:b1:linmaps:def}{linear} terms).
\end{solution}

\begin{solution}{exo:b1:multivar:9}
With $\rho = \sqrt{x^2+y^2}$: $\frac{\partial \rho}{\partial x} =
\frac{x}{\rho}$, so
\[
\frac{\partial f}{\partial x} = \varphi'(\rho)\,\frac{x}{\rho},
\qquad
\frac{\partial^2 f}{\partial x^2}
= \varphi''(\rho)\,\frac{x^2}{\rho^2}
+ \varphi'(\rho)\,\Bigl(\frac{1}{\rho} -
\frac{x^2}{\rho^3}\Bigr),
\]
(quotient and chain rules). Adding the symmetric $y$-expression:
the $\varphi''$ terms collect $\frac{x^2 + y^2}{\rho^2} = 1$, the
$\varphi'$ terms collect $\frac{2}{\rho} - \frac{x^2 +
y^2}{\rho^3} = \frac{1}{\rho}$:
\[
\Delta f = \varphi'' + \frac{\varphi'}{\rho} .
\]
Radial harmonic: solve $\varphi'' + \frac{\varphi'}{\rho} = 0$: the
first-order equation $u' + \frac u\rho = 0$ for $u = \varphi'$
gives $u = \frac{c}{\rho}$ (\cref{thm:b1:diffeq:homogeneous1}),
then $\varphi = c\ln\rho + d$. The radial harmonic functions are
$c\,\ln\sqrt{x^2 + y^2} + d$ --- the logarithmic potential of
\cref{exo:b1:multivar:7} and the constants, nothing else.
\end{solution}

\begin{solution}{exo:b1:multivar:10}
For $z = xy$ at $(1,1)$: partials $y = 1$ and $x = 1$, tangent
plane $z = 1 + (x - 1) + (y - 1) = x + y - 1$. Horizontal
tangent plane means both partials vanish, i.e.\ $\nabla f = 0$:
by \cref{ex:b1:multivar:extrema}, exactly the \omterm{thm:b1:multivar:critical}{critical points}
$(0, 0)$ and $(1, 1)$, with horizontal planes $z = 0$ and $z =
-1$. ``Horizontal tangent plane'' and ``\omterm{thm:b1:multivar:critical}{critical point}'' are the
same notion, seen on the graph and in the formula.
\end{solution}

\begin{solution}{exo:b1:multivar:11}
\begin{enumerate}
  \item For fixed $y$, the one-variable function $x \mapsto
        f(x,y)$ has zero \omterm{def:b1:derivative:def}{derivative} on $\R$, hence is constant
        (\cref{thm:b1:derivative:mvt}): $f(x, y) = f(0, y)$ for
        all $x$: $f$ depends only on $y$.
  \item Let $g(u, v) = f\bigl(\frac{u+v}2, \frac{u-v}2\bigr)$.
        By the chain rule (\cref{thm:b1:multivar:chain}, applied
        in the variable $v$ with $u$ frozen):
        \[
        \frac{\partial g}{\partial v}
        = \frac12\,\frac{\partial f}{\partial x}
        - \frac12\,\frac{\partial f}{\partial y} = 0 .
        \]
        By (1), $g$ depends only on $u$: $g(u, v) = \varphi(u)$
        with $\varphi$ of class $C^1$, and reversing the change
        of variables ($u = x + y$, $v = x - y$),
        \[
        f(x, y) = \varphi(x + y).
        \]
        Conversely every such $f$ satisfies $f_x = f_y =
        \varphi'$: the solutions are exactly the $C^1$ functions
        of $x + y$.
\end{enumerate}
\end{solution}

\begin{solution}{exo:b1:multivar:12}
On the \omterm{def:b1:topology:open}{open} disc, an extremum would be critical: $\nabla f = (y,
x) = 0$ only at $(0,0)$, where $f = 0$; it is a saddle ($f(\pm
\varepsilon, \pm\varepsilon) = \varepsilon^2 > 0 > -\varepsilon^2
= f(\pm\varepsilon, \mp\varepsilon)$), so no extremum there. By
the admitted extreme value theorem the bounds are attained,
necessarily on the \omterm{def:b1:topology:closure}{boundary} circle. There, with
\cref{exo:b1:multivar:5},
\[
f(\cos t, \sin t) = \cos t\sin t = \frac{\sin 2t}{2}
\in \intcc{-\tfrac12}{\tfrac12},
\]
with maximum $\frac12$ at $t = \frac\pi4, \frac{5\pi}4$ (points
$\pm\frac{1}{\sqrt2}(1,1)$) and minimum $-\frac12$ at $t =
\frac{3\pi}4, \frac{7\pi}4$ (points $\pm\frac{1}{\sqrt2}(1,-1)$).
Global maximum $\frac12$, global minimum $-\frac12$.
\end{solution}

\begin{solution}{pb:b1:multivar:1}
\textbf{1.} Expanding the right-hand side:
\[
r\Bigl(h + \frac srk\Bigr)^{\!2} + \frac{rt - s^2}{r}k^2
= rh^2 + 2shk + \frac{s^2}{r}k^2 + \frac{rt - s^2}{r}k^2
= q(h,k) .
\]
If $rt - s^2 > 0$: both squares carry the factor sign of $r$,
and $q(h,k) = 0$ forces $k = 0$ then $h = 0$: strict sign of $r$
off the origin. If $rt - s^2 < 0$: $q(1, 0) = r$ while
$q\bigl(-\frac sr, 1\bigr) = \frac{rt - s^2}{r}$ has the
opposite sign: both signs occur.

\textbf{2.} If $r = 0 \neq t$, exchange the roles of $h$ and $k$
(the completed square in $k$), with the same conclusions; note
$rt - s^2 = -s^2 < 0$ as soon as $s \neq 0$, and indeed $q =
2shk + tk^2$ takes both signs then ($k$ small against $h$). If
$r = t = 0$: $q = 2shk$, both signs iff $s \neq 0$, i.e.\ iff
$rt - s^2 = -s^2 < 0$. Summary: $q$ takes both signs $\iff rt -
s^2 < 0$; $q$ vanishes only at the origin (definite) $\iff rt -
s^2 > 0$ (which forces $r \neq 0$, since $r = 0$ gives $q(1,0) =
0$).

\textbf{3.} With $f_x = rx + sy + \beta$ and $f_y = sx + ty +
\gamma$, direct expansion gives
\[
f(a + h, b + k) = f(a, b) + h\,f_x(a,b) + k\,f_y(a,b)
+ \tfrac12 q(h, k),
\]
with no higher terms (the function is a \omterm{def:b1:poly:def}{polynomial} of degree
$2$). At a \omterm{thm:b1:multivar:critical}{critical point} the \omterm{def:b1:linmaps:def}{linear} part vanishes: $f(X_0 + H)
- f(X_0) = \frac12 q(H)$ \emph{exactly}, so the sign study of
questions 1--2 classifies: strict global minimum ($rt - s^2 >
0$, $r > 0$), strict global maximum ($rt - s^2 > 0$, $r < 0$),
saddle --- both signs in every \omterm{def:b1:topology:open}{neighborhood} --- when $rt - s^2 <
0$. This proves \cref{met:b1:multivar:monge} for quadratic
functions.

\textbf{4.} The form $q - c(h^2 + k^2)$ has data $r - c$, $s$,
$t - c$. With $c = \frac{rt - s^2}{r + t}$ (note $t > 0$ since
$rt > s^2 \geq 0$ and $r > 0$):
\[
(r - c)(t - c) - s^2 = rt - s^2 - c(r + t) + c^2 = c^2 \geq 0,
\]
and $r - c \geq 0$ ($c \leq r \iff rt - s^2 \leq r^2 + rt$,
true). If $r - c > 0$, the completed square of question 1 shows
$q - c(h^2 + k^2) \geq 0$; if $r - c = 0$, then $s = 0$ (from
$-s^2 = -c^2 \leq \dots$ forcing the displayed quantity
$\geq 0$ with first factor $0$) and the form is $(t - c)k^2 \geq
0$. In both cases $q(h,k) \geq c(h^2 + k^2)$.

\textbf{5.} By questions 3--4, $f(X) = f(X_0) + \frac12 q(X -
X_0) \geq f(X_0) + \frac c2\norm{X - X_0}^2 \to +\infty$ as
$\norm X \to \infty$: $f$ is coercive, and the inequality is
strict for $X \neq X_0$: the \omterm{thm:b1:multivar:critical}{critical point} is the unique global
minimizer. For the cubic $f = x^3 + y^3 - 3xy$ of
\cref{ex:b1:multivar:extrema}, no such conclusion holds: $f(x,0)
= x^3 \to -\infty$, and the local minimum at $(1,1)$ is not
global --- exactness of the quadratic expansion is what failed.

\textbf{6.} Expanding the squared norm:
\[
f(u,v) = u^2\norm{C_1}^2 + 2uv\,\langle C_1, C_2\rangle +
v^2\norm{C_2}^2 - 2u\,\langle C_1, b\rangle - 2v\,\langle C_2,
b\rangle + \norm b^2 ,
\]
a quadratic function of $(u, v)$ whose quadratic part is
$\frac12(ru^2 + 2suv + tv^2)$ with $r = 2\norm{C_1}^2$, $s =
2\langle C_1, C_2\rangle$, $t = 2\norm{C_2}^2$.

\textbf{7.} $rt - s^2 = 4\bigl(\norm{C_1}^2\norm{C_2}^2 -
\langle C_1, C_2\rangle^2\bigr) \geq 0$ by Cauchy--Schwarz
(\cref{thm:b1:euclid:cs}), with equality exactly when $C_1, C_2$
are proportional (or one is zero), i.e.\ when the pair is
linked. So $rt - s^2 > 0 \iff (C_1, C_2)$ \omterm{def:b1:vspaces:free}{free}.

\textbf{8.} $\nabla f = 0$ reads
\[
\norm{C_1}^2 u + \langle C_1, C_2\rangle v = \langle C_1,
b\rangle,
\qquad
\langle C_1, C_2\rangle u + \norm{C_2}^2 v = \langle C_2,
b\rangle,
\]
the \omterm{pb:b1:multivar:1}{normal equations} with the Gram matrix on the left; they say
$\langle uC_1 + vC_2 - b,\ C_i\rangle = 0$ for $i = 1, 2$, i.e.\
$b - p \perp \operatorname{Vect}(C_1, C_2)$ for $p = uC_1 +
vC_2$. By Part I (questions 3, 5) the unique \omterm{thm:b1:multivar:critical}{critical point} is
the unique global minimizer, and by
\cref{thm:b1:euclid:projection} the characterization ``$p \in
F$, $b - p \perp F$'' identifies $p$ as the \omterm{thm:b1:euclid:projection}{orthogonal
projection} of $b$ onto $F = \operatorname{Vect}(C_1, C_2)$.

\textbf{9.} $\norm{b - p}^2 = \norm b^2 - 2\langle b, p\rangle +
\norm p^2$, and $\langle p, b - p\rangle = 0$ gives $\norm p^2 =
\langle p, b\rangle$: the minimum equals $\norm b^2 - \langle p,
b\rangle$. Pythagoras: $\norm b^2 = \norm p^2 + \norm{b - p}^2$
--- the data vector splits into its explained part $p$ and its
residual part $b - p$, \omterm{def:b1:euclid:orthogonal}{orthogonal} to each other.

\textbf{10.} $E(\alpha, \beta) = \norm{\alpha C_1 + \beta C_2 -
b}^2$ with the stated $C_1, C_2, b$; the \omterm{pb:b1:multivar:1}{normal equations} of
question 8 are, entrywise,
\[
n\alpha + \Bigl(\sum x_i\Bigr)\beta = \sum y_i,
\qquad
\Bigl(\sum x_i\Bigr)\alpha + \Bigl(\sum x_i^2\Bigr)\beta = \sum
x_i y_i .
\]

\textbf{11.} Dividing by $n$: $\alpha + \beta\bar x = \bar y$
and $\alpha\bar x + \beta\bigl(v_x + \bar x^2\bigr) = c_{xy} +
\bar x\bar y$. Substituting $\alpha = \bar y - \beta\bar x$ into
the second: $\beta v_x = c_{xy}$, so
\[
\beta = \frac{c_{xy}}{v_x}, \qquad \alpha = \bar y - \beta\bar
x ,
\]
and the first equation says precisely that $(\bar x, \bar y)$
lies on the line.

\textbf{12.} $v_x = \frac1n\sum(x_i - \bar x)^2 \geq 0$, zero
iff every $x_i$ equals $\bar x$. And $rt - s^2 = 4\bigl(n\sum
x_i^2 - (\sum x_i)^2\bigr) = 4n^2 v_x$: question 7's freeness
condition is $v_x > 0$, i.e.\ the $x_i$ not all equal.

\textbf{13.} With $\alpha = \bar y - \beta\bar x$, center the
data ($\tilde x_i = x_i - \bar x$, $\tilde y_i = y_i - \bar y$):
\[
E(\alpha, \beta) = \sum_i(\tilde y_i - \beta\tilde x_i)^2
= n\,v_y - 2\beta\,n\,c_{xy} + \beta^2 n\,v_x ,
\]
minimized at $\beta = c_{xy}/v_x$ with value $E_{\min} =
n\bigl(v_y - \frac{c_{xy}^2}{v_x}\bigr) = n v_y(1 - \rho^2)$.
Since $E_{\min} \geq 0$ and $v_y > 0$: $\rho^2 \leq 1$, and
$\abs\rho = 1$ iff $E_{\min} = 0$, i.e.\ iff every residual
vanishes: the points lie exactly on the line.

\textbf{14.} $n = 4$: $\bar x = \frac64 = 1.5$, $\bar y = 2$,
$\sum x_i^2 = 14$ so $v_x = 3.5 - 2.25 = 1.25$, $\sum x_iy_i =
0 + 1 + 6 + 12 = 19$ so $c_{xy} = 4.75 - 3 = 1.75$. Hence
\[
\beta = \frac{1.75}{1.25} = 1.4,
\qquad
\alpha = 2 - 1.4\times1.5 = -0.1 :
\qquad y = 1.4\,x - 0.1 .
\]
Fitted values $-0.1,\ 1.3,\ 2.7,\ 4.1$; residuals $0.1,\ -0.3,\
0.3,\ -0.1$; $E_{\min} = 0.01 + 0.09 + 0.09 + 0.01 = 0.2$
(check: $v_y = \frac{26}4 - 4 = 2.5$ and $4(2.5 -
\frac{1.75^2}{1.25}) = 4(2.5 - 2.45) = 0.2$).

\textbf{15.} The two \omterm{pb:b1:multivar:1}{normal equations} of question 10 are
exactly $\sum_i\varepsilon_i = 0$ and $\sum_i x_i\varepsilon_i =
0$: the residual vector is \omterm{def:b1:euclid:orthogonal}{orthogonal} to $C_1 = (1, \dots, 1)$
and to $C_2 = (x_i)$ --- to the whole model space. In
particular the optimal line always balances its errors: they
sum to zero.

\textbf{16.} $\frac{\dd}{\dd\alpha}\sum(y_i - \alpha)^2 =
-2\sum(y_i - \alpha) = 0$ gives $\alpha = \bar y$ (and the
second \omterm{def:b1:derivative:def}{derivative} $2n > 0$ makes it the global minimum, the
function being a coercive quadratic in one variable). Minimal
value $\sum(y_i - \bar y)^2 = n v_y$: the variance is the
irreducible quadratic error of a constant model.

\textbf{17.} $\frac{\dd}{\dd\beta}\sum(y_i - \beta x_i)^2 =
-2\sum x_i(y_i - \beta x_i) = 0$ gives $\beta_0 = \frac{\sum
x_iy_i}{\sum x_i^2}$. Writing both slopes over centered
quantities: $\beta_0 = \frac{c_{xy} + \bar x\bar y}{v_x + \bar
x^2}$ equals $\frac{c_{xy}}{v_x}$ iff $c_{xy}\bar x^2 = \bar
x\bar y\,v_x$, i.e.\ iff $\bar x = 0$ (centered abscissas) or
$\bar y = \beta\bar x$ --- the latter meaning $\alpha = 0$: the
full \omterm{pb:b1:multivar:1}{regression line} already passes through the origin.

\textbf{18.} $g(\tau) = \norm{M - P_0}^2 - 2\tau\langle M - P_0,
w\rangle + \tau^2$ is minimal at $\tau^* = \langle M - P_0,
w\rangle$, with value $d(M,D)^2 = \norm{M - P_0}^2 - \langle M -
P_0, w\rangle^2$. In the plane, complete $w$ into an \omterm{def:b1:euclid:orthogonal}{orthonormal}
\omterm{def:b1:vspaces:free}{basis} $(w, n)$ with $n = \frac{(a, b)}{\sqrt{a^2+b^2}}$ (unit
normal of $D$): then $\norm{M - P_0}^2 = \langle M - P_0,
w\rangle^2 + \langle M - P_0, n\rangle^2$, so $d(M, D) =
\abs{\langle M - P_0, n\rangle}$, and with $aP_{0x} + bP_{0y} =
-c$:
\[
d = \frac{\abs{a x_0 + b y_0 + c}}{\sqrt{a^2 + b^2}} .
\]

\textbf{19.} From $E_{\min}/n = v_y - \frac{c_{xy}^2}{v_x}$ and
$\beta^2 v_x = \frac{c_{xy}^2}{v_x}$:
\[
v_y = \beta^2 v_x + \frac{E_{\min}}n :
\]
total variance $=$ variance along the fitted line $+$ residual
variance. (Dividing by $v_y$: $1 = \rho^2 + (1 - \rho^2)$, the
share of variance ``explained'' by the line is $\rho^2$.)

\textbf{20.} The model space is
$\operatorname{Vect}\bigl(C_1, C_2, C_3\bigr)$ with $C_1 = (1)$,
$C_2 = (x_i)$, $C_3 = (x_i^2)$; minimizing $\norm{aC_1 + bC_2 +
cC_3 - b}^2$ leads, exactly as in question 8, to the $3\times3$
Gram system, whose matrix has entries $\langle C_i, C_j\rangle =
\sum_k x_k^{\,i+j-2}$: the displayed moment matrix. It is
invertible iff $(C_1, C_2, C_3)$ is \omterm{def:b1:vspaces:free}{free}
(\cref{exo:b1:euclid:11}); a relation $a + bx_k + cx_k^2 = 0$
for all $k$ makes every $x_k$ a root of one \omterm{def:b1:poly:def}{polynomial} of
degree $\leq 2$, impossible with three distinct values unless $a
= b = c = 0$. Compare the moment matrices of the weekend problem
of \cref{ch:b1:det}, where their determinants were squared
Vandermonde values.

\textbf{21.} First: $r = 2, s = 1, t = 2$, $rt - s^2 = 3 > 0$,
$r > 0$: strict global minimum at the origin. Second: $r = 2, s
= 3, t = 2$, $rt - s^2 = -5 < 0$: saddle. Third: $r = s = t =
2$, $rt - s^2 = 0$: degenerate --- but $x^2 + 2xy + y^2 = (x +
y)^2 \geq 0$ vanishes on the whole line $y = -x$: a
(non-strict) global minimum attained along a line, invisible to
the determinant test.

\textbf{22.} New sums ($n = 5$): $\bar x = 3.2$, $\bar y =
1.6$, $\sum x_iy_i = 19$ so $c_{xy} = 3.8 - 5.12 = -1.32$,
$\sum x_i^2 = 114$ so $v_x = 22.8 - 10.24 = 12.56$. New slope:
\[
\beta = \frac{-1.32}{12.56} \approx -0.105 :
\]
one point turned a clearly increasing trend ($\beta = 1.4$) into
a slightly decreasing one. The squared error charges a residual
$\varepsilon^2$, so a single distant point --- large
$\abs{x_{i} - \bar x}$ \emph{and} large residual --- dominates
both $c_{xy}$ and $v_x$: \omterm{pb:b1:multivar:1}{least squares} is efficient but not
robust.

\textbf{23.} Set $W = \sum w_i$ and define $\bar x_w =
\frac1W\sum w_ix_i$, $\bar y_w$ likewise, $v_x^w = \frac1W\sum
w_ix_i^2 - \bar x_w^2$, $c^w_{xy} = \frac1W\sum w_ix_iy_i - \bar
x_w\bar y_w$. The \omterm{def:b1:logic:map}{map} $\langle u, v\rangle_w = \sum_i w_iu_iv_i$
is an \omterm{def:b1:euclid:def}{inner product} on $\R^n$ ($w_i > 0$ gives definiteness), so
Part II applies verbatim and the \omterm{pb:b1:multivar:1}{normal equations} divide by $W$
into
\[
\alpha + \beta\bar x_w = \bar y_w,
\qquad
\alpha\bar x_w + \beta(v^w_x + \bar x_w^2) = c^w_{xy} + \bar
x_w\bar y_w ,
\]
whence $\beta = c^w_{xy}/v^w_x$ and $\alpha = \bar y_w -
\beta\bar x_w$: the same formulas, with every average weighted.

\textbf{24.} $E_{\min} = \sum\varepsilon_i^2 = 0$ forces every
$\varepsilon_i = 0$: all points exactly on the line; by question
13 this is the case $\abs\rho = 1$. Test: for $(1,1), (2,3),
(3,5)$: $\bar x = 2$, $\bar y = 3$, $v_x = \frac{14}3 - 4 =
\frac23$, $c_{xy} = \frac{22}3 - 6 = \frac43$: $\beta = 2$,
$\alpha = -1$, and indeed $y_i = 2x_i - 1$ for all three points;
$v_y = \frac{35}3 - 9 = \frac83$ and $\rho^2 =
\frac{(4/3)^2}{(2/3)(8/3)} = 1$.

\textbf{25.} (i) For quadratic functions the order-two expansion
is an \emph{identity}, so the sign study of the quadratic form
--- pure algebra, questions 1--2 --- classifies \omterm{thm:b1:multivar:critical}{critical points}
with no admitted Taylor theorem. (ii) The \omterm{pb:b1:multivar:1}{normal equations} say
``residual \omterm{def:b1:euclid:orthogonal}{orthogonal} to the model space'', so the analytic
minimum \emph{is} the \omterm{thm:b1:euclid:projection}{orthogonal projection} of the data vector:
calculus and Euclidean geometry compute the same object. (iii)
The correlation $\rho = c_{xy}/\sqrt{v_xv_y}$ measures the share
$\rho^2$ of the variance explained by the line, and
Cauchy--Schwarz bounds it: $\abs\rho \leq 1$, with equality only
for aligned data. (iv) Bases and freeness (Chapter 18), Gram
and moment matrices and their determinants (Chapters 22--23),
and \omterm{thm:b1:euclid:projection}{orthogonal projection} (Chapter 23) all reappeared as the
working parts of one algorithm. Parts II--III develop the
\emph{method of \omterm{pb:b1:multivar:1}{least squares}} (Legendre, Gauss) through its
\emph{\omterm{pb:b1:multivar:1}{normal equations}}.
\end{solution}



\backmatter
\printindex

\end{document}
